\documentclass[12pt,aspectratio=1610]{beamer}
\geometry{paperwidth=20cm,paperheight=12.5cm}
\usepackage[italian,provide=*]{babel}
\usepackage{lmodern,fontspec}
\IfFontExistsTF{NimbusSans-Regular.otf}{
\setsansfont{NimbusSans-Regular.otf}[BoldFont=NimbusSans-Bold.otf,ItalicFont=NimbusSans-Italic.otf]
\setmainfont{NimbusSans-Regular.otf}[BoldFont=NimbusSans-Bold.otf,ItalicFont=NimbusSans-Italic.otf]
}{\setsansfont{TeX Gyre Heros}\setmainfont{TeX Gyre Heros}}
\IfFontExistsTF{lmmono10-regular.otf}{\setmonofont{lmmono10-regular.otf}[Scale=.88]}{\setmonofont{DejaVu Sans Mono}[Scale=.84]}
\usefonttheme{professionalfonts}
\usepackage{amsmath,amssymb,booktabs,array,tabularx,graphicx,tikz,listings,bookmark}
\usetikzlibrary{calc,arrows.meta,positioning}
\definecolor{Paper}{HTML}{F8F6F1}
\definecolor{Ink}{HTML}{253344}\definecolor{Navy}{HTML}{15283E}
\definecolor{Quantum}{HTML}{5146A5}\definecolor{Classical}{HTML}{087C74}
\definecolor{Measure}{HTML}{A05B10}\definecolor{Rose}{HTML}{9B4564}
\definecolor{Muted}{HTML}{657180}\definecolor{Rule}{HTML}{D6DAD9}
\definecolor{Wash}{HTML}{EAE7F6}\definecolor{Night}{HTML}{112236}
\definecolor{Lilac}{HTML}{ABB5FF}\definecolor{Mint}{HTML}{72D8C8}\definecolor{Gold}{HTML}{EEC184}
\setbeamercolor{background canvas}{bg=Paper}\setbeamercolor{normal text}{fg=Ink}
\setbeamertemplate{navigation symbols}{}\setbeamertemplate{headline}{}\setbeamertemplate{footline}{}
\setbeamersize{text margin left=0pt,text margin right=0pt}\setbeameroption{hide notes}
\newsavebox{\BodyBox}\newsavebox{\SideBox}\newsavebox{\RecallBox}\newsavebox{\TitleBox}
\newcommand{\MeasureBox}[2]{\typeout{GEOBOX:#1:\the\wd#2:\the\ht#2:\the\dp#2}}
\lstset{language=R,basicstyle=\ttfamily\fontsize{10.2}{13.4}\selectfont,
columns=fullflexible,keepspaces=true,showstringspaces=false,breaklines=true,breakatwhitespace=false,
commentstyle=\color{Muted},keywordstyle=\color{Classical},stringstyle=\color{Quantum},
aboveskip=0pt,belowskip=0pt,xleftmargin=0pt,xrightmargin=0pt,tabsize=2}
\hypersetup{colorlinks=true,linkcolor=Classical,urlcolor=Classical,pdfpagemode=UseOutlines,pdfstartview=Fit,bookmarksopen=true,bookmarksopenlevel=0}
\bookmarksetup{open,openlevel=0,depth=2}
\tikzset{arr/.style={-{Stealth[length=2mm]},line width=.8pt,draw=Classical}}
\hypersetup{pdftitle={Campionamento, stima e incertezza},pdfauthor={},pdfsubject={Corso di statistica e data science per Geologia, primo anno; modulo 05}}
\begin{document}
\begin{frame}[plain]\hypertarget{slide1}{}
\bookmark[level=0,page=1]{Campionamento, stima e incertezza}
\begin{tikzpicture}[remember picture,overlay]
\fill[Night] (current page.south west) rectangle (current page.north east);
\draw[Mint,line width=2pt] ([xshift=1.05cm,yshift=-1.35cm]current page.north west)--([xshift=3.0cm,yshift=-1.35cm]current page.north west);
\node[anchor=north west,inner sep=0,text=Mint,font=\bfseries\fontsize{11}{14}\selectfont] at ([xshift=1.05cm,yshift=-1.7cm]current page.north west) {MODULO 05 / 10 · PRIMO ANNO};
\node[anchor=north west,inner sep=0,text=white,text width=15.8cm,align=left,font=\bfseries\fontsize{31}{38}\selectfont] at ([xshift=1.05cm,yshift=-2.65cm]current page.north west) {Campionamento, stima e incertezza};
\node[anchor=north west,inner sep=0,text=Lilac,text width=15.8cm,align=left,font=\fontsize{15}{21}\selectfont] at ([xshift=1.05cm,yshift=-6.4cm]current page.north west) {Statistica e data science per le scienze della Terra};
\node[anchor=north west,inner sep=0,text=white,text width=16.5cm,align=left,font=\fontsize{11}{15}\selectfont] at ([xshift=1.05cm,yshift=-8.6cm]current page.north west) {110 slide · esempi, laboratorio R, esercizi e soluzioni\par\smallskip Dati simulati dichiarati; fonti e provenienza nei materiali allegati.};
\node[anchor=north east,inner sep=0,text=Mint,font=\fontsize{8}{10}\selectfont] at ([xshift=19.05cm,yshift=-.2cm]current page.north west) {\hyperlink{slide3}{Indice}};
\node[anchor=south east,inner sep=0,text=Lilac!45,font=\bfseries\fontsize{64}{68}\selectfont] at ([xshift=18.95cm,yshift=-11.35cm]current page.north west) {05};
\draw[Lilac!30,line width=.35pt] ([xshift=.95cm,yshift=-11.92cm]current page.north west)--([xshift=19.05cm,yshift=-11.92cm]current page.north west);
\node[anchor=north west,inner sep=0pt,text=Lilac,text width=9.0cm,font=\fontsize{7.7}{9.3}\selectfont] at ([xshift=.95cm,yshift=-12.06cm]current page.north west) {Modulo 05 · Stima e incertezza};
\node[anchor=north west,inner sep=0pt,text=Mint,text width=8.2cm,align=right,font=\fontsize{7.3}{9.1}\selectfont] at ([xshift=9.45cm,yshift=-12.06cm]current page.north west) {};
\node[anchor=north east,inner sep=0pt,text=Lilac,font=\fontsize{7.7}{9.3}\selectfont] at ([xshift=19.05cm,yshift=-12.06cm]current page.north west) {1/110};
\end{tikzpicture}\end{frame}
\begin{lrbox}{\TitleBox}\begin{minipage}[t]{18.2cm}\color{Ink}\raggedright\sloppy\fontsize{20.7}{24}\selectfont
\bfseries\color{Navy}Che cosa saprai fare
\end{minipage}\end{lrbox}
\begin{lrbox}{\BodyBox}\begin{minipage}[t]{10.25cm}\color{Ink}\raggedright\sloppy\fontsize{12.4}{16.4}\selectfont
1. Definire popolazione, quantità e unità di campionamento.\par\medskip 2. Distinguere deviazione standard ed errore standard.\par\medskip 3. Interpretare intervalli e bootstrap con i loro limiti.
\end{minipage}\end{lrbox}
\begin{lrbox}{\SideBox}\begin{minipage}[t]{7.45cm}\color{Ink}\raggedright\sloppy\fontsize{12.2}{16.1}\selectfont
Quanto può cambiare una stima in un altro campione?
\end{minipage}\end{lrbox}
\begin{lrbox}{\RecallBox}\begin{minipage}[t]{16.45cm}\color{Ink}\raggedright\sloppy\fontsize{10.1}{13.2}\selectfont
Il criterio è saper motivare e interpretare, non memorizzare una lista di comandi.
\end{minipage}\end{lrbox}
\MeasureBox{2-TitleBox}{\TitleBox}
\MeasureBox{2-BodyBox}{\BodyBox}
\MeasureBox{2-SideBox}{\SideBox}
\MeasureBox{2-RecallBox}{\RecallBox}
\begin{frame}[plain]
\hypertarget{slide2}{}\bookmark[level=1,page=2]{Che cosa saprai fare}
\begin{tikzpicture}[remember picture,overlay]
\node[anchor=north east,inner sep=0pt,text=Classical,font=\fontsize{8}{10}\selectfont] at ([xshift=19.05cm,yshift=-.16cm]current page.north west) {\hyperlink{slide3}{Indice}};
\node[anchor=north west,inner sep=0pt,text=Navy,font=\bfseries] at ([xshift=.95cm,yshift=-.62cm]current page.north west) {\usebox{\TitleBox}};
\draw[Classical,line width=1.3pt] ([xshift=.95cm,yshift=-1.68cm]current page.north west)--([xshift=2.25cm,yshift=-1.68cm]current page.north west);
\node[anchor=north west,inner sep=0pt] at ([xshift=.95cm,yshift=-2.12cm]current page.north west) {\usebox{\BodyBox}};
\node[anchor=north west,inner sep=0pt,text=Quantum,font=\bfseries\fontsize{11.5}{14}\selectfont] at ([xshift=11.6cm,yshift=-2.1cm]current page.north west) {Domanda guida};
\node[anchor=north west,inner sep=0pt] at ([xshift=11.6cm,yshift=-2.75cm]current page.north west) {\usebox{\SideBox}};
\draw[Rule,line width=.5pt] ([xshift=.95cm,yshift=-10.43cm]current page.north west)--([xshift=19.05cm,yshift=-10.43cm]current page.north west);
\node[anchor=north west,inner sep=0pt,text=Classical,font=\bfseries\fontsize{9.4}{12}\selectfont] at ([xshift=.95cm,yshift=-10.66cm]current page.north west) {Richiami};
\node[anchor=north west,inner sep=0pt] at ([xshift=2.6cm,yshift=-10.64cm]current page.north west) {\usebox{\RecallBox}};
\draw[Rule,line width=.35pt] ([xshift=.95cm,yshift=-11.92cm]current page.north west)--([xshift=19.05cm,yshift=-11.92cm]current page.north west);
\node[anchor=north west,inner sep=0pt,text=Muted,text width=9.0cm,font=\fontsize{7.7}{9.3}\selectfont] at ([xshift=.95cm,yshift=-12.06cm]current page.north west) {Modulo 05 · Stima e incertezza};
\node[anchor=north west,inner sep=0pt,text=Classical,text width=8.2cm,align=right,font=\fontsize{7.3}{9.1}\selectfont] at ([xshift=9.45cm,yshift=-12.06cm]current page.north west) {};
\node[anchor=north east,inner sep=0pt,text=Muted,font=\fontsize{7.7}{9.3}\selectfont] at ([xshift=19.05cm,yshift=-12.06cm]current page.north west) {2/110};

\end{tikzpicture}
\note{Il criterio è saper motivare e interpretare, non memorizzare una lista di comandi.}
\end{frame}
\begin{frame}[plain]\hypertarget{slide3}{}\bookmark[level=1,page=3]{Mappa del modulo}\begin{tikzpicture}[remember picture,overlay]
\node[anchor=north west,inner sep=0,text=Navy,font=\bfseries\fontsize{22}{27}\selectfont] at ([xshift=.95cm,yshift=-.65cm]current page.north west) {Mappa del modulo};
\draw[Classical,line width=1.3pt] ([xshift=.95cm,yshift=-1.68cm]current page.north west)--([xshift=2.25cm,yshift=-1.68cm]current page.north west);
\node[anchor=north west,inner sep=0,text=Classical,font=\bfseries\fontsize{13}{16}\selectfont] at ([xshift=0.95cm,yshift=-2.23cm]current page.north west) {01};
\node[anchor=north west,inner sep=0,text=Navy,text width=7.55cm,font=\bfseries\fontsize{12.4}{15.5}\selectfont] at ([xshift=1.67cm,yshift=-2.23cm]current page.north west) {\hyperlink{slide6}{Definire che cosa stimiamo}};
\node[anchor=north west,inner sep=0,text=Muted,font=\fontsize{8.7}{11}\selectfont] at ([xshift=1.67cm,yshift=-2.95cm]current page.north west) {Slide 6–15 · esercizio e soluzione};
\node[anchor=north west,inner sep=0,text=Classical,font=\bfseries\fontsize{13}{16}\selectfont] at ([xshift=0.95cm,yshift=-3.66cm]current page.north west) {02};
\node[anchor=north west,inner sep=0,text=Navy,text width=7.55cm,font=\bfseries\fontsize{12.4}{15.5}\selectfont] at ([xshift=1.67cm,yshift=-3.66cm]current page.north west) {\hyperlink{slide16}{Campionamento e distorsioni}};
\node[anchor=north west,inner sep=0,text=Muted,font=\fontsize{8.7}{11}\selectfont] at ([xshift=1.67cm,yshift=-4.38cm]current page.north west) {Slide 16–25 · esercizio e soluzione};
\node[anchor=north west,inner sep=0,text=Classical,font=\bfseries\fontsize{13}{16}\selectfont] at ([xshift=0.95cm,yshift=-5.09cm]current page.north west) {03};
\node[anchor=north west,inner sep=0,text=Navy,text width=7.55cm,font=\bfseries\fontsize{12.4}{15.5}\selectfont] at ([xshift=1.67cm,yshift=-5.09cm]current page.north west) {\hyperlink{slide26}{Ripetere idealmente il campionamento}};
\node[anchor=north west,inner sep=0,text=Muted,font=\fontsize{8.7}{11}\selectfont] at ([xshift=1.67cm,yshift=-5.81cm]current page.north west) {Slide 26–35 · esercizio e soluzione};
\node[anchor=north west,inner sep=0,text=Classical,font=\bfseries\fontsize{13}{16}\selectfont] at ([xshift=0.95cm,yshift=-6.52cm]current page.north west) {04};
\node[anchor=north west,inner sep=0,text=Navy,text width=7.55cm,font=\bfseries\fontsize{12.4}{15.5}\selectfont] at ([xshift=1.67cm,yshift=-6.52cm]current page.north west) {\hyperlink{slide36}{Legge dei grandi numeri e limite centrale}};
\node[anchor=north west,inner sep=0,text=Muted,font=\fontsize{8.7}{11}\selectfont] at ([xshift=1.67cm,yshift=-7.239999999999999cm]current page.north west) {Slide 36–45 · esercizio e soluzione};
\node[anchor=north west,inner sep=0,text=Classical,font=\bfseries\fontsize{13}{16}\selectfont] at ([xshift=0.95cm,yshift=-7.949999999999999cm]current page.north west) {05};
\node[anchor=north west,inner sep=0,text=Navy,text width=7.55cm,font=\bfseries\fontsize{12.4}{15.5}\selectfont] at ([xshift=1.67cm,yshift=-7.949999999999999cm]current page.north west) {\hyperlink{slide46}{Errore standard}};
\node[anchor=north west,inner sep=0,text=Muted,font=\fontsize{8.7}{11}\selectfont] at ([xshift=1.67cm,yshift=-8.67cm]current page.north west) {Slide 46–55 · esercizio e soluzione};
\node[anchor=north west,inner sep=0,text=Classical,font=\bfseries\fontsize{13}{16}\selectfont] at ([xshift=10.299999999999999cm,yshift=-2.23cm]current page.north west) {06};
\node[anchor=north west,inner sep=0,text=Navy,text width=7.55cm,font=\bfseries\fontsize{12.4}{15.5}\selectfont] at ([xshift=11.02cm,yshift=-2.23cm]current page.north west) {\hyperlink{slide56}{Intervalli di confidenza}};
\node[anchor=north west,inner sep=0,text=Muted,font=\fontsize{8.7}{11}\selectfont] at ([xshift=11.02cm,yshift=-2.95cm]current page.north west) {Slide 56–65 · esercizio e soluzione};
\node[anchor=north west,inner sep=0,text=Classical,font=\bfseries\fontsize{13}{16}\selectfont] at ([xshift=10.299999999999999cm,yshift=-3.66cm]current page.north west) {07};
\node[anchor=north west,inner sep=0,text=Navy,text width=7.55cm,font=\bfseries\fontsize{12.4}{15.5}\selectfont] at ([xshift=11.02cm,yshift=-3.66cm]current page.north west) {\hyperlink{slide66}{Bootstrap: ricampionare il campione}};
\node[anchor=north west,inner sep=0,text=Muted,font=\fontsize{8.7}{11}\selectfont] at ([xshift=11.02cm,yshift=-4.38cm]current page.north west) {Slide 66–75 · esercizio e soluzione};
\node[anchor=north west,inner sep=0,text=Classical,font=\bfseries\fontsize{13}{16}\selectfont] at ([xshift=10.299999999999999cm,yshift=-5.09cm]current page.north west) {08};
\node[anchor=north west,inner sep=0,text=Navy,text width=7.55cm,font=\bfseries\fontsize{12.4}{15.5}\selectfont] at ([xshift=11.02cm,yshift=-5.09cm]current page.north west) {\hyperlink{slide76}{Ricampionare al livello corretto}};
\node[anchor=north west,inner sep=0,text=Muted,font=\fontsize{8.7}{11}\selectfont] at ([xshift=11.02cm,yshift=-5.81cm]current page.north west) {Slide 76–85 · esercizio e soluzione};
\node[anchor=north west,inner sep=0,text=Classical,font=\bfseries\fontsize{13}{16}\selectfont] at ([xshift=10.299999999999999cm,yshift=-6.52cm]current page.north west) {09};
\node[anchor=north west,inner sep=0,text=Navy,text width=7.55cm,font=\bfseries\fontsize{12.4}{15.5}\selectfont] at ([xshift=11.02cm,yshift=-6.52cm]current page.north west) {\hyperlink{slide86}{Precisione e progetto di campionamento}};
\node[anchor=north west,inner sep=0,text=Muted,font=\fontsize{8.7}{11}\selectfont] at ([xshift=11.02cm,yshift=-7.239999999999999cm]current page.north west) {Slide 86–95 · esercizio e soluzione};
\node[anchor=north west,inner sep=0,text=Classical,font=\bfseries\fontsize{13}{16}\selectfont] at ([xshift=10.299999999999999cm,yshift=-7.949999999999999cm]current page.north west) {10};
\node[anchor=north west,inner sep=0,text=Navy,text width=7.55cm,font=\bfseries\fontsize{12.4}{15.5}\selectfont] at ([xshift=11.02cm,yshift=-7.949999999999999cm]current page.north west) {\hyperlink{slide96}{Laboratorio di stima}};
\node[anchor=north west,inner sep=0,text=Muted,font=\fontsize{8.7}{11}\selectfont] at ([xshift=11.02cm,yshift=-8.67cm]current page.north west) {Slide 96–105 · esercizio e soluzione};
\node[anchor=north west,inner sep=0,text=Muted,text width=18cm,font=\fontsize{10}{13}\selectfont] at ([xshift=.95cm,yshift=-10.55cm]current page.north west) {I titoli sono cliccabili. Ogni sezione contiene concetti, esempi, un esercizio e una soluzione commentata.};
\draw[Rule,line width=.35pt] ([xshift=.95cm,yshift=-11.92cm]current page.north west)--([xshift=19.05cm,yshift=-11.92cm]current page.north west);
\node[anchor=north west,inner sep=0pt,text=Muted,text width=9.0cm,font=\fontsize{7.7}{9.3}\selectfont] at ([xshift=.95cm,yshift=-12.06cm]current page.north west) {Percorso e navigazione};
\node[anchor=north west,inner sep=0pt,text=Classical,text width=8.2cm,align=right,font=\fontsize{7.3}{9.1}\selectfont] at ([xshift=9.45cm,yshift=-12.06cm]current page.north west) {};
\node[anchor=north east,inner sep=0pt,text=Muted,font=\fontsize{7.7}{9.3}\selectfont] at ([xshift=19.05cm,yshift=-12.06cm]current page.north west) {3/110};
\end{tikzpicture}\end{frame}
\begin{lrbox}{\TitleBox}\begin{minipage}[t]{18.2cm}\color{Ink}\raggedright\sloppy\fontsize{20.7}{24}\selectfont
\bfseries\color{Navy}Dati, fonti e materiali
\end{minipage}\end{lrbox}
\begin{lrbox}{\BodyBox}\begin{minipage}[t]{10.25cm}\color{Ink}\raggedright\sloppy\fontsize{12.4}{16.4}\selectfont
Quasi tutti gli esempi sono simulazioni originali. Il file di 24 provini riprende il modulo 1. I 160 fossili e le 90 rocce sono nuovi casi didattici, non i dataset reali delle presentazioni iniziali.\par\medskip Ogni file ha un dizionario delle variabili e una nota di provenienza. Le figure sono originali e i risultati numerici di riferimento sono verificati in Python. Gli script R sono predisposti per il laboratorio.
\end{minipage}\end{lrbox}
\begin{lrbox}{\SideBox}\begin{minipage}[t]{7.45cm}\color{Ink}\raggedright\sloppy\fontsize{12.2}{16.1}\selectfont
\noindent{\color{Classical}\bfseries 01}\hspace{.35cm}Dati e dizionario\par\vspace{.22cm}\noindent{\color{Classical}\bfseries 02}\hspace{.35cm}Script R commentato\par\vspace{.22cm}\noindent{\color{Classical}\bfseries 03}\hspace{.35cm}Guida docente\par\vspace{.22cm}\noindent{\color{Classical}\bfseries 04}\hspace{.35cm}Scheda esercizi\par\vspace{.22cm}\noindent{\color{Classical}\bfseries 05}\hspace{.35cm}Fonti e risultati\par\vspace{.22cm}
\end{minipage}\end{lrbox}
\begin{lrbox}{\RecallBox}\begin{minipage}[t]{16.45cm}\color{Ink}\raggedright\sloppy\fontsize{10.1}{13.2}\selectfont
Distinguere sempre osservazioni, simulazioni e trasformazioni didattiche.
\end{minipage}\end{lrbox}
\MeasureBox{4-TitleBox}{\TitleBox}
\MeasureBox{4-BodyBox}{\BodyBox}
\MeasureBox{4-SideBox}{\SideBox}
\MeasureBox{4-RecallBox}{\RecallBox}
\begin{frame}[plain]
\hypertarget{slide4}{}\bookmark[level=1,page=4]{Dati, fonti e materiali}
\begin{tikzpicture}[remember picture,overlay]
\node[anchor=north east,inner sep=0pt,text=Classical,font=\fontsize{8}{10}\selectfont] at ([xshift=19.05cm,yshift=-.16cm]current page.north west) {\hyperlink{slide3}{Indice}};
\node[anchor=north west,inner sep=0pt,text=Navy,font=\bfseries] at ([xshift=.95cm,yshift=-.62cm]current page.north west) {\usebox{\TitleBox}};
\draw[Classical,line width=1.3pt] ([xshift=.95cm,yshift=-1.68cm]current page.north west)--([xshift=2.25cm,yshift=-1.68cm]current page.north west);
\node[anchor=north west,inner sep=0pt] at ([xshift=.95cm,yshift=-2.12cm]current page.north west) {\usebox{\BodyBox}};
\node[anchor=north west,inner sep=0pt,text=Quantum,font=\bfseries\fontsize{11.5}{14}\selectfont] at ([xshift=11.6cm,yshift=-2.1cm]current page.north west) {Passaggi da seguire};
\node[anchor=north west,inner sep=0pt] at ([xshift=11.6cm,yshift=-2.75cm]current page.north west) {\usebox{\SideBox}};
\draw[Rule,line width=.5pt] ([xshift=.95cm,yshift=-10.43cm]current page.north west)--([xshift=19.05cm,yshift=-10.43cm]current page.north west);
\node[anchor=north west,inner sep=0pt,text=Classical,font=\bfseries\fontsize{9.4}{12}\selectfont] at ([xshift=.95cm,yshift=-10.66cm]current page.north west) {Richiami};
\node[anchor=north west,inner sep=0pt] at ([xshift=2.6cm,yshift=-10.64cm]current page.north west) {\usebox{\RecallBox}};
\draw[Rule,line width=.35pt] ([xshift=.95cm,yshift=-11.92cm]current page.north west)--([xshift=19.05cm,yshift=-11.92cm]current page.north west);
\node[anchor=north west,inner sep=0pt,text=Muted,text width=9.0cm,font=\fontsize{7.7}{9.3}\selectfont] at ([xshift=.95cm,yshift=-12.06cm]current page.north west) {Modulo 05 · Stima e incertezza};
\node[anchor=north west,inner sep=0pt,text=Classical,text width=8.2cm,align=right,font=\fontsize{7.3}{9.1}\selectfont] at ([xshift=9.45cm,yshift=-12.06cm]current page.north west) {CORSO};
\node[anchor=north east,inner sep=0pt,text=Muted,font=\fontsize{7.7}{9.3}\selectfont] at ([xshift=19.05cm,yshift=-12.06cm]current page.north west) {4/110};

\end{tikzpicture}
\note{Distinguere sempre osservazioni, simulazioni e trasformazioni didattiche.}
\end{frame}
\begin{lrbox}{\TitleBox}\begin{minipage}[t]{18.2cm}\color{Ink}\raggedright\sloppy\fontsize{20.7}{24}\selectfont
\bfseries\color{Navy}Prima di cominciare
\end{minipage}\end{lrbox}
\begin{lrbox}{\BodyBox}\begin{minipage}[t]{10.25cm}\color{Ink}\raggedright\sloppy\fontsize{12.4}{16.4}\selectfont
Distribuzioni, media, variabilità e simulazione del modulo4. Nessun calcolo integrale richiesto.\par\medskip Usa la prima attività per controllare il punto di partenza. I frammenti di codice nelle slide servono a leggere il procedimento: lo script allegato contiene il percorso completo con controlli e salvataggi.
\end{minipage}\end{lrbox}
\begin{lrbox}{\SideBox}\begin{minipage}[t]{7.45cm}\color{Ink}\raggedright\sloppy\fontsize{12.2}{16.1}\selectfont
Nucleo: interpretazioni e procedure essenziali.\par\medskip Laboratorio: esecuzione e modifica dello script.\par\medskip Approfondimenti: indicati esplicitamente, non richiesti come calcolo formale.
\end{minipage}\end{lrbox}
\begin{lrbox}{\RecallBox}\begin{minipage}[t]{16.45cm}\color{Ink}\raggedright\sloppy\fontsize{10.1}{13.2}\selectfont
Lavora in una cartella del modulo e conserva intatti i file distribuiti.
\end{minipage}\end{lrbox}
\MeasureBox{5-TitleBox}{\TitleBox}
\MeasureBox{5-BodyBox}{\BodyBox}
\MeasureBox{5-SideBox}{\SideBox}
\MeasureBox{5-RecallBox}{\RecallBox}
\begin{frame}[plain]
\hypertarget{slide5}{}\bookmark[level=1,page=5]{Prima di cominciare}
\begin{tikzpicture}[remember picture,overlay]
\node[anchor=north east,inner sep=0pt,text=Classical,font=\fontsize{8}{10}\selectfont] at ([xshift=19.05cm,yshift=-.16cm]current page.north west) {\hyperlink{slide3}{Indice}};
\node[anchor=north west,inner sep=0pt,text=Navy,font=\bfseries] at ([xshift=.95cm,yshift=-.62cm]current page.north west) {\usebox{\TitleBox}};
\draw[Classical,line width=1.3pt] ([xshift=.95cm,yshift=-1.68cm]current page.north west)--([xshift=2.25cm,yshift=-1.68cm]current page.north west);
\node[anchor=north west,inner sep=0pt] at ([xshift=.95cm,yshift=-2.12cm]current page.north west) {\usebox{\BodyBox}};
\node[anchor=north west,inner sep=0pt,text=Quantum,font=\bfseries\fontsize{11.5}{14}\selectfont] at ([xshift=11.6cm,yshift=-2.1cm]current page.north west) {Esempio e interpretazione};
\node[anchor=north west,inner sep=0pt] at ([xshift=11.6cm,yshift=-2.75cm]current page.north west) {\usebox{\SideBox}};
\draw[Rule,line width=.5pt] ([xshift=.95cm,yshift=-10.43cm]current page.north west)--([xshift=19.05cm,yshift=-10.43cm]current page.north west);
\node[anchor=north west,inner sep=0pt,text=Classical,font=\bfseries\fontsize{9.4}{12}\selectfont] at ([xshift=.95cm,yshift=-10.66cm]current page.north west) {Richiami};
\node[anchor=north west,inner sep=0pt] at ([xshift=2.6cm,yshift=-10.64cm]current page.north west) {\usebox{\RecallBox}};
\draw[Rule,line width=.35pt] ([xshift=.95cm,yshift=-11.92cm]current page.north west)--([xshift=19.05cm,yshift=-11.92cm]current page.north west);
\node[anchor=north west,inner sep=0pt,text=Muted,text width=9.0cm,font=\fontsize{7.7}{9.3}\selectfont] at ([xshift=.95cm,yshift=-12.06cm]current page.north west) {Modulo 05 · Stima e incertezza};
\node[anchor=north west,inner sep=0pt,text=Classical,text width=8.2cm,align=right,font=\fontsize{7.3}{9.1}\selectfont] at ([xshift=9.45cm,yshift=-12.06cm]current page.north west) {};
\node[anchor=north east,inner sep=0pt,text=Muted,font=\fontsize{7.7}{9.3}\selectfont] at ([xshift=19.05cm,yshift=-12.06cm]current page.north west) {5/110};

\end{tikzpicture}
\note{Lavora in una cartella del modulo e conserva intatti i file distribuiti.}
\end{frame}
\begin{frame}[plain]\hypertarget{slide6}{}
\bookmark[level=0,page=6]{Definire che cosa stimiamo}
\begin{tikzpicture}[remember picture,overlay]
\fill[Night] (current page.south west) rectangle (current page.north east);
\draw[Mint,line width=2pt] ([xshift=1.05cm,yshift=-1.35cm]current page.north west)--([xshift=3.0cm,yshift=-1.35cm]current page.north west);
\node[anchor=north west,inner sep=0,text=Mint,font=\bfseries\fontsize{11}{14}\selectfont] at ([xshift=1.05cm,yshift=-1.7cm]current page.north west) {SEZIONE 01};
\node[anchor=north west,inner sep=0,text=white,text width=15.8cm,align=left,font=\bfseries\fontsize{31}{38}\selectfont] at ([xshift=1.05cm,yshift=-2.65cm]current page.north west) {Definire che cosa stimiamo};
\node[anchor=north west,inner sep=0,text=Lilac,text width=15.8cm,align=left,font=\fontsize{15}{21}\selectfont] at ([xshift=1.05cm,yshift=-6.4cm]current page.north west) {Quale quantità della popolazione vogliamo conoscere?};
\node[anchor=north west,inner sep=0,text=white,text width=16.5cm,align=left,font=\fontsize{11}{15}\selectfont] at ([xshift=1.05cm,yshift=-8.6cm]current page.north west) {Comprendere il concetto · leggere un esempio · applicare e discutere};
\node[anchor=north east,inner sep=0,text=Mint,font=\fontsize{8}{10}\selectfont] at ([xshift=19.05cm,yshift=-.2cm]current page.north west) {\hyperlink{slide3}{Indice}};
\node[anchor=south east,inner sep=0,text=Lilac!45,font=\bfseries\fontsize{64}{68}\selectfont] at ([xshift=18.95cm,yshift=-11.35cm]current page.north west) {01};
\draw[Lilac!30,line width=.35pt] ([xshift=.95cm,yshift=-11.92cm]current page.north west)--([xshift=19.05cm,yshift=-11.92cm]current page.north west);
\node[anchor=north west,inner sep=0pt,text=Lilac,text width=9.0cm,font=\fontsize{7.7}{9.3}\selectfont] at ([xshift=.95cm,yshift=-12.06cm]current page.north west) {01 · Definire che cosa stimiamo};
\node[anchor=north west,inner sep=0pt,text=Mint,text width=8.2cm,align=right,font=\fontsize{7.3}{9.1}\selectfont] at ([xshift=9.45cm,yshift=-12.06cm]current page.north west) {\href{https://openintrostat.github.io/ims/}{IMS} · CORSO};
\node[anchor=north east,inner sep=0pt,text=Lilac,font=\fontsize{7.7}{9.3}\selectfont] at ([xshift=19.05cm,yshift=-12.06cm]current page.north west) {6/110};
\end{tikzpicture}\end{frame}
\begin{lrbox}{\TitleBox}\begin{minipage}[t]{18.2cm}\color{Ink}\raggedright\sloppy\fontsize{20.7}{24}\selectfont
\bfseries\color{Navy}Dalla descrizione alla stima
\end{minipage}\end{lrbox}
\begin{lrbox}{\BodyBox}\begin{minipage}[t]{10.25cm}\color{Ink}\raggedright\sloppy\fontsize{12.4}{16.4}\selectfont
La media del file descrive esattamente i valori presenti, salvo errori di calcolo o misura. Quando la utilizziamo per parlare di una popolazione più ampia diventa una stima. Questo passaggio richiede un legame giustificato tra campione e popolazione.
\end{minipage}\end{lrbox}
\begin{lrbox}{\SideBox}\begin{minipage}[t]{7.45cm}\color{Ink}\raggedright\sloppy\fontsize{12.2}{16.1}\selectfont
Descrizione: diametro medio dei 160 fossili simulati.\par\smallskip Stima: media del modello generatore idealizzato, se il campionamento è coerente.
\end{minipage}\end{lrbox}
\begin{lrbox}{\RecallBox}\begin{minipage}[t]{16.45cm}\color{Ink}\raggedright\sloppy\fontsize{10.1}{13.2}\selectfont
La stessa cifra può avere un ruolo descrittivo o inferenziale: cambia la domanda.
\end{minipage}\end{lrbox}
\MeasureBox{7-TitleBox}{\TitleBox}
\MeasureBox{7-BodyBox}{\BodyBox}
\MeasureBox{7-SideBox}{\SideBox}
\MeasureBox{7-RecallBox}{\RecallBox}
\begin{frame}[plain]
\hypertarget{slide7}{}\bookmark[level=1,page=7]{Dalla descrizione alla stima}
\begin{tikzpicture}[remember picture,overlay]
\node[anchor=north east,inner sep=0pt,text=Classical,font=\fontsize{8}{10}\selectfont] at ([xshift=19.05cm,yshift=-.16cm]current page.north west) {\hyperlink{slide3}{Indice}};
\node[anchor=north west,inner sep=0pt,text=Navy,font=\bfseries] at ([xshift=.95cm,yshift=-.62cm]current page.north west) {\usebox{\TitleBox}};
\draw[Classical,line width=1.3pt] ([xshift=.95cm,yshift=-1.68cm]current page.north west)--([xshift=2.25cm,yshift=-1.68cm]current page.north west);
\node[anchor=north west,inner sep=0pt] at ([xshift=.95cm,yshift=-2.12cm]current page.north west) {\usebox{\BodyBox}};
\node[anchor=north west,inner sep=0pt,text=Quantum,font=\bfseries\fontsize{11.5}{14}\selectfont] at ([xshift=11.6cm,yshift=-2.1cm]current page.north west) {Esempio e interpretazione};
\node[anchor=north west,inner sep=0pt] at ([xshift=11.6cm,yshift=-2.75cm]current page.north west) {\usebox{\SideBox}};
\draw[Rule,line width=.5pt] ([xshift=.95cm,yshift=-10.43cm]current page.north west)--([xshift=19.05cm,yshift=-10.43cm]current page.north west);
\node[anchor=north west,inner sep=0pt,text=Classical,font=\bfseries\fontsize{9.4}{12}\selectfont] at ([xshift=.95cm,yshift=-10.66cm]current page.north west) {Richiami};
\node[anchor=north west,inner sep=0pt] at ([xshift=2.6cm,yshift=-10.64cm]current page.north west) {\usebox{\RecallBox}};
\draw[Rule,line width=.35pt] ([xshift=.95cm,yshift=-11.92cm]current page.north west)--([xshift=19.05cm,yshift=-11.92cm]current page.north west);
\node[anchor=north west,inner sep=0pt,text=Muted,text width=9.0cm,font=\fontsize{7.7}{9.3}\selectfont] at ([xshift=.95cm,yshift=-12.06cm]current page.north west) {01 · Definire che cosa stimiamo};
\node[anchor=north west,inner sep=0pt,text=Classical,text width=8.2cm,align=right,font=\fontsize{7.3}{9.1}\selectfont] at ([xshift=9.45cm,yshift=-12.06cm]current page.north west) {\href{https://openintrostat.github.io/ims/}{IMS} · CORSO};
\node[anchor=north east,inner sep=0pt,text=Muted,font=\fontsize{7.7}{9.3}\selectfont] at ([xshift=19.05cm,yshift=-12.06cm]current page.north west) {7/110};

\end{tikzpicture}
\note{Idea della slide: La media del file descrive esattamente i valori presenti, salvo errori di calcolo o misura. Quando la utilizziamo per parlare di una popolazione più ampia diventa una stima. Questo passaggio richiede un legame giustificato tra campione e popolazione.\par\medskip Punto da far emergere: La stessa cifra può avere un ruolo descrittivo o inferenziale: cambia la domanda.\par\medskip Fonti della sezione: IMS, CORSO}
\end{frame}
\begin{lrbox}{\TitleBox}\begin{minipage}[t]{18.2cm}\color{Ink}\raggedright\sloppy\fontsize{20.7}{24}\selectfont
\bfseries\color{Navy}Parametro e stimatore
\end{minipage}\end{lrbox}
\begin{lrbox}{\BodyBox}\begin{minipage}[t]{10.25cm}\color{Ink}\raggedright\sloppy\fontsize{12.4}{16.4}\selectfont
Il parametro è la quantità ignota della popolazione; lo stimatore è la regola che calcoliamo sul campione. La stima è il valore numerico ottenuto. Distinguere i tre termini permette di capire perché lo stesso procedimento può dare risultati diversi in campioni diversi.
\end{minipage}\end{lrbox}
\begin{lrbox}{\SideBox}\begin{minipage}[t]{7.45cm}\color{Ink}\raggedright\sloppy\fontsize{12.2}{16.1}\selectfont
Parametro: \(\mu\).\par\smallskip Stimatore: media campionaria.\par\smallskip Stima: 6,355 \(\mu\)m nel file didattico.
\end{minipage}\end{lrbox}
\begin{lrbox}{\RecallBox}\begin{minipage}[t]{16.45cm}\color{Ink}\raggedright\sloppy\fontsize{10.1}{13.2}\selectfont
Una stima non diventa un parametro noto perché è riportata con molti decimali.
\end{minipage}\end{lrbox}
\MeasureBox{8-TitleBox}{\TitleBox}
\MeasureBox{8-BodyBox}{\BodyBox}
\MeasureBox{8-SideBox}{\SideBox}
\MeasureBox{8-RecallBox}{\RecallBox}
\begin{frame}[plain]
\hypertarget{slide8}{}\bookmark[level=1,page=8]{Parametro e stimatore}
\begin{tikzpicture}[remember picture,overlay]
\node[anchor=north east,inner sep=0pt,text=Classical,font=\fontsize{8}{10}\selectfont] at ([xshift=19.05cm,yshift=-.16cm]current page.north west) {\hyperlink{slide3}{Indice}};
\node[anchor=north west,inner sep=0pt,text=Navy,font=\bfseries] at ([xshift=.95cm,yshift=-.62cm]current page.north west) {\usebox{\TitleBox}};
\draw[Classical,line width=1.3pt] ([xshift=.95cm,yshift=-1.68cm]current page.north west)--([xshift=2.25cm,yshift=-1.68cm]current page.north west);
\node[anchor=north west,inner sep=0pt] at ([xshift=.95cm,yshift=-2.12cm]current page.north west) {\usebox{\BodyBox}};
\node[anchor=north west,inner sep=0pt,text=Quantum,font=\bfseries\fontsize{11.5}{14}\selectfont] at ([xshift=11.6cm,yshift=-2.1cm]current page.north west) {Esempio e interpretazione};
\node[anchor=north west,inner sep=0pt] at ([xshift=11.6cm,yshift=-2.75cm]current page.north west) {\usebox{\SideBox}};
\draw[Rule,line width=.5pt] ([xshift=.95cm,yshift=-10.43cm]current page.north west)--([xshift=19.05cm,yshift=-10.43cm]current page.north west);
\node[anchor=north west,inner sep=0pt,text=Classical,font=\bfseries\fontsize{9.4}{12}\selectfont] at ([xshift=.95cm,yshift=-10.66cm]current page.north west) {Richiami};
\node[anchor=north west,inner sep=0pt] at ([xshift=2.6cm,yshift=-10.64cm]current page.north west) {\usebox{\RecallBox}};
\draw[Rule,line width=.35pt] ([xshift=.95cm,yshift=-11.92cm]current page.north west)--([xshift=19.05cm,yshift=-11.92cm]current page.north west);
\node[anchor=north west,inner sep=0pt,text=Muted,text width=9.0cm,font=\fontsize{7.7}{9.3}\selectfont] at ([xshift=.95cm,yshift=-12.06cm]current page.north west) {01 · Definire che cosa stimiamo};
\node[anchor=north west,inner sep=0pt,text=Classical,text width=8.2cm,align=right,font=\fontsize{7.3}{9.1}\selectfont] at ([xshift=9.45cm,yshift=-12.06cm]current page.north west) {\href{https://openintrostat.github.io/ims/}{IMS} · CORSO};
\node[anchor=north east,inner sep=0pt,text=Muted,font=\fontsize{7.7}{9.3}\selectfont] at ([xshift=19.05cm,yshift=-12.06cm]current page.north west) {8/110};

\end{tikzpicture}
\note{Idea della slide: Il parametro è la quantità ignota della popolazione; lo stimatore è la regola che calcoliamo sul campione. La stima è il valore numerico ottenuto. Distinguere i tre termini permette di capire perché lo stesso procedimento può dare risultati diversi in campioni diversi.\par\medskip Punto da far emergere: Una stima non diventa un parametro noto perché è riportata con molti decimali.\par\medskip Fonti della sezione: IMS, CORSO}
\end{frame}
\begin{lrbox}{\TitleBox}\begin{minipage}[t]{18.2cm}\color{Ink}\raggedright\sloppy\fontsize{20.7}{24}\selectfont
\bfseries\color{Navy}Definire la popolazione
\end{minipage}\end{lrbox}
\begin{lrbox}{\BodyBox}\begin{minipage}[t]{10.25cm}\color{Ink}\raggedright\sloppy\fontsize{12.4}{16.4}\selectfont
La popolazione deve essere delimitata per luogo, tempo, tipo di oggetto e protocollo. Tutti i fossili di un livello non coincidono con tutti i fossili di una formazione. Una stima può essere precisa per una popolazione ristretta e non trasferibile a un dominio più ampio.
\end{minipage}\end{lrbox}
\begin{lrbox}{\SideBox}\begin{minipage}[t]{7.45cm}\color{Ink}\raggedright\sloppy\fontsize{12.2}{16.1}\selectfont
«Individui osservabili con questo metodo, in questi campioni e in questo intervallo» è più controllabile di «tutti i fossili».
\end{minipage}\end{lrbox}
\begin{lrbox}{\RecallBox}\begin{minipage}[t]{16.45cm}\color{Ink}\raggedright\sloppy\fontsize{10.1}{13.2}\selectfont
La validità esterna nasce dal disegno dello studio, non dalla sola dimensione del campione.
\end{minipage}\end{lrbox}
\MeasureBox{9-TitleBox}{\TitleBox}
\MeasureBox{9-BodyBox}{\BodyBox}
\MeasureBox{9-SideBox}{\SideBox}
\MeasureBox{9-RecallBox}{\RecallBox}
\begin{frame}[plain]
\hypertarget{slide9}{}\bookmark[level=1,page=9]{Definire la popolazione}
\begin{tikzpicture}[remember picture,overlay]
\node[anchor=north east,inner sep=0pt,text=Classical,font=\fontsize{8}{10}\selectfont] at ([xshift=19.05cm,yshift=-.16cm]current page.north west) {\hyperlink{slide3}{Indice}};
\node[anchor=north west,inner sep=0pt,text=Navy,font=\bfseries] at ([xshift=.95cm,yshift=-.62cm]current page.north west) {\usebox{\TitleBox}};
\draw[Classical,line width=1.3pt] ([xshift=.95cm,yshift=-1.68cm]current page.north west)--([xshift=2.25cm,yshift=-1.68cm]current page.north west);
\node[anchor=north west,inner sep=0pt] at ([xshift=.95cm,yshift=-2.12cm]current page.north west) {\usebox{\BodyBox}};
\node[anchor=north west,inner sep=0pt,text=Quantum,font=\bfseries\fontsize{11.5}{14}\selectfont] at ([xshift=11.6cm,yshift=-2.1cm]current page.north west) {Esempio e interpretazione};
\node[anchor=north west,inner sep=0pt] at ([xshift=11.6cm,yshift=-2.75cm]current page.north west) {\usebox{\SideBox}};
\draw[Rule,line width=.5pt] ([xshift=.95cm,yshift=-10.43cm]current page.north west)--([xshift=19.05cm,yshift=-10.43cm]current page.north west);
\node[anchor=north west,inner sep=0pt,text=Classical,font=\bfseries\fontsize{9.4}{12}\selectfont] at ([xshift=.95cm,yshift=-10.66cm]current page.north west) {Richiami};
\node[anchor=north west,inner sep=0pt] at ([xshift=2.6cm,yshift=-10.64cm]current page.north west) {\usebox{\RecallBox}};
\draw[Rule,line width=.35pt] ([xshift=.95cm,yshift=-11.92cm]current page.north west)--([xshift=19.05cm,yshift=-11.92cm]current page.north west);
\node[anchor=north west,inner sep=0pt,text=Muted,text width=9.0cm,font=\fontsize{7.7}{9.3}\selectfont] at ([xshift=.95cm,yshift=-12.06cm]current page.north west) {01 · Definire che cosa stimiamo};
\node[anchor=north west,inner sep=0pt,text=Classical,text width=8.2cm,align=right,font=\fontsize{7.3}{9.1}\selectfont] at ([xshift=9.45cm,yshift=-12.06cm]current page.north west) {\href{https://openintrostat.github.io/ims/}{IMS} · CORSO};
\node[anchor=north east,inner sep=0pt,text=Muted,font=\fontsize{7.7}{9.3}\selectfont] at ([xshift=19.05cm,yshift=-12.06cm]current page.north west) {9/110};

\end{tikzpicture}
\note{Idea della slide: La popolazione deve essere delimitata per luogo, tempo, tipo di oggetto e protocollo. Tutti i fossili di un livello non coincidono con tutti i fossili di una formazione. Una stima può essere precisa per una popolazione ristretta e non trasferibile a un dominio più ampio.\par\medskip Punto da far emergere: La validità esterna nasce dal disegno dello studio, non dalla sola dimensione del campione.\par\medskip Fonti della sezione: IMS, CORSO}
\end{frame}
\begin{lrbox}{\TitleBox}\begin{minipage}[t]{18.2cm}\color{Ink}\raggedright\sloppy\fontsize{20.7}{24}\selectfont
\bfseries\color{Navy}Media, proporzione e differenza
\end{minipage}\end{lrbox}
\begin{lrbox}{\BodyBox}\begin{minipage}[t]{10.25cm}\color{Ink}\raggedright\sloppy\fontsize{12.4}{16.4}\selectfont
Quantità diverse richiedono stimatori e unità diverse. Una proporzione descrive una quota, una media una grandezza, una differenza un contrasto tra gruppi. La scelta deve precedere il calcolo dell'intervallo, altrimenti rischiamo di quantificare bene la domanda sbagliata.
\end{minipage}\end{lrbox}
\begin{lrbox}{\SideBox}\begin{minipage}[t]{7.45cm}\color{Ink}\raggedright\sloppy\fontsize{12.2}{16.1}\selectfont
Media del diametro: \(\mu\)m.\par\smallskip Quota di Spiky: adimensionale.\par\smallskip Differenza Spiky–Rounded: \(\mu\)m.
\end{minipage}\end{lrbox}
\begin{lrbox}{\RecallBox}\begin{minipage}[t]{16.45cm}\color{Ink}\raggedright\sloppy\fontsize{10.1}{13.2}\selectfont
L'incertezza deve riferirsi alla quantità che interessa scientificamente.
\end{minipage}\end{lrbox}
\MeasureBox{10-TitleBox}{\TitleBox}
\MeasureBox{10-BodyBox}{\BodyBox}
\MeasureBox{10-SideBox}{\SideBox}
\MeasureBox{10-RecallBox}{\RecallBox}
\begin{frame}[plain]
\hypertarget{slide10}{}\bookmark[level=1,page=10]{Media, proporzione e differenza}
\begin{tikzpicture}[remember picture,overlay]
\node[anchor=north east,inner sep=0pt,text=Classical,font=\fontsize{8}{10}\selectfont] at ([xshift=19.05cm,yshift=-.16cm]current page.north west) {\hyperlink{slide3}{Indice}};
\node[anchor=north west,inner sep=0pt,text=Navy,font=\bfseries] at ([xshift=.95cm,yshift=-.62cm]current page.north west) {\usebox{\TitleBox}};
\draw[Classical,line width=1.3pt] ([xshift=.95cm,yshift=-1.68cm]current page.north west)--([xshift=2.25cm,yshift=-1.68cm]current page.north west);
\node[anchor=north west,inner sep=0pt] at ([xshift=.95cm,yshift=-2.12cm]current page.north west) {\usebox{\BodyBox}};
\node[anchor=north west,inner sep=0pt,text=Quantum,font=\bfseries\fontsize{11.5}{14}\selectfont] at ([xshift=11.6cm,yshift=-2.1cm]current page.north west) {Esempio e interpretazione};
\node[anchor=north west,inner sep=0pt] at ([xshift=11.6cm,yshift=-2.75cm]current page.north west) {\usebox{\SideBox}};
\draw[Rule,line width=.5pt] ([xshift=.95cm,yshift=-10.43cm]current page.north west)--([xshift=19.05cm,yshift=-10.43cm]current page.north west);
\node[anchor=north west,inner sep=0pt,text=Classical,font=\bfseries\fontsize{9.4}{12}\selectfont] at ([xshift=.95cm,yshift=-10.66cm]current page.north west) {Richiami};
\node[anchor=north west,inner sep=0pt] at ([xshift=2.6cm,yshift=-10.64cm]current page.north west) {\usebox{\RecallBox}};
\draw[Rule,line width=.35pt] ([xshift=.95cm,yshift=-11.92cm]current page.north west)--([xshift=19.05cm,yshift=-11.92cm]current page.north west);
\node[anchor=north west,inner sep=0pt,text=Muted,text width=9.0cm,font=\fontsize{7.7}{9.3}\selectfont] at ([xshift=.95cm,yshift=-12.06cm]current page.north west) {01 · Definire che cosa stimiamo};
\node[anchor=north west,inner sep=0pt,text=Classical,text width=8.2cm,align=right,font=\fontsize{7.3}{9.1}\selectfont] at ([xshift=9.45cm,yshift=-12.06cm]current page.north west) {\href{https://openintrostat.github.io/ims/}{IMS} · CORSO};
\node[anchor=north east,inner sep=0pt,text=Muted,font=\fontsize{7.7}{9.3}\selectfont] at ([xshift=19.05cm,yshift=-12.06cm]current page.north west) {10/110};

\end{tikzpicture}
\note{Idea della slide: Quantità diverse richiedono stimatori e unità diverse. Una proporzione descrive una quota, una media una grandezza, una differenza un contrasto tra gruppi. La scelta deve precedere il calcolo dell'intervallo, altrimenti rischiamo di quantificare bene la domanda sbagliata.\par\medskip Punto da far emergere: L'incertezza deve riferirsi alla quantità che interessa scientificamente.\par\medskip Fonti della sezione: IMS, CORSO}
\end{frame}
\begin{lrbox}{\TitleBox}\begin{minipage}[t]{18.2cm}\color{Ink}\raggedright\sloppy\fontsize{20.7}{24}\selectfont
\bfseries\color{Navy}Unità primaria e misure ripetute
\end{minipage}\end{lrbox}
\begin{lrbox}{\BodyBox}\begin{minipage}[t]{10.25cm}\color{Ink}\raggedright\sloppy\fontsize{12.4}{16.4}\selectfont
Le misure ripetute sullo stesso oggetto aiutano a conoscere quell'oggetto, ma non equivalgono ad altrettanti oggetti selezionati indipendentemente. Prima di stimare l'incertezza individuiamo il livello a cui avviene la selezione e quello a cui esiste dipendenza.
\end{minipage}\end{lrbox}
\begin{lrbox}{\SideBox}\begin{minipage}[t]{7.45cm}\color{Ink}\raggedright\sloppy\fontsize{12.2}{16.1}\selectfont
6 stazioni \(\times\) 4 provini non garantiscono 24 unità indipendenti per un confronto tra ambienti.
\end{minipage}\end{lrbox}
\begin{lrbox}{\RecallBox}\begin{minipage}[t]{16.45cm}\color{Ink}\raggedright\sloppy\fontsize{10.1}{13.2}\selectfont
Contare righe non basta per definire la numerosità inferenziale.
\end{minipage}\end{lrbox}
\MeasureBox{11-TitleBox}{\TitleBox}
\MeasureBox{11-BodyBox}{\BodyBox}
\MeasureBox{11-SideBox}{\SideBox}
\MeasureBox{11-RecallBox}{\RecallBox}
\begin{frame}[plain]
\hypertarget{slide11}{}\bookmark[level=1,page=11]{Unità primaria e misure ripetute}
\begin{tikzpicture}[remember picture,overlay]
\node[anchor=north east,inner sep=0pt,text=Classical,font=\fontsize{8}{10}\selectfont] at ([xshift=19.05cm,yshift=-.16cm]current page.north west) {\hyperlink{slide3}{Indice}};
\node[anchor=north west,inner sep=0pt,text=Navy,font=\bfseries] at ([xshift=.95cm,yshift=-.62cm]current page.north west) {\usebox{\TitleBox}};
\draw[Classical,line width=1.3pt] ([xshift=.95cm,yshift=-1.68cm]current page.north west)--([xshift=2.25cm,yshift=-1.68cm]current page.north west);
\node[anchor=north west,inner sep=0pt] at ([xshift=.95cm,yshift=-2.12cm]current page.north west) {\usebox{\BodyBox}};
\node[anchor=north west,inner sep=0pt,text=Quantum,font=\bfseries\fontsize{11.5}{14}\selectfont] at ([xshift=11.6cm,yshift=-2.1cm]current page.north west) {Esempio e interpretazione};
\node[anchor=north west,inner sep=0pt] at ([xshift=11.6cm,yshift=-2.75cm]current page.north west) {\usebox{\SideBox}};
\draw[Rule,line width=.5pt] ([xshift=.95cm,yshift=-10.43cm]current page.north west)--([xshift=19.05cm,yshift=-10.43cm]current page.north west);
\node[anchor=north west,inner sep=0pt,text=Classical,font=\bfseries\fontsize{9.4}{12}\selectfont] at ([xshift=.95cm,yshift=-10.66cm]current page.north west) {Richiami};
\node[anchor=north west,inner sep=0pt] at ([xshift=2.6cm,yshift=-10.64cm]current page.north west) {\usebox{\RecallBox}};
\draw[Rule,line width=.35pt] ([xshift=.95cm,yshift=-11.92cm]current page.north west)--([xshift=19.05cm,yshift=-11.92cm]current page.north west);
\node[anchor=north west,inner sep=0pt,text=Muted,text width=9.0cm,font=\fontsize{7.7}{9.3}\selectfont] at ([xshift=.95cm,yshift=-12.06cm]current page.north west) {01 · Definire che cosa stimiamo};
\node[anchor=north west,inner sep=0pt,text=Classical,text width=8.2cm,align=right,font=\fontsize{7.3}{9.1}\selectfont] at ([xshift=9.45cm,yshift=-12.06cm]current page.north west) {\href{https://openintrostat.github.io/ims/}{IMS} · CORSO};
\node[anchor=north east,inner sep=0pt,text=Muted,font=\fontsize{7.7}{9.3}\selectfont] at ([xshift=19.05cm,yshift=-12.06cm]current page.north west) {11/110};

\end{tikzpicture}
\note{Idea della slide: Le misure ripetute sullo stesso oggetto aiutano a conoscere quell'oggetto, ma non equivalgono ad altrettanti oggetti selezionati indipendentemente. Prima di stimare l'incertezza individuiamo il livello a cui avviene la selezione e quello a cui esiste dipendenza.\par\medskip Punto da far emergere: Contare righe non basta per definire la numerosità inferenziale.\par\medskip Fonti della sezione: IMS, CORSO}
\end{frame}
\begin{lrbox}{\TitleBox}\begin{minipage}[t]{18.2cm}\color{Ink}\raggedright\sloppy\fontsize{20.7}{24}\selectfont
\bfseries\color{Navy}Stimare non significa spiegare
\end{minipage}\end{lrbox}
\begin{lrbox}{\BodyBox}\begin{minipage}[t]{10.25cm}\color{Ink}\raggedright\sloppy\fontsize{12.4}{16.4}\selectfont
Un intervallo per una media o una differenza quantifica incertezza sotto un modello. Non identifica automaticamente il processo che genera il risultato. Una differenza tra stazioni può essere compatibile con molte spiegazioni geologiche e con limiti del campionamento.
\end{minipage}\end{lrbox}
\begin{lrbox}{\SideBox}\begin{minipage}[t]{7.45cm}\color{Ink}\raggedright\sloppy\fontsize{12.2}{16.1}\selectfont
Una media maggiore in un gruppo può riflettere composizione, selezione, conservazione o misura. La stima non distingue queste possibilità da sola.
\end{minipage}\end{lrbox}
\begin{lrbox}{\RecallBox}\begin{minipage}[t]{16.45cm}\color{Ink}\raggedright\sloppy\fontsize{10.1}{13.2}\selectfont
La precisione statistica non equivale alla certezza della spiegazione.
\end{minipage}\end{lrbox}
\MeasureBox{12-TitleBox}{\TitleBox}
\MeasureBox{12-BodyBox}{\BodyBox}
\MeasureBox{12-SideBox}{\SideBox}
\MeasureBox{12-RecallBox}{\RecallBox}
\begin{frame}[plain]
\hypertarget{slide12}{}\bookmark[level=1,page=12]{Stimare non significa spiegare}
\begin{tikzpicture}[remember picture,overlay]
\node[anchor=north east,inner sep=0pt,text=Classical,font=\fontsize{8}{10}\selectfont] at ([xshift=19.05cm,yshift=-.16cm]current page.north west) {\hyperlink{slide3}{Indice}};
\node[anchor=north west,inner sep=0pt,text=Navy,font=\bfseries] at ([xshift=.95cm,yshift=-.62cm]current page.north west) {\usebox{\TitleBox}};
\draw[Classical,line width=1.3pt] ([xshift=.95cm,yshift=-1.68cm]current page.north west)--([xshift=2.25cm,yshift=-1.68cm]current page.north west);
\node[anchor=north west,inner sep=0pt] at ([xshift=.95cm,yshift=-2.12cm]current page.north west) {\usebox{\BodyBox}};
\node[anchor=north west,inner sep=0pt,text=Quantum,font=\bfseries\fontsize{11.5}{14}\selectfont] at ([xshift=11.6cm,yshift=-2.1cm]current page.north west) {Esempio e interpretazione};
\node[anchor=north west,inner sep=0pt] at ([xshift=11.6cm,yshift=-2.75cm]current page.north west) {\usebox{\SideBox}};
\draw[Rule,line width=.5pt] ([xshift=.95cm,yshift=-10.43cm]current page.north west)--([xshift=19.05cm,yshift=-10.43cm]current page.north west);
\node[anchor=north west,inner sep=0pt,text=Classical,font=\bfseries\fontsize{9.4}{12}\selectfont] at ([xshift=.95cm,yshift=-10.66cm]current page.north west) {Richiami};
\node[anchor=north west,inner sep=0pt] at ([xshift=2.6cm,yshift=-10.64cm]current page.north west) {\usebox{\RecallBox}};
\draw[Rule,line width=.35pt] ([xshift=.95cm,yshift=-11.92cm]current page.north west)--([xshift=19.05cm,yshift=-11.92cm]current page.north west);
\node[anchor=north west,inner sep=0pt,text=Muted,text width=9.0cm,font=\fontsize{7.7}{9.3}\selectfont] at ([xshift=.95cm,yshift=-12.06cm]current page.north west) {01 · Definire che cosa stimiamo};
\node[anchor=north west,inner sep=0pt,text=Classical,text width=8.2cm,align=right,font=\fontsize{7.3}{9.1}\selectfont] at ([xshift=9.45cm,yshift=-12.06cm]current page.north west) {\href{https://openintrostat.github.io/ims/}{IMS} · CORSO};
\node[anchor=north east,inner sep=0pt,text=Muted,font=\fontsize{7.7}{9.3}\selectfont] at ([xshift=19.05cm,yshift=-12.06cm]current page.north west) {12/110};

\end{tikzpicture}
\note{Idea della slide: Un intervallo per una media o una differenza quantifica incertezza sotto un modello. Non identifica automaticamente il processo che genera il risultato. Una differenza tra stazioni può essere compatibile con molte spiegazioni geologiche e con limiti del campionamento.\par\medskip Punto da far emergere: La precisione statistica non equivale alla certezza della spiegazione.\par\medskip Fonti della sezione: IMS, CORSO}
\end{frame}
\begin{lrbox}{\TitleBox}\begin{minipage}[t]{18.2cm}\color{Ink}\raggedright\sloppy\fontsize{20.7}{24}\selectfont
\bfseries\color{Navy}Una domanda abbastanza precisa
\end{minipage}\end{lrbox}
\begin{lrbox}{\BodyBox}\begin{minipage}[t]{10.25cm}\color{Ink}\raggedright\sloppy\fontsize{12.4}{16.4}\selectfont
La domanda deve specificare quantità, popolazione e contrasto. «I fossili sono diversi?» è troppo generica per scegliere una procedura. «Quanto differisce il diametro medio tra due morfotipi nel campionamento definito?» rende espliciti oggetto della stima e unità.
\end{minipage}\end{lrbox}
\begin{lrbox}{\SideBox}\begin{minipage}[t]{7.45cm}\color{Ink}\raggedright\sloppy\fontsize{12.2}{16.1}\selectfont
\noindent{\color{Classical}\bfseries 01}\hspace{.35cm}Popolazione\par\vspace{.22cm}\noindent{\color{Classical}\bfseries 02}\hspace{.35cm}Quantità\par\vspace{.22cm}\noindent{\color{Classical}\bfseries 03}\hspace{.35cm}Stimatore\par\vspace{.22cm}\noindent{\color{Classical}\bfseries 04}\hspace{.35cm}Incertezza\par\vspace{.22cm}\noindent{\color{Classical}\bfseries 05}\hspace{.35cm}Limiti\par\vspace{.22cm}
\end{minipage}\end{lrbox}
\begin{lrbox}{\RecallBox}\begin{minipage}[t]{16.45cm}\color{Ink}\raggedright\sloppy\fontsize{10.1}{13.2}\selectfont
Un buon intervallo inizia con una quantità ben definita.
\end{minipage}\end{lrbox}
\MeasureBox{13-TitleBox}{\TitleBox}
\MeasureBox{13-BodyBox}{\BodyBox}
\MeasureBox{13-SideBox}{\SideBox}
\MeasureBox{13-RecallBox}{\RecallBox}
\begin{frame}[plain]
\hypertarget{slide13}{}\bookmark[level=1,page=13]{Una domanda abbastanza precisa}
\begin{tikzpicture}[remember picture,overlay]
\node[anchor=north east,inner sep=0pt,text=Classical,font=\fontsize{8}{10}\selectfont] at ([xshift=19.05cm,yshift=-.16cm]current page.north west) {\hyperlink{slide3}{Indice}};
\node[anchor=north west,inner sep=0pt,text=Navy,font=\bfseries] at ([xshift=.95cm,yshift=-.62cm]current page.north west) {\usebox{\TitleBox}};
\draw[Classical,line width=1.3pt] ([xshift=.95cm,yshift=-1.68cm]current page.north west)--([xshift=2.25cm,yshift=-1.68cm]current page.north west);
\node[anchor=north west,inner sep=0pt] at ([xshift=.95cm,yshift=-2.12cm]current page.north west) {\usebox{\BodyBox}};
\node[anchor=north west,inner sep=0pt,text=Quantum,font=\bfseries\fontsize{11.5}{14}\selectfont] at ([xshift=11.6cm,yshift=-2.1cm]current page.north west) {Passaggi da seguire};
\node[anchor=north west,inner sep=0pt] at ([xshift=11.6cm,yshift=-2.75cm]current page.north west) {\usebox{\SideBox}};
\draw[Rule,line width=.5pt] ([xshift=.95cm,yshift=-10.43cm]current page.north west)--([xshift=19.05cm,yshift=-10.43cm]current page.north west);
\node[anchor=north west,inner sep=0pt,text=Classical,font=\bfseries\fontsize{9.4}{12}\selectfont] at ([xshift=.95cm,yshift=-10.66cm]current page.north west) {Richiami};
\node[anchor=north west,inner sep=0pt] at ([xshift=2.6cm,yshift=-10.64cm]current page.north west) {\usebox{\RecallBox}};
\draw[Rule,line width=.35pt] ([xshift=.95cm,yshift=-11.92cm]current page.north west)--([xshift=19.05cm,yshift=-11.92cm]current page.north west);
\node[anchor=north west,inner sep=0pt,text=Muted,text width=9.0cm,font=\fontsize{7.7}{9.3}\selectfont] at ([xshift=.95cm,yshift=-12.06cm]current page.north west) {01 · Definire che cosa stimiamo};
\node[anchor=north west,inner sep=0pt,text=Classical,text width=8.2cm,align=right,font=\fontsize{7.3}{9.1}\selectfont] at ([xshift=9.45cm,yshift=-12.06cm]current page.north west) {\href{https://openintrostat.github.io/ims/}{IMS} · CORSO};
\node[anchor=north east,inner sep=0pt,text=Muted,font=\fontsize{7.7}{9.3}\selectfont] at ([xshift=19.05cm,yshift=-12.06cm]current page.north west) {13/110};

\end{tikzpicture}
\note{Idea della slide: La domanda deve specificare quantità, popolazione e contrasto. «I fossili sono diversi?» è troppo generica per scegliere una procedura. «Quanto differisce il diametro medio tra due morfotipi nel campionamento definito?» rende espliciti oggetto della stima e unità.\par\medskip Punto da far emergere: Un buon intervallo inizia con una quantità ben definita.\par\medskip Fonti della sezione: IMS, CORSO}
\end{frame}
\begin{lrbox}{\TitleBox}\begin{minipage}[t]{18.2cm}\color{Ink}\raggedright\sloppy\fontsize{20.7}{24}\selectfont
\bfseries\color{Navy}Attività 01 · applica il concetto
\end{minipage}\end{lrbox}
\begin{lrbox}{\BodyBox}\begin{minipage}[t]{10.25cm}\color{Ink}\raggedright\sloppy\fontsize{12.4}{16.4}\selectfont
Per una campagna con molte misure per stazione, distingui media tra tutti i provini e media tra stazioni. Scrivi quale quantità useresti per descrivere una stazione tipica.
\end{minipage}\end{lrbox}
\begin{lrbox}{\SideBox}\begin{minipage}[t]{7.45cm}\color{Ink}\raggedright\sloppy\fontsize{12.2}{16.1}\selectfont
\noindent{\color{Classical}\bfseries 01}\hspace{.35cm}Individua la domanda\par\vspace{.22cm}\noindent{\color{Classical}\bfseries 02}\hspace{.35cm}Esplicita le ipotesi\par\vspace{.22cm}\noindent{\color{Classical}\bfseries 03}\hspace{.35cm}Calcola o rappresenta\par\vspace{.22cm}\noindent{\color{Classical}\bfseries 04}\hspace{.35cm}Scrivi una conclusione\par\vspace{.22cm}\noindent{\color{Classical}\bfseries 05}\hspace{.35cm}Indica un limite\par\vspace{.22cm}
\end{minipage}\end{lrbox}
\begin{lrbox}{\RecallBox}\begin{minipage}[t]{16.45cm}\color{Ink}\raggedright\sloppy\fontsize{10.1}{13.2}\selectfont
Quale quantità della popolazione vogliamo conoscere?
\end{minipage}\end{lrbox}
\MeasureBox{14-TitleBox}{\TitleBox}
\MeasureBox{14-BodyBox}{\BodyBox}
\MeasureBox{14-SideBox}{\SideBox}
\MeasureBox{14-RecallBox}{\RecallBox}
\begin{frame}[plain]
\hypertarget{slide14}{}\bookmark[level=1,page=14]{Attività 01 · applica il concetto}
\begin{tikzpicture}[remember picture,overlay]
\node[anchor=north east,inner sep=0pt,text=Classical,font=\fontsize{8}{10}\selectfont] at ([xshift=19.05cm,yshift=-.16cm]current page.north west) {\hyperlink{slide3}{Indice}};
\node[anchor=north west,inner sep=0pt,text=Navy,font=\bfseries] at ([xshift=.95cm,yshift=-.62cm]current page.north west) {\usebox{\TitleBox}};
\draw[Classical,line width=1.3pt] ([xshift=.95cm,yshift=-1.68cm]current page.north west)--([xshift=2.25cm,yshift=-1.68cm]current page.north west);
\node[anchor=north west,inner sep=0pt] at ([xshift=.95cm,yshift=-2.12cm]current page.north west) {\usebox{\BodyBox}};
\node[anchor=north west,inner sep=0pt,text=Quantum,font=\bfseries\fontsize{11.5}{14}\selectfont] at ([xshift=11.6cm,yshift=-2.1cm]current page.north west) {Consegna di lavoro};
\node[anchor=north west,inner sep=0pt] at ([xshift=11.6cm,yshift=-2.75cm]current page.north west) {\usebox{\SideBox}};
\draw[Rule,line width=.5pt] ([xshift=.95cm,yshift=-10.43cm]current page.north west)--([xshift=19.05cm,yshift=-10.43cm]current page.north west);
\node[anchor=north west,inner sep=0pt,text=Classical,font=\bfseries\fontsize{9.4}{12}\selectfont] at ([xshift=.95cm,yshift=-10.66cm]current page.north west) {Richiami};
\node[anchor=north west,inner sep=0pt] at ([xshift=2.6cm,yshift=-10.64cm]current page.north west) {\usebox{\RecallBox}};
\draw[Rule,line width=.35pt] ([xshift=.95cm,yshift=-11.92cm]current page.north west)--([xshift=19.05cm,yshift=-11.92cm]current page.north west);
\node[anchor=north west,inner sep=0pt,text=Muted,text width=9.0cm,font=\fontsize{7.7}{9.3}\selectfont] at ([xshift=.95cm,yshift=-12.06cm]current page.north west) {01 · Definire che cosa stimiamo};
\node[anchor=north west,inner sep=0pt,text=Classical,text width=8.2cm,align=right,font=\fontsize{7.3}{9.1}\selectfont] at ([xshift=9.45cm,yshift=-12.06cm]current page.north west) {\href{https://openintrostat.github.io/ims/}{IMS} · CORSO};
\node[anchor=north east,inner sep=0pt,text=Muted,font=\fontsize{7.7}{9.3}\selectfont] at ([xshift=19.05cm,yshift=-12.06cm]current page.north west) {14/110};

\end{tikzpicture}
\note{Consegna: Per una campagna con molte misure per stazione, distingui media tra tutti i provini e media tra stazioni. Scrivi quale quantità useresti per descrivere una stazione tipica.\par\medskip Soluzione per il docente: La media tra provini pesa di più le stazioni con più provini. La media delle medie di stazione assegna uguale peso alle stazioni. Per una stazione tipica, la seconda può essere appropriata, ma il disegno di selezione e la popolazione di riferimento devono essere dichiarati.}
\end{frame}
\begin{lrbox}{\TitleBox}\begin{minipage}[t]{18.2cm}\color{Ink}\raggedright\sloppy\fontsize{20.7}{24}\selectfont
\bfseries\color{Navy}Discussione 01 · soluzione ragionata
\end{minipage}\end{lrbox}
\begin{lrbox}{\BodyBox}\begin{minipage}[t]{10.25cm}\color{Ink}\raggedright\sloppy\fontsize{12.4}{16.4}\selectfont
La media tra provini pesa di più le stazioni con più provini. La media delle medie di stazione assegna uguale peso alle stazioni. Per una stazione tipica, la seconda può essere appropriata, ma il disegno di selezione e la popolazione di riferimento devono essere dichiarati.
\end{minipage}\end{lrbox}
\begin{lrbox}{\SideBox}\begin{minipage}[t]{7.45cm}\color{Ink}\raggedright\sloppy\fontsize{12.2}{16.1}\selectfont
Un buon intervallo inizia con una quantità ben definita.
\end{minipage}\end{lrbox}
\begin{lrbox}{\RecallBox}\begin{minipage}[t]{16.45cm}\color{Ink}\raggedright\sloppy\fontsize{10.1}{13.2}\selectfont
Valuta il collegamento tra domanda, procedimento e conclusione, non soltanto il numero finale.
\end{minipage}\end{lrbox}
\MeasureBox{15-TitleBox}{\TitleBox}
\MeasureBox{15-BodyBox}{\BodyBox}
\MeasureBox{15-SideBox}{\SideBox}
\MeasureBox{15-RecallBox}{\RecallBox}
\begin{frame}[plain]
\hypertarget{slide15}{}\bookmark[level=1,page=15]{Discussione 01 · soluzione ragionata}
\begin{tikzpicture}[remember picture,overlay]
\node[anchor=north east,inner sep=0pt,text=Classical,font=\fontsize{8}{10}\selectfont] at ([xshift=19.05cm,yshift=-.16cm]current page.north west) {\hyperlink{slide3}{Indice}};
\node[anchor=north west,inner sep=0pt,text=Navy,font=\bfseries] at ([xshift=.95cm,yshift=-.62cm]current page.north west) {\usebox{\TitleBox}};
\draw[Classical,line width=1.3pt] ([xshift=.95cm,yshift=-1.68cm]current page.north west)--([xshift=2.25cm,yshift=-1.68cm]current page.north west);
\node[anchor=north west,inner sep=0pt] at ([xshift=.95cm,yshift=-2.12cm]current page.north west) {\usebox{\BodyBox}};
\node[anchor=north west,inner sep=0pt,text=Quantum,font=\bfseries\fontsize{11.5}{14}\selectfont] at ([xshift=11.6cm,yshift=-2.1cm]current page.north west) {Criterio di lettura};
\node[anchor=north west,inner sep=0pt] at ([xshift=11.6cm,yshift=-2.75cm]current page.north west) {\usebox{\SideBox}};
\draw[Rule,line width=.5pt] ([xshift=.95cm,yshift=-10.43cm]current page.north west)--([xshift=19.05cm,yshift=-10.43cm]current page.north west);
\node[anchor=north west,inner sep=0pt,text=Classical,font=\bfseries\fontsize{9.4}{12}\selectfont] at ([xshift=.95cm,yshift=-10.66cm]current page.north west) {Richiami};
\node[anchor=north west,inner sep=0pt] at ([xshift=2.6cm,yshift=-10.64cm]current page.north west) {\usebox{\RecallBox}};
\draw[Rule,line width=.35pt] ([xshift=.95cm,yshift=-11.92cm]current page.north west)--([xshift=19.05cm,yshift=-11.92cm]current page.north west);
\node[anchor=north west,inner sep=0pt,text=Muted,text width=9.0cm,font=\fontsize{7.7}{9.3}\selectfont] at ([xshift=.95cm,yshift=-12.06cm]current page.north west) {01 · Definire che cosa stimiamo};
\node[anchor=north west,inner sep=0pt,text=Classical,text width=8.2cm,align=right,font=\fontsize{7.3}{9.1}\selectfont] at ([xshift=9.45cm,yshift=-12.06cm]current page.north west) {\href{https://openintrostat.github.io/ims/}{IMS} · CORSO};
\node[anchor=north east,inner sep=0pt,text=Muted,font=\fontsize{7.7}{9.3}\selectfont] at ([xshift=19.05cm,yshift=-12.06cm]current page.north west) {15/110};

\end{tikzpicture}
\note{La media tra provini pesa di più le stazioni con più provini. La media delle medie di stazione assegna uguale peso alle stazioni. Per una stazione tipica, la seconda può essere appropriata, ma il disegno di selezione e la popolazione di riferimento devono essere dichiarati.}
\end{frame}
\begin{frame}[plain]\hypertarget{slide16}{}
\bookmark[level=0,page=16]{Campionamento e distorsioni}
\begin{tikzpicture}[remember picture,overlay]
\fill[Night] (current page.south west) rectangle (current page.north east);
\draw[Mint,line width=2pt] ([xshift=1.05cm,yshift=-1.35cm]current page.north west)--([xshift=3.0cm,yshift=-1.35cm]current page.north west);
\node[anchor=north west,inner sep=0,text=Mint,font=\bfseries\fontsize{11}{14}\selectfont] at ([xshift=1.05cm,yshift=-1.7cm]current page.north west) {SEZIONE 02};
\node[anchor=north west,inner sep=0,text=white,text width=15.8cm,align=left,font=\bfseries\fontsize{31}{38}\selectfont] at ([xshift=1.05cm,yshift=-2.65cm]current page.north west) {Campionamento e distorsioni};
\node[anchor=north west,inner sep=0,text=Lilac,text width=15.8cm,align=left,font=\fontsize{15}{21}\selectfont] at ([xshift=1.05cm,yshift=-6.4cm]current page.north west) {Che cosa succede se il campione non rappresenta la popolazione?};
\node[anchor=north west,inner sep=0,text=white,text width=16.5cm,align=left,font=\fontsize{11}{15}\selectfont] at ([xshift=1.05cm,yshift=-8.6cm]current page.north west) {Comprendere il concetto · leggere un esempio · applicare e discutere};
\node[anchor=north east,inner sep=0,text=Mint,font=\fontsize{8}{10}\selectfont] at ([xshift=19.05cm,yshift=-.2cm]current page.north west) {\hyperlink{slide3}{Indice}};
\node[anchor=south east,inner sep=0,text=Lilac!45,font=\bfseries\fontsize{64}{68}\selectfont] at ([xshift=18.95cm,yshift=-11.35cm]current page.north west) {02};
\draw[Lilac!30,line width=.35pt] ([xshift=.95cm,yshift=-11.92cm]current page.north west)--([xshift=19.05cm,yshift=-11.92cm]current page.north west);
\node[anchor=north west,inner sep=0pt,text=Lilac,text width=9.0cm,font=\fontsize{7.7}{9.3}\selectfont] at ([xshift=.95cm,yshift=-12.06cm]current page.north west) {02 · Campionamento e distorsioni};
\node[anchor=north west,inner sep=0pt,text=Mint,text width=8.2cm,align=right,font=\fontsize{7.3}{9.1}\selectfont] at ([xshift=9.45cm,yshift=-12.06cm]current page.north west) {\href{https://openintrostat.github.io/ims/}{IMS} · CORSO};
\node[anchor=north east,inner sep=0pt,text=Lilac,font=\fontsize{7.7}{9.3}\selectfont] at ([xshift=19.05cm,yshift=-12.06cm]current page.north west) {16/110};
\end{tikzpicture}\end{frame}
\begin{lrbox}{\TitleBox}\begin{minipage}[t]{18.2cm}\color{Ink}\raggedright\sloppy\fontsize{20.7}{24}\selectfont
\bfseries\color{Navy}Base di campionamento
\end{minipage}\end{lrbox}
\begin{lrbox}{\BodyBox}\begin{minipage}[t]{10.25cm}\color{Ink}\raggedright\sloppy\fontsize{12.4}{16.4}\selectfont
La base è l'insieme delle unità effettivamente raggiungibili dalla selezione. Se esclude alcune parti della popolazione, una selezione casuale al suo interno non ripara quella mancata copertura. Il limite va descritto prima di interpretare l'intervallo statistico.
\end{minipage}\end{lrbox}
\begin{lrbox}{\SideBox}\begin{minipage}[t]{7.45cm}\color{Ink}\raggedright\sloppy\fontsize{12.2}{16.1}\selectfont
Tutti i campioni presenti in una collezione non coincidono necessariamente con tutti i campioni possibili di un'area.
\end{minipage}\end{lrbox}
\begin{lrbox}{\RecallBox}\begin{minipage}[t]{16.45cm}\color{Ink}\raggedright\sloppy\fontsize{10.1}{13.2}\selectfont
Casuale rispetto a una base incompleta non significa rappresentativo dell'intera popolazione.
\end{minipage}\end{lrbox}
\MeasureBox{17-TitleBox}{\TitleBox}
\MeasureBox{17-BodyBox}{\BodyBox}
\MeasureBox{17-SideBox}{\SideBox}
\MeasureBox{17-RecallBox}{\RecallBox}
\begin{frame}[plain]
\hypertarget{slide17}{}\bookmark[level=1,page=17]{Base di campionamento}
\begin{tikzpicture}[remember picture,overlay]
\node[anchor=north east,inner sep=0pt,text=Classical,font=\fontsize{8}{10}\selectfont] at ([xshift=19.05cm,yshift=-.16cm]current page.north west) {\hyperlink{slide3}{Indice}};
\node[anchor=north west,inner sep=0pt,text=Navy,font=\bfseries] at ([xshift=.95cm,yshift=-.62cm]current page.north west) {\usebox{\TitleBox}};
\draw[Classical,line width=1.3pt] ([xshift=.95cm,yshift=-1.68cm]current page.north west)--([xshift=2.25cm,yshift=-1.68cm]current page.north west);
\node[anchor=north west,inner sep=0pt] at ([xshift=.95cm,yshift=-2.12cm]current page.north west) {\usebox{\BodyBox}};
\node[anchor=north west,inner sep=0pt,text=Quantum,font=\bfseries\fontsize{11.5}{14}\selectfont] at ([xshift=11.6cm,yshift=-2.1cm]current page.north west) {Esempio e interpretazione};
\node[anchor=north west,inner sep=0pt] at ([xshift=11.6cm,yshift=-2.75cm]current page.north west) {\usebox{\SideBox}};
\draw[Rule,line width=.5pt] ([xshift=.95cm,yshift=-10.43cm]current page.north west)--([xshift=19.05cm,yshift=-10.43cm]current page.north west);
\node[anchor=north west,inner sep=0pt,text=Classical,font=\bfseries\fontsize{9.4}{12}\selectfont] at ([xshift=.95cm,yshift=-10.66cm]current page.north west) {Richiami};
\node[anchor=north west,inner sep=0pt] at ([xshift=2.6cm,yshift=-10.64cm]current page.north west) {\usebox{\RecallBox}};
\draw[Rule,line width=.35pt] ([xshift=.95cm,yshift=-11.92cm]current page.north west)--([xshift=19.05cm,yshift=-11.92cm]current page.north west);
\node[anchor=north west,inner sep=0pt,text=Muted,text width=9.0cm,font=\fontsize{7.7}{9.3}\selectfont] at ([xshift=.95cm,yshift=-12.06cm]current page.north west) {02 · Campionamento e distorsioni};
\node[anchor=north west,inner sep=0pt,text=Classical,text width=8.2cm,align=right,font=\fontsize{7.3}{9.1}\selectfont] at ([xshift=9.45cm,yshift=-12.06cm]current page.north west) {\href{https://openintrostat.github.io/ims/}{IMS} · CORSO};
\node[anchor=north east,inner sep=0pt,text=Muted,font=\fontsize{7.7}{9.3}\selectfont] at ([xshift=19.05cm,yshift=-12.06cm]current page.north west) {17/110};

\end{tikzpicture}
\note{Idea della slide: La base è l'insieme delle unità effettivamente raggiungibili dalla selezione. Se esclude alcune parti della popolazione, una selezione casuale al suo interno non ripara quella mancata copertura. Il limite va descritto prima di interpretare l'intervallo statistico.\par\medskip Punto da far emergere: Casuale rispetto a una base incompleta non significa rappresentativo dell'intera popolazione.\par\medskip Fonti della sezione: IMS, CORSO}
\end{frame}
\begin{lrbox}{\TitleBox}\begin{minipage}[t]{18.2cm}\color{Ink}\raggedright\sloppy\fontsize{20.7}{24}\selectfont
\bfseries\color{Navy}Campionamento casuale semplice
\end{minipage}\end{lrbox}
\begin{lrbox}{\BodyBox}\begin{minipage}[t]{10.25cm}\color{Ink}\raggedright\sloppy\fontsize{12.4}{16.4}\selectfont
Nel disegno ideale, ogni unità della base ha la stessa probabilità di essere selezionata. Questo rende più semplice collegare la variabilità del campione all'incertezza delle stime. In pratica vanno documentati base, metodo di selezione e unità non osservate.
\end{minipage}\end{lrbox}
\begin{lrbox}{\SideBox}\begin{minipage}[t]{7.45cm}\color{Ink}\raggedright\sloppy\fontsize{12.2}{16.1}\selectfont
Selezionare con un generatore casuale una lista completa di ID è diverso da scegliere a occhio campioni che sembrano casuali.
\end{minipage}\end{lrbox}
\begin{lrbox}{\RecallBox}\begin{minipage}[t]{16.45cm}\color{Ink}\raggedright\sloppy\fontsize{10.1}{13.2}\selectfont
La casualità è una proprietà del procedimento di selezione, non dell'aspetto del campione.
\end{minipage}\end{lrbox}
\MeasureBox{18-TitleBox}{\TitleBox}
\MeasureBox{18-BodyBox}{\BodyBox}
\MeasureBox{18-SideBox}{\SideBox}
\MeasureBox{18-RecallBox}{\RecallBox}
\begin{frame}[plain]
\hypertarget{slide18}{}\bookmark[level=1,page=18]{Campionamento casuale semplice}
\begin{tikzpicture}[remember picture,overlay]
\node[anchor=north east,inner sep=0pt,text=Classical,font=\fontsize{8}{10}\selectfont] at ([xshift=19.05cm,yshift=-.16cm]current page.north west) {\hyperlink{slide3}{Indice}};
\node[anchor=north west,inner sep=0pt,text=Navy,font=\bfseries] at ([xshift=.95cm,yshift=-.62cm]current page.north west) {\usebox{\TitleBox}};
\draw[Classical,line width=1.3pt] ([xshift=.95cm,yshift=-1.68cm]current page.north west)--([xshift=2.25cm,yshift=-1.68cm]current page.north west);
\node[anchor=north west,inner sep=0pt] at ([xshift=.95cm,yshift=-2.12cm]current page.north west) {\usebox{\BodyBox}};
\node[anchor=north west,inner sep=0pt,text=Quantum,font=\bfseries\fontsize{11.5}{14}\selectfont] at ([xshift=11.6cm,yshift=-2.1cm]current page.north west) {Esempio e interpretazione};
\node[anchor=north west,inner sep=0pt] at ([xshift=11.6cm,yshift=-2.75cm]current page.north west) {\usebox{\SideBox}};
\draw[Rule,line width=.5pt] ([xshift=.95cm,yshift=-10.43cm]current page.north west)--([xshift=19.05cm,yshift=-10.43cm]current page.north west);
\node[anchor=north west,inner sep=0pt,text=Classical,font=\bfseries\fontsize{9.4}{12}\selectfont] at ([xshift=.95cm,yshift=-10.66cm]current page.north west) {Richiami};
\node[anchor=north west,inner sep=0pt] at ([xshift=2.6cm,yshift=-10.64cm]current page.north west) {\usebox{\RecallBox}};
\draw[Rule,line width=.35pt] ([xshift=.95cm,yshift=-11.92cm]current page.north west)--([xshift=19.05cm,yshift=-11.92cm]current page.north west);
\node[anchor=north west,inner sep=0pt,text=Muted,text width=9.0cm,font=\fontsize{7.7}{9.3}\selectfont] at ([xshift=.95cm,yshift=-12.06cm]current page.north west) {02 · Campionamento e distorsioni};
\node[anchor=north west,inner sep=0pt,text=Classical,text width=8.2cm,align=right,font=\fontsize{7.3}{9.1}\selectfont] at ([xshift=9.45cm,yshift=-12.06cm]current page.north west) {\href{https://openintrostat.github.io/ims/}{IMS} · CORSO};
\node[anchor=north east,inner sep=0pt,text=Muted,font=\fontsize{7.7}{9.3}\selectfont] at ([xshift=19.05cm,yshift=-12.06cm]current page.north west) {18/110};

\end{tikzpicture}
\note{Idea della slide: Nel disegno ideale, ogni unità della base ha la stessa probabilità di essere selezionata. Questo rende più semplice collegare la variabilità del campione all'incertezza delle stime. In pratica vanno documentati base, metodo di selezione e unità non osservate.\par\medskip Punto da far emergere: La casualità è una proprietà del procedimento di selezione, non dell'aspetto del campione.\par\medskip Fonti della sezione: IMS, CORSO}
\end{frame}
\begin{lrbox}{\TitleBox}\begin{minipage}[t]{18.2cm}\color{Ink}\raggedright\sloppy\fontsize{20.7}{24}\selectfont
\bfseries\color{Navy}Campionamento stratificato
\end{minipage}\end{lrbox}
\begin{lrbox}{\BodyBox}\begin{minipage}[t]{10.25cm}\color{Ink}\raggedright\sloppy\fontsize{12.4}{16.4}\selectfont
Dividere la popolazione in strati consente di garantire copertura di domini importanti. Se le quote campionate differiscono dalle quote della popolazione, una media complessiva richiede pesi appropriati. La semplice unione delle righe può sovrarappresentare gli strati campionati più intensamente.
\end{minipage}\end{lrbox}
\begin{lrbox}{\SideBox}\begin{minipage}[t]{7.45cm}\color{Ink}\raggedright\sloppy\fontsize{12.2}{16.1}\selectfont
Strato A: 30\% della popolazione.\par\smallskip Strato B: 70\%.\par\smallskip Un campione 50–50 non giustifica automaticamente pesi uguali nella media della popolazione.
\end{minipage}\end{lrbox}
\begin{lrbox}{\RecallBox}\begin{minipage}[t]{16.45cm}\color{Ink}\raggedright\sloppy\fontsize{10.1}{13.2}\selectfont
Separare il beneficio di copertura dalla regola usata per aggregare gli strati.
\end{minipage}\end{lrbox}
\MeasureBox{19-TitleBox}{\TitleBox}
\MeasureBox{19-BodyBox}{\BodyBox}
\MeasureBox{19-SideBox}{\SideBox}
\MeasureBox{19-RecallBox}{\RecallBox}
\begin{frame}[plain]
\hypertarget{slide19}{}\bookmark[level=1,page=19]{Campionamento stratificato}
\begin{tikzpicture}[remember picture,overlay]
\node[anchor=north east,inner sep=0pt,text=Classical,font=\fontsize{8}{10}\selectfont] at ([xshift=19.05cm,yshift=-.16cm]current page.north west) {\hyperlink{slide3}{Indice}};
\node[anchor=north west,inner sep=0pt,text=Navy,font=\bfseries] at ([xshift=.95cm,yshift=-.62cm]current page.north west) {\usebox{\TitleBox}};
\draw[Classical,line width=1.3pt] ([xshift=.95cm,yshift=-1.68cm]current page.north west)--([xshift=2.25cm,yshift=-1.68cm]current page.north west);
\node[anchor=north west,inner sep=0pt] at ([xshift=.95cm,yshift=-2.12cm]current page.north west) {\usebox{\BodyBox}};
\node[anchor=north west,inner sep=0pt,text=Quantum,font=\bfseries\fontsize{11.5}{14}\selectfont] at ([xshift=11.6cm,yshift=-2.1cm]current page.north west) {Esempio e interpretazione};
\node[anchor=north west,inner sep=0pt] at ([xshift=11.6cm,yshift=-2.75cm]current page.north west) {\usebox{\SideBox}};
\draw[Rule,line width=.5pt] ([xshift=.95cm,yshift=-10.43cm]current page.north west)--([xshift=19.05cm,yshift=-10.43cm]current page.north west);
\node[anchor=north west,inner sep=0pt,text=Classical,font=\bfseries\fontsize{9.4}{12}\selectfont] at ([xshift=.95cm,yshift=-10.66cm]current page.north west) {Richiami};
\node[anchor=north west,inner sep=0pt] at ([xshift=2.6cm,yshift=-10.64cm]current page.north west) {\usebox{\RecallBox}};
\draw[Rule,line width=.35pt] ([xshift=.95cm,yshift=-11.92cm]current page.north west)--([xshift=19.05cm,yshift=-11.92cm]current page.north west);
\node[anchor=north west,inner sep=0pt,text=Muted,text width=9.0cm,font=\fontsize{7.7}{9.3}\selectfont] at ([xshift=.95cm,yshift=-12.06cm]current page.north west) {02 · Campionamento e distorsioni};
\node[anchor=north west,inner sep=0pt,text=Classical,text width=8.2cm,align=right,font=\fontsize{7.3}{9.1}\selectfont] at ([xshift=9.45cm,yshift=-12.06cm]current page.north west) {\href{https://openintrostat.github.io/ims/}{IMS} · CORSO};
\node[anchor=north east,inner sep=0pt,text=Muted,font=\fontsize{7.7}{9.3}\selectfont] at ([xshift=19.05cm,yshift=-12.06cm]current page.north west) {19/110};

\end{tikzpicture}
\note{Idea della slide: Dividere la popolazione in strati consente di garantire copertura di domini importanti. Se le quote campionate differiscono dalle quote della popolazione, una media complessiva richiede pesi appropriati. La semplice unione delle righe può sovrarappresentare gli strati campionati più intensamente.\par\medskip Punto da far emergere: Separare il beneficio di copertura dalla regola usata per aggregare gli strati.\par\medskip Fonti della sezione: IMS, CORSO}
\end{frame}
\begin{lrbox}{\TitleBox}\begin{minipage}[t]{18.2cm}\color{Ink}\raggedright\sloppy\fontsize{20.7}{24}\selectfont
\bfseries\color{Navy}Selezione per accessibilità
\end{minipage}\end{lrbox}
\begin{lrbox}{\BodyBox}\begin{minipage}[t]{10.25cm}\color{Ink}\raggedright\sloppy\fontsize{12.4}{16.4}\selectfont
Il campionamento comodo può essere utile per una prima esplorazione, ma limita la generalizzazione. Escludere campioni fragili, difficili da raggiungere o poco riconoscibili può modificare la distribuzione osservata. Questi effetti non sono quantificati automaticamente dall'errore standard.
\end{minipage}\end{lrbox}
\begin{lrbox}{\SideBox}\begin{minipage}[t]{7.45cm}\color{Ink}\raggedright\sloppy\fontsize{12.2}{16.1}\selectfont
Misurare soltanto individui ben conservati può descrivere bene quel sottoinsieme e male l'insieme originario.
\end{minipage}\end{lrbox}
\begin{lrbox}{\RecallBox}\begin{minipage}[t]{16.45cm}\color{Ink}\raggedright\sloppy\fontsize{10.1}{13.2}\selectfont
Un intervallo stretto può circondare una stima sistematicamente spostata.
\end{minipage}\end{lrbox}
\MeasureBox{20-TitleBox}{\TitleBox}
\MeasureBox{20-BodyBox}{\BodyBox}
\MeasureBox{20-SideBox}{\SideBox}
\MeasureBox{20-RecallBox}{\RecallBox}
\begin{frame}[plain]
\hypertarget{slide20}{}\bookmark[level=1,page=20]{Selezione per accessibilità}
\begin{tikzpicture}[remember picture,overlay]
\node[anchor=north east,inner sep=0pt,text=Classical,font=\fontsize{8}{10}\selectfont] at ([xshift=19.05cm,yshift=-.16cm]current page.north west) {\hyperlink{slide3}{Indice}};
\node[anchor=north west,inner sep=0pt,text=Navy,font=\bfseries] at ([xshift=.95cm,yshift=-.62cm]current page.north west) {\usebox{\TitleBox}};
\draw[Classical,line width=1.3pt] ([xshift=.95cm,yshift=-1.68cm]current page.north west)--([xshift=2.25cm,yshift=-1.68cm]current page.north west);
\node[anchor=north west,inner sep=0pt] at ([xshift=.95cm,yshift=-2.12cm]current page.north west) {\usebox{\BodyBox}};
\node[anchor=north west,inner sep=0pt,text=Quantum,font=\bfseries\fontsize{11.5}{14}\selectfont] at ([xshift=11.6cm,yshift=-2.1cm]current page.north west) {Esempio e interpretazione};
\node[anchor=north west,inner sep=0pt] at ([xshift=11.6cm,yshift=-2.75cm]current page.north west) {\usebox{\SideBox}};
\draw[Rule,line width=.5pt] ([xshift=.95cm,yshift=-10.43cm]current page.north west)--([xshift=19.05cm,yshift=-10.43cm]current page.north west);
\node[anchor=north west,inner sep=0pt,text=Classical,font=\bfseries\fontsize{9.4}{12}\selectfont] at ([xshift=.95cm,yshift=-10.66cm]current page.north west) {Richiami};
\node[anchor=north west,inner sep=0pt] at ([xshift=2.6cm,yshift=-10.64cm]current page.north west) {\usebox{\RecallBox}};
\draw[Rule,line width=.35pt] ([xshift=.95cm,yshift=-11.92cm]current page.north west)--([xshift=19.05cm,yshift=-11.92cm]current page.north west);
\node[anchor=north west,inner sep=0pt,text=Muted,text width=9.0cm,font=\fontsize{7.7}{9.3}\selectfont] at ([xshift=.95cm,yshift=-12.06cm]current page.north west) {02 · Campionamento e distorsioni};
\node[anchor=north west,inner sep=0pt,text=Classical,text width=8.2cm,align=right,font=\fontsize{7.3}{9.1}\selectfont] at ([xshift=9.45cm,yshift=-12.06cm]current page.north west) {\href{https://openintrostat.github.io/ims/}{IMS} · CORSO};
\node[anchor=north east,inner sep=0pt,text=Muted,font=\fontsize{7.7}{9.3}\selectfont] at ([xshift=19.05cm,yshift=-12.06cm]current page.north west) {20/110};

\end{tikzpicture}
\note{Idea della slide: Il campionamento comodo può essere utile per una prima esplorazione, ma limita la generalizzazione. Escludere campioni fragili, difficili da raggiungere o poco riconoscibili può modificare la distribuzione osservata. Questi effetti non sono quantificati automaticamente dall'errore standard.\par\medskip Punto da far emergere: Un intervallo stretto può circondare una stima sistematicamente spostata.\par\medskip Fonti della sezione: IMS, CORSO}
\end{frame}
\begin{lrbox}{\TitleBox}\begin{minipage}[t]{18.2cm}\color{Ink}\raggedright\sloppy\fontsize{20.7}{24}\selectfont
\bfseries\color{Navy}Non risposta e perdita di campioni
\end{minipage}\end{lrbox}
\begin{lrbox}{\BodyBox}\begin{minipage}[t]{10.25cm}\color{Ink}\raggedright\sloppy\fontsize{12.4}{16.4}\selectfont
La perdita di osservazioni può dipendere dal valore che avremmo misurato. Se gli individui piccoli sono più difficili da identificare, il campione disponibile può risultare selezionato. Il problema riguarda il processo di osservazione, non soltanto il comando na.rm.
\end{minipage}\end{lrbox}
\begin{lrbox}{\SideBox}\begin{minipage}[t]{7.45cm}\color{Ink}\raggedright\sloppy\fontsize{12.2}{16.1}\selectfont
Registro utile: selezionati, ritrovati, misurabili, classificati e inclusi nell'analisi.
\end{minipage}\end{lrbox}
\begin{lrbox}{\RecallBox}\begin{minipage}[t]{16.45cm}\color{Ink}\raggedright\sloppy\fontsize{10.1}{13.2}\selectfont
Le fasi di perdita devono essere documentate separatamente.
\end{minipage}\end{lrbox}
\MeasureBox{21-TitleBox}{\TitleBox}
\MeasureBox{21-BodyBox}{\BodyBox}
\MeasureBox{21-SideBox}{\SideBox}
\MeasureBox{21-RecallBox}{\RecallBox}
\begin{frame}[plain]
\hypertarget{slide21}{}\bookmark[level=1,page=21]{Non risposta e perdita di campioni}
\begin{tikzpicture}[remember picture,overlay]
\node[anchor=north east,inner sep=0pt,text=Classical,font=\fontsize{8}{10}\selectfont] at ([xshift=19.05cm,yshift=-.16cm]current page.north west) {\hyperlink{slide3}{Indice}};
\node[anchor=north west,inner sep=0pt,text=Navy,font=\bfseries] at ([xshift=.95cm,yshift=-.62cm]current page.north west) {\usebox{\TitleBox}};
\draw[Classical,line width=1.3pt] ([xshift=.95cm,yshift=-1.68cm]current page.north west)--([xshift=2.25cm,yshift=-1.68cm]current page.north west);
\node[anchor=north west,inner sep=0pt] at ([xshift=.95cm,yshift=-2.12cm]current page.north west) {\usebox{\BodyBox}};
\node[anchor=north west,inner sep=0pt,text=Quantum,font=\bfseries\fontsize{11.5}{14}\selectfont] at ([xshift=11.6cm,yshift=-2.1cm]current page.north west) {Esempio e interpretazione};
\node[anchor=north west,inner sep=0pt] at ([xshift=11.6cm,yshift=-2.75cm]current page.north west) {\usebox{\SideBox}};
\draw[Rule,line width=.5pt] ([xshift=.95cm,yshift=-10.43cm]current page.north west)--([xshift=19.05cm,yshift=-10.43cm]current page.north west);
\node[anchor=north west,inner sep=0pt,text=Classical,font=\bfseries\fontsize{9.4}{12}\selectfont] at ([xshift=.95cm,yshift=-10.66cm]current page.north west) {Richiami};
\node[anchor=north west,inner sep=0pt] at ([xshift=2.6cm,yshift=-10.64cm]current page.north west) {\usebox{\RecallBox}};
\draw[Rule,line width=.35pt] ([xshift=.95cm,yshift=-11.92cm]current page.north west)--([xshift=19.05cm,yshift=-11.92cm]current page.north west);
\node[anchor=north west,inner sep=0pt,text=Muted,text width=9.0cm,font=\fontsize{7.7}{9.3}\selectfont] at ([xshift=.95cm,yshift=-12.06cm]current page.north west) {02 · Campionamento e distorsioni};
\node[anchor=north west,inner sep=0pt,text=Classical,text width=8.2cm,align=right,font=\fontsize{7.3}{9.1}\selectfont] at ([xshift=9.45cm,yshift=-12.06cm]current page.north west) {\href{https://openintrostat.github.io/ims/}{IMS} · CORSO};
\node[anchor=north east,inner sep=0pt,text=Muted,font=\fontsize{7.7}{9.3}\selectfont] at ([xshift=19.05cm,yshift=-12.06cm]current page.north west) {21/110};

\end{tikzpicture}
\note{Idea della slide: La perdita di osservazioni può dipendere dal valore che avremmo misurato. Se gli individui piccoli sono più difficili da identificare, il campione disponibile può risultare selezionato. Il problema riguarda il processo di osservazione, non soltanto il comando na.rm.\par\medskip Punto da far emergere: Le fasi di perdita devono essere documentate separatamente.\par\medskip Fonti della sezione: IMS, CORSO}
\end{frame}
\begin{lrbox}{\TitleBox}\begin{minipage}[t]{18.2cm}\color{Ink}\raggedright\sloppy\fontsize{20.7}{24}\selectfont
\bfseries\color{Navy}Pseudoreplicazione
\end{minipage}\end{lrbox}
\begin{lrbox}{\BodyBox}\begin{minipage}[t]{10.25cm}\color{Ink}\raggedright\sloppy\fontsize{12.4}{16.4}\selectfont
Trattare misure dipendenti come repliche indipendenti può produrre un'incertezza troppo piccola. La dipendenza può derivare da stesso campione, stazione o livello. Il problema non si risolve rinominando le righe: occorre riflettere il disegno nella sintesi o nel modello.
\end{minipage}\end{lrbox}
\begin{lrbox}{\SideBox}\begin{minipage}[t]{7.45cm}\color{Ink}\raggedright\sloppy\fontsize{12.2}{16.1}\selectfont
20 letture dello stesso provino non sostituiscono 20 provini distinti estratti dalla popolazione.
\end{minipage}\end{lrbox}
\begin{lrbox}{\RecallBox}\begin{minipage}[t]{16.45cm}\color{Ink}\raggedright\sloppy\fontsize{10.1}{13.2}\selectfont
Il numero utile di unità dipende dalla domanda e dalla struttura del disegno.
\end{minipage}\end{lrbox}
\MeasureBox{22-TitleBox}{\TitleBox}
\MeasureBox{22-BodyBox}{\BodyBox}
\MeasureBox{22-SideBox}{\SideBox}
\MeasureBox{22-RecallBox}{\RecallBox}
\begin{frame}[plain]
\hypertarget{slide22}{}\bookmark[level=1,page=22]{Pseudoreplicazione}
\begin{tikzpicture}[remember picture,overlay]
\node[anchor=north east,inner sep=0pt,text=Classical,font=\fontsize{8}{10}\selectfont] at ([xshift=19.05cm,yshift=-.16cm]current page.north west) {\hyperlink{slide3}{Indice}};
\node[anchor=north west,inner sep=0pt,text=Navy,font=\bfseries] at ([xshift=.95cm,yshift=-.62cm]current page.north west) {\usebox{\TitleBox}};
\draw[Classical,line width=1.3pt] ([xshift=.95cm,yshift=-1.68cm]current page.north west)--([xshift=2.25cm,yshift=-1.68cm]current page.north west);
\node[anchor=north west,inner sep=0pt] at ([xshift=.95cm,yshift=-2.12cm]current page.north west) {\usebox{\BodyBox}};
\node[anchor=north west,inner sep=0pt,text=Quantum,font=\bfseries\fontsize{11.5}{14}\selectfont] at ([xshift=11.6cm,yshift=-2.1cm]current page.north west) {Esempio e interpretazione};
\node[anchor=north west,inner sep=0pt] at ([xshift=11.6cm,yshift=-2.75cm]current page.north west) {\usebox{\SideBox}};
\draw[Rule,line width=.5pt] ([xshift=.95cm,yshift=-10.43cm]current page.north west)--([xshift=19.05cm,yshift=-10.43cm]current page.north west);
\node[anchor=north west,inner sep=0pt,text=Classical,font=\bfseries\fontsize{9.4}{12}\selectfont] at ([xshift=.95cm,yshift=-10.66cm]current page.north west) {Richiami};
\node[anchor=north west,inner sep=0pt] at ([xshift=2.6cm,yshift=-10.64cm]current page.north west) {\usebox{\RecallBox}};
\draw[Rule,line width=.35pt] ([xshift=.95cm,yshift=-11.92cm]current page.north west)--([xshift=19.05cm,yshift=-11.92cm]current page.north west);
\node[anchor=north west,inner sep=0pt,text=Muted,text width=9.0cm,font=\fontsize{7.7}{9.3}\selectfont] at ([xshift=.95cm,yshift=-12.06cm]current page.north west) {02 · Campionamento e distorsioni};
\node[anchor=north west,inner sep=0pt,text=Classical,text width=8.2cm,align=right,font=\fontsize{7.3}{9.1}\selectfont] at ([xshift=9.45cm,yshift=-12.06cm]current page.north west) {\href{https://openintrostat.github.io/ims/}{IMS} · CORSO};
\node[anchor=north east,inner sep=0pt,text=Muted,font=\fontsize{7.7}{9.3}\selectfont] at ([xshift=19.05cm,yshift=-12.06cm]current page.north west) {22/110};

\end{tikzpicture}
\note{Idea della slide: Trattare misure dipendenti come repliche indipendenti può produrre un'incertezza troppo piccola. La dipendenza può derivare da stesso campione, stazione o livello. Il problema non si risolve rinominando le righe: occorre riflettere il disegno nella sintesi o nel modello.\par\medskip Punto da far emergere: Il numero utile di unità dipende dalla domanda e dalla struttura del disegno.\par\medskip Fonti della sezione: IMS, CORSO}
\end{frame}
\begin{lrbox}{\TitleBox}\begin{minipage}[t]{18.2cm}\color{Ink}\raggedright\sloppy\fontsize{20.7}{24}\selectfont
\bfseries\color{Navy}Dichiarare il dominio di validità
\end{minipage}\end{lrbox}
\begin{lrbox}{\BodyBox}\begin{minipage}[t]{10.25cm}\color{Ink}\raggedright\sloppy\fontsize{12.4}{16.4}\selectfont
Quando il disegno non consente un'inferenza ampia, possiamo comunque descrivere il campione e formulare ipotesi. La conclusione deve essere calibrata sul supporto osservativo. Dichiarare il limite è un risultato metodologico, non un fallimento dell'analisi.
\end{minipage}\end{lrbox}
\begin{lrbox}{\SideBox}\begin{minipage}[t]{7.45cm}\color{Ink}\raggedright\sloppy\fontsize{12.2}{16.1}\selectfont
«La stima riguarda i campioni accessibili e misurabili secondo il protocollo adottato».
\end{minipage}\end{lrbox}
\begin{lrbox}{\RecallBox}\begin{minipage}[t]{16.45cm}\color{Ink}\raggedright\sloppy\fontsize{10.1}{13.2}\selectfont
Un'affermazione delimitata è più informativa di una generalizzazione non giustificata.
\end{minipage}\end{lrbox}
\MeasureBox{23-TitleBox}{\TitleBox}
\MeasureBox{23-BodyBox}{\BodyBox}
\MeasureBox{23-SideBox}{\SideBox}
\MeasureBox{23-RecallBox}{\RecallBox}
\begin{frame}[plain]
\hypertarget{slide23}{}\bookmark[level=1,page=23]{Dichiarare il dominio di validità}
\begin{tikzpicture}[remember picture,overlay]
\node[anchor=north east,inner sep=0pt,text=Classical,font=\fontsize{8}{10}\selectfont] at ([xshift=19.05cm,yshift=-.16cm]current page.north west) {\hyperlink{slide3}{Indice}};
\node[anchor=north west,inner sep=0pt,text=Navy,font=\bfseries] at ([xshift=.95cm,yshift=-.62cm]current page.north west) {\usebox{\TitleBox}};
\draw[Classical,line width=1.3pt] ([xshift=.95cm,yshift=-1.68cm]current page.north west)--([xshift=2.25cm,yshift=-1.68cm]current page.north west);
\node[anchor=north west,inner sep=0pt] at ([xshift=.95cm,yshift=-2.12cm]current page.north west) {\usebox{\BodyBox}};
\node[anchor=north west,inner sep=0pt,text=Quantum,font=\bfseries\fontsize{11.5}{14}\selectfont] at ([xshift=11.6cm,yshift=-2.1cm]current page.north west) {Esempio e interpretazione};
\node[anchor=north west,inner sep=0pt] at ([xshift=11.6cm,yshift=-2.75cm]current page.north west) {\usebox{\SideBox}};
\draw[Rule,line width=.5pt] ([xshift=.95cm,yshift=-10.43cm]current page.north west)--([xshift=19.05cm,yshift=-10.43cm]current page.north west);
\node[anchor=north west,inner sep=0pt,text=Classical,font=\bfseries\fontsize{9.4}{12}\selectfont] at ([xshift=.95cm,yshift=-10.66cm]current page.north west) {Richiami};
\node[anchor=north west,inner sep=0pt] at ([xshift=2.6cm,yshift=-10.64cm]current page.north west) {\usebox{\RecallBox}};
\draw[Rule,line width=.35pt] ([xshift=.95cm,yshift=-11.92cm]current page.north west)--([xshift=19.05cm,yshift=-11.92cm]current page.north west);
\node[anchor=north west,inner sep=0pt,text=Muted,text width=9.0cm,font=\fontsize{7.7}{9.3}\selectfont] at ([xshift=.95cm,yshift=-12.06cm]current page.north west) {02 · Campionamento e distorsioni};
\node[anchor=north west,inner sep=0pt,text=Classical,text width=8.2cm,align=right,font=\fontsize{7.3}{9.1}\selectfont] at ([xshift=9.45cm,yshift=-12.06cm]current page.north west) {\href{https://openintrostat.github.io/ims/}{IMS} · CORSO};
\node[anchor=north east,inner sep=0pt,text=Muted,font=\fontsize{7.7}{9.3}\selectfont] at ([xshift=19.05cm,yshift=-12.06cm]current page.north west) {23/110};

\end{tikzpicture}
\note{Idea della slide: Quando il disegno non consente un'inferenza ampia, possiamo comunque descrivere il campione e formulare ipotesi. La conclusione deve essere calibrata sul supporto osservativo. Dichiarare il limite è un risultato metodologico, non un fallimento dell'analisi.\par\medskip Punto da far emergere: Un'affermazione delimitata è più informativa di una generalizzazione non giustificata.\par\medskip Fonti della sezione: IMS, CORSO}
\end{frame}
\begin{lrbox}{\TitleBox}\begin{minipage}[t]{18.2cm}\color{Ink}\raggedright\sloppy\fontsize{20.7}{24}\selectfont
\bfseries\color{Navy}Attività 02 · applica il concetto
\end{minipage}\end{lrbox}
\begin{lrbox}{\BodyBox}\begin{minipage}[t]{10.25cm}\color{Ink}\raggedright\sloppy\fontsize{12.4}{16.4}\selectfont
Devi stimare una proprietà lungo una valle, ma puoi campionare solo vicino a una strada. Un campione dieci volte più grande risolve il problema?
\end{minipage}\end{lrbox}
\begin{lrbox}{\SideBox}\begin{minipage}[t]{7.45cm}\color{Ink}\raggedright\sloppy\fontsize{12.2}{16.1}\selectfont
\noindent{\color{Classical}\bfseries 01}\hspace{.35cm}Individua la domanda\par\vspace{.22cm}\noindent{\color{Classical}\bfseries 02}\hspace{.35cm}Esplicita le ipotesi\par\vspace{.22cm}\noindent{\color{Classical}\bfseries 03}\hspace{.35cm}Calcola o rappresenta\par\vspace{.22cm}\noindent{\color{Classical}\bfseries 04}\hspace{.35cm}Scrivi una conclusione\par\vspace{.22cm}\noindent{\color{Classical}\bfseries 05}\hspace{.35cm}Indica un limite\par\vspace{.22cm}
\end{minipage}\end{lrbox}
\begin{lrbox}{\RecallBox}\begin{minipage}[t]{16.45cm}\color{Ink}\raggedright\sloppy\fontsize{10.1}{13.2}\selectfont
Che cosa succede se il campione non rappresenta la popolazione?
\end{minipage}\end{lrbox}
\MeasureBox{24-TitleBox}{\TitleBox}
\MeasureBox{24-BodyBox}{\BodyBox}
\MeasureBox{24-SideBox}{\SideBox}
\MeasureBox{24-RecallBox}{\RecallBox}
\begin{frame}[plain]
\hypertarget{slide24}{}\bookmark[level=1,page=24]{Attività 02 · applica il concetto}
\begin{tikzpicture}[remember picture,overlay]
\node[anchor=north east,inner sep=0pt,text=Classical,font=\fontsize{8}{10}\selectfont] at ([xshift=19.05cm,yshift=-.16cm]current page.north west) {\hyperlink{slide3}{Indice}};
\node[anchor=north west,inner sep=0pt,text=Navy,font=\bfseries] at ([xshift=.95cm,yshift=-.62cm]current page.north west) {\usebox{\TitleBox}};
\draw[Classical,line width=1.3pt] ([xshift=.95cm,yshift=-1.68cm]current page.north west)--([xshift=2.25cm,yshift=-1.68cm]current page.north west);
\node[anchor=north west,inner sep=0pt] at ([xshift=.95cm,yshift=-2.12cm]current page.north west) {\usebox{\BodyBox}};
\node[anchor=north west,inner sep=0pt,text=Quantum,font=\bfseries\fontsize{11.5}{14}\selectfont] at ([xshift=11.6cm,yshift=-2.1cm]current page.north west) {Consegna di lavoro};
\node[anchor=north west,inner sep=0pt] at ([xshift=11.6cm,yshift=-2.75cm]current page.north west) {\usebox{\SideBox}};
\draw[Rule,line width=.5pt] ([xshift=.95cm,yshift=-10.43cm]current page.north west)--([xshift=19.05cm,yshift=-10.43cm]current page.north west);
\node[anchor=north west,inner sep=0pt,text=Classical,font=\bfseries\fontsize{9.4}{12}\selectfont] at ([xshift=.95cm,yshift=-10.66cm]current page.north west) {Richiami};
\node[anchor=north west,inner sep=0pt] at ([xshift=2.6cm,yshift=-10.64cm]current page.north west) {\usebox{\RecallBox}};
\draw[Rule,line width=.35pt] ([xshift=.95cm,yshift=-11.92cm]current page.north west)--([xshift=19.05cm,yshift=-11.92cm]current page.north west);
\node[anchor=north west,inner sep=0pt,text=Muted,text width=9.0cm,font=\fontsize{7.7}{9.3}\selectfont] at ([xshift=.95cm,yshift=-12.06cm]current page.north west) {02 · Campionamento e distorsioni};
\node[anchor=north west,inner sep=0pt,text=Classical,text width=8.2cm,align=right,font=\fontsize{7.3}{9.1}\selectfont] at ([xshift=9.45cm,yshift=-12.06cm]current page.north west) {\href{https://openintrostat.github.io/ims/}{IMS} · CORSO};
\node[anchor=north east,inner sep=0pt,text=Muted,font=\fontsize{7.7}{9.3}\selectfont] at ([xshift=19.05cm,yshift=-12.06cm]current page.north west) {24/110};

\end{tikzpicture}
\note{Consegna: Devi stimare una proprietà lungo una valle, ma puoi campionare solo vicino a una strada. Un campione dieci volte più grande risolve il problema?\par\medskip Soluzione per il docente: Aumentare n vicino alla strada può ridurre la variabilità casuale della stima per quella zona, ma non corregge l'assenza di altre aree. Serve ridefinire il dominio di validità o modificare il disegno di campionamento. Precisione e rappresentatività sono proprietà diverse.}
\end{frame}
\begin{lrbox}{\TitleBox}\begin{minipage}[t]{18.2cm}\color{Ink}\raggedright\sloppy\fontsize{20.7}{24}\selectfont
\bfseries\color{Navy}Discussione 02 · soluzione ragionata
\end{minipage}\end{lrbox}
\begin{lrbox}{\BodyBox}\begin{minipage}[t]{10.25cm}\color{Ink}\raggedright\sloppy\fontsize{12.4}{16.4}\selectfont
Aumentare n vicino alla strada può ridurre la variabilità casuale della stima per quella zona, ma non corregge l'assenza di altre aree. Serve ridefinire il dominio di validità o modificare il disegno di campionamento. Precisione e rappresentatività sono proprietà diverse.
\end{minipage}\end{lrbox}
\begin{lrbox}{\SideBox}\begin{minipage}[t]{7.45cm}\color{Ink}\raggedright\sloppy\fontsize{12.2}{16.1}\selectfont
Un'affermazione delimitata è più informativa di una generalizzazione non giustificata.
\end{minipage}\end{lrbox}
\begin{lrbox}{\RecallBox}\begin{minipage}[t]{16.45cm}\color{Ink}\raggedright\sloppy\fontsize{10.1}{13.2}\selectfont
Valuta il collegamento tra domanda, procedimento e conclusione, non soltanto il numero finale.
\end{minipage}\end{lrbox}
\MeasureBox{25-TitleBox}{\TitleBox}
\MeasureBox{25-BodyBox}{\BodyBox}
\MeasureBox{25-SideBox}{\SideBox}
\MeasureBox{25-RecallBox}{\RecallBox}
\begin{frame}[plain]
\hypertarget{slide25}{}\bookmark[level=1,page=25]{Discussione 02 · soluzione ragionata}
\begin{tikzpicture}[remember picture,overlay]
\node[anchor=north east,inner sep=0pt,text=Classical,font=\fontsize{8}{10}\selectfont] at ([xshift=19.05cm,yshift=-.16cm]current page.north west) {\hyperlink{slide3}{Indice}};
\node[anchor=north west,inner sep=0pt,text=Navy,font=\bfseries] at ([xshift=.95cm,yshift=-.62cm]current page.north west) {\usebox{\TitleBox}};
\draw[Classical,line width=1.3pt] ([xshift=.95cm,yshift=-1.68cm]current page.north west)--([xshift=2.25cm,yshift=-1.68cm]current page.north west);
\node[anchor=north west,inner sep=0pt] at ([xshift=.95cm,yshift=-2.12cm]current page.north west) {\usebox{\BodyBox}};
\node[anchor=north west,inner sep=0pt,text=Quantum,font=\bfseries\fontsize{11.5}{14}\selectfont] at ([xshift=11.6cm,yshift=-2.1cm]current page.north west) {Criterio di lettura};
\node[anchor=north west,inner sep=0pt] at ([xshift=11.6cm,yshift=-2.75cm]current page.north west) {\usebox{\SideBox}};
\draw[Rule,line width=.5pt] ([xshift=.95cm,yshift=-10.43cm]current page.north west)--([xshift=19.05cm,yshift=-10.43cm]current page.north west);
\node[anchor=north west,inner sep=0pt,text=Classical,font=\bfseries\fontsize{9.4}{12}\selectfont] at ([xshift=.95cm,yshift=-10.66cm]current page.north west) {Richiami};
\node[anchor=north west,inner sep=0pt] at ([xshift=2.6cm,yshift=-10.64cm]current page.north west) {\usebox{\RecallBox}};
\draw[Rule,line width=.35pt] ([xshift=.95cm,yshift=-11.92cm]current page.north west)--([xshift=19.05cm,yshift=-11.92cm]current page.north west);
\node[anchor=north west,inner sep=0pt,text=Muted,text width=9.0cm,font=\fontsize{7.7}{9.3}\selectfont] at ([xshift=.95cm,yshift=-12.06cm]current page.north west) {02 · Campionamento e distorsioni};
\node[anchor=north west,inner sep=0pt,text=Classical,text width=8.2cm,align=right,font=\fontsize{7.3}{9.1}\selectfont] at ([xshift=9.45cm,yshift=-12.06cm]current page.north west) {\href{https://openintrostat.github.io/ims/}{IMS} · CORSO};
\node[anchor=north east,inner sep=0pt,text=Muted,font=\fontsize{7.7}{9.3}\selectfont] at ([xshift=19.05cm,yshift=-12.06cm]current page.north west) {25/110};

\end{tikzpicture}
\note{Aumentare n vicino alla strada può ridurre la variabilità casuale della stima per quella zona, ma non corregge l'assenza di altre aree. Serve ridefinire il dominio di validità o modificare il disegno di campionamento. Precisione e rappresentatività sono proprietà diverse.}
\end{frame}
\begin{frame}[plain]\hypertarget{slide26}{}
\bookmark[level=0,page=26]{Ripetere idealmente il campionamento}
\begin{tikzpicture}[remember picture,overlay]
\fill[Night] (current page.south west) rectangle (current page.north east);
\draw[Mint,line width=2pt] ([xshift=1.05cm,yshift=-1.35cm]current page.north west)--([xshift=3.0cm,yshift=-1.35cm]current page.north west);
\node[anchor=north west,inner sep=0,text=Mint,font=\bfseries\fontsize{11}{14}\selectfont] at ([xshift=1.05cm,yshift=-1.7cm]current page.north west) {SEZIONE 03};
\node[anchor=north west,inner sep=0,text=white,text width=15.8cm,align=left,font=\bfseries\fontsize{31}{38}\selectfont] at ([xshift=1.05cm,yshift=-2.65cm]current page.north west) {Ripetere idealmente il campionamento};
\node[anchor=north west,inner sep=0,text=Lilac,text width=15.8cm,align=left,font=\fontsize{15}{21}\selectfont] at ([xshift=1.05cm,yshift=-6.4cm]current page.north west) {Perché due campioni della stessa popolazione danno medie diverse?};
\node[anchor=north west,inner sep=0,text=white,text width=16.5cm,align=left,font=\fontsize{11}{15}\selectfont] at ([xshift=1.05cm,yshift=-8.6cm]current page.north west) {Comprendere il concetto · leggere un esempio · applicare e discutere};
\node[anchor=north east,inner sep=0,text=Mint,font=\fontsize{8}{10}\selectfont] at ([xshift=19.05cm,yshift=-.2cm]current page.north west) {\hyperlink{slide3}{Indice}};
\node[anchor=south east,inner sep=0,text=Lilac!45,font=\bfseries\fontsize{64}{68}\selectfont] at ([xshift=18.95cm,yshift=-11.35cm]current page.north west) {03};
\draw[Lilac!30,line width=.35pt] ([xshift=.95cm,yshift=-11.92cm]current page.north west)--([xshift=19.05cm,yshift=-11.92cm]current page.north west);
\node[anchor=north west,inner sep=0pt,text=Lilac,text width=9.0cm,font=\fontsize{7.7}{9.3}\selectfont] at ([xshift=.95cm,yshift=-12.06cm]current page.north west) {03 · Ripetere idealmente il campionamento};
\node[anchor=north west,inner sep=0pt,text=Mint,text width=8.2cm,align=right,font=\fontsize{7.3}{9.1}\selectfont] at ([xshift=9.45cm,yshift=-12.06cm]current page.north west) {\href{https://openintrostat.github.io/ims/}{IMS} · \href{https://stat.ethz.ch/R-manual/R-devel/library/stats/html/00Index.html}{RSTATS}};
\node[anchor=north east,inner sep=0pt,text=Lilac,font=\fontsize{7.7}{9.3}\selectfont] at ([xshift=19.05cm,yshift=-12.06cm]current page.north west) {26/110};
\end{tikzpicture}\end{frame}
\begin{lrbox}{\TitleBox}\begin{minipage}[t]{18.2cm}\color{Ink}\raggedright\sloppy\fontsize{20.7}{24}\selectfont
\bfseries\color{Navy}Un campione osservato e molti possibili
\end{minipage}\end{lrbox}
\begin{lrbox}{\BodyBox}\begin{minipage}[t]{10.25cm}\color{Ink}\raggedright\sloppy\fontsize{12.4}{16.4}\selectfont
Il dataset disponibile è una realizzazione del procedimento di campionamento. Se potessimo ripetere lo stesso procedimento, otterremmo altri valori e altre statistiche. L'incertezza di una stima descrive questa variabilità ipotetica sotto le condizioni specificate.
\end{minipage}\end{lrbox}
\begin{lrbox}{\SideBox}\begin{minipage}[t]{7.45cm}\color{Ink}\raggedright\sloppy\fontsize{12.2}{16.1}\selectfont
Un nuovo campione dalla stessa popolazione non è una duplicazione del file esistente. È una nuova selezione di unità.
\end{minipage}\end{lrbox}
\begin{lrbox}{\RecallBox}\begin{minipage}[t]{16.45cm}\color{Ink}\raggedright\sloppy\fontsize{10.1}{13.2}\selectfont
La ripetizione ideale mantiene il procedimento, non gli stessi dati.
\end{minipage}\end{lrbox}
\MeasureBox{27-TitleBox}{\TitleBox}
\MeasureBox{27-BodyBox}{\BodyBox}
\MeasureBox{27-SideBox}{\SideBox}
\MeasureBox{27-RecallBox}{\RecallBox}
\begin{frame}[plain]
\hypertarget{slide27}{}\bookmark[level=1,page=27]{Un campione osservato e molti possibili}
\begin{tikzpicture}[remember picture,overlay]
\node[anchor=north east,inner sep=0pt,text=Classical,font=\fontsize{8}{10}\selectfont] at ([xshift=19.05cm,yshift=-.16cm]current page.north west) {\hyperlink{slide3}{Indice}};
\node[anchor=north west,inner sep=0pt,text=Navy,font=\bfseries] at ([xshift=.95cm,yshift=-.62cm]current page.north west) {\usebox{\TitleBox}};
\draw[Classical,line width=1.3pt] ([xshift=.95cm,yshift=-1.68cm]current page.north west)--([xshift=2.25cm,yshift=-1.68cm]current page.north west);
\node[anchor=north west,inner sep=0pt] at ([xshift=.95cm,yshift=-2.12cm]current page.north west) {\usebox{\BodyBox}};
\node[anchor=north west,inner sep=0pt,text=Quantum,font=\bfseries\fontsize{11.5}{14}\selectfont] at ([xshift=11.6cm,yshift=-2.1cm]current page.north west) {Esempio e interpretazione};
\node[anchor=north west,inner sep=0pt] at ([xshift=11.6cm,yshift=-2.75cm]current page.north west) {\usebox{\SideBox}};
\draw[Rule,line width=.5pt] ([xshift=.95cm,yshift=-10.43cm]current page.north west)--([xshift=19.05cm,yshift=-10.43cm]current page.north west);
\node[anchor=north west,inner sep=0pt,text=Classical,font=\bfseries\fontsize{9.4}{12}\selectfont] at ([xshift=.95cm,yshift=-10.66cm]current page.north west) {Richiami};
\node[anchor=north west,inner sep=0pt] at ([xshift=2.6cm,yshift=-10.64cm]current page.north west) {\usebox{\RecallBox}};
\draw[Rule,line width=.35pt] ([xshift=.95cm,yshift=-11.92cm]current page.north west)--([xshift=19.05cm,yshift=-11.92cm]current page.north west);
\node[anchor=north west,inner sep=0pt,text=Muted,text width=9.0cm,font=\fontsize{7.7}{9.3}\selectfont] at ([xshift=.95cm,yshift=-12.06cm]current page.north west) {03 · Ripetere idealmente il campionamento};
\node[anchor=north west,inner sep=0pt,text=Classical,text width=8.2cm,align=right,font=\fontsize{7.3}{9.1}\selectfont] at ([xshift=9.45cm,yshift=-12.06cm]current page.north west) {\href{https://openintrostat.github.io/ims/}{IMS} · \href{https://stat.ethz.ch/R-manual/R-devel/library/stats/html/00Index.html}{RSTATS}};
\node[anchor=north east,inner sep=0pt,text=Muted,font=\fontsize{7.7}{9.3}\selectfont] at ([xshift=19.05cm,yshift=-12.06cm]current page.north west) {27/110};

\end{tikzpicture}
\note{Idea della slide: Il dataset disponibile è una realizzazione del procedimento di campionamento. Se potessimo ripetere lo stesso procedimento, otterremmo altri valori e altre statistiche. L'incertezza di una stima descrive questa variabilità ipotetica sotto le condizioni specificate.\par\medskip Punto da far emergere: La ripetizione ideale mantiene il procedimento, non gli stessi dati.\par\medskip Fonti della sezione: IMS, RSTATS}
\end{frame}
\begin{lrbox}{\TitleBox}\begin{minipage}[t]{18.2cm}\color{Ink}\raggedright\sloppy\fontsize{20.7}{24}\selectfont
\bfseries\color{Navy}Distribuzione campionaria
\end{minipage}\end{lrbox}
\begin{lrbox}{\BodyBox}\begin{minipage}[t]{10.25cm}\color{Ink}\raggedright\sloppy\fontsize{12.4}{16.4}\selectfont
La distribuzione campionaria è la distribuzione di una statistica su campioni ripetuti. Può riguardare media, proporzione o differenza. Non si legge dall'istogramma dei dati individuali: per studiarla servono teoria, simulazione o ricampionamento.
\end{minipage}\end{lrbox}
\begin{lrbox}{\SideBox}\begin{minipage}[t]{7.45cm}\color{Ink}\raggedright\sloppy\fontsize{12.2}{16.1}\selectfont
Dati: diametri dei singoli individui.\par\smallskip Statistica: media del campione.\par\smallskip Distribuzione campionaria: molte medie possibili.
\end{minipage}\end{lrbox}
\begin{lrbox}{\RecallBox}\begin{minipage}[t]{16.45cm}\color{Ink}\raggedright\sloppy\fontsize{10.1}{13.2}\selectfont
Il livello di aggregazione è essenziale per leggere correttamente il grafico.
\end{minipage}\end{lrbox}
\MeasureBox{28-TitleBox}{\TitleBox}
\MeasureBox{28-BodyBox}{\BodyBox}
\MeasureBox{28-SideBox}{\SideBox}
\MeasureBox{28-RecallBox}{\RecallBox}
\begin{frame}[plain]
\hypertarget{slide28}{}\bookmark[level=1,page=28]{Distribuzione campionaria}
\begin{tikzpicture}[remember picture,overlay]
\node[anchor=north east,inner sep=0pt,text=Classical,font=\fontsize{8}{10}\selectfont] at ([xshift=19.05cm,yshift=-.16cm]current page.north west) {\hyperlink{slide3}{Indice}};
\node[anchor=north west,inner sep=0pt,text=Navy,font=\bfseries] at ([xshift=.95cm,yshift=-.62cm]current page.north west) {\usebox{\TitleBox}};
\draw[Classical,line width=1.3pt] ([xshift=.95cm,yshift=-1.68cm]current page.north west)--([xshift=2.25cm,yshift=-1.68cm]current page.north west);
\node[anchor=north west,inner sep=0pt] at ([xshift=.95cm,yshift=-2.12cm]current page.north west) {\usebox{\BodyBox}};
\node[anchor=north west,inner sep=0pt,text=Quantum,font=\bfseries\fontsize{11.5}{14}\selectfont] at ([xshift=11.6cm,yshift=-2.1cm]current page.north west) {Esempio e interpretazione};
\node[anchor=north west,inner sep=0pt] at ([xshift=11.6cm,yshift=-2.75cm]current page.north west) {\usebox{\SideBox}};
\draw[Rule,line width=.5pt] ([xshift=.95cm,yshift=-10.43cm]current page.north west)--([xshift=19.05cm,yshift=-10.43cm]current page.north west);
\node[anchor=north west,inner sep=0pt,text=Classical,font=\bfseries\fontsize{9.4}{12}\selectfont] at ([xshift=.95cm,yshift=-10.66cm]current page.north west) {Richiami};
\node[anchor=north west,inner sep=0pt] at ([xshift=2.6cm,yshift=-10.64cm]current page.north west) {\usebox{\RecallBox}};
\draw[Rule,line width=.35pt] ([xshift=.95cm,yshift=-11.92cm]current page.north west)--([xshift=19.05cm,yshift=-11.92cm]current page.north west);
\node[anchor=north west,inner sep=0pt,text=Muted,text width=9.0cm,font=\fontsize{7.7}{9.3}\selectfont] at ([xshift=.95cm,yshift=-12.06cm]current page.north west) {03 · Ripetere idealmente il campionamento};
\node[anchor=north west,inner sep=0pt,text=Classical,text width=8.2cm,align=right,font=\fontsize{7.3}{9.1}\selectfont] at ([xshift=9.45cm,yshift=-12.06cm]current page.north west) {\href{https://openintrostat.github.io/ims/}{IMS} · \href{https://stat.ethz.ch/R-manual/R-devel/library/stats/html/00Index.html}{RSTATS}};
\node[anchor=north east,inner sep=0pt,text=Muted,font=\fontsize{7.7}{9.3}\selectfont] at ([xshift=19.05cm,yshift=-12.06cm]current page.north west) {28/110};

\end{tikzpicture}
\note{Idea della slide: La distribuzione campionaria è la distribuzione di una statistica su campioni ripetuti. Può riguardare media, proporzione o differenza. Non si legge dall'istogramma dei dati individuali: per studiarla servono teoria, simulazione o ricampionamento.\par\medskip Punto da far emergere: Il livello di aggregazione è essenziale per leggere correttamente il grafico.\par\medskip Fonti della sezione: IMS, RSTATS}
\end{frame}
\begin{lrbox}{\TitleBox}\begin{minipage}[t]{18.2cm}\color{Ink}\raggedright\sloppy\fontsize{20.7}{24}\selectfont
\bfseries\color{Navy}Una simulazione con parametro noto
\end{minipage}\end{lrbox}
\begin{lrbox}{\BodyBox}\begin{minipage}[t]{8.2cm}\color{Ink}\raggedright\sloppy\fontsize{12.4}{16.4}\selectfont
Quando conosciamo il generatore, possiamo produrre molti campioni e confrontare le stime con il parametro vero del modello. È un contesto ideale per distinguere bias e dispersione. Nei dati reali, il parametro non è solitamente disponibile per questo confronto diretto.
\end{minipage}\end{lrbox}
\begin{lrbox}{\SideBox}\begin{minipage}[t]{9.35cm}\color{Ink}\raggedright\sloppy\fontsize{12.2}{16.1}\selectfont
\begin{lstlisting}
set.seed(42)
medie <- replicate(4000,
  mean(rexp(5, rate = 1)))
hist(medie)
\end{lstlisting}
\end{minipage}\end{lrbox}
\begin{lrbox}{\RecallBox}\begin{minipage}[t]{16.45cm}\color{Ink}\raggedright\sloppy\fontsize{10.1}{13.2}\selectfont
La simulazione serve a capire il comportamento della procedura sotto ipotesi note.
\end{minipage}\end{lrbox}
\MeasureBox{29-TitleBox}{\TitleBox}
\MeasureBox{29-BodyBox}{\BodyBox}
\MeasureBox{29-SideBox}{\SideBox}
\MeasureBox{29-RecallBox}{\RecallBox}
\begin{frame}[plain]
\hypertarget{slide29}{}\bookmark[level=1,page=29]{Una simulazione con parametro noto}
\begin{tikzpicture}[remember picture,overlay]
\node[anchor=north east,inner sep=0pt,text=Classical,font=\fontsize{8}{10}\selectfont] at ([xshift=19.05cm,yshift=-.16cm]current page.north west) {\hyperlink{slide3}{Indice}};
\node[anchor=north west,inner sep=0pt,text=Navy,font=\bfseries] at ([xshift=.95cm,yshift=-.62cm]current page.north west) {\usebox{\TitleBox}};
\draw[Classical,line width=1.3pt] ([xshift=.95cm,yshift=-1.68cm]current page.north west)--([xshift=2.25cm,yshift=-1.68cm]current page.north west);
\node[anchor=north west,inner sep=0pt] at ([xshift=.95cm,yshift=-2.12cm]current page.north west) {\usebox{\BodyBox}};
\node[anchor=north west,inner sep=0pt,text=Quantum,font=\bfseries\fontsize{11.5}{14}\selectfont] at ([xshift=9.7cm,yshift=-2.1cm]current page.north west) {Protocollo R};
\node[anchor=north west,inner sep=0pt] at ([xshift=9.7cm,yshift=-2.75cm]current page.north west) {\usebox{\SideBox}};
\draw[Rule,line width=.5pt] ([xshift=.95cm,yshift=-10.43cm]current page.north west)--([xshift=19.05cm,yshift=-10.43cm]current page.north west);
\node[anchor=north west,inner sep=0pt,text=Classical,font=\bfseries\fontsize{9.4}{12}\selectfont] at ([xshift=.95cm,yshift=-10.66cm]current page.north west) {Richiami};
\node[anchor=north west,inner sep=0pt] at ([xshift=2.6cm,yshift=-10.64cm]current page.north west) {\usebox{\RecallBox}};
\draw[Rule,line width=.35pt] ([xshift=.95cm,yshift=-11.92cm]current page.north west)--([xshift=19.05cm,yshift=-11.92cm]current page.north west);
\node[anchor=north west,inner sep=0pt,text=Muted,text width=9.0cm,font=\fontsize{7.7}{9.3}\selectfont] at ([xshift=.95cm,yshift=-12.06cm]current page.north west) {03 · Ripetere idealmente il campionamento};
\node[anchor=north west,inner sep=0pt,text=Classical,text width=8.2cm,align=right,font=\fontsize{7.3}{9.1}\selectfont] at ([xshift=9.45cm,yshift=-12.06cm]current page.north west) {\href{https://openintrostat.github.io/ims/}{IMS} · \href{https://stat.ethz.ch/R-manual/R-devel/library/stats/html/00Index.html}{RSTATS}};
\node[anchor=north east,inner sep=0pt,text=Muted,font=\fontsize{7.7}{9.3}\selectfont] at ([xshift=19.05cm,yshift=-12.06cm]current page.north west) {29/110};

\end{tikzpicture}
\note{Idea della slide: Quando conosciamo il generatore, possiamo produrre molti campioni e confrontare le stime con il parametro vero del modello. È un contesto ideale per distinguere bias e dispersione. Nei dati reali, il parametro non è solitamente disponibile per questo confronto diretto.\par\medskip Punto da far emergere: La simulazione serve a capire il comportamento della procedura sotto ipotesi note.\par\medskip Fonti della sezione: IMS, RSTATS}
\end{frame}
\begin{lrbox}{\TitleBox}\begin{minipage}[t]{18.2cm}\color{Ink}\raggedright\sloppy\fontsize{20.7}{24}\selectfont
\bfseries\color{Navy}Piccoli campioni
\end{minipage}\end{lrbox}
\begin{lrbox}{\BodyBox}\begin{minipage}[t]{7.1cm}\color{Ink}\raggedright\sloppy\fontsize{11.8}{15.8}\selectfont
Con pochi dati, le medie possono restare asimmetriche e relativamente disperse. Una singola osservazione grande può incidere molto. Non esiste un numero magico che renda ogni distribuzione campionaria ben approssimata da una normale.
\end{minipage}\end{lrbox}
\begin{lrbox}{\SideBox}\begin{minipage}[t]{10.4cm}\color{Ink}\raggedright\sloppy\fontsize{12.2}{16.1}\selectfont
\centering\includegraphics[width=10.4cm,height=6.25cm,keepaspectratio]{figure/clt5.pdf}\par\vspace{.22cm}{\raggedright\color{Muted}\fontsize{9.7}{12.5}\selectfont 4.000 campioni simulati, ciascuno di dimensione 5.\par}
\end{minipage}\end{lrbox}
\begin{lrbox}{\RecallBox}\begin{minipage}[t]{16.45cm}\color{Ink}\raggedright\sloppy\fontsize{10.1}{13.2}\selectfont
La qualità dell'approssimazione dipende da distribuzione, statistiche e dipendenza.
\end{minipage}\end{lrbox}
\MeasureBox{30-TitleBox}{\TitleBox}
\MeasureBox{30-BodyBox}{\BodyBox}
\MeasureBox{30-SideBox}{\SideBox}
\MeasureBox{30-RecallBox}{\RecallBox}
\begin{frame}[plain]
\hypertarget{slide30}{}\bookmark[level=1,page=30]{Piccoli campioni}
\begin{tikzpicture}[remember picture,overlay]
\node[anchor=north east,inner sep=0pt,text=Classical,font=\fontsize{8}{10}\selectfont] at ([xshift=19.05cm,yshift=-.16cm]current page.north west) {\hyperlink{slide3}{Indice}};
\node[anchor=north west,inner sep=0pt,text=Navy,font=\bfseries] at ([xshift=.95cm,yshift=-.62cm]current page.north west) {\usebox{\TitleBox}};
\draw[Classical,line width=1.3pt] ([xshift=.95cm,yshift=-1.68cm]current page.north west)--([xshift=2.25cm,yshift=-1.68cm]current page.north west);
\node[anchor=north west,inner sep=0pt] at ([xshift=.95cm,yshift=-2.12cm]current page.north west) {\usebox{\BodyBox}};
\node[anchor=north west,inner sep=0pt,text=Quantum,font=\bfseries\fontsize{11.5}{14}\selectfont] at ([xshift=8.65cm,yshift=-2.1cm]current page.north west) {Figura commentata};
\node[anchor=north west,inner sep=0pt] at ([xshift=8.65cm,yshift=-2.75cm]current page.north west) {\usebox{\SideBox}};
\draw[Rule,line width=.5pt] ([xshift=.95cm,yshift=-10.43cm]current page.north west)--([xshift=19.05cm,yshift=-10.43cm]current page.north west);
\node[anchor=north west,inner sep=0pt,text=Classical,font=\bfseries\fontsize{9.4}{12}\selectfont] at ([xshift=.95cm,yshift=-10.66cm]current page.north west) {Richiami};
\node[anchor=north west,inner sep=0pt] at ([xshift=2.6cm,yshift=-10.64cm]current page.north west) {\usebox{\RecallBox}};
\draw[Rule,line width=.35pt] ([xshift=.95cm,yshift=-11.92cm]current page.north west)--([xshift=19.05cm,yshift=-11.92cm]current page.north west);
\node[anchor=north west,inner sep=0pt,text=Muted,text width=9.0cm,font=\fontsize{7.7}{9.3}\selectfont] at ([xshift=.95cm,yshift=-12.06cm]current page.north west) {03 · Ripetere idealmente il campionamento};
\node[anchor=north west,inner sep=0pt,text=Classical,text width=8.2cm,align=right,font=\fontsize{7.3}{9.1}\selectfont] at ([xshift=9.45cm,yshift=-12.06cm]current page.north west) {\href{https://openintrostat.github.io/ims/}{IMS} · \href{https://stat.ethz.ch/R-manual/R-devel/library/stats/html/00Index.html}{RSTATS}};
\node[anchor=north east,inner sep=0pt,text=Muted,font=\fontsize{7.7}{9.3}\selectfont] at ([xshift=19.05cm,yshift=-12.06cm]current page.north west) {30/110};

\end{tikzpicture}
\note{Idea della slide: Con pochi dati, le medie possono restare asimmetriche e relativamente disperse. Una singola osservazione grande può incidere molto. Non esiste un numero magico che renda ogni distribuzione campionaria ben approssimata da una normale.\par\medskip Punto da far emergere: La qualità dell'approssimazione dipende da distribuzione, statistiche e dipendenza.\par\medskip Fonti della sezione: IMS, RSTATS}
\end{frame}
\begin{lrbox}{\TitleBox}\begin{minipage}[t]{18.2cm}\color{Ink}\raggedright\sloppy\fontsize{20.7}{24}\selectfont
\bfseries\color{Navy}Campioni più grandi
\end{minipage}\end{lrbox}
\begin{lrbox}{\BodyBox}\begin{minipage}[t]{7.1cm}\color{Ink}\raggedright\sloppy\fontsize{11.8}{15.8}\selectfont
Aumentando n nel modello di dati indipendenti e identicamente distribuiti con varianza finita, le medie si concentrano e diventano più vicine a una forma normale. Questo non significa che i dati originali diventino normali. La trasformazione riguarda la statistica calcolata su un intero campione.
\end{minipage}\end{lrbox}
\begin{lrbox}{\SideBox}\begin{minipage}[t]{10.4cm}\color{Ink}\raggedright\sloppy\fontsize{12.2}{16.1}\selectfont
\centering\includegraphics[width=10.4cm,height=6.25cm,keepaspectratio]{figure/clt40.pdf}\par\vspace{.22cm}{\raggedright\color{Muted}\fontsize{9.7}{12.5}\selectfont Stesso generatore, campioni di dimensione 40.\par}
\end{minipage}\end{lrbox}
\begin{lrbox}{\RecallBox}\begin{minipage}[t]{16.45cm}\color{Ink}\raggedright\sloppy\fontsize{10.1}{13.2}\selectfont
Non applicare conclusioni sulla media alla forma delle osservazioni individuali.
\end{minipage}\end{lrbox}
\MeasureBox{31-TitleBox}{\TitleBox}
\MeasureBox{31-BodyBox}{\BodyBox}
\MeasureBox{31-SideBox}{\SideBox}
\MeasureBox{31-RecallBox}{\RecallBox}
\begin{frame}[plain]
\hypertarget{slide31}{}\bookmark[level=1,page=31]{Campioni più grandi}
\begin{tikzpicture}[remember picture,overlay]
\node[anchor=north east,inner sep=0pt,text=Classical,font=\fontsize{8}{10}\selectfont] at ([xshift=19.05cm,yshift=-.16cm]current page.north west) {\hyperlink{slide3}{Indice}};
\node[anchor=north west,inner sep=0pt,text=Navy,font=\bfseries] at ([xshift=.95cm,yshift=-.62cm]current page.north west) {\usebox{\TitleBox}};
\draw[Classical,line width=1.3pt] ([xshift=.95cm,yshift=-1.68cm]current page.north west)--([xshift=2.25cm,yshift=-1.68cm]current page.north west);
\node[anchor=north west,inner sep=0pt] at ([xshift=.95cm,yshift=-2.12cm]current page.north west) {\usebox{\BodyBox}};
\node[anchor=north west,inner sep=0pt,text=Quantum,font=\bfseries\fontsize{11.5}{14}\selectfont] at ([xshift=8.65cm,yshift=-2.1cm]current page.north west) {Figura commentata};
\node[anchor=north west,inner sep=0pt] at ([xshift=8.65cm,yshift=-2.75cm]current page.north west) {\usebox{\SideBox}};
\draw[Rule,line width=.5pt] ([xshift=.95cm,yshift=-10.43cm]current page.north west)--([xshift=19.05cm,yshift=-10.43cm]current page.north west);
\node[anchor=north west,inner sep=0pt,text=Classical,font=\bfseries\fontsize{9.4}{12}\selectfont] at ([xshift=.95cm,yshift=-10.66cm]current page.north west) {Richiami};
\node[anchor=north west,inner sep=0pt] at ([xshift=2.6cm,yshift=-10.64cm]current page.north west) {\usebox{\RecallBox}};
\draw[Rule,line width=.35pt] ([xshift=.95cm,yshift=-11.92cm]current page.north west)--([xshift=19.05cm,yshift=-11.92cm]current page.north west);
\node[anchor=north west,inner sep=0pt,text=Muted,text width=9.0cm,font=\fontsize{7.7}{9.3}\selectfont] at ([xshift=.95cm,yshift=-12.06cm]current page.north west) {03 · Ripetere idealmente il campionamento};
\node[anchor=north west,inner sep=0pt,text=Classical,text width=8.2cm,align=right,font=\fontsize{7.3}{9.1}\selectfont] at ([xshift=9.45cm,yshift=-12.06cm]current page.north west) {\href{https://openintrostat.github.io/ims/}{IMS} · \href{https://stat.ethz.ch/R-manual/R-devel/library/stats/html/00Index.html}{RSTATS}};
\node[anchor=north east,inner sep=0pt,text=Muted,font=\fontsize{7.7}{9.3}\selectfont] at ([xshift=19.05cm,yshift=-12.06cm]current page.north west) {31/110};

\end{tikzpicture}
\note{Idea della slide: Aumentando n nel modello di dati indipendenti e identicamente distribuiti con varianza finita, le medie si concentrano e diventano più vicine a una forma normale. Questo non significa che i dati originali diventino normali. La trasformazione riguarda la statistica calcolata su un intero campione.\par\medskip Punto da far emergere: Non applicare conclusioni sulla media alla forma delle osservazioni individuali.\par\medskip Fonti della sezione: IMS, RSTATS}
\end{frame}
\begin{lrbox}{\TitleBox}\begin{minipage}[t]{18.2cm}\color{Ink}\raggedright\sloppy\fontsize{20.7}{24}\selectfont
\bfseries\color{Navy}Bias dello stimatore
\end{minipage}\end{lrbox}
\begin{lrbox}{\BodyBox}\begin{minipage}[t]{10.25cm}\color{Ink}\raggedright\sloppy\fontsize{12.4}{16.4}\selectfont
Il bias è la differenza tra il valore atteso dello stimatore e il parametro. Una procedura può essere molto stabile ma centrata sul valore sbagliato. Anche un disegno di selezione distorto può produrre questo comportamento, pur usando una formula corretta.
\end{minipage}\end{lrbox}
\begin{lrbox}{\SideBox}\begin{minipage}[t]{7.45cm}\color{Ink}\raggedright\sloppy\fontsize{12.2}{16.1}\selectfont
Una media stimata sempre su una zona più densa può convergere alla media di quella zona, non a quella dell'intera area.
\end{minipage}\end{lrbox}
\begin{lrbox}{\RecallBox}\begin{minipage}[t]{16.45cm}\color{Ink}\raggedright\sloppy\fontsize{10.1}{13.2}\selectfont
La variabilità ridotta non elimina uno spostamento sistematico.
\end{minipage}\end{lrbox}
\MeasureBox{32-TitleBox}{\TitleBox}
\MeasureBox{32-BodyBox}{\BodyBox}
\MeasureBox{32-SideBox}{\SideBox}
\MeasureBox{32-RecallBox}{\RecallBox}
\begin{frame}[plain]
\hypertarget{slide32}{}\bookmark[level=1,page=32]{Bias dello stimatore}
\begin{tikzpicture}[remember picture,overlay]
\node[anchor=north east,inner sep=0pt,text=Classical,font=\fontsize{8}{10}\selectfont] at ([xshift=19.05cm,yshift=-.16cm]current page.north west) {\hyperlink{slide3}{Indice}};
\node[anchor=north west,inner sep=0pt,text=Navy,font=\bfseries] at ([xshift=.95cm,yshift=-.62cm]current page.north west) {\usebox{\TitleBox}};
\draw[Classical,line width=1.3pt] ([xshift=.95cm,yshift=-1.68cm]current page.north west)--([xshift=2.25cm,yshift=-1.68cm]current page.north west);
\node[anchor=north west,inner sep=0pt] at ([xshift=.95cm,yshift=-2.12cm]current page.north west) {\usebox{\BodyBox}};
\node[anchor=north west,inner sep=0pt,text=Quantum,font=\bfseries\fontsize{11.5}{14}\selectfont] at ([xshift=11.6cm,yshift=-2.1cm]current page.north west) {Esempio e interpretazione};
\node[anchor=north west,inner sep=0pt] at ([xshift=11.6cm,yshift=-2.75cm]current page.north west) {\usebox{\SideBox}};
\draw[Rule,line width=.5pt] ([xshift=.95cm,yshift=-10.43cm]current page.north west)--([xshift=19.05cm,yshift=-10.43cm]current page.north west);
\node[anchor=north west,inner sep=0pt,text=Classical,font=\bfseries\fontsize{9.4}{12}\selectfont] at ([xshift=.95cm,yshift=-10.66cm]current page.north west) {Richiami};
\node[anchor=north west,inner sep=0pt] at ([xshift=2.6cm,yshift=-10.64cm]current page.north west) {\usebox{\RecallBox}};
\draw[Rule,line width=.35pt] ([xshift=.95cm,yshift=-11.92cm]current page.north west)--([xshift=19.05cm,yshift=-11.92cm]current page.north west);
\node[anchor=north west,inner sep=0pt,text=Muted,text width=9.0cm,font=\fontsize{7.7}{9.3}\selectfont] at ([xshift=.95cm,yshift=-12.06cm]current page.north west) {03 · Ripetere idealmente il campionamento};
\node[anchor=north west,inner sep=0pt,text=Classical,text width=8.2cm,align=right,font=\fontsize{7.3}{9.1}\selectfont] at ([xshift=9.45cm,yshift=-12.06cm]current page.north west) {\href{https://openintrostat.github.io/ims/}{IMS} · \href{https://stat.ethz.ch/R-manual/R-devel/library/stats/html/00Index.html}{RSTATS}};
\node[anchor=north east,inner sep=0pt,text=Muted,font=\fontsize{7.7}{9.3}\selectfont] at ([xshift=19.05cm,yshift=-12.06cm]current page.north west) {32/110};

\end{tikzpicture}
\note{Idea della slide: Il bias è la differenza tra il valore atteso dello stimatore e il parametro. Una procedura può essere molto stabile ma centrata sul valore sbagliato. Anche un disegno di selezione distorto può produrre questo comportamento, pur usando una formula corretta.\par\medskip Punto da far emergere: La variabilità ridotta non elimina uno spostamento sistematico.\par\medskip Fonti della sezione: IMS, RSTATS}
\end{frame}
\begin{lrbox}{\TitleBox}\begin{minipage}[t]{18.2cm}\color{Ink}\raggedright\sloppy\fontsize{20.7}{24}\selectfont
\bfseries\color{Navy}Precisione e accuratezza della stima
\end{minipage}\end{lrbox}
\begin{lrbox}{\BodyBox}\begin{minipage}[t]{10.25cm}\color{Ink}\raggedright\sloppy\fontsize{12.4}{16.4}\selectfont
La precisione riguarda la dispersione delle stime ripetute; l'accuratezza considera anche la vicinanza al parametro. Un aumento di n può migliorare la precisione senza correggere il bias. La qualità di una stima richiede entrambi i ragionamenti.
\end{minipage}\end{lrbox}
\begin{lrbox}{\SideBox}\begin{minipage}[t]{7.45cm}\color{Ink}\raggedright\sloppy\fontsize{12.2}{16.1}\selectfont
\fontsize{10.3}{13.2}\selectfont\setlength{\tabcolsep}{4pt}\renewcommand{\arraystretch}{1.35}\begin{tabularx}{7.45cm}{@{}>{\raggedright\arraybackslash}X >{\raggedright\arraybackslash}X >{\raggedright\arraybackslash}X@{}}\toprule
\textbf{Procedura} & \textbf{Dispersione} & \textbf{Centro}\\
\midrule
A & piccola & sbagliato\\
B & moderata & corretto\\
\bottomrule\end{tabularx}
\end{minipage}\end{lrbox}
\begin{lrbox}{\RecallBox}\begin{minipage}[t]{16.45cm}\color{Ink}\raggedright\sloppy\fontsize{10.1}{13.2}\selectfont
Una stima numericamente stabile può essere scientificamente poco affidabile.
\end{minipage}\end{lrbox}
\MeasureBox{33-TitleBox}{\TitleBox}
\MeasureBox{33-BodyBox}{\BodyBox}
\MeasureBox{33-SideBox}{\SideBox}
\MeasureBox{33-RecallBox}{\RecallBox}
\begin{frame}[plain]
\hypertarget{slide33}{}\bookmark[level=1,page=33]{Precisione e accuratezza della stima}
\begin{tikzpicture}[remember picture,overlay]
\node[anchor=north east,inner sep=0pt,text=Classical,font=\fontsize{8}{10}\selectfont] at ([xshift=19.05cm,yshift=-.16cm]current page.north west) {\hyperlink{slide3}{Indice}};
\node[anchor=north west,inner sep=0pt,text=Navy,font=\bfseries] at ([xshift=.95cm,yshift=-.62cm]current page.north west) {\usebox{\TitleBox}};
\draw[Classical,line width=1.3pt] ([xshift=.95cm,yshift=-1.68cm]current page.north west)--([xshift=2.25cm,yshift=-1.68cm]current page.north west);
\node[anchor=north west,inner sep=0pt] at ([xshift=.95cm,yshift=-2.12cm]current page.north west) {\usebox{\BodyBox}};
\node[anchor=north west,inner sep=0pt,text=Quantum,font=\bfseries\fontsize{11.5}{14}\selectfont] at ([xshift=11.6cm,yshift=-2.1cm]current page.north west) {Dati e confronto};
\node[anchor=north west,inner sep=0pt] at ([xshift=11.6cm,yshift=-2.75cm]current page.north west) {\usebox{\SideBox}};
\draw[Rule,line width=.5pt] ([xshift=.95cm,yshift=-10.43cm]current page.north west)--([xshift=19.05cm,yshift=-10.43cm]current page.north west);
\node[anchor=north west,inner sep=0pt,text=Classical,font=\bfseries\fontsize{9.4}{12}\selectfont] at ([xshift=.95cm,yshift=-10.66cm]current page.north west) {Richiami};
\node[anchor=north west,inner sep=0pt] at ([xshift=2.6cm,yshift=-10.64cm]current page.north west) {\usebox{\RecallBox}};
\draw[Rule,line width=.35pt] ([xshift=.95cm,yshift=-11.92cm]current page.north west)--([xshift=19.05cm,yshift=-11.92cm]current page.north west);
\node[anchor=north west,inner sep=0pt,text=Muted,text width=9.0cm,font=\fontsize{7.7}{9.3}\selectfont] at ([xshift=.95cm,yshift=-12.06cm]current page.north west) {03 · Ripetere idealmente il campionamento};
\node[anchor=north west,inner sep=0pt,text=Classical,text width=8.2cm,align=right,font=\fontsize{7.3}{9.1}\selectfont] at ([xshift=9.45cm,yshift=-12.06cm]current page.north west) {\href{https://openintrostat.github.io/ims/}{IMS} · \href{https://stat.ethz.ch/R-manual/R-devel/library/stats/html/00Index.html}{RSTATS}};
\node[anchor=north east,inner sep=0pt,text=Muted,font=\fontsize{7.7}{9.3}\selectfont] at ([xshift=19.05cm,yshift=-12.06cm]current page.north west) {33/110};

\end{tikzpicture}
\note{Idea della slide: La precisione riguarda la dispersione delle stime ripetute; l'accuratezza considera anche la vicinanza al parametro. Un aumento di n può migliorare la precisione senza correggere il bias. La qualità di una stima richiede entrambi i ragionamenti.\par\medskip Punto da far emergere: Una stima numericamente stabile può essere scientificamente poco affidabile.\par\medskip Fonti della sezione: IMS, RSTATS}
\end{frame}
\begin{lrbox}{\TitleBox}\begin{minipage}[t]{18.2cm}\color{Ink}\raggedright\sloppy\fontsize{20.7}{24}\selectfont
\bfseries\color{Navy}Attività 03 · applica il concetto
\end{minipage}\end{lrbox}
\begin{lrbox}{\BodyBox}\begin{minipage}[t]{10.25cm}\color{Ink}\raggedright\sloppy\fontsize{12.4}{16.4}\selectfont
Genera campioni di 5 e 40 osservazioni da una distribuzione esponenziale con media 1. Confronta le distribuzioni delle medie su molte repliche.
\end{minipage}\end{lrbox}
\begin{lrbox}{\SideBox}\begin{minipage}[t]{7.45cm}\color{Ink}\raggedright\sloppy\fontsize{12.2}{16.1}\selectfont
\noindent{\color{Classical}\bfseries 01}\hspace{.35cm}Individua la domanda\par\vspace{.22cm}\noindent{\color{Classical}\bfseries 02}\hspace{.35cm}Esplicita le ipotesi\par\vspace{.22cm}\noindent{\color{Classical}\bfseries 03}\hspace{.35cm}Calcola o rappresenta\par\vspace{.22cm}\noindent{\color{Classical}\bfseries 04}\hspace{.35cm}Scrivi una conclusione\par\vspace{.22cm}\noindent{\color{Classical}\bfseries 05}\hspace{.35cm}Indica un limite\par\vspace{.22cm}
\end{minipage}\end{lrbox}
\begin{lrbox}{\RecallBox}\begin{minipage}[t]{16.45cm}\color{Ink}\raggedright\sloppy\fontsize{10.1}{13.2}\selectfont
Perché due campioni della stessa popolazione danno medie diverse?
\end{minipage}\end{lrbox}
\MeasureBox{34-TitleBox}{\TitleBox}
\MeasureBox{34-BodyBox}{\BodyBox}
\MeasureBox{34-SideBox}{\SideBox}
\MeasureBox{34-RecallBox}{\RecallBox}
\begin{frame}[plain]
\hypertarget{slide34}{}\bookmark[level=1,page=34]{Attività 03 · applica il concetto}
\begin{tikzpicture}[remember picture,overlay]
\node[anchor=north east,inner sep=0pt,text=Classical,font=\fontsize{8}{10}\selectfont] at ([xshift=19.05cm,yshift=-.16cm]current page.north west) {\hyperlink{slide3}{Indice}};
\node[anchor=north west,inner sep=0pt,text=Navy,font=\bfseries] at ([xshift=.95cm,yshift=-.62cm]current page.north west) {\usebox{\TitleBox}};
\draw[Classical,line width=1.3pt] ([xshift=.95cm,yshift=-1.68cm]current page.north west)--([xshift=2.25cm,yshift=-1.68cm]current page.north west);
\node[anchor=north west,inner sep=0pt] at ([xshift=.95cm,yshift=-2.12cm]current page.north west) {\usebox{\BodyBox}};
\node[anchor=north west,inner sep=0pt,text=Quantum,font=\bfseries\fontsize{11.5}{14}\selectfont] at ([xshift=11.6cm,yshift=-2.1cm]current page.north west) {Consegna di lavoro};
\node[anchor=north west,inner sep=0pt] at ([xshift=11.6cm,yshift=-2.75cm]current page.north west) {\usebox{\SideBox}};
\draw[Rule,line width=.5pt] ([xshift=.95cm,yshift=-10.43cm]current page.north west)--([xshift=19.05cm,yshift=-10.43cm]current page.north west);
\node[anchor=north west,inner sep=0pt,text=Classical,font=\bfseries\fontsize{9.4}{12}\selectfont] at ([xshift=.95cm,yshift=-10.66cm]current page.north west) {Richiami};
\node[anchor=north west,inner sep=0pt] at ([xshift=2.6cm,yshift=-10.64cm]current page.north west) {\usebox{\RecallBox}};
\draw[Rule,line width=.35pt] ([xshift=.95cm,yshift=-11.92cm]current page.north west)--([xshift=19.05cm,yshift=-11.92cm]current page.north west);
\node[anchor=north west,inner sep=0pt,text=Muted,text width=9.0cm,font=\fontsize{7.7}{9.3}\selectfont] at ([xshift=.95cm,yshift=-12.06cm]current page.north west) {03 · Ripetere idealmente il campionamento};
\node[anchor=north west,inner sep=0pt,text=Classical,text width=8.2cm,align=right,font=\fontsize{7.3}{9.1}\selectfont] at ([xshift=9.45cm,yshift=-12.06cm]current page.north west) {\href{https://openintrostat.github.io/ims/}{IMS} · \href{https://stat.ethz.ch/R-manual/R-devel/library/stats/html/00Index.html}{RSTATS}};
\node[anchor=north east,inner sep=0pt,text=Muted,font=\fontsize{7.7}{9.3}\selectfont] at ([xshift=19.05cm,yshift=-12.06cm]current page.north west) {34/110};

\end{tikzpicture}
\note{Consegna: Genera campioni di 5 e 40 osservazioni da una distribuzione esponenziale con media 1. Confronta le distribuzioni delle medie su molte repliche.\par\medskip Soluzione per il docente: Entrambe sono centrate vicino a 1, ma con n=40 le medie sono meno disperse e più simmetriche. La distribuzione di origine resta asimmetrica. È la distribuzione della media, non quella dei dati, a cambiare con n.}
\end{frame}
\begin{lrbox}{\TitleBox}\begin{minipage}[t]{18.2cm}\color{Ink}\raggedright\sloppy\fontsize{20.7}{24}\selectfont
\bfseries\color{Navy}Discussione 03 · soluzione ragionata
\end{minipage}\end{lrbox}
\begin{lrbox}{\BodyBox}\begin{minipage}[t]{10.25cm}\color{Ink}\raggedright\sloppy\fontsize{12.4}{16.4}\selectfont
Entrambe sono centrate vicino a 1, ma con n=40 le medie sono meno disperse e più simmetriche. La distribuzione di origine resta asimmetrica. È la distribuzione della media, non quella dei dati, a cambiare con n.
\end{minipage}\end{lrbox}
\begin{lrbox}{\SideBox}\begin{minipage}[t]{7.45cm}\color{Ink}\raggedright\sloppy\fontsize{12.2}{16.1}\selectfont
Una stima numericamente stabile può essere scientificamente poco affidabile.
\end{minipage}\end{lrbox}
\begin{lrbox}{\RecallBox}\begin{minipage}[t]{16.45cm}\color{Ink}\raggedright\sloppy\fontsize{10.1}{13.2}\selectfont
Valuta il collegamento tra domanda, procedimento e conclusione, non soltanto il numero finale.
\end{minipage}\end{lrbox}
\MeasureBox{35-TitleBox}{\TitleBox}
\MeasureBox{35-BodyBox}{\BodyBox}
\MeasureBox{35-SideBox}{\SideBox}
\MeasureBox{35-RecallBox}{\RecallBox}
\begin{frame}[plain]
\hypertarget{slide35}{}\bookmark[level=1,page=35]{Discussione 03 · soluzione ragionata}
\begin{tikzpicture}[remember picture,overlay]
\node[anchor=north east,inner sep=0pt,text=Classical,font=\fontsize{8}{10}\selectfont] at ([xshift=19.05cm,yshift=-.16cm]current page.north west) {\hyperlink{slide3}{Indice}};
\node[anchor=north west,inner sep=0pt,text=Navy,font=\bfseries] at ([xshift=.95cm,yshift=-.62cm]current page.north west) {\usebox{\TitleBox}};
\draw[Classical,line width=1.3pt] ([xshift=.95cm,yshift=-1.68cm]current page.north west)--([xshift=2.25cm,yshift=-1.68cm]current page.north west);
\node[anchor=north west,inner sep=0pt] at ([xshift=.95cm,yshift=-2.12cm]current page.north west) {\usebox{\BodyBox}};
\node[anchor=north west,inner sep=0pt,text=Quantum,font=\bfseries\fontsize{11.5}{14}\selectfont] at ([xshift=11.6cm,yshift=-2.1cm]current page.north west) {Criterio di lettura};
\node[anchor=north west,inner sep=0pt] at ([xshift=11.6cm,yshift=-2.75cm]current page.north west) {\usebox{\SideBox}};
\draw[Rule,line width=.5pt] ([xshift=.95cm,yshift=-10.43cm]current page.north west)--([xshift=19.05cm,yshift=-10.43cm]current page.north west);
\node[anchor=north west,inner sep=0pt,text=Classical,font=\bfseries\fontsize{9.4}{12}\selectfont] at ([xshift=.95cm,yshift=-10.66cm]current page.north west) {Richiami};
\node[anchor=north west,inner sep=0pt] at ([xshift=2.6cm,yshift=-10.64cm]current page.north west) {\usebox{\RecallBox}};
\draw[Rule,line width=.35pt] ([xshift=.95cm,yshift=-11.92cm]current page.north west)--([xshift=19.05cm,yshift=-11.92cm]current page.north west);
\node[anchor=north west,inner sep=0pt,text=Muted,text width=9.0cm,font=\fontsize{7.7}{9.3}\selectfont] at ([xshift=.95cm,yshift=-12.06cm]current page.north west) {03 · Ripetere idealmente il campionamento};
\node[anchor=north west,inner sep=0pt,text=Classical,text width=8.2cm,align=right,font=\fontsize{7.3}{9.1}\selectfont] at ([xshift=9.45cm,yshift=-12.06cm]current page.north west) {\href{https://openintrostat.github.io/ims/}{IMS} · \href{https://stat.ethz.ch/R-manual/R-devel/library/stats/html/00Index.html}{RSTATS}};
\node[anchor=north east,inner sep=0pt,text=Muted,font=\fontsize{7.7}{9.3}\selectfont] at ([xshift=19.05cm,yshift=-12.06cm]current page.north west) {35/110};

\end{tikzpicture}
\note{Entrambe sono centrate vicino a 1, ma con n=40 le medie sono meno disperse e più simmetriche. La distribuzione di origine resta asimmetrica. È la distribuzione della media, non quella dei dati, a cambiare con n.}
\end{frame}
\begin{frame}[plain]\hypertarget{slide36}{}
\bookmark[level=0,page=36]{Legge dei grandi numeri e limite centrale}
\begin{tikzpicture}[remember picture,overlay]
\fill[Night] (current page.south west) rectangle (current page.north east);
\draw[Mint,line width=2pt] ([xshift=1.05cm,yshift=-1.35cm]current page.north west)--([xshift=3.0cm,yshift=-1.35cm]current page.north west);
\node[anchor=north west,inner sep=0,text=Mint,font=\bfseries\fontsize{11}{14}\selectfont] at ([xshift=1.05cm,yshift=-1.7cm]current page.north west) {SEZIONE 04};
\node[anchor=north west,inner sep=0,text=white,text width=15.8cm,align=left,font=\bfseries\fontsize{31}{38}\selectfont] at ([xshift=1.05cm,yshift=-2.65cm]current page.north west) {Legge dei grandi numeri e limite centrale};
\node[anchor=north west,inner sep=0,text=Lilac,text width=15.8cm,align=left,font=\fontsize{15}{21}\selectfont] at ([xshift=1.05cm,yshift=-6.4cm]current page.north west) {Quale risultato riguarda una media che si stabilizza e quale la sua distribuzione?};
\node[anchor=north west,inner sep=0,text=white,text width=16.5cm,align=left,font=\fontsize{11}{15}\selectfont] at ([xshift=1.05cm,yshift=-8.6cm]current page.north west) {Comprendere il concetto · leggere un esempio · applicare e discutere};
\node[anchor=north east,inner sep=0,text=Mint,font=\fontsize{8}{10}\selectfont] at ([xshift=19.05cm,yshift=-.2cm]current page.north west) {\hyperlink{slide3}{Indice}};
\node[anchor=south east,inner sep=0,text=Lilac!45,font=\bfseries\fontsize{64}{68}\selectfont] at ([xshift=18.95cm,yshift=-11.35cm]current page.north west) {04};
\draw[Lilac!30,line width=.35pt] ([xshift=.95cm,yshift=-11.92cm]current page.north west)--([xshift=19.05cm,yshift=-11.92cm]current page.north west);
\node[anchor=north west,inner sep=0pt,text=Lilac,text width=9.0cm,font=\fontsize{7.7}{9.3}\selectfont] at ([xshift=.95cm,yshift=-12.06cm]current page.north west) {04 · Legge dei grandi numeri e limite centrale};
\node[anchor=north west,inner sep=0pt,text=Mint,text width=8.2cm,align=right,font=\fontsize{7.3}{9.1}\selectfont] at ([xshift=9.45cm,yshift=-12.06cm]current page.north west) {\href{https://openintrostat.github.io/ims/}{IMS} · \href{https://www.openintro.org/go?id=stat\_os4\_probability\_chapter}{PROB}};
\node[anchor=north east,inner sep=0pt,text=Lilac,font=\fontsize{7.7}{9.3}\selectfont] at ([xshift=19.05cm,yshift=-12.06cm]current page.north west) {36/110};
\end{tikzpicture}\end{frame}
\begin{lrbox}{\TitleBox}\begin{minipage}[t]{18.2cm}\color{Ink}\raggedright\sloppy\fontsize{20.7}{24}\selectfont
\bfseries\color{Navy}La media cumulativa
\end{minipage}\end{lrbox}
\begin{lrbox}{\BodyBox}\begin{minipage}[t]{7.1cm}\color{Ink}\raggedright\sloppy\fontsize{11.8}{15.8}\selectfont
Aggiungendo osservazioni, la media cumulativa può muoversi verso il parametro senza farlo in modo monotono. Un nuovo estremo può spostarla temporaneamente. Il grafico di una sequenza mostra il comportamento intuitivo, ma non sostituisce le condizioni del teorema.
\end{minipage}\end{lrbox}
\begin{lrbox}{\SideBox}\begin{minipage}[t]{10.4cm}\color{Ink}\raggedright\sloppy\fontsize{12.2}{16.1}\selectfont
\centering\includegraphics[width=10.4cm,height=6.25cm,keepaspectratio]{figure/lln.pdf}\par\vspace{.22cm}{\raggedright\color{Muted}\fontsize{9.7}{12.5}\selectfont Il generatore esponenziale ha media uno; la traiettoria è una sola realizzazione.\par}
\end{minipage}\end{lrbox}
\begin{lrbox}{\RecallBox}\begin{minipage}[t]{16.45cm}\color{Ink}\raggedright\sloppy\fontsize{10.1}{13.2}\selectfont
Stabilizzazione non significa che ogni osservazione aggiuntiva migliori la stima.
\end{minipage}\end{lrbox}
\MeasureBox{37-TitleBox}{\TitleBox}
\MeasureBox{37-BodyBox}{\BodyBox}
\MeasureBox{37-SideBox}{\SideBox}
\MeasureBox{37-RecallBox}{\RecallBox}
\begin{frame}[plain]
\hypertarget{slide37}{}\bookmark[level=1,page=37]{La media cumulativa}
\begin{tikzpicture}[remember picture,overlay]
\node[anchor=north east,inner sep=0pt,text=Classical,font=\fontsize{8}{10}\selectfont] at ([xshift=19.05cm,yshift=-.16cm]current page.north west) {\hyperlink{slide3}{Indice}};
\node[anchor=north west,inner sep=0pt,text=Navy,font=\bfseries] at ([xshift=.95cm,yshift=-.62cm]current page.north west) {\usebox{\TitleBox}};
\draw[Classical,line width=1.3pt] ([xshift=.95cm,yshift=-1.68cm]current page.north west)--([xshift=2.25cm,yshift=-1.68cm]current page.north west);
\node[anchor=north west,inner sep=0pt] at ([xshift=.95cm,yshift=-2.12cm]current page.north west) {\usebox{\BodyBox}};
\node[anchor=north west,inner sep=0pt,text=Quantum,font=\bfseries\fontsize{11.5}{14}\selectfont] at ([xshift=8.65cm,yshift=-2.1cm]current page.north west) {Figura commentata};
\node[anchor=north west,inner sep=0pt] at ([xshift=8.65cm,yshift=-2.75cm]current page.north west) {\usebox{\SideBox}};
\draw[Rule,line width=.5pt] ([xshift=.95cm,yshift=-10.43cm]current page.north west)--([xshift=19.05cm,yshift=-10.43cm]current page.north west);
\node[anchor=north west,inner sep=0pt,text=Classical,font=\bfseries\fontsize{9.4}{12}\selectfont] at ([xshift=.95cm,yshift=-10.66cm]current page.north west) {Richiami};
\node[anchor=north west,inner sep=0pt] at ([xshift=2.6cm,yshift=-10.64cm]current page.north west) {\usebox{\RecallBox}};
\draw[Rule,line width=.35pt] ([xshift=.95cm,yshift=-11.92cm]current page.north west)--([xshift=19.05cm,yshift=-11.92cm]current page.north west);
\node[anchor=north west,inner sep=0pt,text=Muted,text width=9.0cm,font=\fontsize{7.7}{9.3}\selectfont] at ([xshift=.95cm,yshift=-12.06cm]current page.north west) {04 · Legge dei grandi numeri e limite centrale};
\node[anchor=north west,inner sep=0pt,text=Classical,text width=8.2cm,align=right,font=\fontsize{7.3}{9.1}\selectfont] at ([xshift=9.45cm,yshift=-12.06cm]current page.north west) {\href{https://openintrostat.github.io/ims/}{IMS} · \href{https://www.openintro.org/go?id=stat\_os4\_probability\_chapter}{PROB}};
\node[anchor=north east,inner sep=0pt,text=Muted,font=\fontsize{7.7}{9.3}\selectfont] at ([xshift=19.05cm,yshift=-12.06cm]current page.north west) {37/110};

\end{tikzpicture}
\note{Idea della slide: Aggiungendo osservazioni, la media cumulativa può muoversi verso il parametro senza farlo in modo monotono. Un nuovo estremo può spostarla temporaneamente. Il grafico di una sequenza mostra il comportamento intuitivo, ma non sostituisce le condizioni del teorema.\par\medskip Punto da far emergere: Stabilizzazione non significa che ogni osservazione aggiuntiva migliori la stima.\par\medskip Fonti della sezione: IMS, PROB}
\end{frame}
\begin{lrbox}{\TitleBox}\begin{minipage}[t]{18.2cm}\color{Ink}\raggedright\sloppy\fontsize{20.7}{24}\selectfont
\bfseries\color{Navy}Legge dei grandi numeri
\end{minipage}\end{lrbox}
\begin{lrbox}{\BodyBox}\begin{minipage}[t]{10.25cm}\color{Ink}\raggedright\sloppy\fontsize{12.4}{16.4}\selectfont
Per osservazioni indipendenti e identicamente distribuite con valore atteso finito, la media campionaria si avvicina al valore atteso al crescere di n. È una condizione sufficiente introduttiva. Il teorema non corregge errori di selezione o definizioni incoerenti della popolazione.
\end{minipage}\end{lrbox}
\begin{lrbox}{\SideBox}\begin{minipage}[t]{7.45cm}\color{Ink}\raggedright\sloppy\fontsize{12.2}{16.1}\selectfont
$\bar X_n\longrightarrow\mu$\par\medskip Il significato della freccia viene illustrato attraverso simulazioni, senza dimostrazione formale.
\end{minipage}\end{lrbox}
\begin{lrbox}{\RecallBox}\begin{minipage}[t]{16.45cm}\color{Ink}\raggedright\sloppy\fontsize{10.1}{13.2}\selectfont
Il limite riguarda il modello di campionamento adottato.
\end{minipage}\end{lrbox}
\MeasureBox{38-TitleBox}{\TitleBox}
\MeasureBox{38-BodyBox}{\BodyBox}
\MeasureBox{38-SideBox}{\SideBox}
\MeasureBox{38-RecallBox}{\RecallBox}
\begin{frame}[plain]
\hypertarget{slide38}{}\bookmark[level=1,page=38]{Legge dei grandi numeri}
\begin{tikzpicture}[remember picture,overlay]
\node[anchor=north east,inner sep=0pt,text=Classical,font=\fontsize{8}{10}\selectfont] at ([xshift=19.05cm,yshift=-.16cm]current page.north west) {\hyperlink{slide3}{Indice}};
\node[anchor=north west,inner sep=0pt,text=Navy,font=\bfseries] at ([xshift=.95cm,yshift=-.62cm]current page.north west) {\usebox{\TitleBox}};
\draw[Classical,line width=1.3pt] ([xshift=.95cm,yshift=-1.68cm]current page.north west)--([xshift=2.25cm,yshift=-1.68cm]current page.north west);
\node[anchor=north west,inner sep=0pt] at ([xshift=.95cm,yshift=-2.12cm]current page.north west) {\usebox{\BodyBox}};
\node[anchor=north west,inner sep=0pt,text=Quantum,font=\bfseries\fontsize{11.5}{14}\selectfont] at ([xshift=11.6cm,yshift=-2.1cm]current page.north west) {Esempio e interpretazione};
\node[anchor=north west,inner sep=0pt] at ([xshift=11.6cm,yshift=-2.75cm]current page.north west) {\usebox{\SideBox}};
\draw[Rule,line width=.5pt] ([xshift=.95cm,yshift=-10.43cm]current page.north west)--([xshift=19.05cm,yshift=-10.43cm]current page.north west);
\node[anchor=north west,inner sep=0pt,text=Classical,font=\bfseries\fontsize{9.4}{12}\selectfont] at ([xshift=.95cm,yshift=-10.66cm]current page.north west) {Richiami};
\node[anchor=north west,inner sep=0pt] at ([xshift=2.6cm,yshift=-10.64cm]current page.north west) {\usebox{\RecallBox}};
\draw[Rule,line width=.35pt] ([xshift=.95cm,yshift=-11.92cm]current page.north west)--([xshift=19.05cm,yshift=-11.92cm]current page.north west);
\node[anchor=north west,inner sep=0pt,text=Muted,text width=9.0cm,font=\fontsize{7.7}{9.3}\selectfont] at ([xshift=.95cm,yshift=-12.06cm]current page.north west) {04 · Legge dei grandi numeri e limite centrale};
\node[anchor=north west,inner sep=0pt,text=Classical,text width=8.2cm,align=right,font=\fontsize{7.3}{9.1}\selectfont] at ([xshift=9.45cm,yshift=-12.06cm]current page.north west) {\href{https://openintrostat.github.io/ims/}{IMS} · \href{https://www.openintro.org/go?id=stat\_os4\_probability\_chapter}{PROB}};
\node[anchor=north east,inner sep=0pt,text=Muted,font=\fontsize{7.7}{9.3}\selectfont] at ([xshift=19.05cm,yshift=-12.06cm]current page.north west) {38/110};

\end{tikzpicture}
\note{Idea della slide: Per osservazioni indipendenti e identicamente distribuite con valore atteso finito, la media campionaria si avvicina al valore atteso al crescere di n. È una condizione sufficiente introduttiva. Il teorema non corregge errori di selezione o definizioni incoerenti della popolazione.\par\medskip Punto da far emergere: Il limite riguarda il modello di campionamento adottato.\par\medskip Fonti della sezione: IMS, PROB}
\end{frame}
\begin{lrbox}{\TitleBox}\begin{minipage}[t]{18.2cm}\color{Ink}\raggedright\sloppy\fontsize{20.7}{24}\selectfont
\bfseries\color{Navy}Una popolazione sbagliata resta sbagliata
\end{minipage}\end{lrbox}
\begin{lrbox}{\BodyBox}\begin{minipage}[t]{10.25cm}\color{Ink}\raggedright\sloppy\fontsize{12.4}{16.4}\selectfont
Se il campione proviene sistematicamente da un sottoinsieme, la media può stabilizzarsi benissimo attorno alla media di quel sottoinsieme. Questo rende visibile perché «più dati» non equivale a «meno bias». Occorre migliorare il disegno, non soltanto la numerosità.
\end{minipage}\end{lrbox}
\begin{lrbox}{\SideBox}\begin{minipage}[t]{7.45cm}\color{Ink}\raggedright\sloppy\fontsize{12.2}{16.1}\selectfont
Misure sempre vicino alla stessa sorgente non rappresentano automaticamente l'intero acquifero.
\end{minipage}\end{lrbox}
\begin{lrbox}{\RecallBox}\begin{minipage}[t]{16.45cm}\color{Ink}\raggedright\sloppy\fontsize{10.1}{13.2}\selectfont
La legge dei grandi numeri non è una garanzia universale di verità.
\end{minipage}\end{lrbox}
\MeasureBox{39-TitleBox}{\TitleBox}
\MeasureBox{39-BodyBox}{\BodyBox}
\MeasureBox{39-SideBox}{\SideBox}
\MeasureBox{39-RecallBox}{\RecallBox}
\begin{frame}[plain]
\hypertarget{slide39}{}\bookmark[level=1,page=39]{Una popolazione sbagliata resta sbagliata}
\begin{tikzpicture}[remember picture,overlay]
\node[anchor=north east,inner sep=0pt,text=Classical,font=\fontsize{8}{10}\selectfont] at ([xshift=19.05cm,yshift=-.16cm]current page.north west) {\hyperlink{slide3}{Indice}};
\node[anchor=north west,inner sep=0pt,text=Navy,font=\bfseries] at ([xshift=.95cm,yshift=-.62cm]current page.north west) {\usebox{\TitleBox}};
\draw[Classical,line width=1.3pt] ([xshift=.95cm,yshift=-1.68cm]current page.north west)--([xshift=2.25cm,yshift=-1.68cm]current page.north west);
\node[anchor=north west,inner sep=0pt] at ([xshift=.95cm,yshift=-2.12cm]current page.north west) {\usebox{\BodyBox}};
\node[anchor=north west,inner sep=0pt,text=Quantum,font=\bfseries\fontsize{11.5}{14}\selectfont] at ([xshift=11.6cm,yshift=-2.1cm]current page.north west) {Esempio e interpretazione};
\node[anchor=north west,inner sep=0pt] at ([xshift=11.6cm,yshift=-2.75cm]current page.north west) {\usebox{\SideBox}};
\draw[Rule,line width=.5pt] ([xshift=.95cm,yshift=-10.43cm]current page.north west)--([xshift=19.05cm,yshift=-10.43cm]current page.north west);
\node[anchor=north west,inner sep=0pt,text=Classical,font=\bfseries\fontsize{9.4}{12}\selectfont] at ([xshift=.95cm,yshift=-10.66cm]current page.north west) {Richiami};
\node[anchor=north west,inner sep=0pt] at ([xshift=2.6cm,yshift=-10.64cm]current page.north west) {\usebox{\RecallBox}};
\draw[Rule,line width=.35pt] ([xshift=.95cm,yshift=-11.92cm]current page.north west)--([xshift=19.05cm,yshift=-11.92cm]current page.north west);
\node[anchor=north west,inner sep=0pt,text=Muted,text width=9.0cm,font=\fontsize{7.7}{9.3}\selectfont] at ([xshift=.95cm,yshift=-12.06cm]current page.north west) {04 · Legge dei grandi numeri e limite centrale};
\node[anchor=north west,inner sep=0pt,text=Classical,text width=8.2cm,align=right,font=\fontsize{7.3}{9.1}\selectfont] at ([xshift=9.45cm,yshift=-12.06cm]current page.north west) {\href{https://openintrostat.github.io/ims/}{IMS} · \href{https://www.openintro.org/go?id=stat\_os4\_probability\_chapter}{PROB}};
\node[anchor=north east,inner sep=0pt,text=Muted,font=\fontsize{7.7}{9.3}\selectfont] at ([xshift=19.05cm,yshift=-12.06cm]current page.north west) {39/110};

\end{tikzpicture}
\note{Idea della slide: Se il campione proviene sistematicamente da un sottoinsieme, la media può stabilizzarsi benissimo attorno alla media di quel sottoinsieme. Questo rende visibile perché «più dati» non equivale a «meno bias». Occorre migliorare il disegno, non soltanto la numerosità.\par\medskip Punto da far emergere: La legge dei grandi numeri non è una garanzia universale di verità.\par\medskip Fonti della sezione: IMS, PROB}
\end{frame}
\begin{lrbox}{\TitleBox}\begin{minipage}[t]{18.2cm}\color{Ink}\raggedright\sloppy\fontsize{20.7}{24}\selectfont
\bfseries\color{Navy}Teorema del limite centrale
\end{minipage}\end{lrbox}
\begin{lrbox}{\BodyBox}\begin{minipage}[t]{10.25cm}\color{Ink}\raggedright\sloppy\fontsize{12.4}{16.4}\selectfont
Nel modello di dati indipendenti e identicamente distribuiti con media e varianza finite, la media standardizzata tende a una distribuzione normale. La velocità dell'approssimazione dipende dalla forma di origine. La normalità è uno strumento approssimato per la statistica, non una trasformazione fisica dei dati.
\end{minipage}\end{lrbox}
\begin{lrbox}{\SideBox}\begin{minipage}[t]{7.45cm}\color{Ink}\raggedright\sloppy\fontsize{12.2}{16.1}\selectfont
$\dfrac{\bar X-\mu}{\sigma/\sqrt n}\ \approx\ N(0,1)$\par\medskip L'approssimazione è riferita a campioni ripetuti abbastanza grandi per il caso considerato.
\end{minipage}\end{lrbox}
\begin{lrbox}{\RecallBox}\begin{minipage}[t]{16.45cm}\color{Ink}\raggedright\sloppy\fontsize{10.1}{13.2}\selectfont
Non esiste una soglia universale n=30 valida per ogni problema.
\end{minipage}\end{lrbox}
\MeasureBox{40-TitleBox}{\TitleBox}
\MeasureBox{40-BodyBox}{\BodyBox}
\MeasureBox{40-SideBox}{\SideBox}
\MeasureBox{40-RecallBox}{\RecallBox}
\begin{frame}[plain]
\hypertarget{slide40}{}\bookmark[level=1,page=40]{Teorema del limite centrale}
\begin{tikzpicture}[remember picture,overlay]
\node[anchor=north east,inner sep=0pt,text=Classical,font=\fontsize{8}{10}\selectfont] at ([xshift=19.05cm,yshift=-.16cm]current page.north west) {\hyperlink{slide3}{Indice}};
\node[anchor=north west,inner sep=0pt,text=Navy,font=\bfseries] at ([xshift=.95cm,yshift=-.62cm]current page.north west) {\usebox{\TitleBox}};
\draw[Classical,line width=1.3pt] ([xshift=.95cm,yshift=-1.68cm]current page.north west)--([xshift=2.25cm,yshift=-1.68cm]current page.north west);
\node[anchor=north west,inner sep=0pt] at ([xshift=.95cm,yshift=-2.12cm]current page.north west) {\usebox{\BodyBox}};
\node[anchor=north west,inner sep=0pt,text=Quantum,font=\bfseries\fontsize{11.5}{14}\selectfont] at ([xshift=11.6cm,yshift=-2.1cm]current page.north west) {Esempio e interpretazione};
\node[anchor=north west,inner sep=0pt] at ([xshift=11.6cm,yshift=-2.75cm]current page.north west) {\usebox{\SideBox}};
\draw[Rule,line width=.5pt] ([xshift=.95cm,yshift=-10.43cm]current page.north west)--([xshift=19.05cm,yshift=-10.43cm]current page.north west);
\node[anchor=north west,inner sep=0pt,text=Classical,font=\bfseries\fontsize{9.4}{12}\selectfont] at ([xshift=.95cm,yshift=-10.66cm]current page.north west) {Richiami};
\node[anchor=north west,inner sep=0pt] at ([xshift=2.6cm,yshift=-10.64cm]current page.north west) {\usebox{\RecallBox}};
\draw[Rule,line width=.35pt] ([xshift=.95cm,yshift=-11.92cm]current page.north west)--([xshift=19.05cm,yshift=-11.92cm]current page.north west);
\node[anchor=north west,inner sep=0pt,text=Muted,text width=9.0cm,font=\fontsize{7.7}{9.3}\selectfont] at ([xshift=.95cm,yshift=-12.06cm]current page.north west) {04 · Legge dei grandi numeri e limite centrale};
\node[anchor=north west,inner sep=0pt,text=Classical,text width=8.2cm,align=right,font=\fontsize{7.3}{9.1}\selectfont] at ([xshift=9.45cm,yshift=-12.06cm]current page.north west) {\href{https://openintrostat.github.io/ims/}{IMS} · \href{https://www.openintro.org/go?id=stat\_os4\_probability\_chapter}{PROB}};
\node[anchor=north east,inner sep=0pt,text=Muted,font=\fontsize{7.7}{9.3}\selectfont] at ([xshift=19.05cm,yshift=-12.06cm]current page.north west) {40/110};

\end{tikzpicture}
\note{Idea della slide: Nel modello di dati indipendenti e identicamente distribuiti con media e varianza finite, la media standardizzata tende a una distribuzione normale. La velocità dell'approssimazione dipende dalla forma di origine. La normalità è uno strumento approssimato per la statistica, non una trasformazione fisica dei dati.\par\medskip Punto da far emergere: Non esiste una soglia universale n=30 valida per ogni problema.\par\medskip Fonti della sezione: IMS, PROB}
\end{frame}
\begin{lrbox}{\TitleBox}\begin{minipage}[t]{18.2cm}\color{Ink}\raggedright\sloppy\fontsize{20.7}{24}\selectfont
\bfseries\color{Navy}Che cosa deve avere varianza finita
\end{minipage}\end{lrbox}
\begin{lrbox}{\BodyBox}\begin{minipage}[t]{10.25cm}\color{Ink}\raggedright\sloppy\fontsize{12.4}{16.4}\selectfont
La condizione riguarda la distribuzione che genera le osservazioni. Code estremamente pesanti possono rendere inadeguata la versione elementare del teorema. Nei dati geologici reali verifichiamo inoltre estremi, selezione e dipendenza, senza assumere che una formula risolva tutto.
\end{minipage}\end{lrbox}
\begin{lrbox}{\SideBox}\begin{minipage}[t]{7.45cm}\color{Ink}\raggedright\sloppy\fontsize{12.2}{16.1}\selectfont
L'esempio esponenziale ha varianza finita e permette di osservare la convergenza delle medie. Altri generatori possono comportarsi diversamente.
\end{minipage}\end{lrbox}
\begin{lrbox}{\RecallBox}\begin{minipage}[t]{16.45cm}\color{Ink}\raggedright\sloppy\fontsize{10.1}{13.2}\selectfont
Le condizioni teoriche non sono dettagli decorativi della formula.
\end{minipage}\end{lrbox}
\MeasureBox{41-TitleBox}{\TitleBox}
\MeasureBox{41-BodyBox}{\BodyBox}
\MeasureBox{41-SideBox}{\SideBox}
\MeasureBox{41-RecallBox}{\RecallBox}
\begin{frame}[plain]
\hypertarget{slide41}{}\bookmark[level=1,page=41]{Che cosa deve avere varianza finita}
\begin{tikzpicture}[remember picture,overlay]
\node[anchor=north east,inner sep=0pt,text=Classical,font=\fontsize{8}{10}\selectfont] at ([xshift=19.05cm,yshift=-.16cm]current page.north west) {\hyperlink{slide3}{Indice}};
\node[anchor=north west,inner sep=0pt,text=Navy,font=\bfseries] at ([xshift=.95cm,yshift=-.62cm]current page.north west) {\usebox{\TitleBox}};
\draw[Classical,line width=1.3pt] ([xshift=.95cm,yshift=-1.68cm]current page.north west)--([xshift=2.25cm,yshift=-1.68cm]current page.north west);
\node[anchor=north west,inner sep=0pt] at ([xshift=.95cm,yshift=-2.12cm]current page.north west) {\usebox{\BodyBox}};
\node[anchor=north west,inner sep=0pt,text=Quantum,font=\bfseries\fontsize{11.5}{14}\selectfont] at ([xshift=11.6cm,yshift=-2.1cm]current page.north west) {Esempio e interpretazione};
\node[anchor=north west,inner sep=0pt] at ([xshift=11.6cm,yshift=-2.75cm]current page.north west) {\usebox{\SideBox}};
\draw[Rule,line width=.5pt] ([xshift=.95cm,yshift=-10.43cm]current page.north west)--([xshift=19.05cm,yshift=-10.43cm]current page.north west);
\node[anchor=north west,inner sep=0pt,text=Classical,font=\bfseries\fontsize{9.4}{12}\selectfont] at ([xshift=.95cm,yshift=-10.66cm]current page.north west) {Richiami};
\node[anchor=north west,inner sep=0pt] at ([xshift=2.6cm,yshift=-10.64cm]current page.north west) {\usebox{\RecallBox}};
\draw[Rule,line width=.35pt] ([xshift=.95cm,yshift=-11.92cm]current page.north west)--([xshift=19.05cm,yshift=-11.92cm]current page.north west);
\node[anchor=north west,inner sep=0pt,text=Muted,text width=9.0cm,font=\fontsize{7.7}{9.3}\selectfont] at ([xshift=.95cm,yshift=-12.06cm]current page.north west) {04 · Legge dei grandi numeri e limite centrale};
\node[anchor=north west,inner sep=0pt,text=Classical,text width=8.2cm,align=right,font=\fontsize{7.3}{9.1}\selectfont] at ([xshift=9.45cm,yshift=-12.06cm]current page.north west) {\href{https://openintrostat.github.io/ims/}{IMS} · \href{https://www.openintro.org/go?id=stat\_os4\_probability\_chapter}{PROB}};
\node[anchor=north east,inner sep=0pt,text=Muted,font=\fontsize{7.7}{9.3}\selectfont] at ([xshift=19.05cm,yshift=-12.06cm]current page.north west) {41/110};

\end{tikzpicture}
\note{Idea della slide: La condizione riguarda la distribuzione che genera le osservazioni. Code estremamente pesanti possono rendere inadeguata la versione elementare del teorema. Nei dati geologici reali verifichiamo inoltre estremi, selezione e dipendenza, senza assumere che una formula risolva tutto.\par\medskip Punto da far emergere: Le condizioni teoriche non sono dettagli decorativi della formula.\par\medskip Fonti della sezione: IMS, PROB}
\end{frame}
\begin{lrbox}{\TitleBox}\begin{minipage}[t]{18.2cm}\color{Ink}\raggedright\sloppy\fontsize{20.7}{24}\selectfont
\bfseries\color{Navy}Dipendenza e numerosità effettiva
\end{minipage}\end{lrbox}
\begin{lrbox}{\BodyBox}\begin{minipage}[t]{10.25cm}\color{Ink}\raggedright\sloppy\fontsize{12.4}{16.4}\selectfont
Osservazioni positivamente correlate condividono informazione e possono produrre medie meno precise di quanto suggerisca il solo n. Le formule elementari per dati indipendenti non vanno trasferite automaticamente a serie temporali o campioni concentrati nello spazio.
\end{minipage}\end{lrbox}
\begin{lrbox}{\SideBox}\begin{minipage}[t]{7.45cm}\color{Ink}\raggedright\sloppy\fontsize{12.2}{16.1}\selectfont
100 misure consecutive in una serie molto persistente possono contenere meno informazione nuova di 100 misure indipendenti.
\end{minipage}\end{lrbox}
\begin{lrbox}{\RecallBox}\begin{minipage}[t]{16.45cm}\color{Ink}\raggedright\sloppy\fontsize{10.1}{13.2}\selectfont
La numerosità effettiva non si identifica contando le righe del file.
\end{minipage}\end{lrbox}
\MeasureBox{42-TitleBox}{\TitleBox}
\MeasureBox{42-BodyBox}{\BodyBox}
\MeasureBox{42-SideBox}{\SideBox}
\MeasureBox{42-RecallBox}{\RecallBox}
\begin{frame}[plain]
\hypertarget{slide42}{}\bookmark[level=1,page=42]{Dipendenza e numerosità effettiva}
\begin{tikzpicture}[remember picture,overlay]
\node[anchor=north east,inner sep=0pt,text=Classical,font=\fontsize{8}{10}\selectfont] at ([xshift=19.05cm,yshift=-.16cm]current page.north west) {\hyperlink{slide3}{Indice}};
\node[anchor=north west,inner sep=0pt,text=Navy,font=\bfseries] at ([xshift=.95cm,yshift=-.62cm]current page.north west) {\usebox{\TitleBox}};
\draw[Classical,line width=1.3pt] ([xshift=.95cm,yshift=-1.68cm]current page.north west)--([xshift=2.25cm,yshift=-1.68cm]current page.north west);
\node[anchor=north west,inner sep=0pt] at ([xshift=.95cm,yshift=-2.12cm]current page.north west) {\usebox{\BodyBox}};
\node[anchor=north west,inner sep=0pt,text=Quantum,font=\bfseries\fontsize{11.5}{14}\selectfont] at ([xshift=11.6cm,yshift=-2.1cm]current page.north west) {Esempio e interpretazione};
\node[anchor=north west,inner sep=0pt] at ([xshift=11.6cm,yshift=-2.75cm]current page.north west) {\usebox{\SideBox}};
\draw[Rule,line width=.5pt] ([xshift=.95cm,yshift=-10.43cm]current page.north west)--([xshift=19.05cm,yshift=-10.43cm]current page.north west);
\node[anchor=north west,inner sep=0pt,text=Classical,font=\bfseries\fontsize{9.4}{12}\selectfont] at ([xshift=.95cm,yshift=-10.66cm]current page.north west) {Richiami};
\node[anchor=north west,inner sep=0pt] at ([xshift=2.6cm,yshift=-10.64cm]current page.north west) {\usebox{\RecallBox}};
\draw[Rule,line width=.35pt] ([xshift=.95cm,yshift=-11.92cm]current page.north west)--([xshift=19.05cm,yshift=-11.92cm]current page.north west);
\node[anchor=north west,inner sep=0pt,text=Muted,text width=9.0cm,font=\fontsize{7.7}{9.3}\selectfont] at ([xshift=.95cm,yshift=-12.06cm]current page.north west) {04 · Legge dei grandi numeri e limite centrale};
\node[anchor=north west,inner sep=0pt,text=Classical,text width=8.2cm,align=right,font=\fontsize{7.3}{9.1}\selectfont] at ([xshift=9.45cm,yshift=-12.06cm]current page.north west) {\href{https://openintrostat.github.io/ims/}{IMS} · \href{https://www.openintro.org/go?id=stat\_os4\_probability\_chapter}{PROB}};
\node[anchor=north east,inner sep=0pt,text=Muted,font=\fontsize{7.7}{9.3}\selectfont] at ([xshift=19.05cm,yshift=-12.06cm]current page.north west) {42/110};

\end{tikzpicture}
\note{Idea della slide: Osservazioni positivamente correlate condividono informazione e possono produrre medie meno precise di quanto suggerisca il solo n. Le formule elementari per dati indipendenti non vanno trasferite automaticamente a serie temporali o campioni concentrati nello spazio.\par\medskip Punto da far emergere: La numerosità effettiva non si identifica contando le righe del file.\par\medskip Fonti della sezione: IMS, PROB}
\end{frame}
\begin{lrbox}{\TitleBox}\begin{minipage}[t]{18.2cm}\color{Ink}\raggedright\sloppy\fontsize{20.7}{24}\selectfont
\bfseries\color{Navy}Due risultati, due domande
\end{minipage}\end{lrbox}
\begin{lrbox}{\BodyBox}\begin{minipage}[t]{10.25cm}\color{Ink}\raggedright\sloppy\fontsize{12.4}{16.4}\selectfont
La legge dei grandi numeri risponde all'idea di stabilizzazione. Il limite centrale aiuta a descrivere la forma dell'incertezza della media. Sono collegati ma non intercambiabili; comprenderli separatamente chiarisce l'origine delle formule dei successivi intervalli.
\end{minipage}\end{lrbox}
\begin{lrbox}{\SideBox}\begin{minipage}[t]{7.45cm}\color{Ink}\raggedright\sloppy\fontsize{12.2}{16.1}\selectfont
Stabilizzazione: dove tende la media?\par\smallskip Limite centrale: come variano le medie standardizzate?
\end{minipage}\end{lrbox}
\begin{lrbox}{\RecallBox}\begin{minipage}[t]{16.45cm}\color{Ink}\raggedright\sloppy\fontsize{10.1}{13.2}\selectfont
Tenere distinti centro, dispersione e forma della distribuzione campionaria.
\end{minipage}\end{lrbox}
\MeasureBox{43-TitleBox}{\TitleBox}
\MeasureBox{43-BodyBox}{\BodyBox}
\MeasureBox{43-SideBox}{\SideBox}
\MeasureBox{43-RecallBox}{\RecallBox}
\begin{frame}[plain]
\hypertarget{slide43}{}\bookmark[level=1,page=43]{Due risultati, due domande}
\begin{tikzpicture}[remember picture,overlay]
\node[anchor=north east,inner sep=0pt,text=Classical,font=\fontsize{8}{10}\selectfont] at ([xshift=19.05cm,yshift=-.16cm]current page.north west) {\hyperlink{slide3}{Indice}};
\node[anchor=north west,inner sep=0pt,text=Navy,font=\bfseries] at ([xshift=.95cm,yshift=-.62cm]current page.north west) {\usebox{\TitleBox}};
\draw[Classical,line width=1.3pt] ([xshift=.95cm,yshift=-1.68cm]current page.north west)--([xshift=2.25cm,yshift=-1.68cm]current page.north west);
\node[anchor=north west,inner sep=0pt] at ([xshift=.95cm,yshift=-2.12cm]current page.north west) {\usebox{\BodyBox}};
\node[anchor=north west,inner sep=0pt,text=Quantum,font=\bfseries\fontsize{11.5}{14}\selectfont] at ([xshift=11.6cm,yshift=-2.1cm]current page.north west) {Esempio e interpretazione};
\node[anchor=north west,inner sep=0pt] at ([xshift=11.6cm,yshift=-2.75cm]current page.north west) {\usebox{\SideBox}};
\draw[Rule,line width=.5pt] ([xshift=.95cm,yshift=-10.43cm]current page.north west)--([xshift=19.05cm,yshift=-10.43cm]current page.north west);
\node[anchor=north west,inner sep=0pt,text=Classical,font=\bfseries\fontsize{9.4}{12}\selectfont] at ([xshift=.95cm,yshift=-10.66cm]current page.north west) {Richiami};
\node[anchor=north west,inner sep=0pt] at ([xshift=2.6cm,yshift=-10.64cm]current page.north west) {\usebox{\RecallBox}};
\draw[Rule,line width=.35pt] ([xshift=.95cm,yshift=-11.92cm]current page.north west)--([xshift=19.05cm,yshift=-11.92cm]current page.north west);
\node[anchor=north west,inner sep=0pt,text=Muted,text width=9.0cm,font=\fontsize{7.7}{9.3}\selectfont] at ([xshift=.95cm,yshift=-12.06cm]current page.north west) {04 · Legge dei grandi numeri e limite centrale};
\node[anchor=north west,inner sep=0pt,text=Classical,text width=8.2cm,align=right,font=\fontsize{7.3}{9.1}\selectfont] at ([xshift=9.45cm,yshift=-12.06cm]current page.north west) {\href{https://openintrostat.github.io/ims/}{IMS} · \href{https://www.openintro.org/go?id=stat\_os4\_probability\_chapter}{PROB}};
\node[anchor=north east,inner sep=0pt,text=Muted,font=\fontsize{7.7}{9.3}\selectfont] at ([xshift=19.05cm,yshift=-12.06cm]current page.north west) {43/110};

\end{tikzpicture}
\note{Idea della slide: La legge dei grandi numeri risponde all'idea di stabilizzazione. Il limite centrale aiuta a descrivere la forma dell'incertezza della media. Sono collegati ma non intercambiabili; comprenderli separatamente chiarisce l'origine delle formule dei successivi intervalli.\par\medskip Punto da far emergere: Tenere distinti centro, dispersione e forma della distribuzione campionaria.\par\medskip Fonti della sezione: IMS, PROB}
\end{frame}
\begin{lrbox}{\TitleBox}\begin{minipage}[t]{18.2cm}\color{Ink}\raggedright\sloppy\fontsize{20.7}{24}\selectfont
\bfseries\color{Navy}Attività 04 · applica il concetto
\end{minipage}\end{lrbox}
\begin{lrbox}{\BodyBox}\begin{minipage}[t]{10.25cm}\color{Ink}\raggedright\sloppy\fontsize{12.4}{16.4}\selectfont
Spiega perché «con molti dati tutto diventa normale» è una frase sbagliata. Distingui legge dei grandi numeri e teorema del limite centrale.
\end{minipage}\end{lrbox}
\begin{lrbox}{\SideBox}\begin{minipage}[t]{7.45cm}\color{Ink}\raggedright\sloppy\fontsize{12.2}{16.1}\selectfont
\noindent{\color{Classical}\bfseries 01}\hspace{.35cm}Individua la domanda\par\vspace{.22cm}\noindent{\color{Classical}\bfseries 02}\hspace{.35cm}Esplicita le ipotesi\par\vspace{.22cm}\noindent{\color{Classical}\bfseries 03}\hspace{.35cm}Calcola o rappresenta\par\vspace{.22cm}\noindent{\color{Classical}\bfseries 04}\hspace{.35cm}Scrivi una conclusione\par\vspace{.22cm}\noindent{\color{Classical}\bfseries 05}\hspace{.35cm}Indica un limite\par\vspace{.22cm}
\end{minipage}\end{lrbox}
\begin{lrbox}{\RecallBox}\begin{minipage}[t]{16.45cm}\color{Ink}\raggedright\sloppy\fontsize{10.1}{13.2}\selectfont
Quale risultato riguarda una media che si stabilizza e quale la sua distribuzione?
\end{minipage}\end{lrbox}
\MeasureBox{44-TitleBox}{\TitleBox}
\MeasureBox{44-BodyBox}{\BodyBox}
\MeasureBox{44-SideBox}{\SideBox}
\MeasureBox{44-RecallBox}{\RecallBox}
\begin{frame}[plain]
\hypertarget{slide44}{}\bookmark[level=1,page=44]{Attività 04 · applica il concetto}
\begin{tikzpicture}[remember picture,overlay]
\node[anchor=north east,inner sep=0pt,text=Classical,font=\fontsize{8}{10}\selectfont] at ([xshift=19.05cm,yshift=-.16cm]current page.north west) {\hyperlink{slide3}{Indice}};
\node[anchor=north west,inner sep=0pt,text=Navy,font=\bfseries] at ([xshift=.95cm,yshift=-.62cm]current page.north west) {\usebox{\TitleBox}};
\draw[Classical,line width=1.3pt] ([xshift=.95cm,yshift=-1.68cm]current page.north west)--([xshift=2.25cm,yshift=-1.68cm]current page.north west);
\node[anchor=north west,inner sep=0pt] at ([xshift=.95cm,yshift=-2.12cm]current page.north west) {\usebox{\BodyBox}};
\node[anchor=north west,inner sep=0pt,text=Quantum,font=\bfseries\fontsize{11.5}{14}\selectfont] at ([xshift=11.6cm,yshift=-2.1cm]current page.north west) {Consegna di lavoro};
\node[anchor=north west,inner sep=0pt] at ([xshift=11.6cm,yshift=-2.75cm]current page.north west) {\usebox{\SideBox}};
\draw[Rule,line width=.5pt] ([xshift=.95cm,yshift=-10.43cm]current page.north west)--([xshift=19.05cm,yshift=-10.43cm]current page.north west);
\node[anchor=north west,inner sep=0pt,text=Classical,font=\bfseries\fontsize{9.4}{12}\selectfont] at ([xshift=.95cm,yshift=-10.66cm]current page.north west) {Richiami};
\node[anchor=north west,inner sep=0pt] at ([xshift=2.6cm,yshift=-10.64cm]current page.north west) {\usebox{\RecallBox}};
\draw[Rule,line width=.35pt] ([xshift=.95cm,yshift=-11.92cm]current page.north west)--([xshift=19.05cm,yshift=-11.92cm]current page.north west);
\node[anchor=north west,inner sep=0pt,text=Muted,text width=9.0cm,font=\fontsize{7.7}{9.3}\selectfont] at ([xshift=.95cm,yshift=-12.06cm]current page.north west) {04 · Legge dei grandi numeri e limite centrale};
\node[anchor=north west,inner sep=0pt,text=Classical,text width=8.2cm,align=right,font=\fontsize{7.3}{9.1}\selectfont] at ([xshift=9.45cm,yshift=-12.06cm]current page.north west) {\href{https://openintrostat.github.io/ims/}{IMS} · \href{https://www.openintro.org/go?id=stat\_os4\_probability\_chapter}{PROB}};
\node[anchor=north east,inner sep=0pt,text=Muted,font=\fontsize{7.7}{9.3}\selectfont] at ([xshift=19.05cm,yshift=-12.06cm]current page.north west) {44/110};

\end{tikzpicture}
\note{Consegna: Spiega perché «con molti dati tutto diventa normale» è una frase sbagliata. Distingui legge dei grandi numeri e teorema del limite centrale.\par\medskip Soluzione per il docente: La legge dei grandi numeri riguarda la convergenza della media sotto condizioni appropriate. Il limite centrale riguarda la distribuzione della media standardizzata su campioni ripetuti. La distribuzione dei dati può restare asimmetrica; dipendenza forte o varianza infinita richiedono altre condizioni e risultati.}
\end{frame}
\begin{lrbox}{\TitleBox}\begin{minipage}[t]{18.2cm}\color{Ink}\raggedright\sloppy\fontsize{20.7}{24}\selectfont
\bfseries\color{Navy}Discussione 04 · soluzione ragionata
\end{minipage}\end{lrbox}
\begin{lrbox}{\BodyBox}\begin{minipage}[t]{10.25cm}\color{Ink}\raggedright\sloppy\fontsize{12.4}{16.4}\selectfont
La legge dei grandi numeri riguarda la convergenza della media sotto condizioni appropriate. Il limite centrale riguarda la distribuzione della media standardizzata su campioni ripetuti. La distribuzione dei dati può restare asimmetrica; dipendenza forte o varianza infinita richiedono altre condizioni e risultati.
\end{minipage}\end{lrbox}
\begin{lrbox}{\SideBox}\begin{minipage}[t]{7.45cm}\color{Ink}\raggedright\sloppy\fontsize{12.2}{16.1}\selectfont
Tenere distinti centro, dispersione e forma della distribuzione campionaria.
\end{minipage}\end{lrbox}
\begin{lrbox}{\RecallBox}\begin{minipage}[t]{16.45cm}\color{Ink}\raggedright\sloppy\fontsize{10.1}{13.2}\selectfont
Valuta il collegamento tra domanda, procedimento e conclusione, non soltanto il numero finale.
\end{minipage}\end{lrbox}
\MeasureBox{45-TitleBox}{\TitleBox}
\MeasureBox{45-BodyBox}{\BodyBox}
\MeasureBox{45-SideBox}{\SideBox}
\MeasureBox{45-RecallBox}{\RecallBox}
\begin{frame}[plain]
\hypertarget{slide45}{}\bookmark[level=1,page=45]{Discussione 04 · soluzione ragionata}
\begin{tikzpicture}[remember picture,overlay]
\node[anchor=north east,inner sep=0pt,text=Classical,font=\fontsize{8}{10}\selectfont] at ([xshift=19.05cm,yshift=-.16cm]current page.north west) {\hyperlink{slide3}{Indice}};
\node[anchor=north west,inner sep=0pt,text=Navy,font=\bfseries] at ([xshift=.95cm,yshift=-.62cm]current page.north west) {\usebox{\TitleBox}};
\draw[Classical,line width=1.3pt] ([xshift=.95cm,yshift=-1.68cm]current page.north west)--([xshift=2.25cm,yshift=-1.68cm]current page.north west);
\node[anchor=north west,inner sep=0pt] at ([xshift=.95cm,yshift=-2.12cm]current page.north west) {\usebox{\BodyBox}};
\node[anchor=north west,inner sep=0pt,text=Quantum,font=\bfseries\fontsize{11.5}{14}\selectfont] at ([xshift=11.6cm,yshift=-2.1cm]current page.north west) {Criterio di lettura};
\node[anchor=north west,inner sep=0pt] at ([xshift=11.6cm,yshift=-2.75cm]current page.north west) {\usebox{\SideBox}};
\draw[Rule,line width=.5pt] ([xshift=.95cm,yshift=-10.43cm]current page.north west)--([xshift=19.05cm,yshift=-10.43cm]current page.north west);
\node[anchor=north west,inner sep=0pt,text=Classical,font=\bfseries\fontsize{9.4}{12}\selectfont] at ([xshift=.95cm,yshift=-10.66cm]current page.north west) {Richiami};
\node[anchor=north west,inner sep=0pt] at ([xshift=2.6cm,yshift=-10.64cm]current page.north west) {\usebox{\RecallBox}};
\draw[Rule,line width=.35pt] ([xshift=.95cm,yshift=-11.92cm]current page.north west)--([xshift=19.05cm,yshift=-11.92cm]current page.north west);
\node[anchor=north west,inner sep=0pt,text=Muted,text width=9.0cm,font=\fontsize{7.7}{9.3}\selectfont] at ([xshift=.95cm,yshift=-12.06cm]current page.north west) {04 · Legge dei grandi numeri e limite centrale};
\node[anchor=north west,inner sep=0pt,text=Classical,text width=8.2cm,align=right,font=\fontsize{7.3}{9.1}\selectfont] at ([xshift=9.45cm,yshift=-12.06cm]current page.north west) {\href{https://openintrostat.github.io/ims/}{IMS} · \href{https://www.openintro.org/go?id=stat\_os4\_probability\_chapter}{PROB}};
\node[anchor=north east,inner sep=0pt,text=Muted,font=\fontsize{7.7}{9.3}\selectfont] at ([xshift=19.05cm,yshift=-12.06cm]current page.north west) {45/110};

\end{tikzpicture}
\note{La legge dei grandi numeri riguarda la convergenza della media sotto condizioni appropriate. Il limite centrale riguarda la distribuzione della media standardizzata su campioni ripetuti. La distribuzione dei dati può restare asimmetrica; dipendenza forte o varianza infinita richiedono altre condizioni e risultati.}
\end{frame}
\begin{frame}[plain]\hypertarget{slide46}{}
\bookmark[level=0,page=46]{Errore standard}
\begin{tikzpicture}[remember picture,overlay]
\fill[Night] (current page.south west) rectangle (current page.north east);
\draw[Mint,line width=2pt] ([xshift=1.05cm,yshift=-1.35cm]current page.north west)--([xshift=3.0cm,yshift=-1.35cm]current page.north west);
\node[anchor=north west,inner sep=0,text=Mint,font=\bfseries\fontsize{11}{14}\selectfont] at ([xshift=1.05cm,yshift=-1.7cm]current page.north west) {SEZIONE 05};
\node[anchor=north west,inner sep=0,text=white,text width=15.8cm,align=left,font=\bfseries\fontsize{31}{38}\selectfont] at ([xshift=1.05cm,yshift=-2.65cm]current page.north west) {Errore standard};
\node[anchor=north west,inner sep=0,text=Lilac,text width=15.8cm,align=left,font=\fontsize{15}{21}\selectfont] at ([xshift=1.05cm,yshift=-6.4cm]current page.north west) {Quanto è incerta la media, rispetto a quanto variano gli individui?};
\node[anchor=north west,inner sep=0,text=white,text width=16.5cm,align=left,font=\fontsize{11}{15}\selectfont] at ([xshift=1.05cm,yshift=-8.6cm]current page.north west) {Comprendere il concetto · leggere un esempio · applicare e discutere};
\node[anchor=north east,inner sep=0,text=Mint,font=\fontsize{8}{10}\selectfont] at ([xshift=19.05cm,yshift=-.2cm]current page.north west) {\hyperlink{slide3}{Indice}};
\node[anchor=south east,inner sep=0,text=Lilac!45,font=\bfseries\fontsize{64}{68}\selectfont] at ([xshift=18.95cm,yshift=-11.35cm]current page.north west) {05};
\draw[Lilac!30,line width=.35pt] ([xshift=.95cm,yshift=-11.92cm]current page.north west)--([xshift=19.05cm,yshift=-11.92cm]current page.north west);
\node[anchor=north west,inner sep=0pt,text=Lilac,text width=9.0cm,font=\fontsize{7.7}{9.3}\selectfont] at ([xshift=.95cm,yshift=-12.06cm]current page.north west) {05 · Errore standard};
\node[anchor=north west,inner sep=0pt,text=Mint,text width=8.2cm,align=right,font=\fontsize{7.3}{9.1}\selectfont] at ([xshift=9.45cm,yshift=-12.06cm]current page.north west) {\href{https://openintrostat.github.io/ims/}{IMS} · \href{https://stat.ethz.ch/R-manual/R-devel/library/stats/html/00Index.html}{RSTATS}};
\node[anchor=north east,inner sep=0pt,text=Lilac,font=\fontsize{7.7}{9.3}\selectfont] at ([xshift=19.05cm,yshift=-12.06cm]current page.north west) {46/110};
\end{tikzpicture}\end{frame}
\begin{lrbox}{\TitleBox}\begin{minipage}[t]{18.2cm}\color{Ink}\raggedright\sloppy\fontsize{20.7}{24}\selectfont
\bfseries\color{Navy}Deviazione standard dei dati
\end{minipage}\end{lrbox}
\begin{lrbox}{\BodyBox}\begin{minipage}[t]{10.25cm}\color{Ink}\raggedright\sloppy\fontsize{12.4}{16.4}\selectfont
La deviazione standard descrive quanto differiscono gli individui nel campione. È una proprietà della distribuzione osservata. Non tende automaticamente a zero aggiungendo dati, perché la popolazione può restare eterogenea anche quando la sua media è stimata con precisione.
\end{minipage}\end{lrbox}
\begin{lrbox}{\SideBox}\begin{minipage}[t]{7.45cm}\color{Ink}\raggedright\sloppy\fontsize{12.2}{16.1}\selectfont
Nel dataset dei fossili: s circa 1,001 \(\mu\)m.\par\smallskip Questa è dispersione individuale, non incertezza della media.
\end{minipage}\end{lrbox}
\begin{lrbox}{\RecallBox}\begin{minipage}[t]{16.45cm}\color{Ink}\raggedright\sloppy\fontsize{10.1}{13.2}\selectfont
Non sostituire la deviazione standard con l'errore standard per far apparire i dati meno dispersi.
\end{minipage}\end{lrbox}
\MeasureBox{47-TitleBox}{\TitleBox}
\MeasureBox{47-BodyBox}{\BodyBox}
\MeasureBox{47-SideBox}{\SideBox}
\MeasureBox{47-RecallBox}{\RecallBox}
\begin{frame}[plain]
\hypertarget{slide47}{}\bookmark[level=1,page=47]{Deviazione standard dei dati}
\begin{tikzpicture}[remember picture,overlay]
\node[anchor=north east,inner sep=0pt,text=Classical,font=\fontsize{8}{10}\selectfont] at ([xshift=19.05cm,yshift=-.16cm]current page.north west) {\hyperlink{slide3}{Indice}};
\node[anchor=north west,inner sep=0pt,text=Navy,font=\bfseries] at ([xshift=.95cm,yshift=-.62cm]current page.north west) {\usebox{\TitleBox}};
\draw[Classical,line width=1.3pt] ([xshift=.95cm,yshift=-1.68cm]current page.north west)--([xshift=2.25cm,yshift=-1.68cm]current page.north west);
\node[anchor=north west,inner sep=0pt] at ([xshift=.95cm,yshift=-2.12cm]current page.north west) {\usebox{\BodyBox}};
\node[anchor=north west,inner sep=0pt,text=Quantum,font=\bfseries\fontsize{11.5}{14}\selectfont] at ([xshift=11.6cm,yshift=-2.1cm]current page.north west) {Esempio e interpretazione};
\node[anchor=north west,inner sep=0pt] at ([xshift=11.6cm,yshift=-2.75cm]current page.north west) {\usebox{\SideBox}};
\draw[Rule,line width=.5pt] ([xshift=.95cm,yshift=-10.43cm]current page.north west)--([xshift=19.05cm,yshift=-10.43cm]current page.north west);
\node[anchor=north west,inner sep=0pt,text=Classical,font=\bfseries\fontsize{9.4}{12}\selectfont] at ([xshift=.95cm,yshift=-10.66cm]current page.north west) {Richiami};
\node[anchor=north west,inner sep=0pt] at ([xshift=2.6cm,yshift=-10.64cm]current page.north west) {\usebox{\RecallBox}};
\draw[Rule,line width=.35pt] ([xshift=.95cm,yshift=-11.92cm]current page.north west)--([xshift=19.05cm,yshift=-11.92cm]current page.north west);
\node[anchor=north west,inner sep=0pt,text=Muted,text width=9.0cm,font=\fontsize{7.7}{9.3}\selectfont] at ([xshift=.95cm,yshift=-12.06cm]current page.north west) {05 · Errore standard};
\node[anchor=north west,inner sep=0pt,text=Classical,text width=8.2cm,align=right,font=\fontsize{7.3}{9.1}\selectfont] at ([xshift=9.45cm,yshift=-12.06cm]current page.north west) {\href{https://openintrostat.github.io/ims/}{IMS} · \href{https://stat.ethz.ch/R-manual/R-devel/library/stats/html/00Index.html}{RSTATS}};
\node[anchor=north east,inner sep=0pt,text=Muted,font=\fontsize{7.7}{9.3}\selectfont] at ([xshift=19.05cm,yshift=-12.06cm]current page.north west) {47/110};

\end{tikzpicture}
\note{Idea della slide: La deviazione standard descrive quanto differiscono gli individui nel campione. È una proprietà della distribuzione osservata. Non tende automaticamente a zero aggiungendo dati, perché la popolazione può restare eterogenea anche quando la sua media è stimata con precisione.\par\medskip Punto da far emergere: Non sostituire la deviazione standard con l'errore standard per far apparire i dati meno dispersi.\par\medskip Fonti della sezione: IMS, RSTATS}
\end{frame}
\begin{lrbox}{\TitleBox}\begin{minipage}[t]{18.2cm}\color{Ink}\raggedright\sloppy\fontsize{20.7}{24}\selectfont
\bfseries\color{Navy}Errore standard della media
\end{minipage}\end{lrbox}
\begin{lrbox}{\BodyBox}\begin{minipage}[t]{10.25cm}\color{Ink}\raggedright\sloppy\fontsize{12.4}{16.4}\selectfont
L'errore standard è la deviazione standard della distribuzione campionaria dello stimatore. Per una media di dati indipendenti con varianza comune, vale \(\sigma\)/\(\sqrt{}\)n; sostituendo \(\sigma\) con s otteniamo una stima. La formula deriva dal modello, non dal formato della tabella.
\end{minipage}\end{lrbox}
\begin{lrbox}{\SideBox}\begin{minipage}[t]{7.45cm}\color{Ink}\raggedright\sloppy\fontsize{12.2}{16.1}\selectfont
$SE(\bar X)=\sigma/\sqrt n$\par\medskip Stima: $\widehat{SE}=s/\sqrt n$.
\end{minipage}\end{lrbox}
\begin{lrbox}{\RecallBox}\begin{minipage}[t]{16.45cm}\color{Ink}\raggedright\sloppy\fontsize{10.1}{13.2}\selectfont
La formula semplice richiede che la struttura del campionamento sia compatibile.
\end{minipage}\end{lrbox}
\MeasureBox{48-TitleBox}{\TitleBox}
\MeasureBox{48-BodyBox}{\BodyBox}
\MeasureBox{48-SideBox}{\SideBox}
\MeasureBox{48-RecallBox}{\RecallBox}
\begin{frame}[plain]
\hypertarget{slide48}{}\bookmark[level=1,page=48]{Errore standard della media}
\begin{tikzpicture}[remember picture,overlay]
\node[anchor=north east,inner sep=0pt,text=Classical,font=\fontsize{8}{10}\selectfont] at ([xshift=19.05cm,yshift=-.16cm]current page.north west) {\hyperlink{slide3}{Indice}};
\node[anchor=north west,inner sep=0pt,text=Navy,font=\bfseries] at ([xshift=.95cm,yshift=-.62cm]current page.north west) {\usebox{\TitleBox}};
\draw[Classical,line width=1.3pt] ([xshift=.95cm,yshift=-1.68cm]current page.north west)--([xshift=2.25cm,yshift=-1.68cm]current page.north west);
\node[anchor=north west,inner sep=0pt] at ([xshift=.95cm,yshift=-2.12cm]current page.north west) {\usebox{\BodyBox}};
\node[anchor=north west,inner sep=0pt,text=Quantum,font=\bfseries\fontsize{11.5}{14}\selectfont] at ([xshift=11.6cm,yshift=-2.1cm]current page.north west) {Esempio e interpretazione};
\node[anchor=north west,inner sep=0pt] at ([xshift=11.6cm,yshift=-2.75cm]current page.north west) {\usebox{\SideBox}};
\draw[Rule,line width=.5pt] ([xshift=.95cm,yshift=-10.43cm]current page.north west)--([xshift=19.05cm,yshift=-10.43cm]current page.north west);
\node[anchor=north west,inner sep=0pt,text=Classical,font=\bfseries\fontsize{9.4}{12}\selectfont] at ([xshift=.95cm,yshift=-10.66cm]current page.north west) {Richiami};
\node[anchor=north west,inner sep=0pt] at ([xshift=2.6cm,yshift=-10.64cm]current page.north west) {\usebox{\RecallBox}};
\draw[Rule,line width=.35pt] ([xshift=.95cm,yshift=-11.92cm]current page.north west)--([xshift=19.05cm,yshift=-11.92cm]current page.north west);
\node[anchor=north west,inner sep=0pt,text=Muted,text width=9.0cm,font=\fontsize{7.7}{9.3}\selectfont] at ([xshift=.95cm,yshift=-12.06cm]current page.north west) {05 · Errore standard};
\node[anchor=north west,inner sep=0pt,text=Classical,text width=8.2cm,align=right,font=\fontsize{7.3}{9.1}\selectfont] at ([xshift=9.45cm,yshift=-12.06cm]current page.north west) {\href{https://openintrostat.github.io/ims/}{IMS} · \href{https://stat.ethz.ch/R-manual/R-devel/library/stats/html/00Index.html}{RSTATS}};
\node[anchor=north east,inner sep=0pt,text=Muted,font=\fontsize{7.7}{9.3}\selectfont] at ([xshift=19.05cm,yshift=-12.06cm]current page.north west) {48/110};

\end{tikzpicture}
\note{Idea della slide: L'errore standard è la deviazione standard della distribuzione campionaria dello stimatore. Per una media di dati indipendenti con varianza comune, vale \(\sigma\)/\(\sqrt{}\)n; sostituendo \(\sigma\) con s otteniamo una stima. La formula deriva dal modello, non dal formato della tabella.\par\medskip Punto da far emergere: La formula semplice richiede che la struttura del campionamento sia compatibile.\par\medskip Fonti della sezione: IMS, RSTATS}
\end{frame}
\begin{lrbox}{\TitleBox}\begin{minipage}[t]{18.2cm}\color{Ink}\raggedright\sloppy\fontsize{20.7}{24}\selectfont
\bfseries\color{Navy}Calcolo sul caso didattico
\end{minipage}\end{lrbox}
\begin{lrbox}{\BodyBox}\begin{minipage}[t]{8.2cm}\color{Ink}\raggedright\sloppy\fontsize{12.4}{16.4}\selectfont
Applicando la formula elementare ai 160 diametri otteniamo circa 0,079 \(\mu\)m. È un modello di lavoro per una miscela campionata senza quote fisse. Il generatore ha invece 80 individui per morfotipo: per quel disegno va stimata la variabilità entro i gruppi, mantenendo le quote.
\end{minipage}\end{lrbox}
\begin{lrbox}{\SideBox}\begin{minipage}[t]{9.35cm}\color{Ink}\raggedright\sloppy\fontsize{12.2}{16.1}\selectfont
\begin{lstlisting}
x <- d$diametro_tot_um
sd(x) / sqrt(length(x))
# Circa 0.079 µm
\end{lstlisting}
\end{minipage}\end{lrbox}
\begin{lrbox}{\RecallBox}\begin{minipage}[t]{16.45cm}\color{Ink}\raggedright\sloppy\fontsize{10.1}{13.2}\selectfont
La formula elementare non riproduce automaticamente un disegno con quote fisse o gruppi dipendenti.
\end{minipage}\end{lrbox}
\MeasureBox{49-TitleBox}{\TitleBox}
\MeasureBox{49-BodyBox}{\BodyBox}
\MeasureBox{49-SideBox}{\SideBox}
\MeasureBox{49-RecallBox}{\RecallBox}
\begin{frame}[plain]
\hypertarget{slide49}{}\bookmark[level=1,page=49]{Calcolo sul caso didattico}
\begin{tikzpicture}[remember picture,overlay]
\node[anchor=north east,inner sep=0pt,text=Classical,font=\fontsize{8}{10}\selectfont] at ([xshift=19.05cm,yshift=-.16cm]current page.north west) {\hyperlink{slide3}{Indice}};
\node[anchor=north west,inner sep=0pt,text=Navy,font=\bfseries] at ([xshift=.95cm,yshift=-.62cm]current page.north west) {\usebox{\TitleBox}};
\draw[Classical,line width=1.3pt] ([xshift=.95cm,yshift=-1.68cm]current page.north west)--([xshift=2.25cm,yshift=-1.68cm]current page.north west);
\node[anchor=north west,inner sep=0pt] at ([xshift=.95cm,yshift=-2.12cm]current page.north west) {\usebox{\BodyBox}};
\node[anchor=north west,inner sep=0pt,text=Quantum,font=\bfseries\fontsize{11.5}{14}\selectfont] at ([xshift=9.7cm,yshift=-2.1cm]current page.north west) {Protocollo R};
\node[anchor=north west,inner sep=0pt] at ([xshift=9.7cm,yshift=-2.75cm]current page.north west) {\usebox{\SideBox}};
\draw[Rule,line width=.5pt] ([xshift=.95cm,yshift=-10.43cm]current page.north west)--([xshift=19.05cm,yshift=-10.43cm]current page.north west);
\node[anchor=north west,inner sep=0pt,text=Classical,font=\bfseries\fontsize{9.4}{12}\selectfont] at ([xshift=.95cm,yshift=-10.66cm]current page.north west) {Richiami};
\node[anchor=north west,inner sep=0pt] at ([xshift=2.6cm,yshift=-10.64cm]current page.north west) {\usebox{\RecallBox}};
\draw[Rule,line width=.35pt] ([xshift=.95cm,yshift=-11.92cm]current page.north west)--([xshift=19.05cm,yshift=-11.92cm]current page.north west);
\node[anchor=north west,inner sep=0pt,text=Muted,text width=9.0cm,font=\fontsize{7.7}{9.3}\selectfont] at ([xshift=.95cm,yshift=-12.06cm]current page.north west) {05 · Errore standard};
\node[anchor=north west,inner sep=0pt,text=Classical,text width=8.2cm,align=right,font=\fontsize{7.3}{9.1}\selectfont] at ([xshift=9.45cm,yshift=-12.06cm]current page.north west) {\href{https://openintrostat.github.io/ims/}{IMS} · \href{https://stat.ethz.ch/R-manual/R-devel/library/stats/html/00Index.html}{RSTATS}};
\node[anchor=north east,inner sep=0pt,text=Muted,font=\fontsize{7.7}{9.3}\selectfont] at ([xshift=19.05cm,yshift=-12.06cm]current page.north west) {49/110};

\end{tikzpicture}
\note{Idea della slide: Applicando la formula elementare ai 160 diametri otteniamo circa 0,079 \(\mu\)m. È un modello di lavoro per una miscela campionata senza quote fisse. Il generatore ha invece 80 individui per morfotipo: per quel disegno va stimata la variabilità entro i gruppi, mantenendo le quote.\par\medskip Punto da far emergere: La formula elementare non riproduce automaticamente un disegno con quote fisse o gruppi dipendenti.\par\medskip Fonti della sezione: IMS, RSTATS}
\end{frame}
\begin{lrbox}{\TitleBox}\begin{minipage}[t]{18.2cm}\color{Ink}\raggedright\sloppy\fontsize{20.7}{24}\selectfont
\bfseries\color{Navy}La legge della radice quadrata
\end{minipage}\end{lrbox}
\begin{lrbox}{\BodyBox}\begin{minipage}[t]{7.1cm}\color{Ink}\raggedright\sloppy\fontsize{11.8}{15.8}\selectfont
La precisione cresce con la radice della numerosità, non in proporzione diretta. Dimezzare l'errore standard richiede circa quattro volte le unità indipendenti, a parità di dispersione. Questo aiuta a valutare costi e benefici di una campagna di misura.
\end{minipage}\end{lrbox}
\begin{lrbox}{\SideBox}\begin{minipage}[t]{10.4cm}\color{Ink}\raggedright\sloppy\fontsize{12.2}{16.1}\selectfont
\centering\includegraphics[width=10.4cm,height=6.25cm,keepaspectratio]{figure/se_n.pdf}\par\vspace{.22cm}{\raggedright\color{Muted}\fontsize{9.7}{12.5}\selectfont Relazione teorica con deviazione standard della popolazione fissata a uno.\par}
\end{minipage}\end{lrbox}
\begin{lrbox}{\RecallBox}\begin{minipage}[t]{16.45cm}\color{Ink}\raggedright\sloppy\fontsize{10.1}{13.2}\selectfont
Le unità aggiuntive devono apportare informazione, non essere duplicazioni dipendenti.
\end{minipage}\end{lrbox}
\MeasureBox{50-TitleBox}{\TitleBox}
\MeasureBox{50-BodyBox}{\BodyBox}
\MeasureBox{50-SideBox}{\SideBox}
\MeasureBox{50-RecallBox}{\RecallBox}
\begin{frame}[plain]
\hypertarget{slide50}{}\bookmark[level=1,page=50]{La legge della radice quadrata}
\begin{tikzpicture}[remember picture,overlay]
\node[anchor=north east,inner sep=0pt,text=Classical,font=\fontsize{8}{10}\selectfont] at ([xshift=19.05cm,yshift=-.16cm]current page.north west) {\hyperlink{slide3}{Indice}};
\node[anchor=north west,inner sep=0pt,text=Navy,font=\bfseries] at ([xshift=.95cm,yshift=-.62cm]current page.north west) {\usebox{\TitleBox}};
\draw[Classical,line width=1.3pt] ([xshift=.95cm,yshift=-1.68cm]current page.north west)--([xshift=2.25cm,yshift=-1.68cm]current page.north west);
\node[anchor=north west,inner sep=0pt] at ([xshift=.95cm,yshift=-2.12cm]current page.north west) {\usebox{\BodyBox}};
\node[anchor=north west,inner sep=0pt,text=Quantum,font=\bfseries\fontsize{11.5}{14}\selectfont] at ([xshift=8.65cm,yshift=-2.1cm]current page.north west) {Figura commentata};
\node[anchor=north west,inner sep=0pt] at ([xshift=8.65cm,yshift=-2.75cm]current page.north west) {\usebox{\SideBox}};
\draw[Rule,line width=.5pt] ([xshift=.95cm,yshift=-10.43cm]current page.north west)--([xshift=19.05cm,yshift=-10.43cm]current page.north west);
\node[anchor=north west,inner sep=0pt,text=Classical,font=\bfseries\fontsize{9.4}{12}\selectfont] at ([xshift=.95cm,yshift=-10.66cm]current page.north west) {Richiami};
\node[anchor=north west,inner sep=0pt] at ([xshift=2.6cm,yshift=-10.64cm]current page.north west) {\usebox{\RecallBox}};
\draw[Rule,line width=.35pt] ([xshift=.95cm,yshift=-11.92cm]current page.north west)--([xshift=19.05cm,yshift=-11.92cm]current page.north west);
\node[anchor=north west,inner sep=0pt,text=Muted,text width=9.0cm,font=\fontsize{7.7}{9.3}\selectfont] at ([xshift=.95cm,yshift=-12.06cm]current page.north west) {05 · Errore standard};
\node[anchor=north west,inner sep=0pt,text=Classical,text width=8.2cm,align=right,font=\fontsize{7.3}{9.1}\selectfont] at ([xshift=9.45cm,yshift=-12.06cm]current page.north west) {\href{https://openintrostat.github.io/ims/}{IMS} · \href{https://stat.ethz.ch/R-manual/R-devel/library/stats/html/00Index.html}{RSTATS}};
\node[anchor=north east,inner sep=0pt,text=Muted,font=\fontsize{7.7}{9.3}\selectfont] at ([xshift=19.05cm,yshift=-12.06cm]current page.north west) {50/110};

\end{tikzpicture}
\note{Idea della slide: La precisione cresce con la radice della numerosità, non in proporzione diretta. Dimezzare l'errore standard richiede circa quattro volte le unità indipendenti, a parità di dispersione. Questo aiuta a valutare costi e benefici di una campagna di misura.\par\medskip Punto da far emergere: Le unità aggiuntive devono apportare informazione, non essere duplicazioni dipendenti.\par\medskip Fonti della sezione: IMS, RSTATS}
\end{frame}
\begin{lrbox}{\TitleBox}\begin{minipage}[t]{18.2cm}\color{Ink}\raggedright\sloppy\fontsize{20.7}{24}\selectfont
\bfseries\color{Navy}Errore standard di una proporzione
\end{minipage}\end{lrbox}
\begin{lrbox}{\BodyBox}\begin{minipage}[t]{10.25cm}\color{Ink}\raggedright\sloppy\fontsize{12.4}{16.4}\selectfont
Per una proporzione binomiale, la dispersione teorica è \(\sqrt{}\)[p(1–p)/n]. La sostituzione di p con la frequenza osservata è una stima semplice, ma gli intervalli costruiti ingenuamente possono funzionare male con campioni piccoli o quote vicine a zero e uno.
\end{minipage}\end{lrbox}
\begin{lrbox}{\SideBox}\begin{minipage}[t]{7.45cm}\color{Ink}\raggedright\sloppy\fontsize{12.2}{16.1}\selectfont
$SE(\hat p)=\sqrt{p(1-p)/n}$\par\medskip Per gli intervalli useremo procedure appropriate, non una regola unica per ogni p.
\end{minipage}\end{lrbox}
\begin{lrbox}{\RecallBox}\begin{minipage}[t]{16.45cm}\color{Ink}\raggedright\sloppy\fontsize{10.1}{13.2}\selectfont
Un errore standard stimato pari a zero per 0 successi non significa certezza assoluta.
\end{minipage}\end{lrbox}
\MeasureBox{51-TitleBox}{\TitleBox}
\MeasureBox{51-BodyBox}{\BodyBox}
\MeasureBox{51-SideBox}{\SideBox}
\MeasureBox{51-RecallBox}{\RecallBox}
\begin{frame}[plain]
\hypertarget{slide51}{}\bookmark[level=1,page=51]{Errore standard di una proporzione}
\begin{tikzpicture}[remember picture,overlay]
\node[anchor=north east,inner sep=0pt,text=Classical,font=\fontsize{8}{10}\selectfont] at ([xshift=19.05cm,yshift=-.16cm]current page.north west) {\hyperlink{slide3}{Indice}};
\node[anchor=north west,inner sep=0pt,text=Navy,font=\bfseries] at ([xshift=.95cm,yshift=-.62cm]current page.north west) {\usebox{\TitleBox}};
\draw[Classical,line width=1.3pt] ([xshift=.95cm,yshift=-1.68cm]current page.north west)--([xshift=2.25cm,yshift=-1.68cm]current page.north west);
\node[anchor=north west,inner sep=0pt] at ([xshift=.95cm,yshift=-2.12cm]current page.north west) {\usebox{\BodyBox}};
\node[anchor=north west,inner sep=0pt,text=Quantum,font=\bfseries\fontsize{11.5}{14}\selectfont] at ([xshift=11.6cm,yshift=-2.1cm]current page.north west) {Esempio e interpretazione};
\node[anchor=north west,inner sep=0pt] at ([xshift=11.6cm,yshift=-2.75cm]current page.north west) {\usebox{\SideBox}};
\draw[Rule,line width=.5pt] ([xshift=.95cm,yshift=-10.43cm]current page.north west)--([xshift=19.05cm,yshift=-10.43cm]current page.north west);
\node[anchor=north west,inner sep=0pt,text=Classical,font=\bfseries\fontsize{9.4}{12}\selectfont] at ([xshift=.95cm,yshift=-10.66cm]current page.north west) {Richiami};
\node[anchor=north west,inner sep=0pt] at ([xshift=2.6cm,yshift=-10.64cm]current page.north west) {\usebox{\RecallBox}};
\draw[Rule,line width=.35pt] ([xshift=.95cm,yshift=-11.92cm]current page.north west)--([xshift=19.05cm,yshift=-11.92cm]current page.north west);
\node[anchor=north west,inner sep=0pt,text=Muted,text width=9.0cm,font=\fontsize{7.7}{9.3}\selectfont] at ([xshift=.95cm,yshift=-12.06cm]current page.north west) {05 · Errore standard};
\node[anchor=north west,inner sep=0pt,text=Classical,text width=8.2cm,align=right,font=\fontsize{7.3}{9.1}\selectfont] at ([xshift=9.45cm,yshift=-12.06cm]current page.north west) {\href{https://openintrostat.github.io/ims/}{IMS} · \href{https://stat.ethz.ch/R-manual/R-devel/library/stats/html/00Index.html}{RSTATS}};
\node[anchor=north east,inner sep=0pt,text=Muted,font=\fontsize{7.7}{9.3}\selectfont] at ([xshift=19.05cm,yshift=-12.06cm]current page.north west) {51/110};

\end{tikzpicture}
\note{Idea della slide: Per una proporzione binomiale, la dispersione teorica è \(\sqrt{}\)[p(1–p)/n]. La sostituzione di p con la frequenza osservata è una stima semplice, ma gli intervalli costruiti ingenuamente possono funzionare male con campioni piccoli o quote vicine a zero e uno.\par\medskip Punto da far emergere: Un errore standard stimato pari a zero per 0 successi non significa certezza assoluta.\par\medskip Fonti della sezione: IMS, RSTATS}
\end{frame}
\begin{lrbox}{\TitleBox}\begin{minipage}[t]{18.2cm}\color{Ink}\raggedright\sloppy\fontsize{20.7}{24}\selectfont
\bfseries\color{Navy}Errore standard e dimensione del campione
\end{minipage}\end{lrbox}
\begin{lrbox}{\BodyBox}\begin{minipage}[t]{10.25cm}\color{Ink}\raggedright\sloppy\fontsize{12.4}{16.4}\selectfont
La presenza di molti dati non garantisce un piccolo errore in ogni contesto. Dipendenza, pesi, disegno a gruppi e variabilità elevata modificano il calcolo. Anche un intervallo molto stretto può non includere errori sistematici o incertezza del protocollo.
\end{minipage}\end{lrbox}
\begin{lrbox}{\SideBox}\begin{minipage}[t]{7.45cm}\color{Ink}\raggedright\sloppy\fontsize{12.2}{16.1}\selectfont
Domande: le unità sono indipendenti? tutte pesano allo stesso modo? il protocollo ha bias?
\end{minipage}\end{lrbox}
\begin{lrbox}{\RecallBox}\begin{minipage}[t]{16.45cm}\color{Ink}\raggedright\sloppy\fontsize{10.1}{13.2}\selectfont
L'errore standard quantifica una componente dell'incertezza sotto un modello.
\end{minipage}\end{lrbox}
\MeasureBox{52-TitleBox}{\TitleBox}
\MeasureBox{52-BodyBox}{\BodyBox}
\MeasureBox{52-SideBox}{\SideBox}
\MeasureBox{52-RecallBox}{\RecallBox}
\begin{frame}[plain]
\hypertarget{slide52}{}\bookmark[level=1,page=52]{Errore standard e dimensione del campione}
\begin{tikzpicture}[remember picture,overlay]
\node[anchor=north east,inner sep=0pt,text=Classical,font=\fontsize{8}{10}\selectfont] at ([xshift=19.05cm,yshift=-.16cm]current page.north west) {\hyperlink{slide3}{Indice}};
\node[anchor=north west,inner sep=0pt,text=Navy,font=\bfseries] at ([xshift=.95cm,yshift=-.62cm]current page.north west) {\usebox{\TitleBox}};
\draw[Classical,line width=1.3pt] ([xshift=.95cm,yshift=-1.68cm]current page.north west)--([xshift=2.25cm,yshift=-1.68cm]current page.north west);
\node[anchor=north west,inner sep=0pt] at ([xshift=.95cm,yshift=-2.12cm]current page.north west) {\usebox{\BodyBox}};
\node[anchor=north west,inner sep=0pt,text=Quantum,font=\bfseries\fontsize{11.5}{14}\selectfont] at ([xshift=11.6cm,yshift=-2.1cm]current page.north west) {Esempio e interpretazione};
\node[anchor=north west,inner sep=0pt] at ([xshift=11.6cm,yshift=-2.75cm]current page.north west) {\usebox{\SideBox}};
\draw[Rule,line width=.5pt] ([xshift=.95cm,yshift=-10.43cm]current page.north west)--([xshift=19.05cm,yshift=-10.43cm]current page.north west);
\node[anchor=north west,inner sep=0pt,text=Classical,font=\bfseries\fontsize{9.4}{12}\selectfont] at ([xshift=.95cm,yshift=-10.66cm]current page.north west) {Richiami};
\node[anchor=north west,inner sep=0pt] at ([xshift=2.6cm,yshift=-10.64cm]current page.north west) {\usebox{\RecallBox}};
\draw[Rule,line width=.35pt] ([xshift=.95cm,yshift=-11.92cm]current page.north west)--([xshift=19.05cm,yshift=-11.92cm]current page.north west);
\node[anchor=north west,inner sep=0pt,text=Muted,text width=9.0cm,font=\fontsize{7.7}{9.3}\selectfont] at ([xshift=.95cm,yshift=-12.06cm]current page.north west) {05 · Errore standard};
\node[anchor=north west,inner sep=0pt,text=Classical,text width=8.2cm,align=right,font=\fontsize{7.3}{9.1}\selectfont] at ([xshift=9.45cm,yshift=-12.06cm]current page.north west) {\href{https://openintrostat.github.io/ims/}{IMS} · \href{https://stat.ethz.ch/R-manual/R-devel/library/stats/html/00Index.html}{RSTATS}};
\node[anchor=north east,inner sep=0pt,text=Muted,font=\fontsize{7.7}{9.3}\selectfont] at ([xshift=19.05cm,yshift=-12.06cm]current page.north west) {52/110};

\end{tikzpicture}
\note{Idea della slide: La presenza di molti dati non garantisce un piccolo errore in ogni contesto. Dipendenza, pesi, disegno a gruppi e variabilità elevata modificano il calcolo. Anche un intervallo molto stretto può non includere errori sistematici o incertezza del protocollo.\par\medskip Punto da far emergere: L'errore standard quantifica una componente dell'incertezza sotto un modello.\par\medskip Fonti della sezione: IMS, RSTATS}
\end{frame}
\begin{lrbox}{\TitleBox}\begin{minipage}[t]{18.2cm}\color{Ink}\raggedright\sloppy\fontsize{20.7}{24}\selectfont
\bfseries\color{Navy}Etichettare le barre di errore
\end{minipage}\end{lrbox}
\begin{lrbox}{\BodyBox}\begin{minipage}[t]{10.25cm}\color{Ink}\raggedright\sloppy\fontsize{12.4}{16.4}\selectfont
Barre chiamate genericamente «errore» sono ambigue. Possono rappresentare deviazione standard, errore standard, intervallo di confidenza o intervallo di previsione. La scelta deve essere scritta nella didascalia, perché queste barre rispondono a domande diverse.
\end{minipage}\end{lrbox}
\begin{lrbox}{\SideBox}\begin{minipage}[t]{7.45cm}\color{Ink}\raggedright\sloppy\fontsize{12.2}{16.1}\selectfont
«Media \(\pm\) deviazione standard»: eterogeneità.\par\smallskip «Media e IC 95\%»: incertezza della media.
\end{minipage}\end{lrbox}
\begin{lrbox}{\RecallBox}\begin{minipage}[t]{16.45cm}\color{Ink}\raggedright\sloppy\fontsize{10.1}{13.2}\selectfont
Non confrontare barre di tipo diverso come se avessero lo stesso significato.
\end{minipage}\end{lrbox}
\MeasureBox{53-TitleBox}{\TitleBox}
\MeasureBox{53-BodyBox}{\BodyBox}
\MeasureBox{53-SideBox}{\SideBox}
\MeasureBox{53-RecallBox}{\RecallBox}
\begin{frame}[plain]
\hypertarget{slide53}{}\bookmark[level=1,page=53]{Etichettare le barre di errore}
\begin{tikzpicture}[remember picture,overlay]
\node[anchor=north east,inner sep=0pt,text=Classical,font=\fontsize{8}{10}\selectfont] at ([xshift=19.05cm,yshift=-.16cm]current page.north west) {\hyperlink{slide3}{Indice}};
\node[anchor=north west,inner sep=0pt,text=Navy,font=\bfseries] at ([xshift=.95cm,yshift=-.62cm]current page.north west) {\usebox{\TitleBox}};
\draw[Classical,line width=1.3pt] ([xshift=.95cm,yshift=-1.68cm]current page.north west)--([xshift=2.25cm,yshift=-1.68cm]current page.north west);
\node[anchor=north west,inner sep=0pt] at ([xshift=.95cm,yshift=-2.12cm]current page.north west) {\usebox{\BodyBox}};
\node[anchor=north west,inner sep=0pt,text=Quantum,font=\bfseries\fontsize{11.5}{14}\selectfont] at ([xshift=11.6cm,yshift=-2.1cm]current page.north west) {Esempio e interpretazione};
\node[anchor=north west,inner sep=0pt] at ([xshift=11.6cm,yshift=-2.75cm]current page.north west) {\usebox{\SideBox}};
\draw[Rule,line width=.5pt] ([xshift=.95cm,yshift=-10.43cm]current page.north west)--([xshift=19.05cm,yshift=-10.43cm]current page.north west);
\node[anchor=north west,inner sep=0pt,text=Classical,font=\bfseries\fontsize{9.4}{12}\selectfont] at ([xshift=.95cm,yshift=-10.66cm]current page.north west) {Richiami};
\node[anchor=north west,inner sep=0pt] at ([xshift=2.6cm,yshift=-10.64cm]current page.north west) {\usebox{\RecallBox}};
\draw[Rule,line width=.35pt] ([xshift=.95cm,yshift=-11.92cm]current page.north west)--([xshift=19.05cm,yshift=-11.92cm]current page.north west);
\node[anchor=north west,inner sep=0pt,text=Muted,text width=9.0cm,font=\fontsize{7.7}{9.3}\selectfont] at ([xshift=.95cm,yshift=-12.06cm]current page.north west) {05 · Errore standard};
\node[anchor=north west,inner sep=0pt,text=Classical,text width=8.2cm,align=right,font=\fontsize{7.3}{9.1}\selectfont] at ([xshift=9.45cm,yshift=-12.06cm]current page.north west) {\href{https://openintrostat.github.io/ims/}{IMS} · \href{https://stat.ethz.ch/R-manual/R-devel/library/stats/html/00Index.html}{RSTATS}};
\node[anchor=north east,inner sep=0pt,text=Muted,font=\fontsize{7.7}{9.3}\selectfont] at ([xshift=19.05cm,yshift=-12.06cm]current page.north west) {53/110};

\end{tikzpicture}
\note{Idea della slide: Barre chiamate genericamente «errore» sono ambigue. Possono rappresentare deviazione standard, errore standard, intervallo di confidenza o intervallo di previsione. La scelta deve essere scritta nella didascalia, perché queste barre rispondono a domande diverse.\par\medskip Punto da far emergere: Non confrontare barre di tipo diverso come se avessero lo stesso significato.\par\medskip Fonti della sezione: IMS, RSTATS}
\end{frame}
\begin{lrbox}{\TitleBox}\begin{minipage}[t]{18.2cm}\color{Ink}\raggedright\sloppy\fontsize{20.7}{24}\selectfont
\bfseries\color{Navy}Attività 05 · applica il concetto
\end{minipage}\end{lrbox}
\begin{lrbox}{\BodyBox}\begin{minipage}[t]{10.25cm}\color{Ink}\raggedright\sloppy\fontsize{12.4}{16.4}\selectfont
Due campioni hanno deviazione standard 1 \(\mu\)m e numerosità 25 e 100. Calcola l'errore standard ideale della media e spiega che cosa cambia.
\end{minipage}\end{lrbox}
\begin{lrbox}{\SideBox}\begin{minipage}[t]{7.45cm}\color{Ink}\raggedright\sloppy\fontsize{12.2}{16.1}\selectfont
\noindent{\color{Classical}\bfseries 01}\hspace{.35cm}Individua la domanda\par\vspace{.22cm}\noindent{\color{Classical}\bfseries 02}\hspace{.35cm}Esplicita le ipotesi\par\vspace{.22cm}\noindent{\color{Classical}\bfseries 03}\hspace{.35cm}Calcola o rappresenta\par\vspace{.22cm}\noindent{\color{Classical}\bfseries 04}\hspace{.35cm}Scrivi una conclusione\par\vspace{.22cm}\noindent{\color{Classical}\bfseries 05}\hspace{.35cm}Indica un limite\par\vspace{.22cm}
\end{minipage}\end{lrbox}
\begin{lrbox}{\RecallBox}\begin{minipage}[t]{16.45cm}\color{Ink}\raggedright\sloppy\fontsize{10.1}{13.2}\selectfont
Quanto è incerta la media, rispetto a quanto variano gli individui?
\end{minipage}\end{lrbox}
\MeasureBox{54-TitleBox}{\TitleBox}
\MeasureBox{54-BodyBox}{\BodyBox}
\MeasureBox{54-SideBox}{\SideBox}
\MeasureBox{54-RecallBox}{\RecallBox}
\begin{frame}[plain]
\hypertarget{slide54}{}\bookmark[level=1,page=54]{Attività 05 · applica il concetto}
\begin{tikzpicture}[remember picture,overlay]
\node[anchor=north east,inner sep=0pt,text=Classical,font=\fontsize{8}{10}\selectfont] at ([xshift=19.05cm,yshift=-.16cm]current page.north west) {\hyperlink{slide3}{Indice}};
\node[anchor=north west,inner sep=0pt,text=Navy,font=\bfseries] at ([xshift=.95cm,yshift=-.62cm]current page.north west) {\usebox{\TitleBox}};
\draw[Classical,line width=1.3pt] ([xshift=.95cm,yshift=-1.68cm]current page.north west)--([xshift=2.25cm,yshift=-1.68cm]current page.north west);
\node[anchor=north west,inner sep=0pt] at ([xshift=.95cm,yshift=-2.12cm]current page.north west) {\usebox{\BodyBox}};
\node[anchor=north west,inner sep=0pt,text=Quantum,font=\bfseries\fontsize{11.5}{14}\selectfont] at ([xshift=11.6cm,yshift=-2.1cm]current page.north west) {Consegna di lavoro};
\node[anchor=north west,inner sep=0pt] at ([xshift=11.6cm,yshift=-2.75cm]current page.north west) {\usebox{\SideBox}};
\draw[Rule,line width=.5pt] ([xshift=.95cm,yshift=-10.43cm]current page.north west)--([xshift=19.05cm,yshift=-10.43cm]current page.north west);
\node[anchor=north west,inner sep=0pt,text=Classical,font=\bfseries\fontsize{9.4}{12}\selectfont] at ([xshift=.95cm,yshift=-10.66cm]current page.north west) {Richiami};
\node[anchor=north west,inner sep=0pt] at ([xshift=2.6cm,yshift=-10.64cm]current page.north west) {\usebox{\RecallBox}};
\draw[Rule,line width=.35pt] ([xshift=.95cm,yshift=-11.92cm]current page.north west)--([xshift=19.05cm,yshift=-11.92cm]current page.north west);
\node[anchor=north west,inner sep=0pt,text=Muted,text width=9.0cm,font=\fontsize{7.7}{9.3}\selectfont] at ([xshift=.95cm,yshift=-12.06cm]current page.north west) {05 · Errore standard};
\node[anchor=north west,inner sep=0pt,text=Classical,text width=8.2cm,align=right,font=\fontsize{7.3}{9.1}\selectfont] at ([xshift=9.45cm,yshift=-12.06cm]current page.north west) {\href{https://openintrostat.github.io/ims/}{IMS} · \href{https://stat.ethz.ch/R-manual/R-devel/library/stats/html/00Index.html}{RSTATS}};
\node[anchor=north east,inner sep=0pt,text=Muted,font=\fontsize{7.7}{9.3}\selectfont] at ([xshift=19.05cm,yshift=-12.06cm]current page.north west) {54/110};

\end{tikzpicture}
\note{Consegna: Due campioni hanno deviazione standard 1 \(\mu\)m e numerosità 25 e 100. Calcola l'errore standard ideale della media e spiega che cosa cambia.\par\medskip Soluzione per il docente: Gli errori standard stimati sono 1/5 = 0,20 \(\mu\)m e 1/10 = 0,10 \(\mu\)m. Quadruplicare n dimezza l'errore standard nel modello indipendente. La dispersione dei singoli individui resta circa 1 \(\mu\)m: non si è dimezzata la variabilità naturale.}
\end{frame}
\begin{lrbox}{\TitleBox}\begin{minipage}[t]{18.2cm}\color{Ink}\raggedright\sloppy\fontsize{20.7}{24}\selectfont
\bfseries\color{Navy}Discussione 05 · soluzione ragionata
\end{minipage}\end{lrbox}
\begin{lrbox}{\BodyBox}\begin{minipage}[t]{10.25cm}\color{Ink}\raggedright\sloppy\fontsize{12.4}{16.4}\selectfont
Gli errori standard stimati sono 1/5 = 0,20 \(\mu\)m e 1/10 = 0,10 \(\mu\)m. Quadruplicare n dimezza l'errore standard nel modello indipendente. La dispersione dei singoli individui resta circa 1 \(\mu\)m: non si è dimezzata la variabilità naturale.
\end{minipage}\end{lrbox}
\begin{lrbox}{\SideBox}\begin{minipage}[t]{7.45cm}\color{Ink}\raggedright\sloppy\fontsize{12.2}{16.1}\selectfont
Non confrontare barre di tipo diverso come se avessero lo stesso significato.
\end{minipage}\end{lrbox}
\begin{lrbox}{\RecallBox}\begin{minipage}[t]{16.45cm}\color{Ink}\raggedright\sloppy\fontsize{10.1}{13.2}\selectfont
Valuta il collegamento tra domanda, procedimento e conclusione, non soltanto il numero finale.
\end{minipage}\end{lrbox}
\MeasureBox{55-TitleBox}{\TitleBox}
\MeasureBox{55-BodyBox}{\BodyBox}
\MeasureBox{55-SideBox}{\SideBox}
\MeasureBox{55-RecallBox}{\RecallBox}
\begin{frame}[plain]
\hypertarget{slide55}{}\bookmark[level=1,page=55]{Discussione 05 · soluzione ragionata}
\begin{tikzpicture}[remember picture,overlay]
\node[anchor=north east,inner sep=0pt,text=Classical,font=\fontsize{8}{10}\selectfont] at ([xshift=19.05cm,yshift=-.16cm]current page.north west) {\hyperlink{slide3}{Indice}};
\node[anchor=north west,inner sep=0pt,text=Navy,font=\bfseries] at ([xshift=.95cm,yshift=-.62cm]current page.north west) {\usebox{\TitleBox}};
\draw[Classical,line width=1.3pt] ([xshift=.95cm,yshift=-1.68cm]current page.north west)--([xshift=2.25cm,yshift=-1.68cm]current page.north west);
\node[anchor=north west,inner sep=0pt] at ([xshift=.95cm,yshift=-2.12cm]current page.north west) {\usebox{\BodyBox}};
\node[anchor=north west,inner sep=0pt,text=Quantum,font=\bfseries\fontsize{11.5}{14}\selectfont] at ([xshift=11.6cm,yshift=-2.1cm]current page.north west) {Criterio di lettura};
\node[anchor=north west,inner sep=0pt] at ([xshift=11.6cm,yshift=-2.75cm]current page.north west) {\usebox{\SideBox}};
\draw[Rule,line width=.5pt] ([xshift=.95cm,yshift=-10.43cm]current page.north west)--([xshift=19.05cm,yshift=-10.43cm]current page.north west);
\node[anchor=north west,inner sep=0pt,text=Classical,font=\bfseries\fontsize{9.4}{12}\selectfont] at ([xshift=.95cm,yshift=-10.66cm]current page.north west) {Richiami};
\node[anchor=north west,inner sep=0pt] at ([xshift=2.6cm,yshift=-10.64cm]current page.north west) {\usebox{\RecallBox}};
\draw[Rule,line width=.35pt] ([xshift=.95cm,yshift=-11.92cm]current page.north west)--([xshift=19.05cm,yshift=-11.92cm]current page.north west);
\node[anchor=north west,inner sep=0pt,text=Muted,text width=9.0cm,font=\fontsize{7.7}{9.3}\selectfont] at ([xshift=.95cm,yshift=-12.06cm]current page.north west) {05 · Errore standard};
\node[anchor=north west,inner sep=0pt,text=Classical,text width=8.2cm,align=right,font=\fontsize{7.3}{9.1}\selectfont] at ([xshift=9.45cm,yshift=-12.06cm]current page.north west) {\href{https://openintrostat.github.io/ims/}{IMS} · \href{https://stat.ethz.ch/R-manual/R-devel/library/stats/html/00Index.html}{RSTATS}};
\node[anchor=north east,inner sep=0pt,text=Muted,font=\fontsize{7.7}{9.3}\selectfont] at ([xshift=19.05cm,yshift=-12.06cm]current page.north west) {55/110};

\end{tikzpicture}
\note{Gli errori standard stimati sono 1/5 = 0,20 \(\mu\)m e 1/10 = 0,10 \(\mu\)m. Quadruplicare n dimezza l'errore standard nel modello indipendente. La dispersione dei singoli individui resta circa 1 \(\mu\)m: non si è dimezzata la variabilità naturale.}
\end{frame}
\begin{frame}[plain]\hypertarget{slide56}{}
\bookmark[level=0,page=56]{Intervalli di confidenza}
\begin{tikzpicture}[remember picture,overlay]
\fill[Night] (current page.south west) rectangle (current page.north east);
\draw[Mint,line width=2pt] ([xshift=1.05cm,yshift=-1.35cm]current page.north west)--([xshift=3.0cm,yshift=-1.35cm]current page.north west);
\node[anchor=north west,inner sep=0,text=Mint,font=\bfseries\fontsize{11}{14}\selectfont] at ([xshift=1.05cm,yshift=-1.7cm]current page.north west) {SEZIONE 06};
\node[anchor=north west,inner sep=0,text=white,text width=15.8cm,align=left,font=\bfseries\fontsize{31}{38}\selectfont] at ([xshift=1.05cm,yshift=-2.65cm]current page.north west) {Intervalli di confidenza};
\node[anchor=north west,inner sep=0,text=Lilac,text width=15.8cm,align=left,font=\fontsize{15}{21}\selectfont] at ([xshift=1.05cm,yshift=-6.4cm]current page.north west) {Che cosa significa il 95\% in un intervallo?};
\node[anchor=north west,inner sep=0,text=white,text width=16.5cm,align=left,font=\fontsize{11}{15}\selectfont] at ([xshift=1.05cm,yshift=-8.6cm]current page.north west) {Comprendere il concetto · leggere un esempio · applicare e discutere};
\node[anchor=north east,inner sep=0,text=Mint,font=\fontsize{8}{10}\selectfont] at ([xshift=19.05cm,yshift=-.2cm]current page.north west) {\hyperlink{slide3}{Indice}};
\node[anchor=south east,inner sep=0,text=Lilac!45,font=\bfseries\fontsize{64}{68}\selectfont] at ([xshift=18.95cm,yshift=-11.35cm]current page.north west) {06};
\draw[Lilac!30,line width=.35pt] ([xshift=.95cm,yshift=-11.92cm]current page.north west)--([xshift=19.05cm,yshift=-11.92cm]current page.north west);
\node[anchor=north west,inner sep=0pt,text=Lilac,text width=9.0cm,font=\fontsize{7.7}{9.3}\selectfont] at ([xshift=.95cm,yshift=-12.06cm]current page.north west) {06 · Intervalli di confidenza};
\node[anchor=north west,inner sep=0pt,text=Mint,text width=8.2cm,align=right,font=\fontsize{7.3}{9.1}\selectfont] at ([xshift=9.45cm,yshift=-12.06cm]current page.north west) {\href{https://openintrostat.github.io/ims/}{IMS} · \href{https://stat.ethz.ch/R-manual/R-devel/library/stats/html/t.test.html}{RTTEST}};
\node[anchor=north east,inner sep=0pt,text=Lilac,font=\fontsize{7.7}{9.3}\selectfont] at ([xshift=19.05cm,yshift=-12.06cm]current page.north west) {56/110};
\end{tikzpicture}\end{frame}
\begin{lrbox}{\TitleBox}\begin{minipage}[t]{18.2cm}\color{Ink}\raggedright\sloppy\fontsize{20.7}{24}\selectfont
\bfseries\color{Navy}Stima puntuale e intervallo
\end{minipage}\end{lrbox}
\begin{lrbox}{\BodyBox}\begin{minipage}[t]{10.25cm}\color{Ink}\raggedright\sloppy\fontsize{12.4}{16.4}\selectfont
Una stima puntuale dà un numero; un intervallo rende visibile quanto la procedura potrebbe variare tra campioni. L'intervallo non include necessariamente ogni fonte di errore: dipende da modello, selezione e misura. La sua utilità sta nella quantificazione esplicita di quella incertezza.
\end{minipage}\end{lrbox}
\begin{lrbox}{\SideBox}\begin{minipage}[t]{7.45cm}\color{Ink}\raggedright\sloppy\fontsize{12.2}{16.1}\selectfont
Media stimata + intervallo + popolazione + ipotesi formano un risultato leggibile.
\end{minipage}\end{lrbox}
\begin{lrbox}{\RecallBox}\begin{minipage}[t]{16.45cm}\color{Ink}\raggedright\sloppy\fontsize{10.1}{13.2}\selectfont
Il numero centrale senza incertezza può apparire più preciso di quanto sia.
\end{minipage}\end{lrbox}
\MeasureBox{57-TitleBox}{\TitleBox}
\MeasureBox{57-BodyBox}{\BodyBox}
\MeasureBox{57-SideBox}{\SideBox}
\MeasureBox{57-RecallBox}{\RecallBox}
\begin{frame}[plain]
\hypertarget{slide57}{}\bookmark[level=1,page=57]{Stima puntuale e intervallo}
\begin{tikzpicture}[remember picture,overlay]
\node[anchor=north east,inner sep=0pt,text=Classical,font=\fontsize{8}{10}\selectfont] at ([xshift=19.05cm,yshift=-.16cm]current page.north west) {\hyperlink{slide3}{Indice}};
\node[anchor=north west,inner sep=0pt,text=Navy,font=\bfseries] at ([xshift=.95cm,yshift=-.62cm]current page.north west) {\usebox{\TitleBox}};
\draw[Classical,line width=1.3pt] ([xshift=.95cm,yshift=-1.68cm]current page.north west)--([xshift=2.25cm,yshift=-1.68cm]current page.north west);
\node[anchor=north west,inner sep=0pt] at ([xshift=.95cm,yshift=-2.12cm]current page.north west) {\usebox{\BodyBox}};
\node[anchor=north west,inner sep=0pt,text=Quantum,font=\bfseries\fontsize{11.5}{14}\selectfont] at ([xshift=11.6cm,yshift=-2.1cm]current page.north west) {Esempio e interpretazione};
\node[anchor=north west,inner sep=0pt] at ([xshift=11.6cm,yshift=-2.75cm]current page.north west) {\usebox{\SideBox}};
\draw[Rule,line width=.5pt] ([xshift=.95cm,yshift=-10.43cm]current page.north west)--([xshift=19.05cm,yshift=-10.43cm]current page.north west);
\node[anchor=north west,inner sep=0pt,text=Classical,font=\bfseries\fontsize{9.4}{12}\selectfont] at ([xshift=.95cm,yshift=-10.66cm]current page.north west) {Richiami};
\node[anchor=north west,inner sep=0pt] at ([xshift=2.6cm,yshift=-10.64cm]current page.north west) {\usebox{\RecallBox}};
\draw[Rule,line width=.35pt] ([xshift=.95cm,yshift=-11.92cm]current page.north west)--([xshift=19.05cm,yshift=-11.92cm]current page.north west);
\node[anchor=north west,inner sep=0pt,text=Muted,text width=9.0cm,font=\fontsize{7.7}{9.3}\selectfont] at ([xshift=.95cm,yshift=-12.06cm]current page.north west) {06 · Intervalli di confidenza};
\node[anchor=north west,inner sep=0pt,text=Classical,text width=8.2cm,align=right,font=\fontsize{7.3}{9.1}\selectfont] at ([xshift=9.45cm,yshift=-12.06cm]current page.north west) {\href{https://openintrostat.github.io/ims/}{IMS} · \href{https://stat.ethz.ch/R-manual/R-devel/library/stats/html/t.test.html}{RTTEST}};
\node[anchor=north east,inner sep=0pt,text=Muted,font=\fontsize{7.7}{9.3}\selectfont] at ([xshift=19.05cm,yshift=-12.06cm]current page.north west) {57/110};

\end{tikzpicture}
\note{Idea della slide: Una stima puntuale dà un numero; un intervallo rende visibile quanto la procedura potrebbe variare tra campioni. L'intervallo non include necessariamente ogni fonte di errore: dipende da modello, selezione e misura. La sua utilità sta nella quantificazione esplicita di quella incertezza.\par\medskip Punto da far emergere: Il numero centrale senza incertezza può apparire più preciso di quanto sia.\par\medskip Fonti della sezione: IMS, RTTEST}
\end{frame}
\begin{lrbox}{\TitleBox}\begin{minipage}[t]{18.2cm}\color{Ink}\raggedright\sloppy\fontsize{20.7}{24}\selectfont
\bfseries\color{Navy}La costruzione ripetuta
\end{minipage}\end{lrbox}
\begin{lrbox}{\BodyBox}\begin{minipage}[t]{7.1cm}\color{Ink}\raggedright\sloppy\fontsize{11.8}{15.8}\selectfont
Immaginiamo molti campioni e un intervallo costruito con la stessa regola per ciascuno. Alcuni intervalli non conterranno il parametro. La copertura è una proprietà della procedura nel modello, non una garanzia assoluta per il singolo intervallo ottenuto.
\end{minipage}\end{lrbox}
\begin{lrbox}{\SideBox}\begin{minipage}[t]{10.4cm}\color{Ink}\raggedright\sloppy\fontsize{12.2}{16.1}\selectfont
\centering\includegraphics[width=10.4cm,height=6.25cm,keepaspectratio]{figure/coverage.pdf}\par\vspace{.22cm}{\raggedright\color{Muted}\fontsize{9.7}{12.5}\selectfont Intervalli ottenuti da campioni normali simulati; linea verticale: media del generatore.\par}
\end{minipage}\end{lrbox}
\begin{lrbox}{\RecallBox}\begin{minipage}[t]{16.45cm}\color{Ink}\raggedright\sloppy\fontsize{10.1}{13.2}\selectfont
Non sappiamo quale singolo intervallo sia corretto quando il parametro reale è ignoto.
\end{minipage}\end{lrbox}
\MeasureBox{58-TitleBox}{\TitleBox}
\MeasureBox{58-BodyBox}{\BodyBox}
\MeasureBox{58-SideBox}{\SideBox}
\MeasureBox{58-RecallBox}{\RecallBox}
\begin{frame}[plain]
\hypertarget{slide58}{}\bookmark[level=1,page=58]{La costruzione ripetuta}
\begin{tikzpicture}[remember picture,overlay]
\node[anchor=north east,inner sep=0pt,text=Classical,font=\fontsize{8}{10}\selectfont] at ([xshift=19.05cm,yshift=-.16cm]current page.north west) {\hyperlink{slide3}{Indice}};
\node[anchor=north west,inner sep=0pt,text=Navy,font=\bfseries] at ([xshift=.95cm,yshift=-.62cm]current page.north west) {\usebox{\TitleBox}};
\draw[Classical,line width=1.3pt] ([xshift=.95cm,yshift=-1.68cm]current page.north west)--([xshift=2.25cm,yshift=-1.68cm]current page.north west);
\node[anchor=north west,inner sep=0pt] at ([xshift=.95cm,yshift=-2.12cm]current page.north west) {\usebox{\BodyBox}};
\node[anchor=north west,inner sep=0pt,text=Quantum,font=\bfseries\fontsize{11.5}{14}\selectfont] at ([xshift=8.65cm,yshift=-2.1cm]current page.north west) {Figura commentata};
\node[anchor=north west,inner sep=0pt] at ([xshift=8.65cm,yshift=-2.75cm]current page.north west) {\usebox{\SideBox}};
\draw[Rule,line width=.5pt] ([xshift=.95cm,yshift=-10.43cm]current page.north west)--([xshift=19.05cm,yshift=-10.43cm]current page.north west);
\node[anchor=north west,inner sep=0pt,text=Classical,font=\bfseries\fontsize{9.4}{12}\selectfont] at ([xshift=.95cm,yshift=-10.66cm]current page.north west) {Richiami};
\node[anchor=north west,inner sep=0pt] at ([xshift=2.6cm,yshift=-10.64cm]current page.north west) {\usebox{\RecallBox}};
\draw[Rule,line width=.35pt] ([xshift=.95cm,yshift=-11.92cm]current page.north west)--([xshift=19.05cm,yshift=-11.92cm]current page.north west);
\node[anchor=north west,inner sep=0pt,text=Muted,text width=9.0cm,font=\fontsize{7.7}{9.3}\selectfont] at ([xshift=.95cm,yshift=-12.06cm]current page.north west) {06 · Intervalli di confidenza};
\node[anchor=north west,inner sep=0pt,text=Classical,text width=8.2cm,align=right,font=\fontsize{7.3}{9.1}\selectfont] at ([xshift=9.45cm,yshift=-12.06cm]current page.north west) {\href{https://openintrostat.github.io/ims/}{IMS} · \href{https://stat.ethz.ch/R-manual/R-devel/library/stats/html/t.test.html}{RTTEST}};
\node[anchor=north east,inner sep=0pt,text=Muted,font=\fontsize{7.7}{9.3}\selectfont] at ([xshift=19.05cm,yshift=-12.06cm]current page.north west) {58/110};

\end{tikzpicture}
\note{Idea della slide: Immaginiamo molti campioni e un intervallo costruito con la stessa regola per ciascuno. Alcuni intervalli non conterranno il parametro. La copertura è una proprietà della procedura nel modello, non una garanzia assoluta per il singolo intervallo ottenuto.\par\medskip Punto da far emergere: Non sappiamo quale singolo intervallo sia corretto quando il parametro reale è ignoto.\par\medskip Fonti della sezione: IMS, RTTEST}
\end{frame}
\begin{lrbox}{\TitleBox}\begin{minipage}[t]{18.2cm}\color{Ink}\raggedright\sloppy\fontsize{20.7}{24}\selectfont
\bfseries\color{Navy}Il parametro resta fisso
\end{minipage}\end{lrbox}
\begin{lrbox}{\BodyBox}\begin{minipage}[t]{10.25cm}\color{Ink}\raggedright\sloppy\fontsize{12.4}{16.4}\selectfont
Nell'interpretazione frequentista il parametro è fisso, mentre il campione e l'intervallo sono casuali prima dell'osservazione. Dopo aver osservato i dati, l'intervallo contiene o non contiene quel parametro. Questa distinzione evita interpretazioni probabilistiche non giustificate dalla procedura.
\end{minipage}\end{lrbox}
\begin{lrbox}{\SideBox}\begin{minipage}[t]{7.45cm}\color{Ink}\raggedright\sloppy\fontsize{12.2}{16.1}\selectfont
Il 95\% riguarda il comportamento del metodo su ripetizioni ideali, non una distribuzione del parametro costruita automaticamente.
\end{minipage}\end{lrbox}
\begin{lrbox}{\RecallBox}\begin{minipage}[t]{16.45cm}\color{Ink}\raggedright\sloppy\fontsize{10.1}{13.2}\selectfont
Un intervallo bayesiano richiede un'altra impostazione e non viene introdotto qui come sinonimo.
\end{minipage}\end{lrbox}
\MeasureBox{59-TitleBox}{\TitleBox}
\MeasureBox{59-BodyBox}{\BodyBox}
\MeasureBox{59-SideBox}{\SideBox}
\MeasureBox{59-RecallBox}{\RecallBox}
\begin{frame}[plain]
\hypertarget{slide59}{}\bookmark[level=1,page=59]{Il parametro resta fisso}
\begin{tikzpicture}[remember picture,overlay]
\node[anchor=north east,inner sep=0pt,text=Classical,font=\fontsize{8}{10}\selectfont] at ([xshift=19.05cm,yshift=-.16cm]current page.north west) {\hyperlink{slide3}{Indice}};
\node[anchor=north west,inner sep=0pt,text=Navy,font=\bfseries] at ([xshift=.95cm,yshift=-.62cm]current page.north west) {\usebox{\TitleBox}};
\draw[Classical,line width=1.3pt] ([xshift=.95cm,yshift=-1.68cm]current page.north west)--([xshift=2.25cm,yshift=-1.68cm]current page.north west);
\node[anchor=north west,inner sep=0pt] at ([xshift=.95cm,yshift=-2.12cm]current page.north west) {\usebox{\BodyBox}};
\node[anchor=north west,inner sep=0pt,text=Quantum,font=\bfseries\fontsize{11.5}{14}\selectfont] at ([xshift=11.6cm,yshift=-2.1cm]current page.north west) {Esempio e interpretazione};
\node[anchor=north west,inner sep=0pt] at ([xshift=11.6cm,yshift=-2.75cm]current page.north west) {\usebox{\SideBox}};
\draw[Rule,line width=.5pt] ([xshift=.95cm,yshift=-10.43cm]current page.north west)--([xshift=19.05cm,yshift=-10.43cm]current page.north west);
\node[anchor=north west,inner sep=0pt,text=Classical,font=\bfseries\fontsize{9.4}{12}\selectfont] at ([xshift=.95cm,yshift=-10.66cm]current page.north west) {Richiami};
\node[anchor=north west,inner sep=0pt] at ([xshift=2.6cm,yshift=-10.64cm]current page.north west) {\usebox{\RecallBox}};
\draw[Rule,line width=.35pt] ([xshift=.95cm,yshift=-11.92cm]current page.north west)--([xshift=19.05cm,yshift=-11.92cm]current page.north west);
\node[anchor=north west,inner sep=0pt,text=Muted,text width=9.0cm,font=\fontsize{7.7}{9.3}\selectfont] at ([xshift=.95cm,yshift=-12.06cm]current page.north west) {06 · Intervalli di confidenza};
\node[anchor=north west,inner sep=0pt,text=Classical,text width=8.2cm,align=right,font=\fontsize{7.3}{9.1}\selectfont] at ([xshift=9.45cm,yshift=-12.06cm]current page.north west) {\href{https://openintrostat.github.io/ims/}{IMS} · \href{https://stat.ethz.ch/R-manual/R-devel/library/stats/html/t.test.html}{RTTEST}};
\node[anchor=north east,inner sep=0pt,text=Muted,font=\fontsize{7.7}{9.3}\selectfont] at ([xshift=19.05cm,yshift=-12.06cm]current page.north west) {59/110};

\end{tikzpicture}
\note{Idea della slide: Nell'interpretazione frequentista il parametro è fisso, mentre il campione e l'intervallo sono casuali prima dell'osservazione. Dopo aver osservato i dati, l'intervallo contiene o non contiene quel parametro. Questa distinzione evita interpretazioni probabilistiche non giustificate dalla procedura.\par\medskip Punto da far emergere: Un intervallo bayesiano richiede un'altra impostazione e non viene introdotto qui come sinonimo.\par\medskip Fonti della sezione: IMS, RTTEST}
\end{frame}
\begin{lrbox}{\TitleBox}\begin{minipage}[t]{18.2cm}\color{Ink}\raggedright\sloppy\fontsize{20.7}{24}\selectfont
\bfseries\color{Navy}Intervallo t per una media
\end{minipage}\end{lrbox}
\begin{lrbox}{\BodyBox}\begin{minipage}[t]{10.25cm}\color{Ink}\raggedright\sloppy\fontsize{12.4}{16.4}\selectfont
Per dati indipendenti compatibili con le condizioni del modello, l'intervallo usa la media, l'errore standard e un quantile t. Nei piccoli campioni la normalità delle osservazioni è rilevante per l'esattezza del metodo; con campioni ampi la robustezza dipende dal caso.
\end{minipage}\end{lrbox}
\begin{lrbox}{\SideBox}\begin{minipage}[t]{7.45cm}\color{Ink}\raggedright\sloppy\fontsize{12.2}{16.1}\selectfont
$\bar x\pm t_{0.975,n-1}\,s/\sqrt n$\par\medskip In R: t.test(x)\$conf.int.
\end{minipage}\end{lrbox}
\begin{lrbox}{\RecallBox}\begin{minipage}[t]{16.45cm}\color{Ink}\raggedright\sloppy\fontsize{10.1}{13.2}\selectfont
La formula non corregge automaticamente outlier influenti, bias o dipendenza.
\end{minipage}\end{lrbox}
\MeasureBox{60-TitleBox}{\TitleBox}
\MeasureBox{60-BodyBox}{\BodyBox}
\MeasureBox{60-SideBox}{\SideBox}
\MeasureBox{60-RecallBox}{\RecallBox}
\begin{frame}[plain]
\hypertarget{slide60}{}\bookmark[level=1,page=60]{Intervallo t per una media}
\begin{tikzpicture}[remember picture,overlay]
\node[anchor=north east,inner sep=0pt,text=Classical,font=\fontsize{8}{10}\selectfont] at ([xshift=19.05cm,yshift=-.16cm]current page.north west) {\hyperlink{slide3}{Indice}};
\node[anchor=north west,inner sep=0pt,text=Navy,font=\bfseries] at ([xshift=.95cm,yshift=-.62cm]current page.north west) {\usebox{\TitleBox}};
\draw[Classical,line width=1.3pt] ([xshift=.95cm,yshift=-1.68cm]current page.north west)--([xshift=2.25cm,yshift=-1.68cm]current page.north west);
\node[anchor=north west,inner sep=0pt] at ([xshift=.95cm,yshift=-2.12cm]current page.north west) {\usebox{\BodyBox}};
\node[anchor=north west,inner sep=0pt,text=Quantum,font=\bfseries\fontsize{11.5}{14}\selectfont] at ([xshift=11.6cm,yshift=-2.1cm]current page.north west) {Esempio e interpretazione};
\node[anchor=north west,inner sep=0pt] at ([xshift=11.6cm,yshift=-2.75cm]current page.north west) {\usebox{\SideBox}};
\draw[Rule,line width=.5pt] ([xshift=.95cm,yshift=-10.43cm]current page.north west)--([xshift=19.05cm,yshift=-10.43cm]current page.north west);
\node[anchor=north west,inner sep=0pt,text=Classical,font=\bfseries\fontsize{9.4}{12}\selectfont] at ([xshift=.95cm,yshift=-10.66cm]current page.north west) {Richiami};
\node[anchor=north west,inner sep=0pt] at ([xshift=2.6cm,yshift=-10.64cm]current page.north west) {\usebox{\RecallBox}};
\draw[Rule,line width=.35pt] ([xshift=.95cm,yshift=-11.92cm]current page.north west)--([xshift=19.05cm,yshift=-11.92cm]current page.north west);
\node[anchor=north west,inner sep=0pt,text=Muted,text width=9.0cm,font=\fontsize{7.7}{9.3}\selectfont] at ([xshift=.95cm,yshift=-12.06cm]current page.north west) {06 · Intervalli di confidenza};
\node[anchor=north west,inner sep=0pt,text=Classical,text width=8.2cm,align=right,font=\fontsize{7.3}{9.1}\selectfont] at ([xshift=9.45cm,yshift=-12.06cm]current page.north west) {\href{https://openintrostat.github.io/ims/}{IMS} · \href{https://stat.ethz.ch/R-manual/R-devel/library/stats/html/t.test.html}{RTTEST}};
\node[anchor=north east,inner sep=0pt,text=Muted,font=\fontsize{7.7}{9.3}\selectfont] at ([xshift=19.05cm,yshift=-12.06cm]current page.north west) {60/110};

\end{tikzpicture}
\note{Idea della slide: Per dati indipendenti compatibili con le condizioni del modello, l'intervallo usa la media, l'errore standard e un quantile t. Nei piccoli campioni la normalità delle osservazioni è rilevante per l'esattezza del metodo; con campioni ampi la robustezza dipende dal caso.\par\medskip Punto da far emergere: La formula non corregge automaticamente outlier influenti, bias o dipendenza.\par\medskip Fonti della sezione: IMS, RTTEST}
\end{frame}
\begin{lrbox}{\TitleBox}\begin{minipage}[t]{18.2cm}\color{Ink}\raggedright\sloppy\fontsize{20.7}{24}\selectfont
\bfseries\color{Navy}Livello di confidenza e ampiezza
\end{minipage}\end{lrbox}
\begin{lrbox}{\BodyBox}\begin{minipage}[t]{10.25cm}\color{Ink}\raggedright\sloppy\fontsize{12.4}{16.4}\selectfont
A parità di dati, un livello più alto produce in genere un intervallo più ampio. Stiamo chiedendo alla procedura una maggiore copertura, non rendendo i dati più precisi. Il livello va scelto e dichiarato senza modificarlo a posteriori per ottenere il messaggio desiderato.
\end{minipage}\end{lrbox}
\begin{lrbox}{\SideBox}\begin{minipage}[t]{7.45cm}\color{Ink}\raggedright\sloppy\fontsize{12.2}{16.1}\selectfont
IC 90\%: più stretto.\par\smallskip IC 99\%: più ampio.\par\smallskip Stessi dati, diversa regola di copertura.
\end{minipage}\end{lrbox}
\begin{lrbox}{\RecallBox}\begin{minipage}[t]{16.45cm}\color{Ink}\raggedright\sloppy\fontsize{10.1}{13.2}\selectfont
Un intervallo largo può essere la rappresentazione onesta dell'informazione disponibile.
\end{minipage}\end{lrbox}
\MeasureBox{61-TitleBox}{\TitleBox}
\MeasureBox{61-BodyBox}{\BodyBox}
\MeasureBox{61-SideBox}{\SideBox}
\MeasureBox{61-RecallBox}{\RecallBox}
\begin{frame}[plain]
\hypertarget{slide61}{}\bookmark[level=1,page=61]{Livello di confidenza e ampiezza}
\begin{tikzpicture}[remember picture,overlay]
\node[anchor=north east,inner sep=0pt,text=Classical,font=\fontsize{8}{10}\selectfont] at ([xshift=19.05cm,yshift=-.16cm]current page.north west) {\hyperlink{slide3}{Indice}};
\node[anchor=north west,inner sep=0pt,text=Navy,font=\bfseries] at ([xshift=.95cm,yshift=-.62cm]current page.north west) {\usebox{\TitleBox}};
\draw[Classical,line width=1.3pt] ([xshift=.95cm,yshift=-1.68cm]current page.north west)--([xshift=2.25cm,yshift=-1.68cm]current page.north west);
\node[anchor=north west,inner sep=0pt] at ([xshift=.95cm,yshift=-2.12cm]current page.north west) {\usebox{\BodyBox}};
\node[anchor=north west,inner sep=0pt,text=Quantum,font=\bfseries\fontsize{11.5}{14}\selectfont] at ([xshift=11.6cm,yshift=-2.1cm]current page.north west) {Esempio e interpretazione};
\node[anchor=north west,inner sep=0pt] at ([xshift=11.6cm,yshift=-2.75cm]current page.north west) {\usebox{\SideBox}};
\draw[Rule,line width=.5pt] ([xshift=.95cm,yshift=-10.43cm]current page.north west)--([xshift=19.05cm,yshift=-10.43cm]current page.north west);
\node[anchor=north west,inner sep=0pt,text=Classical,font=\bfseries\fontsize{9.4}{12}\selectfont] at ([xshift=.95cm,yshift=-10.66cm]current page.north west) {Richiami};
\node[anchor=north west,inner sep=0pt] at ([xshift=2.6cm,yshift=-10.64cm]current page.north west) {\usebox{\RecallBox}};
\draw[Rule,line width=.35pt] ([xshift=.95cm,yshift=-11.92cm]current page.north west)--([xshift=19.05cm,yshift=-11.92cm]current page.north west);
\node[anchor=north west,inner sep=0pt,text=Muted,text width=9.0cm,font=\fontsize{7.7}{9.3}\selectfont] at ([xshift=.95cm,yshift=-12.06cm]current page.north west) {06 · Intervalli di confidenza};
\node[anchor=north west,inner sep=0pt,text=Classical,text width=8.2cm,align=right,font=\fontsize{7.3}{9.1}\selectfont] at ([xshift=9.45cm,yshift=-12.06cm]current page.north west) {\href{https://openintrostat.github.io/ims/}{IMS} · \href{https://stat.ethz.ch/R-manual/R-devel/library/stats/html/t.test.html}{RTTEST}};
\node[anchor=north east,inner sep=0pt,text=Muted,font=\fontsize{7.7}{9.3}\selectfont] at ([xshift=19.05cm,yshift=-12.06cm]current page.north west) {61/110};

\end{tikzpicture}
\note{Idea della slide: A parità di dati, un livello più alto produce in genere un intervallo più ampio. Stiamo chiedendo alla procedura una maggiore copertura, non rendendo i dati più precisi. Il livello va scelto e dichiarato senza modificarlo a posteriori per ottenere il messaggio desiderato.\par\medskip Punto da far emergere: Un intervallo largo può essere la rappresentazione onesta dell'informazione disponibile.\par\medskip Fonti della sezione: IMS, RTTEST}
\end{frame}
\begin{lrbox}{\TitleBox}\begin{minipage}[t]{18.2cm}\color{Ink}\raggedright\sloppy\fontsize{20.7}{24}\selectfont
\bfseries\color{Navy}Intervallo della media e individui
\end{minipage}\end{lrbox}
\begin{lrbox}{\BodyBox}\begin{minipage}[t]{10.25cm}\color{Ink}\raggedright\sloppy\fontsize{12.4}{16.4}\selectfont
L'intervallo della media può essere stretto anche quando gli individui sono molto variabili. Un intervallo per una nuova osservazione deve includere anche la dispersione individuale e risulta più ampio. Questa distinzione tornerà nel modulo di regressione.
\end{minipage}\end{lrbox}
\begin{lrbox}{\SideBox}\begin{minipage}[t]{7.45cm}\color{Ink}\raggedright\sloppy\fontsize{12.2}{16.1}\selectfont
Con n grande conosciamo bene la media, ma non per questo possiamo prevedere con precisione il diametro di ogni individuo.
\end{minipage}\end{lrbox}
\begin{lrbox}{\RecallBox}\begin{minipage}[t]{16.45cm}\color{Ink}\raggedright\sloppy\fontsize{10.1}{13.2}\selectfont
Confidenza sul parametro e previsione del singolo valore non sono la stessa domanda.
\end{minipage}\end{lrbox}
\MeasureBox{62-TitleBox}{\TitleBox}
\MeasureBox{62-BodyBox}{\BodyBox}
\MeasureBox{62-SideBox}{\SideBox}
\MeasureBox{62-RecallBox}{\RecallBox}
\begin{frame}[plain]
\hypertarget{slide62}{}\bookmark[level=1,page=62]{Intervallo della media e individui}
\begin{tikzpicture}[remember picture,overlay]
\node[anchor=north east,inner sep=0pt,text=Classical,font=\fontsize{8}{10}\selectfont] at ([xshift=19.05cm,yshift=-.16cm]current page.north west) {\hyperlink{slide3}{Indice}};
\node[anchor=north west,inner sep=0pt,text=Navy,font=\bfseries] at ([xshift=.95cm,yshift=-.62cm]current page.north west) {\usebox{\TitleBox}};
\draw[Classical,line width=1.3pt] ([xshift=.95cm,yshift=-1.68cm]current page.north west)--([xshift=2.25cm,yshift=-1.68cm]current page.north west);
\node[anchor=north west,inner sep=0pt] at ([xshift=.95cm,yshift=-2.12cm]current page.north west) {\usebox{\BodyBox}};
\node[anchor=north west,inner sep=0pt,text=Quantum,font=\bfseries\fontsize{11.5}{14}\selectfont] at ([xshift=11.6cm,yshift=-2.1cm]current page.north west) {Esempio e interpretazione};
\node[anchor=north west,inner sep=0pt] at ([xshift=11.6cm,yshift=-2.75cm]current page.north west) {\usebox{\SideBox}};
\draw[Rule,line width=.5pt] ([xshift=.95cm,yshift=-10.43cm]current page.north west)--([xshift=19.05cm,yshift=-10.43cm]current page.north west);
\node[anchor=north west,inner sep=0pt,text=Classical,font=\bfseries\fontsize{9.4}{12}\selectfont] at ([xshift=.95cm,yshift=-10.66cm]current page.north west) {Richiami};
\node[anchor=north west,inner sep=0pt] at ([xshift=2.6cm,yshift=-10.64cm]current page.north west) {\usebox{\RecallBox}};
\draw[Rule,line width=.35pt] ([xshift=.95cm,yshift=-11.92cm]current page.north west)--([xshift=19.05cm,yshift=-11.92cm]current page.north west);
\node[anchor=north west,inner sep=0pt,text=Muted,text width=9.0cm,font=\fontsize{7.7}{9.3}\selectfont] at ([xshift=.95cm,yshift=-12.06cm]current page.north west) {06 · Intervalli di confidenza};
\node[anchor=north west,inner sep=0pt,text=Classical,text width=8.2cm,align=right,font=\fontsize{7.3}{9.1}\selectfont] at ([xshift=9.45cm,yshift=-12.06cm]current page.north west) {\href{https://openintrostat.github.io/ims/}{IMS} · \href{https://stat.ethz.ch/R-manual/R-devel/library/stats/html/t.test.html}{RTTEST}};
\node[anchor=north east,inner sep=0pt,text=Muted,font=\fontsize{7.7}{9.3}\selectfont] at ([xshift=19.05cm,yshift=-12.06cm]current page.north west) {62/110};

\end{tikzpicture}
\note{Idea della slide: L'intervallo della media può essere stretto anche quando gli individui sono molto variabili. Un intervallo per una nuova osservazione deve includere anche la dispersione individuale e risulta più ampio. Questa distinzione tornerà nel modulo di regressione.\par\medskip Punto da far emergere: Confidenza sul parametro e previsione del singolo valore non sono la stessa domanda.\par\medskip Fonti della sezione: IMS, RTTEST}
\end{frame}
\begin{lrbox}{\TitleBox}\begin{minipage}[t]{18.2cm}\color{Ink}\raggedright\sloppy\fontsize{20.7}{24}\selectfont
\bfseries\color{Navy}Intervalli per differenze
\end{minipage}\end{lrbox}
\begin{lrbox}{\BodyBox}\begin{minipage}[t]{10.25cm}\color{Ink}\raggedright\sloppy\fontsize{12.4}{16.4}\selectfont
Quando stimiamo un contrasto, il valore nullo spesso è zero. Un intervallo che include zero mostra compatibilità con assenza di differenza e con altri valori nel suo intervallo. Non dimostra equivalenza; un intervallo ampio può includere anche effetti scientificamente importanti.
\end{minipage}\end{lrbox}
\begin{lrbox}{\SideBox}\begin{minipage}[t]{7.45cm}\color{Ink}\raggedright\sloppy\fontsize{12.2}{16.1}\selectfont
Prima di leggere l'intervallo, definire direzione del contrasto e scala: Spiky–Rounded in \(\mu\)m, non il contrario.
\end{minipage}\end{lrbox}
\begin{lrbox}{\RecallBox}\begin{minipage}[t]{16.45cm}\color{Ink}\raggedright\sloppy\fontsize{10.1}{13.2}\selectfont
«Non conclusivo» è diverso da «nessuna differenza esiste».
\end{minipage}\end{lrbox}
\MeasureBox{63-TitleBox}{\TitleBox}
\MeasureBox{63-BodyBox}{\BodyBox}
\MeasureBox{63-SideBox}{\SideBox}
\MeasureBox{63-RecallBox}{\RecallBox}
\begin{frame}[plain]
\hypertarget{slide63}{}\bookmark[level=1,page=63]{Intervalli per differenze}
\begin{tikzpicture}[remember picture,overlay]
\node[anchor=north east,inner sep=0pt,text=Classical,font=\fontsize{8}{10}\selectfont] at ([xshift=19.05cm,yshift=-.16cm]current page.north west) {\hyperlink{slide3}{Indice}};
\node[anchor=north west,inner sep=0pt,text=Navy,font=\bfseries] at ([xshift=.95cm,yshift=-.62cm]current page.north west) {\usebox{\TitleBox}};
\draw[Classical,line width=1.3pt] ([xshift=.95cm,yshift=-1.68cm]current page.north west)--([xshift=2.25cm,yshift=-1.68cm]current page.north west);
\node[anchor=north west,inner sep=0pt] at ([xshift=.95cm,yshift=-2.12cm]current page.north west) {\usebox{\BodyBox}};
\node[anchor=north west,inner sep=0pt,text=Quantum,font=\bfseries\fontsize{11.5}{14}\selectfont] at ([xshift=11.6cm,yshift=-2.1cm]current page.north west) {Esempio e interpretazione};
\node[anchor=north west,inner sep=0pt] at ([xshift=11.6cm,yshift=-2.75cm]current page.north west) {\usebox{\SideBox}};
\draw[Rule,line width=.5pt] ([xshift=.95cm,yshift=-10.43cm]current page.north west)--([xshift=19.05cm,yshift=-10.43cm]current page.north west);
\node[anchor=north west,inner sep=0pt,text=Classical,font=\bfseries\fontsize{9.4}{12}\selectfont] at ([xshift=.95cm,yshift=-10.66cm]current page.north west) {Richiami};
\node[anchor=north west,inner sep=0pt] at ([xshift=2.6cm,yshift=-10.64cm]current page.north west) {\usebox{\RecallBox}};
\draw[Rule,line width=.35pt] ([xshift=.95cm,yshift=-11.92cm]current page.north west)--([xshift=19.05cm,yshift=-11.92cm]current page.north west);
\node[anchor=north west,inner sep=0pt,text=Muted,text width=9.0cm,font=\fontsize{7.7}{9.3}\selectfont] at ([xshift=.95cm,yshift=-12.06cm]current page.north west) {06 · Intervalli di confidenza};
\node[anchor=north west,inner sep=0pt,text=Classical,text width=8.2cm,align=right,font=\fontsize{7.3}{9.1}\selectfont] at ([xshift=9.45cm,yshift=-12.06cm]current page.north west) {\href{https://openintrostat.github.io/ims/}{IMS} · \href{https://stat.ethz.ch/R-manual/R-devel/library/stats/html/t.test.html}{RTTEST}};
\node[anchor=north east,inner sep=0pt,text=Muted,font=\fontsize{7.7}{9.3}\selectfont] at ([xshift=19.05cm,yshift=-12.06cm]current page.north west) {63/110};

\end{tikzpicture}
\note{Idea della slide: Quando stimiamo un contrasto, il valore nullo spesso è zero. Un intervallo che include zero mostra compatibilità con assenza di differenza e con altri valori nel suo intervallo. Non dimostra equivalenza; un intervallo ampio può includere anche effetti scientificamente importanti.\par\medskip Punto da far emergere: «Non conclusivo» è diverso da «nessuna differenza esiste».\par\medskip Fonti della sezione: IMS, RTTEST}
\end{frame}
\begin{lrbox}{\TitleBox}\begin{minipage}[t]{18.2cm}\color{Ink}\raggedright\sloppy\fontsize{20.7}{24}\selectfont
\bfseries\color{Navy}Attività 06 · applica il concetto
\end{minipage}\end{lrbox}
\begin{lrbox}{\BodyBox}\begin{minipage}[t]{10.25cm}\color{Ink}\raggedright\sloppy\fontsize{12.4}{16.4}\selectfont
Uno studio riporta una media di 6,35 \(\mu\)m con IC 95\% da 6,20 a 6,51. Quale interpretazione è corretta? Riguarda il 95\% dei fossili?
\end{minipage}\end{lrbox}
\begin{lrbox}{\SideBox}\begin{minipage}[t]{7.45cm}\color{Ink}\raggedright\sloppy\fontsize{12.2}{16.1}\selectfont
\noindent{\color{Classical}\bfseries 01}\hspace{.35cm}Individua la domanda\par\vspace{.22cm}\noindent{\color{Classical}\bfseries 02}\hspace{.35cm}Esplicita le ipotesi\par\vspace{.22cm}\noindent{\color{Classical}\bfseries 03}\hspace{.35cm}Calcola o rappresenta\par\vspace{.22cm}\noindent{\color{Classical}\bfseries 04}\hspace{.35cm}Scrivi una conclusione\par\vspace{.22cm}\noindent{\color{Classical}\bfseries 05}\hspace{.35cm}Indica un limite\par\vspace{.22cm}
\end{minipage}\end{lrbox}
\begin{lrbox}{\RecallBox}\begin{minipage}[t]{16.45cm}\color{Ink}\raggedright\sloppy\fontsize{10.1}{13.2}\selectfont
Che cosa significa il 95\% in un intervallo?
\end{minipage}\end{lrbox}
\MeasureBox{64-TitleBox}{\TitleBox}
\MeasureBox{64-BodyBox}{\BodyBox}
\MeasureBox{64-SideBox}{\SideBox}
\MeasureBox{64-RecallBox}{\RecallBox}
\begin{frame}[plain]
\hypertarget{slide64}{}\bookmark[level=1,page=64]{Attività 06 · applica il concetto}
\begin{tikzpicture}[remember picture,overlay]
\node[anchor=north east,inner sep=0pt,text=Classical,font=\fontsize{8}{10}\selectfont] at ([xshift=19.05cm,yshift=-.16cm]current page.north west) {\hyperlink{slide3}{Indice}};
\node[anchor=north west,inner sep=0pt,text=Navy,font=\bfseries] at ([xshift=.95cm,yshift=-.62cm]current page.north west) {\usebox{\TitleBox}};
\draw[Classical,line width=1.3pt] ([xshift=.95cm,yshift=-1.68cm]current page.north west)--([xshift=2.25cm,yshift=-1.68cm]current page.north west);
\node[anchor=north west,inner sep=0pt] at ([xshift=.95cm,yshift=-2.12cm]current page.north west) {\usebox{\BodyBox}};
\node[anchor=north west,inner sep=0pt,text=Quantum,font=\bfseries\fontsize{11.5}{14}\selectfont] at ([xshift=11.6cm,yshift=-2.1cm]current page.north west) {Consegna di lavoro};
\node[anchor=north west,inner sep=0pt] at ([xshift=11.6cm,yshift=-2.75cm]current page.north west) {\usebox{\SideBox}};
\draw[Rule,line width=.5pt] ([xshift=.95cm,yshift=-10.43cm]current page.north west)--([xshift=19.05cm,yshift=-10.43cm]current page.north west);
\node[anchor=north west,inner sep=0pt,text=Classical,font=\bfseries\fontsize{9.4}{12}\selectfont] at ([xshift=.95cm,yshift=-10.66cm]current page.north west) {Richiami};
\node[anchor=north west,inner sep=0pt] at ([xshift=2.6cm,yshift=-10.64cm]current page.north west) {\usebox{\RecallBox}};
\draw[Rule,line width=.35pt] ([xshift=.95cm,yshift=-11.92cm]current page.north west)--([xshift=19.05cm,yshift=-11.92cm]current page.north west);
\node[anchor=north west,inner sep=0pt,text=Muted,text width=9.0cm,font=\fontsize{7.7}{9.3}\selectfont] at ([xshift=.95cm,yshift=-12.06cm]current page.north west) {06 · Intervalli di confidenza};
\node[anchor=north west,inner sep=0pt,text=Classical,text width=8.2cm,align=right,font=\fontsize{7.3}{9.1}\selectfont] at ([xshift=9.45cm,yshift=-12.06cm]current page.north west) {\href{https://openintrostat.github.io/ims/}{IMS} · \href{https://stat.ethz.ch/R-manual/R-devel/library/stats/html/t.test.html}{RTTEST}};
\node[anchor=north east,inner sep=0pt,text=Muted,font=\fontsize{7.7}{9.3}\selectfont] at ([xshift=19.05cm,yshift=-12.06cm]current page.north west) {64/110};

\end{tikzpicture}
\note{Consegna: Uno studio riporta una media di 6,35 \(\mu\)m con IC 95\% da 6,20 a 6,51. Quale interpretazione è corretta? Riguarda il 95\% dei fossili?\par\medskip Soluzione per il docente: L'intervallo riguarda il parametro medio, non la quota di individui. Nella lettura frequentista, una procedura con copertura 95\% contiene il parametro nel 95\% circa delle ripetizioni ideali, sotto le ipotesi. Non assegna automaticamente probabilità 95\% al parametro dentro l'intervallo già osservato.}
\end{frame}
\begin{lrbox}{\TitleBox}\begin{minipage}[t]{18.2cm}\color{Ink}\raggedright\sloppy\fontsize{20.7}{24}\selectfont
\bfseries\color{Navy}Discussione 06 · soluzione ragionata
\end{minipage}\end{lrbox}
\begin{lrbox}{\BodyBox}\begin{minipage}[t]{10.25cm}\color{Ink}\raggedright\sloppy\fontsize{12.4}{16.4}\selectfont
L'intervallo riguarda il parametro medio, non la quota di individui. Nella lettura frequentista, una procedura con copertura 95\% contiene il parametro nel 95\% circa delle ripetizioni ideali, sotto le ipotesi. Non assegna automaticamente probabilità 95\% al parametro dentro l'intervallo già osservato.
\end{minipage}\end{lrbox}
\begin{lrbox}{\SideBox}\begin{minipage}[t]{7.45cm}\color{Ink}\raggedright\sloppy\fontsize{12.2}{16.1}\selectfont
«Non conclusivo» è diverso da «nessuna differenza esiste».
\end{minipage}\end{lrbox}
\begin{lrbox}{\RecallBox}\begin{minipage}[t]{16.45cm}\color{Ink}\raggedright\sloppy\fontsize{10.1}{13.2}\selectfont
Valuta il collegamento tra domanda, procedimento e conclusione, non soltanto il numero finale.
\end{minipage}\end{lrbox}
\MeasureBox{65-TitleBox}{\TitleBox}
\MeasureBox{65-BodyBox}{\BodyBox}
\MeasureBox{65-SideBox}{\SideBox}
\MeasureBox{65-RecallBox}{\RecallBox}
\begin{frame}[plain]
\hypertarget{slide65}{}\bookmark[level=1,page=65]{Discussione 06 · soluzione ragionata}
\begin{tikzpicture}[remember picture,overlay]
\node[anchor=north east,inner sep=0pt,text=Classical,font=\fontsize{8}{10}\selectfont] at ([xshift=19.05cm,yshift=-.16cm]current page.north west) {\hyperlink{slide3}{Indice}};
\node[anchor=north west,inner sep=0pt,text=Navy,font=\bfseries] at ([xshift=.95cm,yshift=-.62cm]current page.north west) {\usebox{\TitleBox}};
\draw[Classical,line width=1.3pt] ([xshift=.95cm,yshift=-1.68cm]current page.north west)--([xshift=2.25cm,yshift=-1.68cm]current page.north west);
\node[anchor=north west,inner sep=0pt] at ([xshift=.95cm,yshift=-2.12cm]current page.north west) {\usebox{\BodyBox}};
\node[anchor=north west,inner sep=0pt,text=Quantum,font=\bfseries\fontsize{11.5}{14}\selectfont] at ([xshift=11.6cm,yshift=-2.1cm]current page.north west) {Criterio di lettura};
\node[anchor=north west,inner sep=0pt] at ([xshift=11.6cm,yshift=-2.75cm]current page.north west) {\usebox{\SideBox}};
\draw[Rule,line width=.5pt] ([xshift=.95cm,yshift=-10.43cm]current page.north west)--([xshift=19.05cm,yshift=-10.43cm]current page.north west);
\node[anchor=north west,inner sep=0pt,text=Classical,font=\bfseries\fontsize{9.4}{12}\selectfont] at ([xshift=.95cm,yshift=-10.66cm]current page.north west) {Richiami};
\node[anchor=north west,inner sep=0pt] at ([xshift=2.6cm,yshift=-10.64cm]current page.north west) {\usebox{\RecallBox}};
\draw[Rule,line width=.35pt] ([xshift=.95cm,yshift=-11.92cm]current page.north west)--([xshift=19.05cm,yshift=-11.92cm]current page.north west);
\node[anchor=north west,inner sep=0pt,text=Muted,text width=9.0cm,font=\fontsize{7.7}{9.3}\selectfont] at ([xshift=.95cm,yshift=-12.06cm]current page.north west) {06 · Intervalli di confidenza};
\node[anchor=north west,inner sep=0pt,text=Classical,text width=8.2cm,align=right,font=\fontsize{7.3}{9.1}\selectfont] at ([xshift=9.45cm,yshift=-12.06cm]current page.north west) {\href{https://openintrostat.github.io/ims/}{IMS} · \href{https://stat.ethz.ch/R-manual/R-devel/library/stats/html/t.test.html}{RTTEST}};
\node[anchor=north east,inner sep=0pt,text=Muted,font=\fontsize{7.7}{9.3}\selectfont] at ([xshift=19.05cm,yshift=-12.06cm]current page.north west) {65/110};

\end{tikzpicture}
\note{L'intervallo riguarda il parametro medio, non la quota di individui. Nella lettura frequentista, una procedura con copertura 95\% contiene il parametro nel 95\% circa delle ripetizioni ideali, sotto le ipotesi. Non assegna automaticamente probabilità 95\% al parametro dentro l'intervallo già osservato.}
\end{frame}
\begin{frame}[plain]\hypertarget{slide66}{}
\bookmark[level=0,page=66]{Bootstrap: ricampionare il campione}
\begin{tikzpicture}[remember picture,overlay]
\fill[Night] (current page.south west) rectangle (current page.north east);
\draw[Mint,line width=2pt] ([xshift=1.05cm,yshift=-1.35cm]current page.north west)--([xshift=3.0cm,yshift=-1.35cm]current page.north west);
\node[anchor=north west,inner sep=0,text=Mint,font=\bfseries\fontsize{11}{14}\selectfont] at ([xshift=1.05cm,yshift=-1.7cm]current page.north west) {SEZIONE 07};
\node[anchor=north west,inner sep=0,text=white,text width=15.8cm,align=left,font=\bfseries\fontsize{31}{38}\selectfont] at ([xshift=1.05cm,yshift=-2.65cm]current page.north west) {Bootstrap: ricampionare il campione};
\node[anchor=north west,inner sep=0,text=Lilac,text width=15.8cm,align=left,font=\fontsize{15}{21}\selectfont] at ([xshift=1.05cm,yshift=-6.4cm]current page.north west) {Come stimiamo una distribuzione campionaria senza conoscere la popolazione?};
\node[anchor=north west,inner sep=0,text=white,text width=16.5cm,align=left,font=\fontsize{11}{15}\selectfont] at ([xshift=1.05cm,yshift=-8.6cm]current page.north west) {Comprendere il concetto · leggere un esempio · applicare e discutere};
\node[anchor=north east,inner sep=0,text=Mint,font=\fontsize{8}{10}\selectfont] at ([xshift=19.05cm,yshift=-.2cm]current page.north west) {\hyperlink{slide3}{Indice}};
\node[anchor=south east,inner sep=0,text=Lilac!45,font=\bfseries\fontsize{64}{68}\selectfont] at ([xshift=18.95cm,yshift=-11.35cm]current page.north west) {07};
\draw[Lilac!30,line width=.35pt] ([xshift=.95cm,yshift=-11.92cm]current page.north west)--([xshift=19.05cm,yshift=-11.92cm]current page.north west);
\node[anchor=north west,inner sep=0pt,text=Lilac,text width=9.0cm,font=\fontsize{7.7}{9.3}\selectfont] at ([xshift=.95cm,yshift=-12.06cm]current page.north west) {07 · Bootstrap: ricampionare il campione};
\node[anchor=north west,inner sep=0pt,text=Mint,text width=8.2cm,align=right,font=\fontsize{7.3}{9.1}\selectfont] at ([xshift=9.45cm,yshift=-12.06cm]current page.north west) {\href{https://openintrostat.github.io/ims/}{IMS} · \href{https://cran.r-project.org/doc/manuals/r-release/R-intro.html}{RBASE}};
\node[anchor=north east,inner sep=0pt,text=Lilac,font=\fontsize{7.7}{9.3}\selectfont] at ([xshift=19.05cm,yshift=-12.06cm]current page.north west) {66/110};
\end{tikzpicture}\end{frame}
\begin{lrbox}{\TitleBox}\begin{minipage}[t]{18.2cm}\color{Ink}\raggedright\sloppy\fontsize{20.7}{24}\selectfont
\bfseries\color{Navy}L'idea del bootstrap non parametrico
\end{minipage}\end{lrbox}
\begin{lrbox}{\BodyBox}\begin{minipage}[t]{10.25cm}\color{Ink}\raggedright\sloppy\fontsize{12.4}{16.4}\selectfont
Il bootstrap usa la distribuzione empirica del campione come approssimazione della popolazione e ricampiona da essa. Ripetendo il calcolo della statistica, otteniamo una distribuzione di repliche. La validità dipende dalla rappresentatività del campione e dalla struttura di dipendenza rispettata.
\end{minipage}\end{lrbox}
\begin{lrbox}{\SideBox}\begin{minipage}[t]{7.45cm}\color{Ink}\raggedright\sloppy\fontsize{12.2}{16.1}\selectfont
\noindent{\color{Classical}\bfseries 01}\hspace{.35cm}Campione osservato\par\vspace{.22cm}\noindent{\color{Classical}\bfseries 02}\hspace{.35cm}Ricampionamento con reinserimento\par\vspace{.22cm}\noindent{\color{Classical}\bfseries 03}\hspace{.35cm}Statistica della replica\par\vspace{.22cm}\noindent{\color{Classical}\bfseries 04}\hspace{.35cm}Distribuzione delle repliche\par\vspace{.22cm}
\end{minipage}\end{lrbox}
\begin{lrbox}{\RecallBox}\begin{minipage}[t]{16.45cm}\color{Ink}\raggedright\sloppy\fontsize{10.1}{13.2}\selectfont
Il bootstrap non crea osservazioni nuove indipendenti dal campione iniziale.
\end{minipage}\end{lrbox}
\MeasureBox{67-TitleBox}{\TitleBox}
\MeasureBox{67-BodyBox}{\BodyBox}
\MeasureBox{67-SideBox}{\SideBox}
\MeasureBox{67-RecallBox}{\RecallBox}
\begin{frame}[plain]
\hypertarget{slide67}{}\bookmark[level=1,page=67]{L'idea del bootstrap non parametrico}
\begin{tikzpicture}[remember picture,overlay]
\node[anchor=north east,inner sep=0pt,text=Classical,font=\fontsize{8}{10}\selectfont] at ([xshift=19.05cm,yshift=-.16cm]current page.north west) {\hyperlink{slide3}{Indice}};
\node[anchor=north west,inner sep=0pt,text=Navy,font=\bfseries] at ([xshift=.95cm,yshift=-.62cm]current page.north west) {\usebox{\TitleBox}};
\draw[Classical,line width=1.3pt] ([xshift=.95cm,yshift=-1.68cm]current page.north west)--([xshift=2.25cm,yshift=-1.68cm]current page.north west);
\node[anchor=north west,inner sep=0pt] at ([xshift=.95cm,yshift=-2.12cm]current page.north west) {\usebox{\BodyBox}};
\node[anchor=north west,inner sep=0pt,text=Quantum,font=\bfseries\fontsize{11.5}{14}\selectfont] at ([xshift=11.6cm,yshift=-2.1cm]current page.north west) {Passaggi da seguire};
\node[anchor=north west,inner sep=0pt] at ([xshift=11.6cm,yshift=-2.75cm]current page.north west) {\usebox{\SideBox}};
\draw[Rule,line width=.5pt] ([xshift=.95cm,yshift=-10.43cm]current page.north west)--([xshift=19.05cm,yshift=-10.43cm]current page.north west);
\node[anchor=north west,inner sep=0pt,text=Classical,font=\bfseries\fontsize{9.4}{12}\selectfont] at ([xshift=.95cm,yshift=-10.66cm]current page.north west) {Richiami};
\node[anchor=north west,inner sep=0pt] at ([xshift=2.6cm,yshift=-10.64cm]current page.north west) {\usebox{\RecallBox}};
\draw[Rule,line width=.35pt] ([xshift=.95cm,yshift=-11.92cm]current page.north west)--([xshift=19.05cm,yshift=-11.92cm]current page.north west);
\node[anchor=north west,inner sep=0pt,text=Muted,text width=9.0cm,font=\fontsize{7.7}{9.3}\selectfont] at ([xshift=.95cm,yshift=-12.06cm]current page.north west) {07 · Bootstrap: ricampionare il campione};
\node[anchor=north west,inner sep=0pt,text=Classical,text width=8.2cm,align=right,font=\fontsize{7.3}{9.1}\selectfont] at ([xshift=9.45cm,yshift=-12.06cm]current page.north west) {\href{https://openintrostat.github.io/ims/}{IMS} · \href{https://cran.r-project.org/doc/manuals/r-release/R-intro.html}{RBASE}};
\node[anchor=north east,inner sep=0pt,text=Muted,font=\fontsize{7.7}{9.3}\selectfont] at ([xshift=19.05cm,yshift=-12.06cm]current page.north west) {67/110};

\end{tikzpicture}
\note{Idea della slide: Il bootstrap usa la distribuzione empirica del campione come approssimazione della popolazione e ricampiona da essa. Ripetendo il calcolo della statistica, otteniamo una distribuzione di repliche. La validità dipende dalla rappresentatività del campione e dalla struttura di dipendenza rispettata.\par\medskip Punto da far emergere: Il bootstrap non crea osservazioni nuove indipendenti dal campione iniziale.\par\medskip Fonti della sezione: IMS, RBASE}
\end{frame}
\begin{lrbox}{\TitleBox}\begin{minipage}[t]{18.2cm}\color{Ink}\raggedright\sloppy\fontsize{20.7}{24}\selectfont
\bfseries\color{Navy}Perché con reinserimento
\end{minipage}\end{lrbox}
\begin{lrbox}{\BodyBox}\begin{minipage}[t]{8.2cm}\color{Ink}\raggedright\sloppy\fontsize{12.4}{16.4}\selectfont
Il reinserimento permette che una replica contenga ripetizioni e ometta alcune unità, pur mantenendo la stessa dimensione. Questo produce variazione della statistica. Permutare soltanto l'ordine dei valori non cambierebbe una media o una mediana.
\end{minipage}\end{lrbox}
\begin{lrbox}{\SideBox}\begin{minipage}[t]{9.35cm}\color{Ink}\raggedright\sloppy\fontsize{12.2}{16.1}\selectfont
\begin{lstlisting}
x <- c(4, 5, 7, 8)
sample(x, length(x), replace = TRUE)
\end{lstlisting}
\end{minipage}\end{lrbox}
\begin{lrbox}{\RecallBox}\begin{minipage}[t]{16.45cm}\color{Ink}\raggedright\sloppy\fontsize{10.1}{13.2}\selectfont
Bootstrap e permutazione usano ripetizioni diverse per rispondere a domande diverse.
\end{minipage}\end{lrbox}
\MeasureBox{68-TitleBox}{\TitleBox}
\MeasureBox{68-BodyBox}{\BodyBox}
\MeasureBox{68-SideBox}{\SideBox}
\MeasureBox{68-RecallBox}{\RecallBox}
\begin{frame}[plain]
\hypertarget{slide68}{}\bookmark[level=1,page=68]{Perché con reinserimento}
\begin{tikzpicture}[remember picture,overlay]
\node[anchor=north east,inner sep=0pt,text=Classical,font=\fontsize{8}{10}\selectfont] at ([xshift=19.05cm,yshift=-.16cm]current page.north west) {\hyperlink{slide3}{Indice}};
\node[anchor=north west,inner sep=0pt,text=Navy,font=\bfseries] at ([xshift=.95cm,yshift=-.62cm]current page.north west) {\usebox{\TitleBox}};
\draw[Classical,line width=1.3pt] ([xshift=.95cm,yshift=-1.68cm]current page.north west)--([xshift=2.25cm,yshift=-1.68cm]current page.north west);
\node[anchor=north west,inner sep=0pt] at ([xshift=.95cm,yshift=-2.12cm]current page.north west) {\usebox{\BodyBox}};
\node[anchor=north west,inner sep=0pt,text=Quantum,font=\bfseries\fontsize{11.5}{14}\selectfont] at ([xshift=9.7cm,yshift=-2.1cm]current page.north west) {Protocollo R};
\node[anchor=north west,inner sep=0pt] at ([xshift=9.7cm,yshift=-2.75cm]current page.north west) {\usebox{\SideBox}};
\draw[Rule,line width=.5pt] ([xshift=.95cm,yshift=-10.43cm]current page.north west)--([xshift=19.05cm,yshift=-10.43cm]current page.north west);
\node[anchor=north west,inner sep=0pt,text=Classical,font=\bfseries\fontsize{9.4}{12}\selectfont] at ([xshift=.95cm,yshift=-10.66cm]current page.north west) {Richiami};
\node[anchor=north west,inner sep=0pt] at ([xshift=2.6cm,yshift=-10.64cm]current page.north west) {\usebox{\RecallBox}};
\draw[Rule,line width=.35pt] ([xshift=.95cm,yshift=-11.92cm]current page.north west)--([xshift=19.05cm,yshift=-11.92cm]current page.north west);
\node[anchor=north west,inner sep=0pt,text=Muted,text width=9.0cm,font=\fontsize{7.7}{9.3}\selectfont] at ([xshift=.95cm,yshift=-12.06cm]current page.north west) {07 · Bootstrap: ricampionare il campione};
\node[anchor=north west,inner sep=0pt,text=Classical,text width=8.2cm,align=right,font=\fontsize{7.3}{9.1}\selectfont] at ([xshift=9.45cm,yshift=-12.06cm]current page.north west) {\href{https://openintrostat.github.io/ims/}{IMS} · \href{https://cran.r-project.org/doc/manuals/r-release/R-intro.html}{RBASE}};
\node[anchor=north east,inner sep=0pt,text=Muted,font=\fontsize{7.7}{9.3}\selectfont] at ([xshift=19.05cm,yshift=-12.06cm]current page.north west) {68/110};

\end{tikzpicture}
\note{Idea della slide: Il reinserimento permette che una replica contenga ripetizioni e ometta alcune unità, pur mantenendo la stessa dimensione. Questo produce variazione della statistica. Permutare soltanto l'ordine dei valori non cambierebbe una media o una mediana.\par\medskip Punto da far emergere: Bootstrap e permutazione usano ripetizioni diverse per rispondere a domande diverse.\par\medskip Fonti della sezione: IMS, RBASE}
\end{frame}
\begin{lrbox}{\TitleBox}\begin{minipage}[t]{18.2cm}\color{Ink}\raggedright\sloppy\fontsize{20.7}{24}\selectfont
\bfseries\color{Navy}Una replica e molte repliche
\end{minipage}\end{lrbox}
\begin{lrbox}{\BodyBox}\begin{minipage}[t]{8.2cm}\color{Ink}\raggedright\sloppy\fontsize{12.4}{16.4}\selectfont
Ogni replica genera un campione e restituisce una statistica. Il vettore finale contiene statistiche, non tutte le osservazioni estratte. Il numero di repliche controlla la precisione numerica della distribuzione bootstrap, senza modificare la numerosità del campione reale.
\end{minipage}\end{lrbox}
\begin{lrbox}{\SideBox}\begin{minipage}[t]{9.35cm}\color{Ink}\raggedright\sloppy\fontsize{12.2}{16.1}\selectfont
\begin{lstlisting}
b <- replicate(5000,
  mean(sample(x, length(x),
              replace = TRUE)))
\end{lstlisting}
\end{minipage}\end{lrbox}
\begin{lrbox}{\RecallBox}\begin{minipage}[t]{16.45cm}\color{Ink}\raggedright\sloppy\fontsize{10.1}{13.2}\selectfont
Cinquemila repliche non trasformano quattro dati osservati in ventimila nuove misure.
\end{minipage}\end{lrbox}
\MeasureBox{69-TitleBox}{\TitleBox}
\MeasureBox{69-BodyBox}{\BodyBox}
\MeasureBox{69-SideBox}{\SideBox}
\MeasureBox{69-RecallBox}{\RecallBox}
\begin{frame}[plain]
\hypertarget{slide69}{}\bookmark[level=1,page=69]{Una replica e molte repliche}
\begin{tikzpicture}[remember picture,overlay]
\node[anchor=north east,inner sep=0pt,text=Classical,font=\fontsize{8}{10}\selectfont] at ([xshift=19.05cm,yshift=-.16cm]current page.north west) {\hyperlink{slide3}{Indice}};
\node[anchor=north west,inner sep=0pt,text=Navy,font=\bfseries] at ([xshift=.95cm,yshift=-.62cm]current page.north west) {\usebox{\TitleBox}};
\draw[Classical,line width=1.3pt] ([xshift=.95cm,yshift=-1.68cm]current page.north west)--([xshift=2.25cm,yshift=-1.68cm]current page.north west);
\node[anchor=north west,inner sep=0pt] at ([xshift=.95cm,yshift=-2.12cm]current page.north west) {\usebox{\BodyBox}};
\node[anchor=north west,inner sep=0pt,text=Quantum,font=\bfseries\fontsize{11.5}{14}\selectfont] at ([xshift=9.7cm,yshift=-2.1cm]current page.north west) {Protocollo R};
\node[anchor=north west,inner sep=0pt] at ([xshift=9.7cm,yshift=-2.75cm]current page.north west) {\usebox{\SideBox}};
\draw[Rule,line width=.5pt] ([xshift=.95cm,yshift=-10.43cm]current page.north west)--([xshift=19.05cm,yshift=-10.43cm]current page.north west);
\node[anchor=north west,inner sep=0pt,text=Classical,font=\bfseries\fontsize{9.4}{12}\selectfont] at ([xshift=.95cm,yshift=-10.66cm]current page.north west) {Richiami};
\node[anchor=north west,inner sep=0pt] at ([xshift=2.6cm,yshift=-10.64cm]current page.north west) {\usebox{\RecallBox}};
\draw[Rule,line width=.35pt] ([xshift=.95cm,yshift=-11.92cm]current page.north west)--([xshift=19.05cm,yshift=-11.92cm]current page.north west);
\node[anchor=north west,inner sep=0pt,text=Muted,text width=9.0cm,font=\fontsize{7.7}{9.3}\selectfont] at ([xshift=.95cm,yshift=-12.06cm]current page.north west) {07 · Bootstrap: ricampionare il campione};
\node[anchor=north west,inner sep=0pt,text=Classical,text width=8.2cm,align=right,font=\fontsize{7.3}{9.1}\selectfont] at ([xshift=9.45cm,yshift=-12.06cm]current page.north west) {\href{https://openintrostat.github.io/ims/}{IMS} · \href{https://cran.r-project.org/doc/manuals/r-release/R-intro.html}{RBASE}};
\node[anchor=north east,inner sep=0pt,text=Muted,font=\fontsize{7.7}{9.3}\selectfont] at ([xshift=19.05cm,yshift=-12.06cm]current page.north west) {69/110};

\end{tikzpicture}
\note{Idea della slide: Ogni replica genera un campione e restituisce una statistica. Il vettore finale contiene statistiche, non tutte le osservazioni estratte. Il numero di repliche controlla la precisione numerica della distribuzione bootstrap, senza modificare la numerosità del campione reale.\par\medskip Punto da far emergere: Cinquemila repliche non trasformano quattro dati osservati in ventimila nuove misure.\par\medskip Fonti della sezione: IMS, RBASE}
\end{frame}
\begin{lrbox}{\TitleBox}\begin{minipage}[t]{18.2cm}\color{Ink}\raggedright\sloppy\fontsize{20.7}{24}\selectfont
\bfseries\color{Navy}Bootstrap della media dei fossili
\end{minipage}\end{lrbox}
\begin{lrbox}{\BodyBox}\begin{minipage}[t]{7.1cm}\color{Ink}\raggedright\sloppy\fontsize{11.8}{15.8}\selectfont
Qui ricampioniamo dalla miscela empirica dei 160 diametri: le quote dei morfotipi possono variare tra repliche. Il generatore ne contiene invece 80 per gruppo. Per imitare quel disegno a quote fisse, il laboratorio aggiunge un bootstrap stratificato e confronta le due incertezze.
\end{minipage}\end{lrbox}
\begin{lrbox}{\SideBox}\begin{minipage}[t]{10.4cm}\color{Ink}\raggedright\sloppy\fontsize{12.2}{16.1}\selectfont
\centering\includegraphics[width=10.4cm,height=6.25cm,keepaspectratio]{figure/bootstrap.pdf}\par\vspace{.22cm}{\raggedright\color{Muted}\fontsize{9.7}{12.5}\selectfont 5.000 repliche pooled: miscela empirica, senza vincolo 80/80.\par}
\end{minipage}\end{lrbox}
\begin{lrbox}{\RecallBox}\begin{minipage}[t]{16.45cm}\color{Ink}\raggedright\sloppy\fontsize{10.1}{13.2}\selectfont
I numeri del bootstrap dipendono dal campione e dalla procedura, non soltanto dal seme.
\end{minipage}\end{lrbox}
\MeasureBox{70-TitleBox}{\TitleBox}
\MeasureBox{70-BodyBox}{\BodyBox}
\MeasureBox{70-SideBox}{\SideBox}
\MeasureBox{70-RecallBox}{\RecallBox}
\begin{frame}[plain]
\hypertarget{slide70}{}\bookmark[level=1,page=70]{Bootstrap della media dei fossili}
\begin{tikzpicture}[remember picture,overlay]
\node[anchor=north east,inner sep=0pt,text=Classical,font=\fontsize{8}{10}\selectfont] at ([xshift=19.05cm,yshift=-.16cm]current page.north west) {\hyperlink{slide3}{Indice}};
\node[anchor=north west,inner sep=0pt,text=Navy,font=\bfseries] at ([xshift=.95cm,yshift=-.62cm]current page.north west) {\usebox{\TitleBox}};
\draw[Classical,line width=1.3pt] ([xshift=.95cm,yshift=-1.68cm]current page.north west)--([xshift=2.25cm,yshift=-1.68cm]current page.north west);
\node[anchor=north west,inner sep=0pt] at ([xshift=.95cm,yshift=-2.12cm]current page.north west) {\usebox{\BodyBox}};
\node[anchor=north west,inner sep=0pt,text=Quantum,font=\bfseries\fontsize{11.5}{14}\selectfont] at ([xshift=8.65cm,yshift=-2.1cm]current page.north west) {Figura commentata};
\node[anchor=north west,inner sep=0pt] at ([xshift=8.65cm,yshift=-2.75cm]current page.north west) {\usebox{\SideBox}};
\draw[Rule,line width=.5pt] ([xshift=.95cm,yshift=-10.43cm]current page.north west)--([xshift=19.05cm,yshift=-10.43cm]current page.north west);
\node[anchor=north west,inner sep=0pt,text=Classical,font=\bfseries\fontsize{9.4}{12}\selectfont] at ([xshift=.95cm,yshift=-10.66cm]current page.north west) {Richiami};
\node[anchor=north west,inner sep=0pt] at ([xshift=2.6cm,yshift=-10.64cm]current page.north west) {\usebox{\RecallBox}};
\draw[Rule,line width=.35pt] ([xshift=.95cm,yshift=-11.92cm]current page.north west)--([xshift=19.05cm,yshift=-11.92cm]current page.north west);
\node[anchor=north west,inner sep=0pt,text=Muted,text width=9.0cm,font=\fontsize{7.7}{9.3}\selectfont] at ([xshift=.95cm,yshift=-12.06cm]current page.north west) {07 · Bootstrap: ricampionare il campione};
\node[anchor=north west,inner sep=0pt,text=Classical,text width=8.2cm,align=right,font=\fontsize{7.3}{9.1}\selectfont] at ([xshift=9.45cm,yshift=-12.06cm]current page.north west) {\href{https://openintrostat.github.io/ims/}{IMS} · \href{https://cran.r-project.org/doc/manuals/r-release/R-intro.html}{RBASE}};
\node[anchor=north east,inner sep=0pt,text=Muted,font=\fontsize{7.7}{9.3}\selectfont] at ([xshift=19.05cm,yshift=-12.06cm]current page.north west) {70/110};

\end{tikzpicture}
\note{Idea della slide: Qui ricampioniamo dalla miscela empirica dei 160 diametri: le quote dei morfotipi possono variare tra repliche. Il generatore ne contiene invece 80 per gruppo. Per imitare quel disegno a quote fisse, il laboratorio aggiunge un bootstrap stratificato e confronta le due incertezze.\par\medskip Punto da far emergere: I numeri del bootstrap dipendono dal campione e dalla procedura, non soltanto dal seme.\par\medskip Fonti della sezione: IMS, RBASE}
\end{frame}
\begin{lrbox}{\TitleBox}\begin{minipage}[t]{18.2cm}\color{Ink}\raggedright\sloppy\fontsize{20.7}{24}\selectfont
\bfseries\color{Navy}Intervallo percentile
\end{minipage}\end{lrbox}
\begin{lrbox}{\BodyBox}\begin{minipage}[t]{8.2cm}\color{Ink}\raggedright\sloppy\fontsize{12.4}{16.4}\selectfont
Il metodo percentile usa quantili della distribuzione delle repliche. Per un intervallo centrale al 95\%, prendiamo 2,5\% e 97,5\%. È semplice da comprendere, ma non ha prestazioni uniformi per tutte le statistiche, forme o piccoli campioni.
\end{minipage}\end{lrbox}
\begin{lrbox}{\SideBox}\begin{minipage}[t]{9.35cm}\color{Ink}\raggedright\sloppy\fontsize{12.2}{16.1}\selectfont
\begin{lstlisting}
quantile(b, c(.025, .975))

Nel calcolo di riferimento: circa 6,199–6,515 µm per la media del dataset completo.
\end{lstlisting}
\end{minipage}\end{lrbox}
\begin{lrbox}{\RecallBox}\begin{minipage}[t]{16.45cm}\color{Ink}\raggedright\sloppy\fontsize{10.1}{13.2}\selectfont
Un metodo semplice non deve essere presentato come esatto in ogni situazione.
\end{minipage}\end{lrbox}
\MeasureBox{71-TitleBox}{\TitleBox}
\MeasureBox{71-BodyBox}{\BodyBox}
\MeasureBox{71-SideBox}{\SideBox}
\MeasureBox{71-RecallBox}{\RecallBox}
\begin{frame}[plain]
\hypertarget{slide71}{}\bookmark[level=1,page=71]{Intervallo percentile}
\begin{tikzpicture}[remember picture,overlay]
\node[anchor=north east,inner sep=0pt,text=Classical,font=\fontsize{8}{10}\selectfont] at ([xshift=19.05cm,yshift=-.16cm]current page.north west) {\hyperlink{slide3}{Indice}};
\node[anchor=north west,inner sep=0pt,text=Navy,font=\bfseries] at ([xshift=.95cm,yshift=-.62cm]current page.north west) {\usebox{\TitleBox}};
\draw[Classical,line width=1.3pt] ([xshift=.95cm,yshift=-1.68cm]current page.north west)--([xshift=2.25cm,yshift=-1.68cm]current page.north west);
\node[anchor=north west,inner sep=0pt] at ([xshift=.95cm,yshift=-2.12cm]current page.north west) {\usebox{\BodyBox}};
\node[anchor=north west,inner sep=0pt,text=Quantum,font=\bfseries\fontsize{11.5}{14}\selectfont] at ([xshift=9.7cm,yshift=-2.1cm]current page.north west) {Protocollo R};
\node[anchor=north west,inner sep=0pt] at ([xshift=9.7cm,yshift=-2.75cm]current page.north west) {\usebox{\SideBox}};
\draw[Rule,line width=.5pt] ([xshift=.95cm,yshift=-10.43cm]current page.north west)--([xshift=19.05cm,yshift=-10.43cm]current page.north west);
\node[anchor=north west,inner sep=0pt,text=Classical,font=\bfseries\fontsize{9.4}{12}\selectfont] at ([xshift=.95cm,yshift=-10.66cm]current page.north west) {Richiami};
\node[anchor=north west,inner sep=0pt] at ([xshift=2.6cm,yshift=-10.64cm]current page.north west) {\usebox{\RecallBox}};
\draw[Rule,line width=.35pt] ([xshift=.95cm,yshift=-11.92cm]current page.north west)--([xshift=19.05cm,yshift=-11.92cm]current page.north west);
\node[anchor=north west,inner sep=0pt,text=Muted,text width=9.0cm,font=\fontsize{7.7}{9.3}\selectfont] at ([xshift=.95cm,yshift=-12.06cm]current page.north west) {07 · Bootstrap: ricampionare il campione};
\node[anchor=north west,inner sep=0pt,text=Classical,text width=8.2cm,align=right,font=\fontsize{7.3}{9.1}\selectfont] at ([xshift=9.45cm,yshift=-12.06cm]current page.north west) {\href{https://openintrostat.github.io/ims/}{IMS} · \href{https://cran.r-project.org/doc/manuals/r-release/R-intro.html}{RBASE}};
\node[anchor=north east,inner sep=0pt,text=Muted,font=\fontsize{7.7}{9.3}\selectfont] at ([xshift=19.05cm,yshift=-12.06cm]current page.north west) {71/110};

\end{tikzpicture}
\note{Idea della slide: Il metodo percentile usa quantili della distribuzione delle repliche. Per un intervallo centrale al 95\%, prendiamo 2,5\% e 97,5\%. È semplice da comprendere, ma non ha prestazioni uniformi per tutte le statistiche, forme o piccoli campioni.\par\medskip Punto da far emergere: Un metodo semplice non deve essere presentato come esatto in ogni situazione.\par\medskip Fonti della sezione: IMS, RBASE}
\end{frame}
\begin{lrbox}{\TitleBox}\begin{minipage}[t]{18.2cm}\color{Ink}\raggedright\sloppy\fontsize{20.7}{24}\selectfont
\bfseries\color{Navy}Bootstrap di una mediana
\end{minipage}\end{lrbox}
\begin{lrbox}{\BodyBox}\begin{minipage}[t]{8.2cm}\color{Ink}\raggedright\sloppy\fontsize{12.4}{16.4}\selectfont
Il ricampionamento può studiare statistiche per cui una formula elementare dell'errore standard è meno immediata. Per una mediana, campioni piccoli o molti valori ripetuti possono produrre una distribuzione bootstrap discreta e irregolare. Questo comportamento va mostrato, non nascosto.
\end{minipage}\end{lrbox}
\begin{lrbox}{\SideBox}\begin{minipage}[t]{9.35cm}\color{Ink}\raggedright\sloppy\fontsize{12.2}{16.1}\selectfont
\begin{lstlisting}
b_med <- replicate(5000,
  median(sample(x, length(x),
                replace = TRUE)))
\end{lstlisting}
\end{minipage}\end{lrbox}
\begin{lrbox}{\RecallBox}\begin{minipage}[t]{16.45cm}\color{Ink}\raggedright\sloppy\fontsize{10.1}{13.2}\selectfont
La flessibilità del bootstrap non elimina le difficoltà dei dati poco informativi.
\end{minipage}\end{lrbox}
\MeasureBox{72-TitleBox}{\TitleBox}
\MeasureBox{72-BodyBox}{\BodyBox}
\MeasureBox{72-SideBox}{\SideBox}
\MeasureBox{72-RecallBox}{\RecallBox}
\begin{frame}[plain]
\hypertarget{slide72}{}\bookmark[level=1,page=72]{Bootstrap di una mediana}
\begin{tikzpicture}[remember picture,overlay]
\node[anchor=north east,inner sep=0pt,text=Classical,font=\fontsize{8}{10}\selectfont] at ([xshift=19.05cm,yshift=-.16cm]current page.north west) {\hyperlink{slide3}{Indice}};
\node[anchor=north west,inner sep=0pt,text=Navy,font=\bfseries] at ([xshift=.95cm,yshift=-.62cm]current page.north west) {\usebox{\TitleBox}};
\draw[Classical,line width=1.3pt] ([xshift=.95cm,yshift=-1.68cm]current page.north west)--([xshift=2.25cm,yshift=-1.68cm]current page.north west);
\node[anchor=north west,inner sep=0pt] at ([xshift=.95cm,yshift=-2.12cm]current page.north west) {\usebox{\BodyBox}};
\node[anchor=north west,inner sep=0pt,text=Quantum,font=\bfseries\fontsize{11.5}{14}\selectfont] at ([xshift=9.7cm,yshift=-2.1cm]current page.north west) {Protocollo R};
\node[anchor=north west,inner sep=0pt] at ([xshift=9.7cm,yshift=-2.75cm]current page.north west) {\usebox{\SideBox}};
\draw[Rule,line width=.5pt] ([xshift=.95cm,yshift=-10.43cm]current page.north west)--([xshift=19.05cm,yshift=-10.43cm]current page.north west);
\node[anchor=north west,inner sep=0pt,text=Classical,font=\bfseries\fontsize{9.4}{12}\selectfont] at ([xshift=.95cm,yshift=-10.66cm]current page.north west) {Richiami};
\node[anchor=north west,inner sep=0pt] at ([xshift=2.6cm,yshift=-10.64cm]current page.north west) {\usebox{\RecallBox}};
\draw[Rule,line width=.35pt] ([xshift=.95cm,yshift=-11.92cm]current page.north west)--([xshift=19.05cm,yshift=-11.92cm]current page.north west);
\node[anchor=north west,inner sep=0pt,text=Muted,text width=9.0cm,font=\fontsize{7.7}{9.3}\selectfont] at ([xshift=.95cm,yshift=-12.06cm]current page.north west) {07 · Bootstrap: ricampionare il campione};
\node[anchor=north west,inner sep=0pt,text=Classical,text width=8.2cm,align=right,font=\fontsize{7.3}{9.1}\selectfont] at ([xshift=9.45cm,yshift=-12.06cm]current page.north west) {\href{https://openintrostat.github.io/ims/}{IMS} · \href{https://cran.r-project.org/doc/manuals/r-release/R-intro.html}{RBASE}};
\node[anchor=north east,inner sep=0pt,text=Muted,font=\fontsize{7.7}{9.3}\selectfont] at ([xshift=19.05cm,yshift=-12.06cm]current page.north west) {72/110};

\end{tikzpicture}
\note{Idea della slide: Il ricampionamento può studiare statistiche per cui una formula elementare dell'errore standard è meno immediata. Per una mediana, campioni piccoli o molti valori ripetuti possono produrre una distribuzione bootstrap discreta e irregolare. Questo comportamento va mostrato, non nascosto.\par\medskip Punto da far emergere: La flessibilità del bootstrap non elimina le difficoltà dei dati poco informativi.\par\medskip Fonti della sezione: IMS, RBASE}
\end{frame}
\begin{lrbox}{\TitleBox}\begin{minipage}[t]{18.2cm}\color{Ink}\raggedright\sloppy\fontsize{20.7}{24}\selectfont
\bfseries\color{Navy}Quali limiti non risolve
\end{minipage}\end{lrbox}
\begin{lrbox}{\BodyBox}\begin{minipage}[t]{10.25cm}\color{Ink}\raggedright\sloppy\fontsize{12.4}{16.4}\selectfont
Il bootstrap non inventa valori di popolazione mai rappresentati nel campione e non corregge selezione, censura o misura sbagliata. Se la statistica dipende fortemente dagli estremi, un piccolo campione può dare un quadro incompleto della sua incertezza.
\end{minipage}\end{lrbox}
\begin{lrbox}{\SideBox}\begin{minipage}[t]{7.45cm}\color{Ink}\raggedright\sloppy\fontsize{12.2}{16.1}\selectfont
Un campione che non contiene individui molto grandi non permette al bootstrap empirico di generarne di nuovi oltre il massimo osservato.
\end{minipage}\end{lrbox}
\begin{lrbox}{\RecallBox}\begin{minipage}[t]{16.45cm}\color{Ink}\raggedright\sloppy\fontsize{10.1}{13.2}\selectfont
Ricampionare una base distorta riproduce quella distorsione.
\end{minipage}\end{lrbox}
\MeasureBox{73-TitleBox}{\TitleBox}
\MeasureBox{73-BodyBox}{\BodyBox}
\MeasureBox{73-SideBox}{\SideBox}
\MeasureBox{73-RecallBox}{\RecallBox}
\begin{frame}[plain]
\hypertarget{slide73}{}\bookmark[level=1,page=73]{Quali limiti non risolve}
\begin{tikzpicture}[remember picture,overlay]
\node[anchor=north east,inner sep=0pt,text=Classical,font=\fontsize{8}{10}\selectfont] at ([xshift=19.05cm,yshift=-.16cm]current page.north west) {\hyperlink{slide3}{Indice}};
\node[anchor=north west,inner sep=0pt,text=Navy,font=\bfseries] at ([xshift=.95cm,yshift=-.62cm]current page.north west) {\usebox{\TitleBox}};
\draw[Classical,line width=1.3pt] ([xshift=.95cm,yshift=-1.68cm]current page.north west)--([xshift=2.25cm,yshift=-1.68cm]current page.north west);
\node[anchor=north west,inner sep=0pt] at ([xshift=.95cm,yshift=-2.12cm]current page.north west) {\usebox{\BodyBox}};
\node[anchor=north west,inner sep=0pt,text=Quantum,font=\bfseries\fontsize{11.5}{14}\selectfont] at ([xshift=11.6cm,yshift=-2.1cm]current page.north west) {Esempio e interpretazione};
\node[anchor=north west,inner sep=0pt] at ([xshift=11.6cm,yshift=-2.75cm]current page.north west) {\usebox{\SideBox}};
\draw[Rule,line width=.5pt] ([xshift=.95cm,yshift=-10.43cm]current page.north west)--([xshift=19.05cm,yshift=-10.43cm]current page.north west);
\node[anchor=north west,inner sep=0pt,text=Classical,font=\bfseries\fontsize{9.4}{12}\selectfont] at ([xshift=.95cm,yshift=-10.66cm]current page.north west) {Richiami};
\node[anchor=north west,inner sep=0pt] at ([xshift=2.6cm,yshift=-10.64cm]current page.north west) {\usebox{\RecallBox}};
\draw[Rule,line width=.35pt] ([xshift=.95cm,yshift=-11.92cm]current page.north west)--([xshift=19.05cm,yshift=-11.92cm]current page.north west);
\node[anchor=north west,inner sep=0pt,text=Muted,text width=9.0cm,font=\fontsize{7.7}{9.3}\selectfont] at ([xshift=.95cm,yshift=-12.06cm]current page.north west) {07 · Bootstrap: ricampionare il campione};
\node[anchor=north west,inner sep=0pt,text=Classical,text width=8.2cm,align=right,font=\fontsize{7.3}{9.1}\selectfont] at ([xshift=9.45cm,yshift=-12.06cm]current page.north west) {\href{https://openintrostat.github.io/ims/}{IMS} · \href{https://cran.r-project.org/doc/manuals/r-release/R-intro.html}{RBASE}};
\node[anchor=north east,inner sep=0pt,text=Muted,font=\fontsize{7.7}{9.3}\selectfont] at ([xshift=19.05cm,yshift=-12.06cm]current page.north west) {73/110};

\end{tikzpicture}
\note{Idea della slide: Il bootstrap non inventa valori di popolazione mai rappresentati nel campione e non corregge selezione, censura o misura sbagliata. Se la statistica dipende fortemente dagli estremi, un piccolo campione può dare un quadro incompleto della sua incertezza.\par\medskip Punto da far emergere: Ricampionare una base distorta riproduce quella distorsione.\par\medskip Fonti della sezione: IMS, RBASE}
\end{frame}
\begin{lrbox}{\TitleBox}\begin{minipage}[t]{18.2cm}\color{Ink}\raggedright\sloppy\fontsize{20.7}{24}\selectfont
\bfseries\color{Navy}Attività 07 · applica il concetto
\end{minipage}\end{lrbox}
\begin{lrbox}{\BodyBox}\begin{minipage}[t]{10.25cm}\color{Ink}\raggedright\sloppy\fontsize{12.4}{16.4}\selectfont
Dai valori 4, 5, 7 e 8 costruisci tre campioni bootstrap di quattro elementi. Spiega perché un valore può comparire più volte.
\end{minipage}\end{lrbox}
\begin{lrbox}{\SideBox}\begin{minipage}[t]{7.45cm}\color{Ink}\raggedright\sloppy\fontsize{12.2}{16.1}\selectfont
\noindent{\color{Classical}\bfseries 01}\hspace{.35cm}Individua la domanda\par\vspace{.22cm}\noindent{\color{Classical}\bfseries 02}\hspace{.35cm}Esplicita le ipotesi\par\vspace{.22cm}\noindent{\color{Classical}\bfseries 03}\hspace{.35cm}Calcola o rappresenta\par\vspace{.22cm}\noindent{\color{Classical}\bfseries 04}\hspace{.35cm}Scrivi una conclusione\par\vspace{.22cm}\noindent{\color{Classical}\bfseries 05}\hspace{.35cm}Indica un limite\par\vspace{.22cm}
\end{minipage}\end{lrbox}
\begin{lrbox}{\RecallBox}\begin{minipage}[t]{16.45cm}\color{Ink}\raggedright\sloppy\fontsize{10.1}{13.2}\selectfont
Come stimiamo una distribuzione campionaria senza conoscere la popolazione?
\end{minipage}\end{lrbox}
\MeasureBox{74-TitleBox}{\TitleBox}
\MeasureBox{74-BodyBox}{\BodyBox}
\MeasureBox{74-SideBox}{\SideBox}
\MeasureBox{74-RecallBox}{\RecallBox}
\begin{frame}[plain]
\hypertarget{slide74}{}\bookmark[level=1,page=74]{Attività 07 · applica il concetto}
\begin{tikzpicture}[remember picture,overlay]
\node[anchor=north east,inner sep=0pt,text=Classical,font=\fontsize{8}{10}\selectfont] at ([xshift=19.05cm,yshift=-.16cm]current page.north west) {\hyperlink{slide3}{Indice}};
\node[anchor=north west,inner sep=0pt,text=Navy,font=\bfseries] at ([xshift=.95cm,yshift=-.62cm]current page.north west) {\usebox{\TitleBox}};
\draw[Classical,line width=1.3pt] ([xshift=.95cm,yshift=-1.68cm]current page.north west)--([xshift=2.25cm,yshift=-1.68cm]current page.north west);
\node[anchor=north west,inner sep=0pt] at ([xshift=.95cm,yshift=-2.12cm]current page.north west) {\usebox{\BodyBox}};
\node[anchor=north west,inner sep=0pt,text=Quantum,font=\bfseries\fontsize{11.5}{14}\selectfont] at ([xshift=11.6cm,yshift=-2.1cm]current page.north west) {Consegna di lavoro};
\node[anchor=north west,inner sep=0pt] at ([xshift=11.6cm,yshift=-2.75cm]current page.north west) {\usebox{\SideBox}};
\draw[Rule,line width=.5pt] ([xshift=.95cm,yshift=-10.43cm]current page.north west)--([xshift=19.05cm,yshift=-10.43cm]current page.north west);
\node[anchor=north west,inner sep=0pt,text=Classical,font=\bfseries\fontsize{9.4}{12}\selectfont] at ([xshift=.95cm,yshift=-10.66cm]current page.north west) {Richiami};
\node[anchor=north west,inner sep=0pt] at ([xshift=2.6cm,yshift=-10.64cm]current page.north west) {\usebox{\RecallBox}};
\draw[Rule,line width=.35pt] ([xshift=.95cm,yshift=-11.92cm]current page.north west)--([xshift=19.05cm,yshift=-11.92cm]current page.north west);
\node[anchor=north west,inner sep=0pt,text=Muted,text width=9.0cm,font=\fontsize{7.7}{9.3}\selectfont] at ([xshift=.95cm,yshift=-12.06cm]current page.north west) {07 · Bootstrap: ricampionare il campione};
\node[anchor=north west,inner sep=0pt,text=Classical,text width=8.2cm,align=right,font=\fontsize{7.3}{9.1}\selectfont] at ([xshift=9.45cm,yshift=-12.06cm]current page.north west) {\href{https://openintrostat.github.io/ims/}{IMS} · \href{https://cran.r-project.org/doc/manuals/r-release/R-intro.html}{RBASE}};
\node[anchor=north east,inner sep=0pt,text=Muted,font=\fontsize{7.7}{9.3}\selectfont] at ([xshift=19.05cm,yshift=-12.06cm]current page.north west) {74/110};

\end{tikzpicture}
\note{Consegna: Dai valori 4, 5, 7 e 8 costruisci tre campioni bootstrap di quattro elementi. Spiega perché un valore può comparire più volte.\par\medskip Soluzione per il docente: Ogni replica seleziona quattro elementi con reinserimento dal campione osservato: per esempio 4,4,7,8 oppure 5,5,5,8. La stessa unità può essere selezionata più volte. Senza reinserimento, una replica completa sarebbe solo un riordinamento e la media non cambierebbe.}
\end{frame}
\begin{lrbox}{\TitleBox}\begin{minipage}[t]{18.2cm}\color{Ink}\raggedright\sloppy\fontsize{20.7}{24}\selectfont
\bfseries\color{Navy}Discussione 07 · soluzione ragionata
\end{minipage}\end{lrbox}
\begin{lrbox}{\BodyBox}\begin{minipage}[t]{10.25cm}\color{Ink}\raggedright\sloppy\fontsize{12.4}{16.4}\selectfont
Ogni replica seleziona quattro elementi con reinserimento dal campione osservato: per esempio 4,4,7,8 oppure 5,5,5,8. La stessa unità può essere selezionata più volte. Senza reinserimento, una replica completa sarebbe solo un riordinamento e la media non cambierebbe.
\end{minipage}\end{lrbox}
\begin{lrbox}{\SideBox}\begin{minipage}[t]{7.45cm}\color{Ink}\raggedright\sloppy\fontsize{12.2}{16.1}\selectfont
Ricampionare una base distorta riproduce quella distorsione.
\end{minipage}\end{lrbox}
\begin{lrbox}{\RecallBox}\begin{minipage}[t]{16.45cm}\color{Ink}\raggedright\sloppy\fontsize{10.1}{13.2}\selectfont
Valuta il collegamento tra domanda, procedimento e conclusione, non soltanto il numero finale.
\end{minipage}\end{lrbox}
\MeasureBox{75-TitleBox}{\TitleBox}
\MeasureBox{75-BodyBox}{\BodyBox}
\MeasureBox{75-SideBox}{\SideBox}
\MeasureBox{75-RecallBox}{\RecallBox}
\begin{frame}[plain]
\hypertarget{slide75}{}\bookmark[level=1,page=75]{Discussione 07 · soluzione ragionata}
\begin{tikzpicture}[remember picture,overlay]
\node[anchor=north east,inner sep=0pt,text=Classical,font=\fontsize{8}{10}\selectfont] at ([xshift=19.05cm,yshift=-.16cm]current page.north west) {\hyperlink{slide3}{Indice}};
\node[anchor=north west,inner sep=0pt,text=Navy,font=\bfseries] at ([xshift=.95cm,yshift=-.62cm]current page.north west) {\usebox{\TitleBox}};
\draw[Classical,line width=1.3pt] ([xshift=.95cm,yshift=-1.68cm]current page.north west)--([xshift=2.25cm,yshift=-1.68cm]current page.north west);
\node[anchor=north west,inner sep=0pt] at ([xshift=.95cm,yshift=-2.12cm]current page.north west) {\usebox{\BodyBox}};
\node[anchor=north west,inner sep=0pt,text=Quantum,font=\bfseries\fontsize{11.5}{14}\selectfont] at ([xshift=11.6cm,yshift=-2.1cm]current page.north west) {Criterio di lettura};
\node[anchor=north west,inner sep=0pt] at ([xshift=11.6cm,yshift=-2.75cm]current page.north west) {\usebox{\SideBox}};
\draw[Rule,line width=.5pt] ([xshift=.95cm,yshift=-10.43cm]current page.north west)--([xshift=19.05cm,yshift=-10.43cm]current page.north west);
\node[anchor=north west,inner sep=0pt,text=Classical,font=\bfseries\fontsize{9.4}{12}\selectfont] at ([xshift=.95cm,yshift=-10.66cm]current page.north west) {Richiami};
\node[anchor=north west,inner sep=0pt] at ([xshift=2.6cm,yshift=-10.64cm]current page.north west) {\usebox{\RecallBox}};
\draw[Rule,line width=.35pt] ([xshift=.95cm,yshift=-11.92cm]current page.north west)--([xshift=19.05cm,yshift=-11.92cm]current page.north west);
\node[anchor=north west,inner sep=0pt,text=Muted,text width=9.0cm,font=\fontsize{7.7}{9.3}\selectfont] at ([xshift=.95cm,yshift=-12.06cm]current page.north west) {07 · Bootstrap: ricampionare il campione};
\node[anchor=north west,inner sep=0pt,text=Classical,text width=8.2cm,align=right,font=\fontsize{7.3}{9.1}\selectfont] at ([xshift=9.45cm,yshift=-12.06cm]current page.north west) {\href{https://openintrostat.github.io/ims/}{IMS} · \href{https://cran.r-project.org/doc/manuals/r-release/R-intro.html}{RBASE}};
\node[anchor=north east,inner sep=0pt,text=Muted,font=\fontsize{7.7}{9.3}\selectfont] at ([xshift=19.05cm,yshift=-12.06cm]current page.north west) {75/110};

\end{tikzpicture}
\note{Ogni replica seleziona quattro elementi con reinserimento dal campione osservato: per esempio 4,4,7,8 oppure 5,5,5,8. La stessa unità può essere selezionata più volte. Senza reinserimento, una replica completa sarebbe solo un riordinamento e la media non cambierebbe.}
\end{frame}
\begin{frame}[plain]\hypertarget{slide76}{}
\bookmark[level=0,page=76]{Ricampionare al livello corretto}
\begin{tikzpicture}[remember picture,overlay]
\fill[Night] (current page.south west) rectangle (current page.north east);
\draw[Mint,line width=2pt] ([xshift=1.05cm,yshift=-1.35cm]current page.north west)--([xshift=3.0cm,yshift=-1.35cm]current page.north west);
\node[anchor=north west,inner sep=0,text=Mint,font=\bfseries\fontsize{11}{14}\selectfont] at ([xshift=1.05cm,yshift=-1.7cm]current page.north west) {SEZIONE 08};
\node[anchor=north west,inner sep=0,text=white,text width=15.8cm,align=left,font=\bfseries\fontsize{31}{38}\selectfont] at ([xshift=1.05cm,yshift=-2.65cm]current page.north west) {Ricampionare al livello corretto};
\node[anchor=north west,inner sep=0,text=Lilac,text width=15.8cm,align=left,font=\fontsize{15}{21}\selectfont] at ([xshift=1.05cm,yshift=-6.4cm]current page.north west) {Quale unità deve essere estratta a ogni replica?};
\node[anchor=north west,inner sep=0,text=white,text width=16.5cm,align=left,font=\fontsize{11}{15}\selectfont] at ([xshift=1.05cm,yshift=-8.6cm]current page.north west) {Comprendere il concetto · leggere un esempio · applicare e discutere};
\node[anchor=north east,inner sep=0,text=Mint,font=\fontsize{8}{10}\selectfont] at ([xshift=19.05cm,yshift=-.2cm]current page.north west) {\hyperlink{slide3}{Indice}};
\node[anchor=south east,inner sep=0,text=Lilac!45,font=\bfseries\fontsize{64}{68}\selectfont] at ([xshift=18.95cm,yshift=-11.35cm]current page.north west) {08};
\draw[Lilac!30,line width=.35pt] ([xshift=.95cm,yshift=-11.92cm]current page.north west)--([xshift=19.05cm,yshift=-11.92cm]current page.north west);
\node[anchor=north west,inner sep=0pt,text=Lilac,text width=9.0cm,font=\fontsize{7.7}{9.3}\selectfont] at ([xshift=.95cm,yshift=-12.06cm]current page.north west) {08 · Ricampionare al livello corretto};
\node[anchor=north west,inner sep=0pt,text=Mint,text width=8.2cm,align=right,font=\fontsize{7.3}{9.1}\selectfont] at ([xshift=9.45cm,yshift=-12.06cm]current page.north west) {\href{https://openintrostat.github.io/ims/}{IMS} · CORSO};
\node[anchor=north east,inner sep=0pt,text=Lilac,font=\fontsize{7.7}{9.3}\selectfont] at ([xshift=19.05cm,yshift=-12.06cm]current page.north west) {76/110};
\end{tikzpicture}\end{frame}
\begin{lrbox}{\TitleBox}\begin{minipage}[t]{18.2cm}\color{Ink}\raggedright\sloppy\fontsize{20.7}{24}\selectfont
\bfseries\color{Navy}Unità di ricampionamento
\end{minipage}\end{lrbox}
\begin{lrbox}{\BodyBox}\begin{minipage}[t]{10.25cm}\color{Ink}\raggedright\sloppy\fontsize{12.4}{16.4}\selectfont
Il bootstrap deve riprodurre il livello di variazione rilevante del campionamento. Per dati indipendenti sono le righe; per gruppi dipendenti possono essere interi gruppi. La scelta dipende da come i dati sono stati raccolti e da quale parametro interessa.
\end{minipage}\end{lrbox}
\begin{lrbox}{\SideBox}\begin{minipage}[t]{7.45cm}\color{Ink}\raggedright\sloppy\fontsize{12.2}{16.1}\selectfont
Il disegno non si deduce dal numero di colonne. Va letto nei metadati e nel protocollo.
\end{minipage}\end{lrbox}
\begin{lrbox}{\RecallBox}\begin{minipage}[t]{16.45cm}\color{Ink}\raggedright\sloppy\fontsize{10.1}{13.2}\selectfont
La procedura più breve da scrivere non è necessariamente quella scientificamente appropriata.
\end{minipage}\end{lrbox}
\MeasureBox{77-TitleBox}{\TitleBox}
\MeasureBox{77-BodyBox}{\BodyBox}
\MeasureBox{77-SideBox}{\SideBox}
\MeasureBox{77-RecallBox}{\RecallBox}
\begin{frame}[plain]
\hypertarget{slide77}{}\bookmark[level=1,page=77]{Unità di ricampionamento}
\begin{tikzpicture}[remember picture,overlay]
\node[anchor=north east,inner sep=0pt,text=Classical,font=\fontsize{8}{10}\selectfont] at ([xshift=19.05cm,yshift=-.16cm]current page.north west) {\hyperlink{slide3}{Indice}};
\node[anchor=north west,inner sep=0pt,text=Navy,font=\bfseries] at ([xshift=.95cm,yshift=-.62cm]current page.north west) {\usebox{\TitleBox}};
\draw[Classical,line width=1.3pt] ([xshift=.95cm,yshift=-1.68cm]current page.north west)--([xshift=2.25cm,yshift=-1.68cm]current page.north west);
\node[anchor=north west,inner sep=0pt] at ([xshift=.95cm,yshift=-2.12cm]current page.north west) {\usebox{\BodyBox}};
\node[anchor=north west,inner sep=0pt,text=Quantum,font=\bfseries\fontsize{11.5}{14}\selectfont] at ([xshift=11.6cm,yshift=-2.1cm]current page.north west) {Esempio e interpretazione};
\node[anchor=north west,inner sep=0pt] at ([xshift=11.6cm,yshift=-2.75cm]current page.north west) {\usebox{\SideBox}};
\draw[Rule,line width=.5pt] ([xshift=.95cm,yshift=-10.43cm]current page.north west)--([xshift=19.05cm,yshift=-10.43cm]current page.north west);
\node[anchor=north west,inner sep=0pt,text=Classical,font=\bfseries\fontsize{9.4}{12}\selectfont] at ([xshift=.95cm,yshift=-10.66cm]current page.north west) {Richiami};
\node[anchor=north west,inner sep=0pt] at ([xshift=2.6cm,yshift=-10.64cm]current page.north west) {\usebox{\RecallBox}};
\draw[Rule,line width=.35pt] ([xshift=.95cm,yshift=-11.92cm]current page.north west)--([xshift=19.05cm,yshift=-11.92cm]current page.north west);
\node[anchor=north west,inner sep=0pt,text=Muted,text width=9.0cm,font=\fontsize{7.7}{9.3}\selectfont] at ([xshift=.95cm,yshift=-12.06cm]current page.north west) {08 · Ricampionare al livello corretto};
\node[anchor=north west,inner sep=0pt,text=Classical,text width=8.2cm,align=right,font=\fontsize{7.3}{9.1}\selectfont] at ([xshift=9.45cm,yshift=-12.06cm]current page.north west) {\href{https://openintrostat.github.io/ims/}{IMS} · CORSO};
\node[anchor=north east,inner sep=0pt,text=Muted,font=\fontsize{7.7}{9.3}\selectfont] at ([xshift=19.05cm,yshift=-12.06cm]current page.north west) {77/110};

\end{tikzpicture}
\note{Idea della slide: Il bootstrap deve riprodurre il livello di variazione rilevante del campionamento. Per dati indipendenti sono le righe; per gruppi dipendenti possono essere interi gruppi. La scelta dipende da come i dati sono stati raccolti e da quale parametro interessa.\par\medskip Punto da far emergere: La procedura più breve da scrivere non è necessariamente quella scientificamente appropriata.\par\medskip Fonti della sezione: IMS, CORSO}
\end{frame}
\begin{lrbox}{\TitleBox}\begin{minipage}[t]{18.2cm}\color{Ink}\raggedright\sloppy\fontsize{20.7}{24}\selectfont
\bfseries\color{Navy}Due gruppi indipendenti
\end{minipage}\end{lrbox}
\begin{lrbox}{\BodyBox}\begin{minipage}[t]{8.2cm}\color{Ink}\raggedright\sloppy\fontsize{12.4}{16.4}\selectfont
Per stimare una differenza tra medie di gruppi indipendenti possiamo ricampionare separatamente entro ciascun gruppo. Ogni replica conserva le rispettive numerosità. La differenza viene calcolata nello stesso ordine, evitando cambi di segno tra script e relazione.
\end{minipage}\end{lrbox}
\begin{lrbox}{\SideBox}\begin{minipage}[t]{9.35cm}\color{Ink}\raggedright\sloppy\fontsize{12.2}{16.1}\selectfont
\begin{lstlisting}
diff_b <- mean(sample(A, length(A), TRUE)) -
          mean(sample(B, length(B), TRUE))
\end{lstlisting}
\end{minipage}\end{lrbox}
\begin{lrbox}{\RecallBox}\begin{minipage}[t]{16.45cm}\color{Ink}\raggedright\sloppy\fontsize{10.1}{13.2}\selectfont
Il contrasto A–B deve rimanere definito in tutte le repliche e nei grafici.
\end{minipage}\end{lrbox}
\MeasureBox{78-TitleBox}{\TitleBox}
\MeasureBox{78-BodyBox}{\BodyBox}
\MeasureBox{78-SideBox}{\SideBox}
\MeasureBox{78-RecallBox}{\RecallBox}
\begin{frame}[plain]
\hypertarget{slide78}{}\bookmark[level=1,page=78]{Due gruppi indipendenti}
\begin{tikzpicture}[remember picture,overlay]
\node[anchor=north east,inner sep=0pt,text=Classical,font=\fontsize{8}{10}\selectfont] at ([xshift=19.05cm,yshift=-.16cm]current page.north west) {\hyperlink{slide3}{Indice}};
\node[anchor=north west,inner sep=0pt,text=Navy,font=\bfseries] at ([xshift=.95cm,yshift=-.62cm]current page.north west) {\usebox{\TitleBox}};
\draw[Classical,line width=1.3pt] ([xshift=.95cm,yshift=-1.68cm]current page.north west)--([xshift=2.25cm,yshift=-1.68cm]current page.north west);
\node[anchor=north west,inner sep=0pt] at ([xshift=.95cm,yshift=-2.12cm]current page.north west) {\usebox{\BodyBox}};
\node[anchor=north west,inner sep=0pt,text=Quantum,font=\bfseries\fontsize{11.5}{14}\selectfont] at ([xshift=9.7cm,yshift=-2.1cm]current page.north west) {Protocollo R};
\node[anchor=north west,inner sep=0pt] at ([xshift=9.7cm,yshift=-2.75cm]current page.north west) {\usebox{\SideBox}};
\draw[Rule,line width=.5pt] ([xshift=.95cm,yshift=-10.43cm]current page.north west)--([xshift=19.05cm,yshift=-10.43cm]current page.north west);
\node[anchor=north west,inner sep=0pt,text=Classical,font=\bfseries\fontsize{9.4}{12}\selectfont] at ([xshift=.95cm,yshift=-10.66cm]current page.north west) {Richiami};
\node[anchor=north west,inner sep=0pt] at ([xshift=2.6cm,yshift=-10.64cm]current page.north west) {\usebox{\RecallBox}};
\draw[Rule,line width=.35pt] ([xshift=.95cm,yshift=-11.92cm]current page.north west)--([xshift=19.05cm,yshift=-11.92cm]current page.north west);
\node[anchor=north west,inner sep=0pt,text=Muted,text width=9.0cm,font=\fontsize{7.7}{9.3}\selectfont] at ([xshift=.95cm,yshift=-12.06cm]current page.north west) {08 · Ricampionare al livello corretto};
\node[anchor=north west,inner sep=0pt,text=Classical,text width=8.2cm,align=right,font=\fontsize{7.3}{9.1}\selectfont] at ([xshift=9.45cm,yshift=-12.06cm]current page.north west) {\href{https://openintrostat.github.io/ims/}{IMS} · CORSO};
\node[anchor=north east,inner sep=0pt,text=Muted,font=\fontsize{7.7}{9.3}\selectfont] at ([xshift=19.05cm,yshift=-12.06cm]current page.north west) {78/110};

\end{tikzpicture}
\note{Idea della slide: Per stimare una differenza tra medie di gruppi indipendenti possiamo ricampionare separatamente entro ciascun gruppo. Ogni replica conserva le rispettive numerosità. La differenza viene calcolata nello stesso ordine, evitando cambi di segno tra script e relazione.\par\medskip Punto da far emergere: Il contrasto A–B deve rimanere definito in tutte le repliche e nei grafici.\par\medskip Fonti della sezione: IMS, CORSO}
\end{frame}
\begin{lrbox}{\TitleBox}\begin{minipage}[t]{18.2cm}\color{Ink}\raggedright\sloppy\fontsize{20.7}{24}\selectfont
\bfseries\color{Navy}Coppie di misure
\end{minipage}\end{lrbox}
\begin{lrbox}{\BodyBox}\begin{minipage}[t]{8.2cm}\color{Ink}\raggedright\sloppy\fontsize{12.4}{16.4}\selectfont
Quando le osservazioni sono appaiate, ricampioniamo le coppie oppure le differenze entro coppia. Estrarre prima e dopo separatamente distrugge l'abbinamento e cambia il problema. Il legame tra misure è informazione da preservare, non un ostacolo da rimuovere.
\end{minipage}\end{lrbox}
\begin{lrbox}{\SideBox}\begin{minipage}[t]{9.35cm}\color{Ink}\raggedright\sloppy\fontsize{12.2}{16.1}\selectfont
\begin{lstlisting}
delta <- coppie$dopo - coppie$prima
b <- replicate(5000,
 mean(sample(delta, length(delta), TRUE)))
\end{lstlisting}
\end{minipage}\end{lrbox}
\begin{lrbox}{\RecallBox}\begin{minipage}[t]{16.45cm}\color{Ink}\raggedright\sloppy\fontsize{10.1}{13.2}\selectfont
La numerosità è il numero di coppie, non la somma delle due colonne.
\end{minipage}\end{lrbox}
\MeasureBox{79-TitleBox}{\TitleBox}
\MeasureBox{79-BodyBox}{\BodyBox}
\MeasureBox{79-SideBox}{\SideBox}
\MeasureBox{79-RecallBox}{\RecallBox}
\begin{frame}[plain]
\hypertarget{slide79}{}\bookmark[level=1,page=79]{Coppie di misure}
\begin{tikzpicture}[remember picture,overlay]
\node[anchor=north east,inner sep=0pt,text=Classical,font=\fontsize{8}{10}\selectfont] at ([xshift=19.05cm,yshift=-.16cm]current page.north west) {\hyperlink{slide3}{Indice}};
\node[anchor=north west,inner sep=0pt,text=Navy,font=\bfseries] at ([xshift=.95cm,yshift=-.62cm]current page.north west) {\usebox{\TitleBox}};
\draw[Classical,line width=1.3pt] ([xshift=.95cm,yshift=-1.68cm]current page.north west)--([xshift=2.25cm,yshift=-1.68cm]current page.north west);
\node[anchor=north west,inner sep=0pt] at ([xshift=.95cm,yshift=-2.12cm]current page.north west) {\usebox{\BodyBox}};
\node[anchor=north west,inner sep=0pt,text=Quantum,font=\bfseries\fontsize{11.5}{14}\selectfont] at ([xshift=9.7cm,yshift=-2.1cm]current page.north west) {Protocollo R};
\node[anchor=north west,inner sep=0pt] at ([xshift=9.7cm,yshift=-2.75cm]current page.north west) {\usebox{\SideBox}};
\draw[Rule,line width=.5pt] ([xshift=.95cm,yshift=-10.43cm]current page.north west)--([xshift=19.05cm,yshift=-10.43cm]current page.north west);
\node[anchor=north west,inner sep=0pt,text=Classical,font=\bfseries\fontsize{9.4}{12}\selectfont] at ([xshift=.95cm,yshift=-10.66cm]current page.north west) {Richiami};
\node[anchor=north west,inner sep=0pt] at ([xshift=2.6cm,yshift=-10.64cm]current page.north west) {\usebox{\RecallBox}};
\draw[Rule,line width=.35pt] ([xshift=.95cm,yshift=-11.92cm]current page.north west)--([xshift=19.05cm,yshift=-11.92cm]current page.north west);
\node[anchor=north west,inner sep=0pt,text=Muted,text width=9.0cm,font=\fontsize{7.7}{9.3}\selectfont] at ([xshift=.95cm,yshift=-12.06cm]current page.north west) {08 · Ricampionare al livello corretto};
\node[anchor=north west,inner sep=0pt,text=Classical,text width=8.2cm,align=right,font=\fontsize{7.3}{9.1}\selectfont] at ([xshift=9.45cm,yshift=-12.06cm]current page.north west) {\href{https://openintrostat.github.io/ims/}{IMS} · CORSO};
\node[anchor=north east,inner sep=0pt,text=Muted,font=\fontsize{7.7}{9.3}\selectfont] at ([xshift=19.05cm,yshift=-12.06cm]current page.north west) {79/110};

\end{tikzpicture}
\note{Idea della slide: Quando le osservazioni sono appaiate, ricampioniamo le coppie oppure le differenze entro coppia. Estrarre prima e dopo separatamente distrugge l'abbinamento e cambia il problema. Il legame tra misure è informazione da preservare, non un ostacolo da rimuovere.\par\medskip Punto da far emergere: La numerosità è il numero di coppie, non la somma delle due colonne.\par\medskip Fonti della sezione: IMS, CORSO}
\end{frame}
\begin{lrbox}{\TitleBox}\begin{minipage}[t]{18.2cm}\color{Ink}\raggedright\sloppy\fontsize{20.7}{24}\selectfont
\bfseries\color{Navy}Stazioni come gruppi
\end{minipage}\end{lrbox}
\begin{lrbox}{\BodyBox}\begin{minipage}[t]{10.25cm}\color{Ink}\raggedright\sloppy\fontsize{12.4}{16.4}\selectfont
Nel bootstrap per cluster estraiamo stazioni e portiamo con esse le misure associate. Se una stazione è estratta due volte, la replica deve includerla due volte. Questo approccio rappresenta la variabilità tra stazioni, ma richiede un numero adeguato di gruppi.
\end{minipage}\end{lrbox}
\begin{lrbox}{\SideBox}\begin{minipage}[t]{7.45cm}\color{Ink}\raggedright\sloppy\fontsize{12.2}{16.1}\selectfont
\noindent{\color{Classical}\bfseries 01}\hspace{.35cm}Elenco stazioni\par\vspace{.22cm}\noindent{\color{Classical}\bfseries 02}\hspace{.35cm}Estrazione con reinserimento\par\vspace{.22cm}\noindent{\color{Classical}\bfseries 03}\hspace{.35cm}Recupero di tutti i provini\par\vspace{.22cm}\noindent{\color{Classical}\bfseries 04}\hspace{.35cm}Calcolo della statistica\par\vspace{.22cm}
\end{minipage}\end{lrbox}
\begin{lrbox}{\RecallBox}\begin{minipage}[t]{16.45cm}\color{Ink}\raggedright\sloppy\fontsize{10.1}{13.2}\selectfont
Sei stazioni sono poche: il metodo non rende automaticamente robusta l'inferenza.
\end{minipage}\end{lrbox}
\MeasureBox{80-TitleBox}{\TitleBox}
\MeasureBox{80-BodyBox}{\BodyBox}
\MeasureBox{80-SideBox}{\SideBox}
\MeasureBox{80-RecallBox}{\RecallBox}
\begin{frame}[plain]
\hypertarget{slide80}{}\bookmark[level=1,page=80]{Stazioni come gruppi}
\begin{tikzpicture}[remember picture,overlay]
\node[anchor=north east,inner sep=0pt,text=Classical,font=\fontsize{8}{10}\selectfont] at ([xshift=19.05cm,yshift=-.16cm]current page.north west) {\hyperlink{slide3}{Indice}};
\node[anchor=north west,inner sep=0pt,text=Navy,font=\bfseries] at ([xshift=.95cm,yshift=-.62cm]current page.north west) {\usebox{\TitleBox}};
\draw[Classical,line width=1.3pt] ([xshift=.95cm,yshift=-1.68cm]current page.north west)--([xshift=2.25cm,yshift=-1.68cm]current page.north west);
\node[anchor=north west,inner sep=0pt] at ([xshift=.95cm,yshift=-2.12cm]current page.north west) {\usebox{\BodyBox}};
\node[anchor=north west,inner sep=0pt,text=Quantum,font=\bfseries\fontsize{11.5}{14}\selectfont] at ([xshift=11.6cm,yshift=-2.1cm]current page.north west) {Passaggi da seguire};
\node[anchor=north west,inner sep=0pt] at ([xshift=11.6cm,yshift=-2.75cm]current page.north west) {\usebox{\SideBox}};
\draw[Rule,line width=.5pt] ([xshift=.95cm,yshift=-10.43cm]current page.north west)--([xshift=19.05cm,yshift=-10.43cm]current page.north west);
\node[anchor=north west,inner sep=0pt,text=Classical,font=\bfseries\fontsize{9.4}{12}\selectfont] at ([xshift=.95cm,yshift=-10.66cm]current page.north west) {Richiami};
\node[anchor=north west,inner sep=0pt] at ([xshift=2.6cm,yshift=-10.64cm]current page.north west) {\usebox{\RecallBox}};
\draw[Rule,line width=.35pt] ([xshift=.95cm,yshift=-11.92cm]current page.north west)--([xshift=19.05cm,yshift=-11.92cm]current page.north west);
\node[anchor=north west,inner sep=0pt,text=Muted,text width=9.0cm,font=\fontsize{7.7}{9.3}\selectfont] at ([xshift=.95cm,yshift=-12.06cm]current page.north west) {08 · Ricampionare al livello corretto};
\node[anchor=north west,inner sep=0pt,text=Classical,text width=8.2cm,align=right,font=\fontsize{7.3}{9.1}\selectfont] at ([xshift=9.45cm,yshift=-12.06cm]current page.north west) {\href{https://openintrostat.github.io/ims/}{IMS} · CORSO};
\node[anchor=north east,inner sep=0pt,text=Muted,font=\fontsize{7.7}{9.3}\selectfont] at ([xshift=19.05cm,yshift=-12.06cm]current page.north west) {80/110};

\end{tikzpicture}
\note{Idea della slide: Nel bootstrap per cluster estraiamo stazioni e portiamo con esse le misure associate. Se una stazione è estratta due volte, la replica deve includerla due volte. Questo approccio rappresenta la variabilità tra stazioni, ma richiede un numero adeguato di gruppi.\par\medskip Punto da far emergere: Sei stazioni sono poche: il metodo non rende automaticamente robusta l'inferenza.\par\medskip Fonti della sezione: IMS, CORSO}
\end{frame}
\begin{lrbox}{\TitleBox}\begin{minipage}[t]{18.2cm}\color{Ink}\raggedright\sloppy\fontsize{20.7}{24}\selectfont
\bfseries\color{Navy}Stratificazione e quote
\end{minipage}\end{lrbox}
\begin{lrbox}{\BodyBox}\begin{minipage}[t]{10.25cm}\color{Ink}\raggedright\sloppy\fontsize{12.4}{16.4}\selectfont
Se il disegno prevede strati, può essere necessario ricampionare entro ciascuno e mantenere pesi coerenti con il parametro. Mescolare tutte le righe può alterare la composizione degli strati e rappresentare un campionamento diverso da quello effettivamente adottato.
\end{minipage}\end{lrbox}
\begin{lrbox}{\SideBox}\begin{minipage}[t]{7.45cm}\color{Ink}\raggedright\sloppy\fontsize{12.2}{16.1}\selectfont
Esempio: tre stazioni per ciascuno di due affioramenti. Il ricampionamento e l'aggregazione devono rispettare il contrasto tra affioramenti.
\end{minipage}\end{lrbox}
\begin{lrbox}{\RecallBox}\begin{minipage}[t]{16.45cm}\color{Ink}\raggedright\sloppy\fontsize{10.1}{13.2}\selectfont
Un bootstrap ben implementato può essere sbagliato se riproduce il disegno sbagliato.
\end{minipage}\end{lrbox}
\MeasureBox{81-TitleBox}{\TitleBox}
\MeasureBox{81-BodyBox}{\BodyBox}
\MeasureBox{81-SideBox}{\SideBox}
\MeasureBox{81-RecallBox}{\RecallBox}
\begin{frame}[plain]
\hypertarget{slide81}{}\bookmark[level=1,page=81]{Stratificazione e quote}
\begin{tikzpicture}[remember picture,overlay]
\node[anchor=north east,inner sep=0pt,text=Classical,font=\fontsize{8}{10}\selectfont] at ([xshift=19.05cm,yshift=-.16cm]current page.north west) {\hyperlink{slide3}{Indice}};
\node[anchor=north west,inner sep=0pt,text=Navy,font=\bfseries] at ([xshift=.95cm,yshift=-.62cm]current page.north west) {\usebox{\TitleBox}};
\draw[Classical,line width=1.3pt] ([xshift=.95cm,yshift=-1.68cm]current page.north west)--([xshift=2.25cm,yshift=-1.68cm]current page.north west);
\node[anchor=north west,inner sep=0pt] at ([xshift=.95cm,yshift=-2.12cm]current page.north west) {\usebox{\BodyBox}};
\node[anchor=north west,inner sep=0pt,text=Quantum,font=\bfseries\fontsize{11.5}{14}\selectfont] at ([xshift=11.6cm,yshift=-2.1cm]current page.north west) {Esempio e interpretazione};
\node[anchor=north west,inner sep=0pt] at ([xshift=11.6cm,yshift=-2.75cm]current page.north west) {\usebox{\SideBox}};
\draw[Rule,line width=.5pt] ([xshift=.95cm,yshift=-10.43cm]current page.north west)--([xshift=19.05cm,yshift=-10.43cm]current page.north west);
\node[anchor=north west,inner sep=0pt,text=Classical,font=\bfseries\fontsize{9.4}{12}\selectfont] at ([xshift=.95cm,yshift=-10.66cm]current page.north west) {Richiami};
\node[anchor=north west,inner sep=0pt] at ([xshift=2.6cm,yshift=-10.64cm]current page.north west) {\usebox{\RecallBox}};
\draw[Rule,line width=.35pt] ([xshift=.95cm,yshift=-11.92cm]current page.north west)--([xshift=19.05cm,yshift=-11.92cm]current page.north west);
\node[anchor=north west,inner sep=0pt,text=Muted,text width=9.0cm,font=\fontsize{7.7}{9.3}\selectfont] at ([xshift=.95cm,yshift=-12.06cm]current page.north west) {08 · Ricampionare al livello corretto};
\node[anchor=north west,inner sep=0pt,text=Classical,text width=8.2cm,align=right,font=\fontsize{7.3}{9.1}\selectfont] at ([xshift=9.45cm,yshift=-12.06cm]current page.north west) {\href{https://openintrostat.github.io/ims/}{IMS} · CORSO};
\node[anchor=north east,inner sep=0pt,text=Muted,font=\fontsize{7.7}{9.3}\selectfont] at ([xshift=19.05cm,yshift=-12.06cm]current page.north west) {81/110};

\end{tikzpicture}
\note{Idea della slide: Se il disegno prevede strati, può essere necessario ricampionare entro ciascuno e mantenere pesi coerenti con il parametro. Mescolare tutte le righe può alterare la composizione degli strati e rappresentare un campionamento diverso da quello effettivamente adottato.\par\medskip Punto da far emergere: Un bootstrap ben implementato può essere sbagliato se riproduce il disegno sbagliato.\par\medskip Fonti della sezione: IMS, CORSO}
\end{frame}
\begin{lrbox}{\TitleBox}\begin{minipage}[t]{18.2cm}\color{Ink}\raggedright\sloppy\fontsize{20.7}{24}\selectfont
\bfseries\color{Navy}Serie temporali: conservare la dipendenza
\end{minipage}\end{lrbox}
\begin{lrbox}{\BodyBox}\begin{minipage}[t]{10.25cm}\color{Ink}\raggedright\sloppy\fontsize{12.4}{16.4}\selectfont
Ricampionare singoli istanti elimina l'ordine e l'autocorrelazione. Per serie dipendenti esistono ricampionamenti a blocchi e modelli del processo. Nel primo anno il punto fondamentale è riconoscere che il bootstrap delle righe non si trasferisce automaticamente nel tempo.
\end{minipage}\end{lrbox}
\begin{lrbox}{\SideBox}\begin{minipage}[t]{7.45cm}\color{Ink}\raggedright\sloppy\fontsize{12.2}{16.1}\selectfont
Un blocco conserva una parte della sequenza locale; la sua lunghezza introduce una scelta da motivare.
\end{minipage}\end{lrbox}
\begin{lrbox}{\RecallBox}\begin{minipage}[t]{16.45cm}\color{Ink}\raggedright\sloppy\fontsize{10.1}{13.2}\selectfont
L'ordine temporale contiene informazione che non va distrutta inconsapevolmente.
\end{minipage}\end{lrbox}
\MeasureBox{82-TitleBox}{\TitleBox}
\MeasureBox{82-BodyBox}{\BodyBox}
\MeasureBox{82-SideBox}{\SideBox}
\MeasureBox{82-RecallBox}{\RecallBox}
\begin{frame}[plain]
\hypertarget{slide82}{}\bookmark[level=1,page=82]{Serie temporali: conservare la dipendenza}
\begin{tikzpicture}[remember picture,overlay]
\node[anchor=north east,inner sep=0pt,text=Classical,font=\fontsize{8}{10}\selectfont] at ([xshift=19.05cm,yshift=-.16cm]current page.north west) {\hyperlink{slide3}{Indice}};
\node[anchor=north west,inner sep=0pt,text=Navy,font=\bfseries] at ([xshift=.95cm,yshift=-.62cm]current page.north west) {\usebox{\TitleBox}};
\draw[Classical,line width=1.3pt] ([xshift=.95cm,yshift=-1.68cm]current page.north west)--([xshift=2.25cm,yshift=-1.68cm]current page.north west);
\node[anchor=north west,inner sep=0pt] at ([xshift=.95cm,yshift=-2.12cm]current page.north west) {\usebox{\BodyBox}};
\node[anchor=north west,inner sep=0pt,text=Quantum,font=\bfseries\fontsize{11.5}{14}\selectfont] at ([xshift=11.6cm,yshift=-2.1cm]current page.north west) {Esempio e interpretazione};
\node[anchor=north west,inner sep=0pt] at ([xshift=11.6cm,yshift=-2.75cm]current page.north west) {\usebox{\SideBox}};
\draw[Rule,line width=.5pt] ([xshift=.95cm,yshift=-10.43cm]current page.north west)--([xshift=19.05cm,yshift=-10.43cm]current page.north west);
\node[anchor=north west,inner sep=0pt,text=Classical,font=\bfseries\fontsize{9.4}{12}\selectfont] at ([xshift=.95cm,yshift=-10.66cm]current page.north west) {Richiami};
\node[anchor=north west,inner sep=0pt] at ([xshift=2.6cm,yshift=-10.64cm]current page.north west) {\usebox{\RecallBox}};
\draw[Rule,line width=.35pt] ([xshift=.95cm,yshift=-11.92cm]current page.north west)--([xshift=19.05cm,yshift=-11.92cm]current page.north west);
\node[anchor=north west,inner sep=0pt,text=Muted,text width=9.0cm,font=\fontsize{7.7}{9.3}\selectfont] at ([xshift=.95cm,yshift=-12.06cm]current page.north west) {08 · Ricampionare al livello corretto};
\node[anchor=north west,inner sep=0pt,text=Classical,text width=8.2cm,align=right,font=\fontsize{7.3}{9.1}\selectfont] at ([xshift=9.45cm,yshift=-12.06cm]current page.north west) {\href{https://openintrostat.github.io/ims/}{IMS} · CORSO};
\node[anchor=north east,inner sep=0pt,text=Muted,font=\fontsize{7.7}{9.3}\selectfont] at ([xshift=19.05cm,yshift=-12.06cm]current page.north west) {82/110};

\end{tikzpicture}
\note{Idea della slide: Ricampionare singoli istanti elimina l'ordine e l'autocorrelazione. Per serie dipendenti esistono ricampionamenti a blocchi e modelli del processo. Nel primo anno il punto fondamentale è riconoscere che il bootstrap delle righe non si trasferisce automaticamente nel tempo.\par\medskip Punto da far emergere: L'ordine temporale contiene informazione che non va distrutta inconsapevolmente.\par\medskip Fonti della sezione: IMS, CORSO}
\end{frame}
\begin{lrbox}{\TitleBox}\begin{minipage}[t]{18.2cm}\color{Ink}\raggedright\sloppy\fontsize{20.7}{24}\selectfont
\bfseries\color{Navy}Una verifica di sensibilità del disegno
\end{minipage}\end{lrbox}
\begin{lrbox}{\BodyBox}\begin{minipage}[t]{10.25cm}\color{Ink}\raggedright\sloppy\fontsize{12.4}{16.4}\selectfont
Confrontare risultati ottenuti con unità di ricampionamento diverse rende visibile quanto la dipendenza conti. Non scegliamo la procedura che produce l'intervallo più stretto. Scegliamo quella coerente con il disegno e discutiamo la fragilità quando i gruppi sono pochi.
\end{minipage}\end{lrbox}
\begin{lrbox}{\SideBox}\begin{minipage}[t]{7.45cm}\color{Ink}\raggedright\sloppy\fontsize{12.2}{16.1}\selectfont
Domanda per il rapporto: quale aspetto del campionamento è imitato da ogni replica?
\end{minipage}\end{lrbox}
\begin{lrbox}{\RecallBox}\begin{minipage}[t]{16.45cm}\color{Ink}\raggedright\sloppy\fontsize{10.1}{13.2}\selectfont
La precisione apparente non è un criterio di selezione del metodo.
\end{minipage}\end{lrbox}
\MeasureBox{83-TitleBox}{\TitleBox}
\MeasureBox{83-BodyBox}{\BodyBox}
\MeasureBox{83-SideBox}{\SideBox}
\MeasureBox{83-RecallBox}{\RecallBox}
\begin{frame}[plain]
\hypertarget{slide83}{}\bookmark[level=1,page=83]{Una verifica di sensibilità del disegno}
\begin{tikzpicture}[remember picture,overlay]
\node[anchor=north east,inner sep=0pt,text=Classical,font=\fontsize{8}{10}\selectfont] at ([xshift=19.05cm,yshift=-.16cm]current page.north west) {\hyperlink{slide3}{Indice}};
\node[anchor=north west,inner sep=0pt,text=Navy,font=\bfseries] at ([xshift=.95cm,yshift=-.62cm]current page.north west) {\usebox{\TitleBox}};
\draw[Classical,line width=1.3pt] ([xshift=.95cm,yshift=-1.68cm]current page.north west)--([xshift=2.25cm,yshift=-1.68cm]current page.north west);
\node[anchor=north west,inner sep=0pt] at ([xshift=.95cm,yshift=-2.12cm]current page.north west) {\usebox{\BodyBox}};
\node[anchor=north west,inner sep=0pt,text=Quantum,font=\bfseries\fontsize{11.5}{14}\selectfont] at ([xshift=11.6cm,yshift=-2.1cm]current page.north west) {Esempio e interpretazione};
\node[anchor=north west,inner sep=0pt] at ([xshift=11.6cm,yshift=-2.75cm]current page.north west) {\usebox{\SideBox}};
\draw[Rule,line width=.5pt] ([xshift=.95cm,yshift=-10.43cm]current page.north west)--([xshift=19.05cm,yshift=-10.43cm]current page.north west);
\node[anchor=north west,inner sep=0pt,text=Classical,font=\bfseries\fontsize{9.4}{12}\selectfont] at ([xshift=.95cm,yshift=-10.66cm]current page.north west) {Richiami};
\node[anchor=north west,inner sep=0pt] at ([xshift=2.6cm,yshift=-10.64cm]current page.north west) {\usebox{\RecallBox}};
\draw[Rule,line width=.35pt] ([xshift=.95cm,yshift=-11.92cm]current page.north west)--([xshift=19.05cm,yshift=-11.92cm]current page.north west);
\node[anchor=north west,inner sep=0pt,text=Muted,text width=9.0cm,font=\fontsize{7.7}{9.3}\selectfont] at ([xshift=.95cm,yshift=-12.06cm]current page.north west) {08 · Ricampionare al livello corretto};
\node[anchor=north west,inner sep=0pt,text=Classical,text width=8.2cm,align=right,font=\fontsize{7.3}{9.1}\selectfont] at ([xshift=9.45cm,yshift=-12.06cm]current page.north west) {\href{https://openintrostat.github.io/ims/}{IMS} · CORSO};
\node[anchor=north east,inner sep=0pt,text=Muted,font=\fontsize{7.7}{9.3}\selectfont] at ([xshift=19.05cm,yshift=-12.06cm]current page.north west) {83/110};

\end{tikzpicture}
\note{Idea della slide: Confrontare risultati ottenuti con unità di ricampionamento diverse rende visibile quanto la dipendenza conti. Non scegliamo la procedura che produce l'intervallo più stretto. Scegliamo quella coerente con il disegno e discutiamo la fragilità quando i gruppi sono pochi.\par\medskip Punto da far emergere: La precisione apparente non è un criterio di selezione del metodo.\par\medskip Fonti della sezione: IMS, CORSO}
\end{frame}
\begin{lrbox}{\TitleBox}\begin{minipage}[t]{18.2cm}\color{Ink}\raggedright\sloppy\fontsize{20.7}{24}\selectfont
\bfseries\color{Navy}Attività 08 · applica il concetto
\end{minipage}\end{lrbox}
\begin{lrbox}{\BodyBox}\begin{minipage}[t]{10.25cm}\color{Ink}\raggedright\sloppy\fontsize{12.4}{16.4}\selectfont
Hai sei stazioni con quattro provini per stazione. Per descrivere l'incertezza tra stazioni, è corretto estrarre 24 righe indipendentemente? Proponi un'alternativa.
\end{minipage}\end{lrbox}
\begin{lrbox}{\SideBox}\begin{minipage}[t]{7.45cm}\color{Ink}\raggedright\sloppy\fontsize{12.2}{16.1}\selectfont
\noindent{\color{Classical}\bfseries 01}\hspace{.35cm}Individua la domanda\par\vspace{.22cm}\noindent{\color{Classical}\bfseries 02}\hspace{.35cm}Esplicita le ipotesi\par\vspace{.22cm}\noindent{\color{Classical}\bfseries 03}\hspace{.35cm}Calcola o rappresenta\par\vspace{.22cm}\noindent{\color{Classical}\bfseries 04}\hspace{.35cm}Scrivi una conclusione\par\vspace{.22cm}\noindent{\color{Classical}\bfseries 05}\hspace{.35cm}Indica un limite\par\vspace{.22cm}
\end{minipage}\end{lrbox}
\begin{lrbox}{\RecallBox}\begin{minipage}[t]{16.45cm}\color{Ink}\raggedright\sloppy\fontsize{10.1}{13.2}\selectfont
Quale unità deve essere estratta a ogni replica?
\end{minipage}\end{lrbox}
\MeasureBox{84-TitleBox}{\TitleBox}
\MeasureBox{84-BodyBox}{\BodyBox}
\MeasureBox{84-SideBox}{\SideBox}
\MeasureBox{84-RecallBox}{\RecallBox}
\begin{frame}[plain]
\hypertarget{slide84}{}\bookmark[level=1,page=84]{Attività 08 · applica il concetto}
\begin{tikzpicture}[remember picture,overlay]
\node[anchor=north east,inner sep=0pt,text=Classical,font=\fontsize{8}{10}\selectfont] at ([xshift=19.05cm,yshift=-.16cm]current page.north west) {\hyperlink{slide3}{Indice}};
\node[anchor=north west,inner sep=0pt,text=Navy,font=\bfseries] at ([xshift=.95cm,yshift=-.62cm]current page.north west) {\usebox{\TitleBox}};
\draw[Classical,line width=1.3pt] ([xshift=.95cm,yshift=-1.68cm]current page.north west)--([xshift=2.25cm,yshift=-1.68cm]current page.north west);
\node[anchor=north west,inner sep=0pt] at ([xshift=.95cm,yshift=-2.12cm]current page.north west) {\usebox{\BodyBox}};
\node[anchor=north west,inner sep=0pt,text=Quantum,font=\bfseries\fontsize{11.5}{14}\selectfont] at ([xshift=11.6cm,yshift=-2.1cm]current page.north west) {Consegna di lavoro};
\node[anchor=north west,inner sep=0pt] at ([xshift=11.6cm,yshift=-2.75cm]current page.north west) {\usebox{\SideBox}};
\draw[Rule,line width=.5pt] ([xshift=.95cm,yshift=-10.43cm]current page.north west)--([xshift=19.05cm,yshift=-10.43cm]current page.north west);
\node[anchor=north west,inner sep=0pt,text=Classical,font=\bfseries\fontsize{9.4}{12}\selectfont] at ([xshift=.95cm,yshift=-10.66cm]current page.north west) {Richiami};
\node[anchor=north west,inner sep=0pt] at ([xshift=2.6cm,yshift=-10.64cm]current page.north west) {\usebox{\RecallBox}};
\draw[Rule,line width=.35pt] ([xshift=.95cm,yshift=-11.92cm]current page.north west)--([xshift=19.05cm,yshift=-11.92cm]current page.north west);
\node[anchor=north west,inner sep=0pt,text=Muted,text width=9.0cm,font=\fontsize{7.7}{9.3}\selectfont] at ([xshift=.95cm,yshift=-12.06cm]current page.north west) {08 · Ricampionare al livello corretto};
\node[anchor=north west,inner sep=0pt,text=Classical,text width=8.2cm,align=right,font=\fontsize{7.3}{9.1}\selectfont] at ([xshift=9.45cm,yshift=-12.06cm]current page.north west) {\href{https://openintrostat.github.io/ims/}{IMS} · CORSO};
\node[anchor=north east,inner sep=0pt,text=Muted,font=\fontsize{7.7}{9.3}\selectfont] at ([xshift=19.05cm,yshift=-12.06cm]current page.north west) {84/110};

\end{tikzpicture}
\note{Consegna: Hai sei stazioni con quattro provini per stazione. Per descrivere l'incertezza tra stazioni, è corretto estrarre 24 righe indipendentemente? Proponi un'alternativa.\par\medskip Soluzione per il docente: Il ricampionamento delle righe ignora l'eventuale dipendenza interna alla stazione. Un'alternativa è ricampionare le stazioni, conservando insieme i provini di ciascuna stazione estratta. Con sole sei stazioni, l'incertezza resta difficile da stimare e va discussa esplicitamente.}
\end{frame}
\begin{lrbox}{\TitleBox}\begin{minipage}[t]{18.2cm}\color{Ink}\raggedright\sloppy\fontsize{20.7}{24}\selectfont
\bfseries\color{Navy}Discussione 08 · soluzione ragionata
\end{minipage}\end{lrbox}
\begin{lrbox}{\BodyBox}\begin{minipage}[t]{10.25cm}\color{Ink}\raggedright\sloppy\fontsize{12.4}{16.4}\selectfont
Il ricampionamento delle righe ignora l'eventuale dipendenza interna alla stazione. Un'alternativa è ricampionare le stazioni, conservando insieme i provini di ciascuna stazione estratta. Con sole sei stazioni, l'incertezza resta difficile da stimare e va discussa esplicitamente.
\end{minipage}\end{lrbox}
\begin{lrbox}{\SideBox}\begin{minipage}[t]{7.45cm}\color{Ink}\raggedright\sloppy\fontsize{12.2}{16.1}\selectfont
La precisione apparente non è un criterio di selezione del metodo.
\end{minipage}\end{lrbox}
\begin{lrbox}{\RecallBox}\begin{minipage}[t]{16.45cm}\color{Ink}\raggedright\sloppy\fontsize{10.1}{13.2}\selectfont
Valuta il collegamento tra domanda, procedimento e conclusione, non soltanto il numero finale.
\end{minipage}\end{lrbox}
\MeasureBox{85-TitleBox}{\TitleBox}
\MeasureBox{85-BodyBox}{\BodyBox}
\MeasureBox{85-SideBox}{\SideBox}
\MeasureBox{85-RecallBox}{\RecallBox}
\begin{frame}[plain]
\hypertarget{slide85}{}\bookmark[level=1,page=85]{Discussione 08 · soluzione ragionata}
\begin{tikzpicture}[remember picture,overlay]
\node[anchor=north east,inner sep=0pt,text=Classical,font=\fontsize{8}{10}\selectfont] at ([xshift=19.05cm,yshift=-.16cm]current page.north west) {\hyperlink{slide3}{Indice}};
\node[anchor=north west,inner sep=0pt,text=Navy,font=\bfseries] at ([xshift=.95cm,yshift=-.62cm]current page.north west) {\usebox{\TitleBox}};
\draw[Classical,line width=1.3pt] ([xshift=.95cm,yshift=-1.68cm]current page.north west)--([xshift=2.25cm,yshift=-1.68cm]current page.north west);
\node[anchor=north west,inner sep=0pt] at ([xshift=.95cm,yshift=-2.12cm]current page.north west) {\usebox{\BodyBox}};
\node[anchor=north west,inner sep=0pt,text=Quantum,font=\bfseries\fontsize{11.5}{14}\selectfont] at ([xshift=11.6cm,yshift=-2.1cm]current page.north west) {Criterio di lettura};
\node[anchor=north west,inner sep=0pt] at ([xshift=11.6cm,yshift=-2.75cm]current page.north west) {\usebox{\SideBox}};
\draw[Rule,line width=.5pt] ([xshift=.95cm,yshift=-10.43cm]current page.north west)--([xshift=19.05cm,yshift=-10.43cm]current page.north west);
\node[anchor=north west,inner sep=0pt,text=Classical,font=\bfseries\fontsize{9.4}{12}\selectfont] at ([xshift=.95cm,yshift=-10.66cm]current page.north west) {Richiami};
\node[anchor=north west,inner sep=0pt] at ([xshift=2.6cm,yshift=-10.64cm]current page.north west) {\usebox{\RecallBox}};
\draw[Rule,line width=.35pt] ([xshift=.95cm,yshift=-11.92cm]current page.north west)--([xshift=19.05cm,yshift=-11.92cm]current page.north west);
\node[anchor=north west,inner sep=0pt,text=Muted,text width=9.0cm,font=\fontsize{7.7}{9.3}\selectfont] at ([xshift=.95cm,yshift=-12.06cm]current page.north west) {08 · Ricampionare al livello corretto};
\node[anchor=north west,inner sep=0pt,text=Classical,text width=8.2cm,align=right,font=\fontsize{7.3}{9.1}\selectfont] at ([xshift=9.45cm,yshift=-12.06cm]current page.north west) {\href{https://openintrostat.github.io/ims/}{IMS} · CORSO};
\node[anchor=north east,inner sep=0pt,text=Muted,font=\fontsize{7.7}{9.3}\selectfont] at ([xshift=19.05cm,yshift=-12.06cm]current page.north west) {85/110};

\end{tikzpicture}
\note{Il ricampionamento delle righe ignora l'eventuale dipendenza interna alla stazione. Un'alternativa è ricampionare le stazioni, conservando insieme i provini di ciascuna stazione estratta. Con sole sei stazioni, l'incertezza resta difficile da stimare e va discussa esplicitamente.}
\end{frame}
\begin{frame}[plain]\hypertarget{slide86}{}
\bookmark[level=0,page=86]{Precisione e progetto di campionamento}
\begin{tikzpicture}[remember picture,overlay]
\fill[Night] (current page.south west) rectangle (current page.north east);
\draw[Mint,line width=2pt] ([xshift=1.05cm,yshift=-1.35cm]current page.north west)--([xshift=3.0cm,yshift=-1.35cm]current page.north west);
\node[anchor=north west,inner sep=0,text=Mint,font=\bfseries\fontsize{11}{14}\selectfont] at ([xshift=1.05cm,yshift=-1.7cm]current page.north west) {SEZIONE 09};
\node[anchor=north west,inner sep=0,text=white,text width=15.8cm,align=left,font=\bfseries\fontsize{31}{38}\selectfont] at ([xshift=1.05cm,yshift=-2.65cm]current page.north west) {Precisione e progetto di campionamento};
\node[anchor=north west,inner sep=0,text=Lilac,text width=15.8cm,align=left,font=\fontsize{15}{21}\selectfont] at ([xshift=1.05cm,yshift=-6.4cm]current page.north west) {Dove conviene investire nuove misure?};
\node[anchor=north west,inner sep=0,text=white,text width=16.5cm,align=left,font=\fontsize{11}{15}\selectfont] at ([xshift=1.05cm,yshift=-8.6cm]current page.north west) {Comprendere il concetto · leggere un esempio · applicare e discutere};
\node[anchor=north east,inner sep=0,text=Mint,font=\fontsize{8}{10}\selectfont] at ([xshift=19.05cm,yshift=-.2cm]current page.north west) {\hyperlink{slide3}{Indice}};
\node[anchor=south east,inner sep=0,text=Lilac!45,font=\bfseries\fontsize{64}{68}\selectfont] at ([xshift=18.95cm,yshift=-11.35cm]current page.north west) {09};
\draw[Lilac!30,line width=.35pt] ([xshift=.95cm,yshift=-11.92cm]current page.north west)--([xshift=19.05cm,yshift=-11.92cm]current page.north west);
\node[anchor=north west,inner sep=0pt,text=Lilac,text width=9.0cm,font=\fontsize{7.7}{9.3}\selectfont] at ([xshift=.95cm,yshift=-12.06cm]current page.north west) {09 · Precisione e progetto di campionamento};
\node[anchor=north west,inner sep=0pt,text=Mint,text width=8.2cm,align=right,font=\fontsize{7.3}{9.1}\selectfont] at ([xshift=9.45cm,yshift=-12.06cm]current page.north west) {\href{https://openintrostat.github.io/ims/}{IMS} · CORSO};
\node[anchor=north east,inner sep=0pt,text=Lilac,font=\fontsize{7.7}{9.3}\selectfont] at ([xshift=19.05cm,yshift=-12.06cm]current page.north west) {86/110};
\end{tikzpicture}\end{frame}
\begin{lrbox}{\TitleBox}\begin{minipage}[t]{18.2cm}\color{Ink}\raggedright\sloppy\fontsize{20.7}{24}\selectfont
\bfseries\color{Navy}Scegliere una precisione utile
\end{minipage}\end{lrbox}
\begin{lrbox}{\BodyBox}\begin{minipage}[t]{10.25cm}\color{Ink}\raggedright\sloppy\fontsize{12.4}{16.4}\selectfont
Prima di raccogliere dati chiediamo quale ampiezza dell'incertezza sarebbe scientificamente informativa. Una stima più precisa di quanto richieda la domanda può essere costosa senza beneficio sostanziale. Al contrario, un intervallo molto ampio può non distinguere scenari importanti.
\end{minipage}\end{lrbox}
\begin{lrbox}{\SideBox}\begin{minipage}[t]{7.45cm}\color{Ink}\raggedright\sloppy\fontsize{12.2}{16.1}\selectfont
Un obiettivo possibile: stimare una media entro un margine giudicato utile rispetto alla variabilità naturale e al protocollo.
\end{minipage}\end{lrbox}
\begin{lrbox}{\RecallBox}\begin{minipage}[t]{16.45cm}\color{Ink}\raggedright\sloppy\fontsize{10.1}{13.2}\selectfont
La precisione desiderata deve essere espressa nelle unità della domanda.
\end{minipage}\end{lrbox}
\MeasureBox{87-TitleBox}{\TitleBox}
\MeasureBox{87-BodyBox}{\BodyBox}
\MeasureBox{87-SideBox}{\SideBox}
\MeasureBox{87-RecallBox}{\RecallBox}
\begin{frame}[plain]
\hypertarget{slide87}{}\bookmark[level=1,page=87]{Scegliere una precisione utile}
\begin{tikzpicture}[remember picture,overlay]
\node[anchor=north east,inner sep=0pt,text=Classical,font=\fontsize{8}{10}\selectfont] at ([xshift=19.05cm,yshift=-.16cm]current page.north west) {\hyperlink{slide3}{Indice}};
\node[anchor=north west,inner sep=0pt,text=Navy,font=\bfseries] at ([xshift=.95cm,yshift=-.62cm]current page.north west) {\usebox{\TitleBox}};
\draw[Classical,line width=1.3pt] ([xshift=.95cm,yshift=-1.68cm]current page.north west)--([xshift=2.25cm,yshift=-1.68cm]current page.north west);
\node[anchor=north west,inner sep=0pt] at ([xshift=.95cm,yshift=-2.12cm]current page.north west) {\usebox{\BodyBox}};
\node[anchor=north west,inner sep=0pt,text=Quantum,font=\bfseries\fontsize{11.5}{14}\selectfont] at ([xshift=11.6cm,yshift=-2.1cm]current page.north west) {Esempio e interpretazione};
\node[anchor=north west,inner sep=0pt] at ([xshift=11.6cm,yshift=-2.75cm]current page.north west) {\usebox{\SideBox}};
\draw[Rule,line width=.5pt] ([xshift=.95cm,yshift=-10.43cm]current page.north west)--([xshift=19.05cm,yshift=-10.43cm]current page.north west);
\node[anchor=north west,inner sep=0pt,text=Classical,font=\bfseries\fontsize{9.4}{12}\selectfont] at ([xshift=.95cm,yshift=-10.66cm]current page.north west) {Richiami};
\node[anchor=north west,inner sep=0pt] at ([xshift=2.6cm,yshift=-10.64cm]current page.north west) {\usebox{\RecallBox}};
\draw[Rule,line width=.35pt] ([xshift=.95cm,yshift=-11.92cm]current page.north west)--([xshift=19.05cm,yshift=-11.92cm]current page.north west);
\node[anchor=north west,inner sep=0pt,text=Muted,text width=9.0cm,font=\fontsize{7.7}{9.3}\selectfont] at ([xshift=.95cm,yshift=-12.06cm]current page.north west) {09 · Precisione e progetto di campionamento};
\node[anchor=north west,inner sep=0pt,text=Classical,text width=8.2cm,align=right,font=\fontsize{7.3}{9.1}\selectfont] at ([xshift=9.45cm,yshift=-12.06cm]current page.north west) {\href{https://openintrostat.github.io/ims/}{IMS} · CORSO};
\node[anchor=north east,inner sep=0pt,text=Muted,font=\fontsize{7.7}{9.3}\selectfont] at ([xshift=19.05cm,yshift=-12.06cm]current page.north west) {87/110};

\end{tikzpicture}
\note{Idea della slide: Prima di raccogliere dati chiediamo quale ampiezza dell'incertezza sarebbe scientificamente informativa. Una stima più precisa di quanto richieda la domanda può essere costosa senza beneficio sostanziale. Al contrario, un intervallo molto ampio può non distinguere scenari importanti.\par\medskip Punto da far emergere: La precisione desiderata deve essere espressa nelle unità della domanda.\par\medskip Fonti della sezione: IMS, CORSO}
\end{frame}
\begin{lrbox}{\TitleBox}\begin{minipage}[t]{18.2cm}\color{Ink}\raggedright\sloppy\fontsize{20.7}{24}\selectfont
\bfseries\color{Navy}Una formula orientativa per n
\end{minipage}\end{lrbox}
\begin{lrbox}{\BodyBox}\begin{minipage}[t]{10.25cm}\color{Ink}\raggedright\sloppy\fontsize{12.4}{16.4}\selectfont
Nel modello semplice con deviazione standard circa nota, il margine dell'intervallo della media è proporzionale a \(\sigma\)/\(\sqrt{}\)n. Risolvere la relazione per n dà un ordine di grandezza. È una pianificazione preliminare, non una garanzia quando disegno e distribuzione sono più complessi.
\end{minipage}\end{lrbox}
\begin{lrbox}{\SideBox}\begin{minipage}[t]{7.45cm}\color{Ink}\raggedright\sloppy\fontsize{12.2}{16.1}\selectfont
$n\approx(z\,\sigma/E)^2$\par\medskip E è il margine desiderato e z il quantile scelto.
\end{minipage}\end{lrbox}
\begin{lrbox}{\RecallBox}\begin{minipage}[t]{16.45cm}\color{Ink}\raggedright\sloppy\fontsize{10.1}{13.2}\selectfont
Non applicare la formula a dati dipendenti fingendo che ogni misura aggiunga la stessa informazione.
\end{minipage}\end{lrbox}
\MeasureBox{88-TitleBox}{\TitleBox}
\MeasureBox{88-BodyBox}{\BodyBox}
\MeasureBox{88-SideBox}{\SideBox}
\MeasureBox{88-RecallBox}{\RecallBox}
\begin{frame}[plain]
\hypertarget{slide88}{}\bookmark[level=1,page=88]{Una formula orientativa per n}
\begin{tikzpicture}[remember picture,overlay]
\node[anchor=north east,inner sep=0pt,text=Classical,font=\fontsize{8}{10}\selectfont] at ([xshift=19.05cm,yshift=-.16cm]current page.north west) {\hyperlink{slide3}{Indice}};
\node[anchor=north west,inner sep=0pt,text=Navy,font=\bfseries] at ([xshift=.95cm,yshift=-.62cm]current page.north west) {\usebox{\TitleBox}};
\draw[Classical,line width=1.3pt] ([xshift=.95cm,yshift=-1.68cm]current page.north west)--([xshift=2.25cm,yshift=-1.68cm]current page.north west);
\node[anchor=north west,inner sep=0pt] at ([xshift=.95cm,yshift=-2.12cm]current page.north west) {\usebox{\BodyBox}};
\node[anchor=north west,inner sep=0pt,text=Quantum,font=\bfseries\fontsize{11.5}{14}\selectfont] at ([xshift=11.6cm,yshift=-2.1cm]current page.north west) {Esempio e interpretazione};
\node[anchor=north west,inner sep=0pt] at ([xshift=11.6cm,yshift=-2.75cm]current page.north west) {\usebox{\SideBox}};
\draw[Rule,line width=.5pt] ([xshift=.95cm,yshift=-10.43cm]current page.north west)--([xshift=19.05cm,yshift=-10.43cm]current page.north west);
\node[anchor=north west,inner sep=0pt,text=Classical,font=\bfseries\fontsize{9.4}{12}\selectfont] at ([xshift=.95cm,yshift=-10.66cm]current page.north west) {Richiami};
\node[anchor=north west,inner sep=0pt] at ([xshift=2.6cm,yshift=-10.64cm]current page.north west) {\usebox{\RecallBox}};
\draw[Rule,line width=.35pt] ([xshift=.95cm,yshift=-11.92cm]current page.north west)--([xshift=19.05cm,yshift=-11.92cm]current page.north west);
\node[anchor=north west,inner sep=0pt,text=Muted,text width=9.0cm,font=\fontsize{7.7}{9.3}\selectfont] at ([xshift=.95cm,yshift=-12.06cm]current page.north west) {09 · Precisione e progetto di campionamento};
\node[anchor=north west,inner sep=0pt,text=Classical,text width=8.2cm,align=right,font=\fontsize{7.3}{9.1}\selectfont] at ([xshift=9.45cm,yshift=-12.06cm]current page.north west) {\href{https://openintrostat.github.io/ims/}{IMS} · CORSO};
\node[anchor=north east,inner sep=0pt,text=Muted,font=\fontsize{7.7}{9.3}\selectfont] at ([xshift=19.05cm,yshift=-12.06cm]current page.north west) {88/110};

\end{tikzpicture}
\note{Idea della slide: Nel modello semplice con deviazione standard circa nota, il margine dell'intervallo della media è proporzionale a \(\sigma\)/\(\sqrt{}\)n. Risolvere la relazione per n dà un ordine di grandezza. È una pianificazione preliminare, non una garanzia quando disegno e distribuzione sono più complessi.\par\medskip Punto da far emergere: Non applicare la formula a dati dipendenti fingendo che ogni misura aggiunga la stessa informazione.\par\medskip Fonti della sezione: IMS, CORSO}
\end{frame}
\begin{lrbox}{\TitleBox}\begin{minipage}[t]{18.2cm}\color{Ink}\raggedright\sloppy\fontsize{20.7}{24}\selectfont
\bfseries\color{Navy}Una stima pilota della variabilità
\end{minipage}\end{lrbox}
\begin{lrbox}{\BodyBox}\begin{minipage}[t]{10.25cm}\color{Ink}\raggedright\sloppy\fontsize{12.4}{16.4}\selectfont
Un piccolo studio iniziale può aiutare a stimare dispersione, problemi di misura e costi. La deviazione standard pilota è però incerta. Conviene valutare più valori plausibili invece di progettare tutto su una singola cifra apparentemente esatta.
\end{minipage}\end{lrbox}
\begin{lrbox}{\SideBox}\begin{minipage}[t]{7.45cm}\color{Ink}\raggedright\sloppy\fontsize{12.2}{16.1}\selectfont
Calcolare n per una dispersione bassa, centrale e alta aiuta a rendere esplicita la sensibilità del piano.
\end{minipage}\end{lrbox}
\begin{lrbox}{\RecallBox}\begin{minipage}[t]{16.45cm}\color{Ink}\raggedright\sloppy\fontsize{10.1}{13.2}\selectfont
Il pilota serve anche a verificare protocollo e base di campionamento, non soltanto s.
\end{minipage}\end{lrbox}
\MeasureBox{89-TitleBox}{\TitleBox}
\MeasureBox{89-BodyBox}{\BodyBox}
\MeasureBox{89-SideBox}{\SideBox}
\MeasureBox{89-RecallBox}{\RecallBox}
\begin{frame}[plain]
\hypertarget{slide89}{}\bookmark[level=1,page=89]{Una stima pilota della variabilità}
\begin{tikzpicture}[remember picture,overlay]
\node[anchor=north east,inner sep=0pt,text=Classical,font=\fontsize{8}{10}\selectfont] at ([xshift=19.05cm,yshift=-.16cm]current page.north west) {\hyperlink{slide3}{Indice}};
\node[anchor=north west,inner sep=0pt,text=Navy,font=\bfseries] at ([xshift=.95cm,yshift=-.62cm]current page.north west) {\usebox{\TitleBox}};
\draw[Classical,line width=1.3pt] ([xshift=.95cm,yshift=-1.68cm]current page.north west)--([xshift=2.25cm,yshift=-1.68cm]current page.north west);
\node[anchor=north west,inner sep=0pt] at ([xshift=.95cm,yshift=-2.12cm]current page.north west) {\usebox{\BodyBox}};
\node[anchor=north west,inner sep=0pt,text=Quantum,font=\bfseries\fontsize{11.5}{14}\selectfont] at ([xshift=11.6cm,yshift=-2.1cm]current page.north west) {Esempio e interpretazione};
\node[anchor=north west,inner sep=0pt] at ([xshift=11.6cm,yshift=-2.75cm]current page.north west) {\usebox{\SideBox}};
\draw[Rule,line width=.5pt] ([xshift=.95cm,yshift=-10.43cm]current page.north west)--([xshift=19.05cm,yshift=-10.43cm]current page.north west);
\node[anchor=north west,inner sep=0pt,text=Classical,font=\bfseries\fontsize{9.4}{12}\selectfont] at ([xshift=.95cm,yshift=-10.66cm]current page.north west) {Richiami};
\node[anchor=north west,inner sep=0pt] at ([xshift=2.6cm,yshift=-10.64cm]current page.north west) {\usebox{\RecallBox}};
\draw[Rule,line width=.35pt] ([xshift=.95cm,yshift=-11.92cm]current page.north west)--([xshift=19.05cm,yshift=-11.92cm]current page.north west);
\node[anchor=north west,inner sep=0pt,text=Muted,text width=9.0cm,font=\fontsize{7.7}{9.3}\selectfont] at ([xshift=.95cm,yshift=-12.06cm]current page.north west) {09 · Precisione e progetto di campionamento};
\node[anchor=north west,inner sep=0pt,text=Classical,text width=8.2cm,align=right,font=\fontsize{7.3}{9.1}\selectfont] at ([xshift=9.45cm,yshift=-12.06cm]current page.north west) {\href{https://openintrostat.github.io/ims/}{IMS} · CORSO};
\node[anchor=north east,inner sep=0pt,text=Muted,font=\fontsize{7.7}{9.3}\selectfont] at ([xshift=19.05cm,yshift=-12.06cm]current page.north west) {89/110};

\end{tikzpicture}
\note{Idea della slide: Un piccolo studio iniziale può aiutare a stimare dispersione, problemi di misura e costi. La deviazione standard pilota è però incerta. Conviene valutare più valori plausibili invece di progettare tutto su una singola cifra apparentemente esatta.\par\medskip Punto da far emergere: Il pilota serve anche a verificare protocollo e base di campionamento, non soltanto s.\par\medskip Fonti della sezione: IMS, CORSO}
\end{frame}
\begin{lrbox}{\TitleBox}\begin{minipage}[t]{18.2cm}\color{Ink}\raggedright\sloppy\fontsize{20.7}{24}\selectfont
\bfseries\color{Navy}Pesi e obiettivo di aggregazione
\end{minipage}\end{lrbox}
\begin{lrbox}{\BodyBox}\begin{minipage}[t]{10.25cm}\color{Ink}\raggedright\sloppy\fontsize{12.4}{16.4}\selectfont
Se la domanda riguarda una media areale, i pesi possono dipendere dalla rappresentatività delle aree. Se riguarda la media tra stazioni campionate, l'obiettivo è diverso. I pesi non sono dettagli da aggiungere dopo aver visto il risultato.
\end{minipage}\end{lrbox}
\begin{lrbox}{\SideBox}\begin{minipage}[t]{7.45cm}\color{Ink}\raggedright\sloppy\fontsize{12.2}{16.1}\selectfont
Uguale peso a ogni provino; uguale peso a ogni stazione; peso proporzionale all'area. Sono tre possibili parametri differenti.
\end{minipage}\end{lrbox}
\begin{lrbox}{\RecallBox}\begin{minipage}[t]{16.45cm}\color{Ink}\raggedright\sloppy\fontsize{10.1}{13.2}\selectfont
La formula della media deve essere coerente con la popolazione definita.
\end{minipage}\end{lrbox}
\MeasureBox{90-TitleBox}{\TitleBox}
\MeasureBox{90-BodyBox}{\BodyBox}
\MeasureBox{90-SideBox}{\SideBox}
\MeasureBox{90-RecallBox}{\RecallBox}
\begin{frame}[plain]
\hypertarget{slide90}{}\bookmark[level=1,page=90]{Pesi e obiettivo di aggregazione}
\begin{tikzpicture}[remember picture,overlay]
\node[anchor=north east,inner sep=0pt,text=Classical,font=\fontsize{8}{10}\selectfont] at ([xshift=19.05cm,yshift=-.16cm]current page.north west) {\hyperlink{slide3}{Indice}};
\node[anchor=north west,inner sep=0pt,text=Navy,font=\bfseries] at ([xshift=.95cm,yshift=-.62cm]current page.north west) {\usebox{\TitleBox}};
\draw[Classical,line width=1.3pt] ([xshift=.95cm,yshift=-1.68cm]current page.north west)--([xshift=2.25cm,yshift=-1.68cm]current page.north west);
\node[anchor=north west,inner sep=0pt] at ([xshift=.95cm,yshift=-2.12cm]current page.north west) {\usebox{\BodyBox}};
\node[anchor=north west,inner sep=0pt,text=Quantum,font=\bfseries\fontsize{11.5}{14}\selectfont] at ([xshift=11.6cm,yshift=-2.1cm]current page.north west) {Esempio e interpretazione};
\node[anchor=north west,inner sep=0pt] at ([xshift=11.6cm,yshift=-2.75cm]current page.north west) {\usebox{\SideBox}};
\draw[Rule,line width=.5pt] ([xshift=.95cm,yshift=-10.43cm]current page.north west)--([xshift=19.05cm,yshift=-10.43cm]current page.north west);
\node[anchor=north west,inner sep=0pt,text=Classical,font=\bfseries\fontsize{9.4}{12}\selectfont] at ([xshift=.95cm,yshift=-10.66cm]current page.north west) {Richiami};
\node[anchor=north west,inner sep=0pt] at ([xshift=2.6cm,yshift=-10.64cm]current page.north west) {\usebox{\RecallBox}};
\draw[Rule,line width=.35pt] ([xshift=.95cm,yshift=-11.92cm]current page.north west)--([xshift=19.05cm,yshift=-11.92cm]current page.north west);
\node[anchor=north west,inner sep=0pt,text=Muted,text width=9.0cm,font=\fontsize{7.7}{9.3}\selectfont] at ([xshift=.95cm,yshift=-12.06cm]current page.north west) {09 · Precisione e progetto di campionamento};
\node[anchor=north west,inner sep=0pt,text=Classical,text width=8.2cm,align=right,font=\fontsize{7.3}{9.1}\selectfont] at ([xshift=9.45cm,yshift=-12.06cm]current page.north west) {\href{https://openintrostat.github.io/ims/}{IMS} · CORSO};
\node[anchor=north east,inner sep=0pt,text=Muted,font=\fontsize{7.7}{9.3}\selectfont] at ([xshift=19.05cm,yshift=-12.06cm]current page.north west) {90/110};

\end{tikzpicture}
\note{Idea della slide: Se la domanda riguarda una media areale, i pesi possono dipendere dalla rappresentatività delle aree. Se riguarda la media tra stazioni campionate, l'obiettivo è diverso. I pesi non sono dettagli da aggiungere dopo aver visto il risultato.\par\medskip Punto da far emergere: La formula della media deve essere coerente con la popolazione definita.\par\medskip Fonti della sezione: IMS, CORSO}
\end{frame}
\begin{lrbox}{\TitleBox}\begin{minipage}[t]{18.2cm}\color{Ink}\raggedright\sloppy\fontsize{20.7}{24}\selectfont
\bfseries\color{Navy}Bias contro varianza
\end{minipage}\end{lrbox}
\begin{lrbox}{\BodyBox}\begin{minipage}[t]{10.25cm}\color{Ink}\raggedright\sloppy\fontsize{12.4}{16.4}\selectfont
Nuove misure possono ridurre la variabilità campionaria; migliore calibrazione o migliore copertura possono ridurre bias. Le due strategie affrontano errori diversi. Un progetto efficace cerca il limite principale invece di aumentare automaticamente il numero di osservazioni.
\end{minipage}\end{lrbox}
\begin{lrbox}{\SideBox}\begin{minipage}[t]{7.45cm}\color{Ink}\raggedright\sloppy\fontsize{12.2}{16.1}\selectfont
Se il volume è sistematicamente sottostimato, ripetere cento volte la stessa procedura non risolve il problema di accuratezza.
\end{minipage}\end{lrbox}
\begin{lrbox}{\RecallBox}\begin{minipage}[t]{16.45cm}\color{Ink}\raggedright\sloppy\fontsize{10.1}{13.2}\selectfont
Più repliche dello stesso errore non costituiscono una correzione.
\end{minipage}\end{lrbox}
\MeasureBox{91-TitleBox}{\TitleBox}
\MeasureBox{91-BodyBox}{\BodyBox}
\MeasureBox{91-SideBox}{\SideBox}
\MeasureBox{91-RecallBox}{\RecallBox}
\begin{frame}[plain]
\hypertarget{slide91}{}\bookmark[level=1,page=91]{Bias contro varianza}
\begin{tikzpicture}[remember picture,overlay]
\node[anchor=north east,inner sep=0pt,text=Classical,font=\fontsize{8}{10}\selectfont] at ([xshift=19.05cm,yshift=-.16cm]current page.north west) {\hyperlink{slide3}{Indice}};
\node[anchor=north west,inner sep=0pt,text=Navy,font=\bfseries] at ([xshift=.95cm,yshift=-.62cm]current page.north west) {\usebox{\TitleBox}};
\draw[Classical,line width=1.3pt] ([xshift=.95cm,yshift=-1.68cm]current page.north west)--([xshift=2.25cm,yshift=-1.68cm]current page.north west);
\node[anchor=north west,inner sep=0pt] at ([xshift=.95cm,yshift=-2.12cm]current page.north west) {\usebox{\BodyBox}};
\node[anchor=north west,inner sep=0pt,text=Quantum,font=\bfseries\fontsize{11.5}{14}\selectfont] at ([xshift=11.6cm,yshift=-2.1cm]current page.north west) {Esempio e interpretazione};
\node[anchor=north west,inner sep=0pt] at ([xshift=11.6cm,yshift=-2.75cm]current page.north west) {\usebox{\SideBox}};
\draw[Rule,line width=.5pt] ([xshift=.95cm,yshift=-10.43cm]current page.north west)--([xshift=19.05cm,yshift=-10.43cm]current page.north west);
\node[anchor=north west,inner sep=0pt,text=Classical,font=\bfseries\fontsize{9.4}{12}\selectfont] at ([xshift=.95cm,yshift=-10.66cm]current page.north west) {Richiami};
\node[anchor=north west,inner sep=0pt] at ([xshift=2.6cm,yshift=-10.64cm]current page.north west) {\usebox{\RecallBox}};
\draw[Rule,line width=.35pt] ([xshift=.95cm,yshift=-11.92cm]current page.north west)--([xshift=19.05cm,yshift=-11.92cm]current page.north west);
\node[anchor=north west,inner sep=0pt,text=Muted,text width=9.0cm,font=\fontsize{7.7}{9.3}\selectfont] at ([xshift=.95cm,yshift=-12.06cm]current page.north west) {09 · Precisione e progetto di campionamento};
\node[anchor=north west,inner sep=0pt,text=Classical,text width=8.2cm,align=right,font=\fontsize{7.3}{9.1}\selectfont] at ([xshift=9.45cm,yshift=-12.06cm]current page.north west) {\href{https://openintrostat.github.io/ims/}{IMS} · CORSO};
\node[anchor=north east,inner sep=0pt,text=Muted,font=\fontsize{7.7}{9.3}\selectfont] at ([xshift=19.05cm,yshift=-12.06cm]current page.north west) {91/110};

\end{tikzpicture}
\note{Idea della slide: Nuove misure possono ridurre la variabilità campionaria; migliore calibrazione o migliore copertura possono ridurre bias. Le due strategie affrontano errori diversi. Un progetto efficace cerca il limite principale invece di aumentare automaticamente il numero di osservazioni.\par\medskip Punto da far emergere: Più repliche dello stesso errore non costituiscono una correzione.\par\medskip Fonti della sezione: IMS, CORSO}
\end{frame}
\begin{lrbox}{\TitleBox}\begin{minipage}[t]{18.2cm}\color{Ink}\raggedright\sloppy\fontsize{20.7}{24}\selectfont
\bfseries\color{Navy}Incertezza del protocollo
\end{minipage}\end{lrbox}
\begin{lrbox}{\BodyBox}\begin{minipage}[t]{10.25cm}\color{Ink}\raggedright\sloppy\fontsize{12.4}{16.4}\selectfont
Differenze tra operatori, strumenti e condizioni possono contribuire alla variabilità osservata. Un disegno che registra questi fattori permette di distinguerli almeno in parte. Senza metadati, la dispersione finale può mescolare processi naturali e componenti di misura.
\end{minipage}\end{lrbox}
\begin{lrbox}{\SideBox}\begin{minipage}[t]{7.45cm}\color{Ink}\raggedright\sloppy\fontsize{12.2}{16.1}\selectfont
Registrare operatore, strumento, data, calibrazione e condizioni del campione quando rilevanti.
\end{minipage}\end{lrbox}
\begin{lrbox}{\RecallBox}\begin{minipage}[t]{16.45cm}\color{Ink}\raggedright\sloppy\fontsize{10.1}{13.2}\selectfont
La qualità dell'incertezza stimata dipende dalla qualità delle informazioni sul processo di misura.
\end{minipage}\end{lrbox}
\MeasureBox{92-TitleBox}{\TitleBox}
\MeasureBox{92-BodyBox}{\BodyBox}
\MeasureBox{92-SideBox}{\SideBox}
\MeasureBox{92-RecallBox}{\RecallBox}
\begin{frame}[plain]
\hypertarget{slide92}{}\bookmark[level=1,page=92]{Incertezza del protocollo}
\begin{tikzpicture}[remember picture,overlay]
\node[anchor=north east,inner sep=0pt,text=Classical,font=\fontsize{8}{10}\selectfont] at ([xshift=19.05cm,yshift=-.16cm]current page.north west) {\hyperlink{slide3}{Indice}};
\node[anchor=north west,inner sep=0pt,text=Navy,font=\bfseries] at ([xshift=.95cm,yshift=-.62cm]current page.north west) {\usebox{\TitleBox}};
\draw[Classical,line width=1.3pt] ([xshift=.95cm,yshift=-1.68cm]current page.north west)--([xshift=2.25cm,yshift=-1.68cm]current page.north west);
\node[anchor=north west,inner sep=0pt] at ([xshift=.95cm,yshift=-2.12cm]current page.north west) {\usebox{\BodyBox}};
\node[anchor=north west,inner sep=0pt,text=Quantum,font=\bfseries\fontsize{11.5}{14}\selectfont] at ([xshift=11.6cm,yshift=-2.1cm]current page.north west) {Esempio e interpretazione};
\node[anchor=north west,inner sep=0pt] at ([xshift=11.6cm,yshift=-2.75cm]current page.north west) {\usebox{\SideBox}};
\draw[Rule,line width=.5pt] ([xshift=.95cm,yshift=-10.43cm]current page.north west)--([xshift=19.05cm,yshift=-10.43cm]current page.north west);
\node[anchor=north west,inner sep=0pt,text=Classical,font=\bfseries\fontsize{9.4}{12}\selectfont] at ([xshift=.95cm,yshift=-10.66cm]current page.north west) {Richiami};
\node[anchor=north west,inner sep=0pt] at ([xshift=2.6cm,yshift=-10.64cm]current page.north west) {\usebox{\RecallBox}};
\draw[Rule,line width=.35pt] ([xshift=.95cm,yshift=-11.92cm]current page.north west)--([xshift=19.05cm,yshift=-11.92cm]current page.north west);
\node[anchor=north west,inner sep=0pt,text=Muted,text width=9.0cm,font=\fontsize{7.7}{9.3}\selectfont] at ([xshift=.95cm,yshift=-12.06cm]current page.north west) {09 · Precisione e progetto di campionamento};
\node[anchor=north west,inner sep=0pt,text=Classical,text width=8.2cm,align=right,font=\fontsize{7.3}{9.1}\selectfont] at ([xshift=9.45cm,yshift=-12.06cm]current page.north west) {\href{https://openintrostat.github.io/ims/}{IMS} · CORSO};
\node[anchor=north east,inner sep=0pt,text=Muted,font=\fontsize{7.7}{9.3}\selectfont] at ([xshift=19.05cm,yshift=-12.06cm]current page.north west) {92/110};

\end{tikzpicture}
\note{Idea della slide: Differenze tra operatori, strumenti e condizioni possono contribuire alla variabilità osservata. Un disegno che registra questi fattori permette di distinguerli almeno in parte. Senza metadati, la dispersione finale può mescolare processi naturali e componenti di misura.\par\medskip Punto da far emergere: La qualità dell'incertezza stimata dipende dalla qualità delle informazioni sul processo di misura.\par\medskip Fonti della sezione: IMS, CORSO}
\end{frame}
\begin{lrbox}{\TitleBox}\begin{minipage}[t]{18.2cm}\color{Ink}\raggedright\sloppy\fontsize{20.7}{24}\selectfont
\bfseries\color{Navy}Pianificare e poi documentare le deviazioni
\end{minipage}\end{lrbox}
\begin{lrbox}{\BodyBox}\begin{minipage}[t]{10.25cm}\color{Ink}\raggedright\sloppy\fontsize{12.4}{16.4}\selectfont
Il piano può cambiare per difficoltà di campo. Le deviazioni vanno registrate, non nascoste. Un protocollo ideale rimasto sulla carta non descrive il campione effettivamente ottenuto; la relazione deve rendere visibile la differenza.
\end{minipage}\end{lrbox}
\begin{lrbox}{\SideBox}\begin{minipage}[t]{7.45cm}\color{Ink}\raggedright\sloppy\fontsize{12.2}{16.1}\selectfont
Pianificati: 10 stazioni distribuite.\par\smallskip Misurati: 7 stazioni accessibili.\par\smallskip Il dominio di inferenza deve essere rivalutato.
\end{minipage}\end{lrbox}
\begin{lrbox}{\RecallBox}\begin{minipage}[t]{16.45cm}\color{Ink}\raggedright\sloppy\fontsize{10.1}{13.2}\selectfont
La validità si riferisce al disegno eseguito, non soltanto a quello previsto.
\end{minipage}\end{lrbox}
\MeasureBox{93-TitleBox}{\TitleBox}
\MeasureBox{93-BodyBox}{\BodyBox}
\MeasureBox{93-SideBox}{\SideBox}
\MeasureBox{93-RecallBox}{\RecallBox}
\begin{frame}[plain]
\hypertarget{slide93}{}\bookmark[level=1,page=93]{Pianificare e poi documentare le deviazioni}
\begin{tikzpicture}[remember picture,overlay]
\node[anchor=north east,inner sep=0pt,text=Classical,font=\fontsize{8}{10}\selectfont] at ([xshift=19.05cm,yshift=-.16cm]current page.north west) {\hyperlink{slide3}{Indice}};
\node[anchor=north west,inner sep=0pt,text=Navy,font=\bfseries] at ([xshift=.95cm,yshift=-.62cm]current page.north west) {\usebox{\TitleBox}};
\draw[Classical,line width=1.3pt] ([xshift=.95cm,yshift=-1.68cm]current page.north west)--([xshift=2.25cm,yshift=-1.68cm]current page.north west);
\node[anchor=north west,inner sep=0pt] at ([xshift=.95cm,yshift=-2.12cm]current page.north west) {\usebox{\BodyBox}};
\node[anchor=north west,inner sep=0pt,text=Quantum,font=\bfseries\fontsize{11.5}{14}\selectfont] at ([xshift=11.6cm,yshift=-2.1cm]current page.north west) {Esempio e interpretazione};
\node[anchor=north west,inner sep=0pt] at ([xshift=11.6cm,yshift=-2.75cm]current page.north west) {\usebox{\SideBox}};
\draw[Rule,line width=.5pt] ([xshift=.95cm,yshift=-10.43cm]current page.north west)--([xshift=19.05cm,yshift=-10.43cm]current page.north west);
\node[anchor=north west,inner sep=0pt,text=Classical,font=\bfseries\fontsize{9.4}{12}\selectfont] at ([xshift=.95cm,yshift=-10.66cm]current page.north west) {Richiami};
\node[anchor=north west,inner sep=0pt] at ([xshift=2.6cm,yshift=-10.64cm]current page.north west) {\usebox{\RecallBox}};
\draw[Rule,line width=.35pt] ([xshift=.95cm,yshift=-11.92cm]current page.north west)--([xshift=19.05cm,yshift=-11.92cm]current page.north west);
\node[anchor=north west,inner sep=0pt,text=Muted,text width=9.0cm,font=\fontsize{7.7}{9.3}\selectfont] at ([xshift=.95cm,yshift=-12.06cm]current page.north west) {09 · Precisione e progetto di campionamento};
\node[anchor=north west,inner sep=0pt,text=Classical,text width=8.2cm,align=right,font=\fontsize{7.3}{9.1}\selectfont] at ([xshift=9.45cm,yshift=-12.06cm]current page.north west) {\href{https://openintrostat.github.io/ims/}{IMS} · CORSO};
\node[anchor=north east,inner sep=0pt,text=Muted,font=\fontsize{7.7}{9.3}\selectfont] at ([xshift=19.05cm,yshift=-12.06cm]current page.north west) {93/110};

\end{tikzpicture}
\note{Idea della slide: Il piano può cambiare per difficoltà di campo. Le deviazioni vanno registrate, non nascoste. Un protocollo ideale rimasto sulla carta non descrive il campione effettivamente ottenuto; la relazione deve rendere visibile la differenza.\par\medskip Punto da far emergere: La validità si riferisce al disegno eseguito, non soltanto a quello previsto.\par\medskip Fonti della sezione: IMS, CORSO}
\end{frame}
\begin{lrbox}{\TitleBox}\begin{minipage}[t]{18.2cm}\color{Ink}\raggedright\sloppy\fontsize{20.7}{24}\selectfont
\bfseries\color{Navy}Attività 09 · applica il concetto
\end{minipage}\end{lrbox}
\begin{lrbox}{\BodyBox}\begin{minipage}[t]{10.25cm}\color{Ink}\raggedright\sloppy\fontsize{12.4}{16.4}\selectfont
Disponi di risorse per 24 misure aggiuntive: tutte nella stessa stazione oppure distribuite in nuove stazioni. Da che cosa dipende la scelta?
\end{minipage}\end{lrbox}
\begin{lrbox}{\SideBox}\begin{minipage}[t]{7.45cm}\color{Ink}\raggedright\sloppy\fontsize{12.2}{16.1}\selectfont
\noindent{\color{Classical}\bfseries 01}\hspace{.35cm}Individua la domanda\par\vspace{.22cm}\noindent{\color{Classical}\bfseries 02}\hspace{.35cm}Esplicita le ipotesi\par\vspace{.22cm}\noindent{\color{Classical}\bfseries 03}\hspace{.35cm}Calcola o rappresenta\par\vspace{.22cm}\noindent{\color{Classical}\bfseries 04}\hspace{.35cm}Scrivi una conclusione\par\vspace{.22cm}\noindent{\color{Classical}\bfseries 05}\hspace{.35cm}Indica un limite\par\vspace{.22cm}
\end{minipage}\end{lrbox}
\begin{lrbox}{\RecallBox}\begin{minipage}[t]{16.45cm}\color{Ink}\raggedright\sloppy\fontsize{10.1}{13.2}\selectfont
Dove conviene investire nuove misure?
\end{minipage}\end{lrbox}
\MeasureBox{94-TitleBox}{\TitleBox}
\MeasureBox{94-BodyBox}{\BodyBox}
\MeasureBox{94-SideBox}{\SideBox}
\MeasureBox{94-RecallBox}{\RecallBox}
\begin{frame}[plain]
\hypertarget{slide94}{}\bookmark[level=1,page=94]{Attività 09 · applica il concetto}
\begin{tikzpicture}[remember picture,overlay]
\node[anchor=north east,inner sep=0pt,text=Classical,font=\fontsize{8}{10}\selectfont] at ([xshift=19.05cm,yshift=-.16cm]current page.north west) {\hyperlink{slide3}{Indice}};
\node[anchor=north west,inner sep=0pt,text=Navy,font=\bfseries] at ([xshift=.95cm,yshift=-.62cm]current page.north west) {\usebox{\TitleBox}};
\draw[Classical,line width=1.3pt] ([xshift=.95cm,yshift=-1.68cm]current page.north west)--([xshift=2.25cm,yshift=-1.68cm]current page.north west);
\node[anchor=north west,inner sep=0pt] at ([xshift=.95cm,yshift=-2.12cm]current page.north west) {\usebox{\BodyBox}};
\node[anchor=north west,inner sep=0pt,text=Quantum,font=\bfseries\fontsize{11.5}{14}\selectfont] at ([xshift=11.6cm,yshift=-2.1cm]current page.north west) {Consegna di lavoro};
\node[anchor=north west,inner sep=0pt] at ([xshift=11.6cm,yshift=-2.75cm]current page.north west) {\usebox{\SideBox}};
\draw[Rule,line width=.5pt] ([xshift=.95cm,yshift=-10.43cm]current page.north west)--([xshift=19.05cm,yshift=-10.43cm]current page.north west);
\node[anchor=north west,inner sep=0pt,text=Classical,font=\bfseries\fontsize{9.4}{12}\selectfont] at ([xshift=.95cm,yshift=-10.66cm]current page.north west) {Richiami};
\node[anchor=north west,inner sep=0pt] at ([xshift=2.6cm,yshift=-10.64cm]current page.north west) {\usebox{\RecallBox}};
\draw[Rule,line width=.35pt] ([xshift=.95cm,yshift=-11.92cm]current page.north west)--([xshift=19.05cm,yshift=-11.92cm]current page.north west);
\node[anchor=north west,inner sep=0pt,text=Muted,text width=9.0cm,font=\fontsize{7.7}{9.3}\selectfont] at ([xshift=.95cm,yshift=-12.06cm]current page.north west) {09 · Precisione e progetto di campionamento};
\node[anchor=north west,inner sep=0pt,text=Classical,text width=8.2cm,align=right,font=\fontsize{7.3}{9.1}\selectfont] at ([xshift=9.45cm,yshift=-12.06cm]current page.north west) {\href{https://openintrostat.github.io/ims/}{IMS} · CORSO};
\node[anchor=north east,inner sep=0pt,text=Muted,font=\fontsize{7.7}{9.3}\selectfont] at ([xshift=19.05cm,yshift=-12.06cm]current page.north west) {94/110};

\end{tikzpicture}
\note{Consegna: Disponi di risorse per 24 misure aggiuntive: tutte nella stessa stazione oppure distribuite in nuove stazioni. Da che cosa dipende la scelta?\par\medskip Soluzione per il docente: Dipende dalla quantità da stimare e dalla variabilità ai diversi livelli. Per caratterizzare una singola stazione, misure interne possono essere utili; per generalizzare all'area, nuove stazioni possono aggiungere informazione non ridondante. Non esiste una risposta basata soltanto sul conteggio totale.}
\end{frame}
\begin{lrbox}{\TitleBox}\begin{minipage}[t]{18.2cm}\color{Ink}\raggedright\sloppy\fontsize{20.7}{24}\selectfont
\bfseries\color{Navy}Discussione 09 · soluzione ragionata
\end{minipage}\end{lrbox}
\begin{lrbox}{\BodyBox}\begin{minipage}[t]{10.25cm}\color{Ink}\raggedright\sloppy\fontsize{12.4}{16.4}\selectfont
Dipende dalla quantità da stimare e dalla variabilità ai diversi livelli. Per caratterizzare una singola stazione, misure interne possono essere utili; per generalizzare all'area, nuove stazioni possono aggiungere informazione non ridondante. Non esiste una risposta basata soltanto sul conteggio totale.
\end{minipage}\end{lrbox}
\begin{lrbox}{\SideBox}\begin{minipage}[t]{7.45cm}\color{Ink}\raggedright\sloppy\fontsize{12.2}{16.1}\selectfont
La validità si riferisce al disegno eseguito, non soltanto a quello previsto.
\end{minipage}\end{lrbox}
\begin{lrbox}{\RecallBox}\begin{minipage}[t]{16.45cm}\color{Ink}\raggedright\sloppy\fontsize{10.1}{13.2}\selectfont
Valuta il collegamento tra domanda, procedimento e conclusione, non soltanto il numero finale.
\end{minipage}\end{lrbox}
\MeasureBox{95-TitleBox}{\TitleBox}
\MeasureBox{95-BodyBox}{\BodyBox}
\MeasureBox{95-SideBox}{\SideBox}
\MeasureBox{95-RecallBox}{\RecallBox}
\begin{frame}[plain]
\hypertarget{slide95}{}\bookmark[level=1,page=95]{Discussione 09 · soluzione ragionata}
\begin{tikzpicture}[remember picture,overlay]
\node[anchor=north east,inner sep=0pt,text=Classical,font=\fontsize{8}{10}\selectfont] at ([xshift=19.05cm,yshift=-.16cm]current page.north west) {\hyperlink{slide3}{Indice}};
\node[anchor=north west,inner sep=0pt,text=Navy,font=\bfseries] at ([xshift=.95cm,yshift=-.62cm]current page.north west) {\usebox{\TitleBox}};
\draw[Classical,line width=1.3pt] ([xshift=.95cm,yshift=-1.68cm]current page.north west)--([xshift=2.25cm,yshift=-1.68cm]current page.north west);
\node[anchor=north west,inner sep=0pt] at ([xshift=.95cm,yshift=-2.12cm]current page.north west) {\usebox{\BodyBox}};
\node[anchor=north west,inner sep=0pt,text=Quantum,font=\bfseries\fontsize{11.5}{14}\selectfont] at ([xshift=11.6cm,yshift=-2.1cm]current page.north west) {Criterio di lettura};
\node[anchor=north west,inner sep=0pt] at ([xshift=11.6cm,yshift=-2.75cm]current page.north west) {\usebox{\SideBox}};
\draw[Rule,line width=.5pt] ([xshift=.95cm,yshift=-10.43cm]current page.north west)--([xshift=19.05cm,yshift=-10.43cm]current page.north west);
\node[anchor=north west,inner sep=0pt,text=Classical,font=\bfseries\fontsize{9.4}{12}\selectfont] at ([xshift=.95cm,yshift=-10.66cm]current page.north west) {Richiami};
\node[anchor=north west,inner sep=0pt] at ([xshift=2.6cm,yshift=-10.64cm]current page.north west) {\usebox{\RecallBox}};
\draw[Rule,line width=.35pt] ([xshift=.95cm,yshift=-11.92cm]current page.north west)--([xshift=19.05cm,yshift=-11.92cm]current page.north west);
\node[anchor=north west,inner sep=0pt,text=Muted,text width=9.0cm,font=\fontsize{7.7}{9.3}\selectfont] at ([xshift=.95cm,yshift=-12.06cm]current page.north west) {09 · Precisione e progetto di campionamento};
\node[anchor=north west,inner sep=0pt,text=Classical,text width=8.2cm,align=right,font=\fontsize{7.3}{9.1}\selectfont] at ([xshift=9.45cm,yshift=-12.06cm]current page.north west) {\href{https://openintrostat.github.io/ims/}{IMS} · CORSO};
\node[anchor=north east,inner sep=0pt,text=Muted,font=\fontsize{7.7}{9.3}\selectfont] at ([xshift=19.05cm,yshift=-12.06cm]current page.north west) {95/110};

\end{tikzpicture}
\note{Dipende dalla quantità da stimare e dalla variabilità ai diversi livelli. Per caratterizzare una singola stazione, misure interne possono essere utili; per generalizzare all'area, nuove stazioni possono aggiungere informazione non ridondante. Non esiste una risposta basata soltanto sul conteggio totale.}
\end{frame}
\begin{frame}[plain]\hypertarget{slide96}{}
\bookmark[level=0,page=96]{Laboratorio di stima}
\begin{tikzpicture}[remember picture,overlay]
\fill[Night] (current page.south west) rectangle (current page.north east);
\draw[Mint,line width=2pt] ([xshift=1.05cm,yshift=-1.35cm]current page.north west)--([xshift=3.0cm,yshift=-1.35cm]current page.north west);
\node[anchor=north west,inner sep=0,text=Mint,font=\bfseries\fontsize{11}{14}\selectfont] at ([xshift=1.05cm,yshift=-1.7cm]current page.north west) {SEZIONE 10};
\node[anchor=north west,inner sep=0,text=white,text width=15.8cm,align=left,font=\bfseries\fontsize{31}{38}\selectfont] at ([xshift=1.05cm,yshift=-2.65cm]current page.north west) {Laboratorio di stima};
\node[anchor=north west,inner sep=0,text=Lilac,text width=15.8cm,align=left,font=\fontsize{15}{21}\selectfont] at ([xshift=1.05cm,yshift=-6.4cm]current page.north west) {Come presentare un numero con la sua incertezza e i suoi limiti?};
\node[anchor=north west,inner sep=0,text=white,text width=16.5cm,align=left,font=\fontsize{11}{15}\selectfont] at ([xshift=1.05cm,yshift=-8.6cm]current page.north west) {Comprendere il concetto · leggere un esempio · applicare e discutere};
\node[anchor=north east,inner sep=0,text=Mint,font=\fontsize{8}{10}\selectfont] at ([xshift=19.05cm,yshift=-.2cm]current page.north west) {\hyperlink{slide3}{Indice}};
\node[anchor=south east,inner sep=0,text=Lilac!45,font=\bfseries\fontsize{64}{68}\selectfont] at ([xshift=18.95cm,yshift=-11.35cm]current page.north west) {10};
\draw[Lilac!30,line width=.35pt] ([xshift=.95cm,yshift=-11.92cm]current page.north west)--([xshift=19.05cm,yshift=-11.92cm]current page.north west);
\node[anchor=north west,inner sep=0pt,text=Lilac,text width=9.0cm,font=\fontsize{7.7}{9.3}\selectfont] at ([xshift=.95cm,yshift=-12.06cm]current page.north west) {10 · Laboratorio di stima};
\node[anchor=north west,inner sep=0pt,text=Mint,text width=8.2cm,align=right,font=\fontsize{7.3}{9.1}\selectfont] at ([xshift=9.45cm,yshift=-12.06cm]current page.north west) {\href{https://openintrostat.github.io/ims/}{IMS} · \href{https://stat.ethz.ch/R-manual/R-devel/library/stats/html/t.test.html}{RTTEST}};
\node[anchor=north east,inner sep=0pt,text=Lilac,font=\fontsize{7.7}{9.3}\selectfont] at ([xshift=19.05cm,yshift=-12.06cm]current page.north west) {96/110};
\end{tikzpicture}\end{frame}
\begin{lrbox}{\TitleBox}\begin{minipage}[t]{18.2cm}\color{Ink}\raggedright\sloppy\fontsize{20.7}{24}\selectfont
\bfseries\color{Navy}Preparare la quantità da stimare
\end{minipage}\end{lrbox}
\begin{lrbox}{\BodyBox}\begin{minipage}[t]{10.25cm}\color{Ink}\raggedright\sloppy\fontsize{12.4}{16.4}\selectfont
Il laboratorio usa 80 individui indipendenti per ciascuno dei due morfotipi. La media con pesi 1/2 descrive una composizione fissata. Confrontiamo questo disegno con il ricampionamento pooled della miscela. I 24 provini raggruppati per stazione costituiscono un terzo caso, distinto.
\end{minipage}\end{lrbox}
\begin{lrbox}{\SideBox}\begin{minipage}[t]{7.45cm}\color{Ink}\raggedright\sloppy\fontsize{12.2}{16.1}\selectfont
File: fossili\_simulati.csv.\par\smallskip Quantità: media del diametro del modello didattico.\par\smallskip Unità: \(\mu\)m.
\end{minipage}\end{lrbox}
\begin{lrbox}{\RecallBox}\begin{minipage}[t]{16.45cm}\color{Ink}\raggedright\sloppy\fontsize{10.1}{13.2}\selectfont
La scelta del file determina le condizioni della dimostrazione.
\end{minipage}\end{lrbox}
\MeasureBox{97-TitleBox}{\TitleBox}
\MeasureBox{97-BodyBox}{\BodyBox}
\MeasureBox{97-SideBox}{\SideBox}
\MeasureBox{97-RecallBox}{\RecallBox}
\begin{frame}[plain]
\hypertarget{slide97}{}\bookmark[level=1,page=97]{Preparare la quantità da stimare}
\begin{tikzpicture}[remember picture,overlay]
\node[anchor=north east,inner sep=0pt,text=Classical,font=\fontsize{8}{10}\selectfont] at ([xshift=19.05cm,yshift=-.16cm]current page.north west) {\hyperlink{slide3}{Indice}};
\node[anchor=north west,inner sep=0pt,text=Navy,font=\bfseries] at ([xshift=.95cm,yshift=-.62cm]current page.north west) {\usebox{\TitleBox}};
\draw[Classical,line width=1.3pt] ([xshift=.95cm,yshift=-1.68cm]current page.north west)--([xshift=2.25cm,yshift=-1.68cm]current page.north west);
\node[anchor=north west,inner sep=0pt] at ([xshift=.95cm,yshift=-2.12cm]current page.north west) {\usebox{\BodyBox}};
\node[anchor=north west,inner sep=0pt,text=Quantum,font=\bfseries\fontsize{11.5}{14}\selectfont] at ([xshift=11.6cm,yshift=-2.1cm]current page.north west) {Esempio e interpretazione};
\node[anchor=north west,inner sep=0pt] at ([xshift=11.6cm,yshift=-2.75cm]current page.north west) {\usebox{\SideBox}};
\draw[Rule,line width=.5pt] ([xshift=.95cm,yshift=-10.43cm]current page.north west)--([xshift=19.05cm,yshift=-10.43cm]current page.north west);
\node[anchor=north west,inner sep=0pt,text=Classical,font=\bfseries\fontsize{9.4}{12}\selectfont] at ([xshift=.95cm,yshift=-10.66cm]current page.north west) {Richiami};
\node[anchor=north west,inner sep=0pt] at ([xshift=2.6cm,yshift=-10.64cm]current page.north west) {\usebox{\RecallBox}};
\draw[Rule,line width=.35pt] ([xshift=.95cm,yshift=-11.92cm]current page.north west)--([xshift=19.05cm,yshift=-11.92cm]current page.north west);
\node[anchor=north west,inner sep=0pt,text=Muted,text width=9.0cm,font=\fontsize{7.7}{9.3}\selectfont] at ([xshift=.95cm,yshift=-12.06cm]current page.north west) {10 · Laboratorio di stima};
\node[anchor=north west,inner sep=0pt,text=Classical,text width=8.2cm,align=right,font=\fontsize{7.3}{9.1}\selectfont] at ([xshift=9.45cm,yshift=-12.06cm]current page.north west) {\href{https://openintrostat.github.io/ims/}{IMS} · \href{https://stat.ethz.ch/R-manual/R-devel/library/stats/html/t.test.html}{RTTEST}};
\node[anchor=north east,inner sep=0pt,text=Muted,font=\fontsize{7.7}{9.3}\selectfont] at ([xshift=19.05cm,yshift=-12.06cm]current page.north west) {97/110};

\end{tikzpicture}
\note{Idea della slide: Il laboratorio usa 80 individui indipendenti per ciascuno dei due morfotipi. La media con pesi 1/2 descrive una composizione fissata. Confrontiamo questo disegno con il ricampionamento pooled della miscela. I 24 provini raggruppati per stazione costituiscono un terzo caso, distinto.\par\medskip Punto da far emergere: La scelta del file determina le condizioni della dimostrazione.\par\medskip Fonti della sezione: IMS, RTTEST}
\end{frame}
\begin{lrbox}{\TitleBox}\begin{minipage}[t]{18.2cm}\color{Ink}\raggedright\sloppy\fontsize{20.7}{24}\selectfont
\bfseries\color{Navy}Calcolare le tre dispersioni giuste
\end{minipage}\end{lrbox}
\begin{lrbox}{\BodyBox}\begin{minipage}[t]{8.2cm}\color{Ink}\raggedright\sloppy\fontsize{12.4}{16.4}\selectfont
Riportiamo deviazione standard dei dati, errore standard stimato e dispersione delle medie bootstrap. Le ultime due riguardano la media e dovrebbero essere confrontabili nel modello pooled. Per le quote fisse confrontiamo invece errore standard stratificato e relativo bootstrap. La prima quantità riguarda gli individui.
\end{minipage}\end{lrbox}
\begin{lrbox}{\SideBox}\begin{minipage}[t]{9.35cm}\color{Ink}\raggedright\sloppy\fontsize{12.2}{16.1}\selectfont
\begin{lstlisting}
s <- sd(x)
se <- s / sqrt(length(x))
se_boot <- sd(b)
\end{lstlisting}
\end{minipage}\end{lrbox}
\begin{lrbox}{\RecallBox}\begin{minipage}[t]{16.45cm}\color{Ink}\raggedright\sloppy\fontsize{10.1}{13.2}\selectfont
Non usare la parola errore senza indicare quale quantità descrive.
\end{minipage}\end{lrbox}
\MeasureBox{98-TitleBox}{\TitleBox}
\MeasureBox{98-BodyBox}{\BodyBox}
\MeasureBox{98-SideBox}{\SideBox}
\MeasureBox{98-RecallBox}{\RecallBox}
\begin{frame}[plain]
\hypertarget{slide98}{}\bookmark[level=1,page=98]{Calcolare le tre dispersioni giuste}
\begin{tikzpicture}[remember picture,overlay]
\node[anchor=north east,inner sep=0pt,text=Classical,font=\fontsize{8}{10}\selectfont] at ([xshift=19.05cm,yshift=-.16cm]current page.north west) {\hyperlink{slide3}{Indice}};
\node[anchor=north west,inner sep=0pt,text=Navy,font=\bfseries] at ([xshift=.95cm,yshift=-.62cm]current page.north west) {\usebox{\TitleBox}};
\draw[Classical,line width=1.3pt] ([xshift=.95cm,yshift=-1.68cm]current page.north west)--([xshift=2.25cm,yshift=-1.68cm]current page.north west);
\node[anchor=north west,inner sep=0pt] at ([xshift=.95cm,yshift=-2.12cm]current page.north west) {\usebox{\BodyBox}};
\node[anchor=north west,inner sep=0pt,text=Quantum,font=\bfseries\fontsize{11.5}{14}\selectfont] at ([xshift=9.7cm,yshift=-2.1cm]current page.north west) {Protocollo R};
\node[anchor=north west,inner sep=0pt] at ([xshift=9.7cm,yshift=-2.75cm]current page.north west) {\usebox{\SideBox}};
\draw[Rule,line width=.5pt] ([xshift=.95cm,yshift=-10.43cm]current page.north west)--([xshift=19.05cm,yshift=-10.43cm]current page.north west);
\node[anchor=north west,inner sep=0pt,text=Classical,font=\bfseries\fontsize{9.4}{12}\selectfont] at ([xshift=.95cm,yshift=-10.66cm]current page.north west) {Richiami};
\node[anchor=north west,inner sep=0pt] at ([xshift=2.6cm,yshift=-10.64cm]current page.north west) {\usebox{\RecallBox}};
\draw[Rule,line width=.35pt] ([xshift=.95cm,yshift=-11.92cm]current page.north west)--([xshift=19.05cm,yshift=-11.92cm]current page.north west);
\node[anchor=north west,inner sep=0pt,text=Muted,text width=9.0cm,font=\fontsize{7.7}{9.3}\selectfont] at ([xshift=.95cm,yshift=-12.06cm]current page.north west) {10 · Laboratorio di stima};
\node[anchor=north west,inner sep=0pt,text=Classical,text width=8.2cm,align=right,font=\fontsize{7.3}{9.1}\selectfont] at ([xshift=9.45cm,yshift=-12.06cm]current page.north west) {\href{https://openintrostat.github.io/ims/}{IMS} · \href{https://stat.ethz.ch/R-manual/R-devel/library/stats/html/t.test.html}{RTTEST}};
\node[anchor=north east,inner sep=0pt,text=Muted,font=\fontsize{7.7}{9.3}\selectfont] at ([xshift=19.05cm,yshift=-12.06cm]current page.north west) {98/110};

\end{tikzpicture}
\note{Idea della slide: Riportiamo deviazione standard dei dati, errore standard stimato e dispersione delle medie bootstrap. Le ultime due riguardano la media e dovrebbero essere confrontabili nel modello pooled. Per le quote fisse confrontiamo invece errore standard stratificato e relativo bootstrap. La prima quantità riguarda gli individui.\par\medskip Punto da far emergere: Non usare la parola errore senza indicare quale quantità descrive.\par\medskip Fonti della sezione: IMS, RTTEST}
\end{frame}
\begin{lrbox}{\TitleBox}\begin{minipage}[t]{18.2cm}\color{Ink}\raggedright\sloppy\fontsize{20.7}{24}\selectfont
\bfseries\color{Navy}Intervallo t e bootstrap
\end{minipage}\end{lrbox}
\begin{lrbox}{\BodyBox}\begin{minipage}[t]{8.2cm}\color{Ink}\raggedright\sloppy\fontsize{12.4}{16.4}\selectfont
Lo script calcola entrambi gli intervalli e li affianca. La differenza numerica può riflettere approssimazione, asimmetria e variabilità Monte Carlo. Non si sceglie a posteriori quello che sostiene la conclusione preferita; si motiva il metodo rispetto al problema.
\end{minipage}\end{lrbox}
\begin{lrbox}{\SideBox}\begin{minipage}[t]{9.35cm}\color{Ink}\raggedright\sloppy\fontsize{12.2}{16.1}\selectfont
\begin{lstlisting}
t.test(x)$conf.int
quantile(b, c(.025, .975))
\end{lstlisting}
\end{minipage}\end{lrbox}
\begin{lrbox}{\RecallBox}\begin{minipage}[t]{16.45cm}\color{Ink}\raggedright\sloppy\fontsize{10.1}{13.2}\selectfont
Il confronto è un controllo di sensibilità, non una gara per ottenere l'intervallo più stretto.
\end{minipage}\end{lrbox}
\MeasureBox{99-TitleBox}{\TitleBox}
\MeasureBox{99-BodyBox}{\BodyBox}
\MeasureBox{99-SideBox}{\SideBox}
\MeasureBox{99-RecallBox}{\RecallBox}
\begin{frame}[plain]
\hypertarget{slide99}{}\bookmark[level=1,page=99]{Intervallo t e bootstrap}
\begin{tikzpicture}[remember picture,overlay]
\node[anchor=north east,inner sep=0pt,text=Classical,font=\fontsize{8}{10}\selectfont] at ([xshift=19.05cm,yshift=-.16cm]current page.north west) {\hyperlink{slide3}{Indice}};
\node[anchor=north west,inner sep=0pt,text=Navy,font=\bfseries] at ([xshift=.95cm,yshift=-.62cm]current page.north west) {\usebox{\TitleBox}};
\draw[Classical,line width=1.3pt] ([xshift=.95cm,yshift=-1.68cm]current page.north west)--([xshift=2.25cm,yshift=-1.68cm]current page.north west);
\node[anchor=north west,inner sep=0pt] at ([xshift=.95cm,yshift=-2.12cm]current page.north west) {\usebox{\BodyBox}};
\node[anchor=north west,inner sep=0pt,text=Quantum,font=\bfseries\fontsize{11.5}{14}\selectfont] at ([xshift=9.7cm,yshift=-2.1cm]current page.north west) {Protocollo R};
\node[anchor=north west,inner sep=0pt] at ([xshift=9.7cm,yshift=-2.75cm]current page.north west) {\usebox{\SideBox}};
\draw[Rule,line width=.5pt] ([xshift=.95cm,yshift=-10.43cm]current page.north west)--([xshift=19.05cm,yshift=-10.43cm]current page.north west);
\node[anchor=north west,inner sep=0pt,text=Classical,font=\bfseries\fontsize{9.4}{12}\selectfont] at ([xshift=.95cm,yshift=-10.66cm]current page.north west) {Richiami};
\node[anchor=north west,inner sep=0pt] at ([xshift=2.6cm,yshift=-10.64cm]current page.north west) {\usebox{\RecallBox}};
\draw[Rule,line width=.35pt] ([xshift=.95cm,yshift=-11.92cm]current page.north west)--([xshift=19.05cm,yshift=-11.92cm]current page.north west);
\node[anchor=north west,inner sep=0pt,text=Muted,text width=9.0cm,font=\fontsize{7.7}{9.3}\selectfont] at ([xshift=.95cm,yshift=-12.06cm]current page.north west) {10 · Laboratorio di stima};
\node[anchor=north west,inner sep=0pt,text=Classical,text width=8.2cm,align=right,font=\fontsize{7.3}{9.1}\selectfont] at ([xshift=9.45cm,yshift=-12.06cm]current page.north west) {\href{https://openintrostat.github.io/ims/}{IMS} · \href{https://stat.ethz.ch/R-manual/R-devel/library/stats/html/t.test.html}{RTTEST}};
\node[anchor=north east,inner sep=0pt,text=Muted,font=\fontsize{7.7}{9.3}\selectfont] at ([xshift=19.05cm,yshift=-12.06cm]current page.north west) {99/110};

\end{tikzpicture}
\note{Idea della slide: Lo script calcola entrambi gli intervalli e li affianca. La differenza numerica può riflettere approssimazione, asimmetria e variabilità Monte Carlo. Non si sceglie a posteriori quello che sostiene la conclusione preferita; si motiva il metodo rispetto al problema.\par\medskip Punto da far emergere: Il confronto è un controllo di sensibilità, non una gara per ottenere l'intervallo più stretto.\par\medskip Fonti della sezione: IMS, RTTEST}
\end{frame}
\begin{lrbox}{\TitleBox}\begin{minipage}[t]{18.2cm}\color{Ink}\raggedright\sloppy\fontsize{20.7}{24}\selectfont
\bfseries\color{Navy}Rappresentare il ricampionamento
\end{minipage}\end{lrbox}
\begin{lrbox}{\BodyBox}\begin{minipage}[t]{7.1cm}\color{Ink}\raggedright\sloppy\fontsize{11.8}{15.8}\selectfont
L'istogramma delle medie bootstrap deve avere un asse che specifichi la statistica e l'unità. Le linee dei quantili indicano la costruzione percentile. Non si tratta della distribuzione dei fossili: ogni barra contiene repliche di una media.
\end{minipage}\end{lrbox}
\begin{lrbox}{\SideBox}\begin{minipage}[t]{10.4cm}\color{Ink}\raggedright\sloppy\fontsize{12.2}{16.1}\selectfont
\centering\includegraphics[width=10.4cm,height=6.25cm,keepaspectratio]{figure/bootstrap.pdf}\par\vspace{.22cm}{\raggedright\color{Muted}\fontsize{9.7}{12.5}\selectfont Ogni valore nella distribuzione è una media di un campione ricampionato.\par}
\end{minipage}\end{lrbox}
\begin{lrbox}{\RecallBox}\begin{minipage}[t]{16.45cm}\color{Ink}\raggedright\sloppy\fontsize{10.1}{13.2}\selectfont
Nella didascalia distinguere n osservazioni e B repliche.
\end{minipage}\end{lrbox}
\MeasureBox{100-TitleBox}{\TitleBox}
\MeasureBox{100-BodyBox}{\BodyBox}
\MeasureBox{100-SideBox}{\SideBox}
\MeasureBox{100-RecallBox}{\RecallBox}
\begin{frame}[plain]
\hypertarget{slide100}{}\bookmark[level=1,page=100]{Rappresentare il ricampionamento}
\begin{tikzpicture}[remember picture,overlay]
\node[anchor=north east,inner sep=0pt,text=Classical,font=\fontsize{8}{10}\selectfont] at ([xshift=19.05cm,yshift=-.16cm]current page.north west) {\hyperlink{slide3}{Indice}};
\node[anchor=north west,inner sep=0pt,text=Navy,font=\bfseries] at ([xshift=.95cm,yshift=-.62cm]current page.north west) {\usebox{\TitleBox}};
\draw[Classical,line width=1.3pt] ([xshift=.95cm,yshift=-1.68cm]current page.north west)--([xshift=2.25cm,yshift=-1.68cm]current page.north west);
\node[anchor=north west,inner sep=0pt] at ([xshift=.95cm,yshift=-2.12cm]current page.north west) {\usebox{\BodyBox}};
\node[anchor=north west,inner sep=0pt,text=Quantum,font=\bfseries\fontsize{11.5}{14}\selectfont] at ([xshift=8.65cm,yshift=-2.1cm]current page.north west) {Figura commentata};
\node[anchor=north west,inner sep=0pt] at ([xshift=8.65cm,yshift=-2.75cm]current page.north west) {\usebox{\SideBox}};
\draw[Rule,line width=.5pt] ([xshift=.95cm,yshift=-10.43cm]current page.north west)--([xshift=19.05cm,yshift=-10.43cm]current page.north west);
\node[anchor=north west,inner sep=0pt,text=Classical,font=\bfseries\fontsize{9.4}{12}\selectfont] at ([xshift=.95cm,yshift=-10.66cm]current page.north west) {Richiami};
\node[anchor=north west,inner sep=0pt] at ([xshift=2.6cm,yshift=-10.64cm]current page.north west) {\usebox{\RecallBox}};
\draw[Rule,line width=.35pt] ([xshift=.95cm,yshift=-11.92cm]current page.north west)--([xshift=19.05cm,yshift=-11.92cm]current page.north west);
\node[anchor=north west,inner sep=0pt,text=Muted,text width=9.0cm,font=\fontsize{7.7}{9.3}\selectfont] at ([xshift=.95cm,yshift=-12.06cm]current page.north west) {10 · Laboratorio di stima};
\node[anchor=north west,inner sep=0pt,text=Classical,text width=8.2cm,align=right,font=\fontsize{7.3}{9.1}\selectfont] at ([xshift=9.45cm,yshift=-12.06cm]current page.north west) {\href{https://openintrostat.github.io/ims/}{IMS} · \href{https://stat.ethz.ch/R-manual/R-devel/library/stats/html/t.test.html}{RTTEST}};
\node[anchor=north east,inner sep=0pt,text=Muted,font=\fontsize{7.7}{9.3}\selectfont] at ([xshift=19.05cm,yshift=-12.06cm]current page.north west) {100/110};

\end{tikzpicture}
\note{Idea della slide: L'istogramma delle medie bootstrap deve avere un asse che specifichi la statistica e l'unità. Le linee dei quantili indicano la costruzione percentile. Non si tratta della distribuzione dei fossili: ogni barra contiene repliche di una media.\par\medskip Punto da far emergere: Nella didascalia distinguere n osservazioni e B repliche.\par\medskip Fonti della sezione: IMS, RTTEST}
\end{frame}
\begin{lrbox}{\TitleBox}\begin{minipage}[t]{18.2cm}\color{Ink}\raggedright\sloppy\fontsize{20.7}{24}\selectfont
\bfseries\color{Navy}Discussione del campionamento gerarchico
\end{minipage}\end{lrbox}
\begin{lrbox}{\BodyBox}\begin{minipage}[t]{10.25cm}\color{Ink}\raggedright\sloppy\fontsize{12.4}{16.4}\selectfont
Il caso dei sei gruppi del modulo 1 serve a ragionare su cosa cambierebbe con dipendenza. Il laboratorio include un confronto dimostrativo tra ricampionamento di righe e di stazioni. I risultati non devono essere presentati come inferenza affidabile su tutta una formazione geologica.
\end{minipage}\end{lrbox}
\begin{lrbox}{\SideBox}\begin{minipage}[t]{7.45cm}\color{Ink}\raggedright\sloppy\fontsize{12.2}{16.1}\selectfont
Un numero piccolo di stazioni limita ciò che possiamo imparare sulla variabilità tra stazioni, indipendentemente dal numero di repliche bootstrap.
\end{minipage}\end{lrbox}
\begin{lrbox}{\RecallBox}\begin{minipage}[t]{16.45cm}\color{Ink}\raggedright\sloppy\fontsize{10.1}{13.2}\selectfont
La simulazione numerica non aumenta il numero di gruppi osservati.
\end{minipage}\end{lrbox}
\MeasureBox{101-TitleBox}{\TitleBox}
\MeasureBox{101-BodyBox}{\BodyBox}
\MeasureBox{101-SideBox}{\SideBox}
\MeasureBox{101-RecallBox}{\RecallBox}
\begin{frame}[plain]
\hypertarget{slide101}{}\bookmark[level=1,page=101]{Discussione del campionamento gerarchico}
\begin{tikzpicture}[remember picture,overlay]
\node[anchor=north east,inner sep=0pt,text=Classical,font=\fontsize{8}{10}\selectfont] at ([xshift=19.05cm,yshift=-.16cm]current page.north west) {\hyperlink{slide3}{Indice}};
\node[anchor=north west,inner sep=0pt,text=Navy,font=\bfseries] at ([xshift=.95cm,yshift=-.62cm]current page.north west) {\usebox{\TitleBox}};
\draw[Classical,line width=1.3pt] ([xshift=.95cm,yshift=-1.68cm]current page.north west)--([xshift=2.25cm,yshift=-1.68cm]current page.north west);
\node[anchor=north west,inner sep=0pt] at ([xshift=.95cm,yshift=-2.12cm]current page.north west) {\usebox{\BodyBox}};
\node[anchor=north west,inner sep=0pt,text=Quantum,font=\bfseries\fontsize{11.5}{14}\selectfont] at ([xshift=11.6cm,yshift=-2.1cm]current page.north west) {Esempio e interpretazione};
\node[anchor=north west,inner sep=0pt] at ([xshift=11.6cm,yshift=-2.75cm]current page.north west) {\usebox{\SideBox}};
\draw[Rule,line width=.5pt] ([xshift=.95cm,yshift=-10.43cm]current page.north west)--([xshift=19.05cm,yshift=-10.43cm]current page.north west);
\node[anchor=north west,inner sep=0pt,text=Classical,font=\bfseries\fontsize{9.4}{12}\selectfont] at ([xshift=.95cm,yshift=-10.66cm]current page.north west) {Richiami};
\node[anchor=north west,inner sep=0pt] at ([xshift=2.6cm,yshift=-10.64cm]current page.north west) {\usebox{\RecallBox}};
\draw[Rule,line width=.35pt] ([xshift=.95cm,yshift=-11.92cm]current page.north west)--([xshift=19.05cm,yshift=-11.92cm]current page.north west);
\node[anchor=north west,inner sep=0pt,text=Muted,text width=9.0cm,font=\fontsize{7.7}{9.3}\selectfont] at ([xshift=.95cm,yshift=-12.06cm]current page.north west) {10 · Laboratorio di stima};
\node[anchor=north west,inner sep=0pt,text=Classical,text width=8.2cm,align=right,font=\fontsize{7.3}{9.1}\selectfont] at ([xshift=9.45cm,yshift=-12.06cm]current page.north west) {\href{https://openintrostat.github.io/ims/}{IMS} · \href{https://stat.ethz.ch/R-manual/R-devel/library/stats/html/t.test.html}{RTTEST}};
\node[anchor=north east,inner sep=0pt,text=Muted,font=\fontsize{7.7}{9.3}\selectfont] at ([xshift=19.05cm,yshift=-12.06cm]current page.north west) {101/110};

\end{tikzpicture}
\note{Idea della slide: Il caso dei sei gruppi del modulo 1 serve a ragionare su cosa cambierebbe con dipendenza. Il laboratorio include un confronto dimostrativo tra ricampionamento di righe e di stazioni. I risultati non devono essere presentati come inferenza affidabile su tutta una formazione geologica.\par\medskip Punto da far emergere: La simulazione numerica non aumenta il numero di gruppi osservati.\par\medskip Fonti della sezione: IMS, RTTEST}
\end{frame}
\begin{lrbox}{\TitleBox}\begin{minipage}[t]{18.2cm}\color{Ink}\raggedright\sloppy\fontsize{20.7}{24}\selectfont
\bfseries\color{Navy}Una conclusione completa
\end{minipage}\end{lrbox}
\begin{lrbox}{\BodyBox}\begin{minipage}[t]{10.25cm}\color{Ink}\raggedright\sloppy\fontsize{12.4}{16.4}\selectfont
La conclusione indica stima, intervallo, unità, popolazione di riferimento, procedura e limite principale. Deve evitare frasi come «il 95\% dei fossili cade nell'intervallo della media». L'oggetto della stima va ripetuto quando potrebbe essere confuso.
\end{minipage}\end{lrbox}
\begin{lrbox}{\SideBox}\begin{minipage}[t]{7.45cm}\color{Ink}\raggedright\sloppy\fontsize{12.2}{16.1}\selectfont
«La media del diametro nel caso simulato è circa 6,355 \(\mu\)m; l'intervallo quantifica l'incertezza della media con quote 80/80; l'intervallo deve specificare se rispetta queste quote».
\end{minipage}\end{lrbox}
\begin{lrbox}{\RecallBox}\begin{minipage}[t]{16.45cm}\color{Ink}\raggedright\sloppy\fontsize{10.1}{13.2}\selectfont
Una buona interpretazione vale più di molte cifre senza condizioni.
\end{minipage}\end{lrbox}
\MeasureBox{102-TitleBox}{\TitleBox}
\MeasureBox{102-BodyBox}{\BodyBox}
\MeasureBox{102-SideBox}{\SideBox}
\MeasureBox{102-RecallBox}{\RecallBox}
\begin{frame}[plain]
\hypertarget{slide102}{}\bookmark[level=1,page=102]{Una conclusione completa}
\begin{tikzpicture}[remember picture,overlay]
\node[anchor=north east,inner sep=0pt,text=Classical,font=\fontsize{8}{10}\selectfont] at ([xshift=19.05cm,yshift=-.16cm]current page.north west) {\hyperlink{slide3}{Indice}};
\node[anchor=north west,inner sep=0pt,text=Navy,font=\bfseries] at ([xshift=.95cm,yshift=-.62cm]current page.north west) {\usebox{\TitleBox}};
\draw[Classical,line width=1.3pt] ([xshift=.95cm,yshift=-1.68cm]current page.north west)--([xshift=2.25cm,yshift=-1.68cm]current page.north west);
\node[anchor=north west,inner sep=0pt] at ([xshift=.95cm,yshift=-2.12cm]current page.north west) {\usebox{\BodyBox}};
\node[anchor=north west,inner sep=0pt,text=Quantum,font=\bfseries\fontsize{11.5}{14}\selectfont] at ([xshift=11.6cm,yshift=-2.1cm]current page.north west) {Esempio e interpretazione};
\node[anchor=north west,inner sep=0pt] at ([xshift=11.6cm,yshift=-2.75cm]current page.north west) {\usebox{\SideBox}};
\draw[Rule,line width=.5pt] ([xshift=.95cm,yshift=-10.43cm]current page.north west)--([xshift=19.05cm,yshift=-10.43cm]current page.north west);
\node[anchor=north west,inner sep=0pt,text=Classical,font=\bfseries\fontsize{9.4}{12}\selectfont] at ([xshift=.95cm,yshift=-10.66cm]current page.north west) {Richiami};
\node[anchor=north west,inner sep=0pt] at ([xshift=2.6cm,yshift=-10.64cm]current page.north west) {\usebox{\RecallBox}};
\draw[Rule,line width=.35pt] ([xshift=.95cm,yshift=-11.92cm]current page.north west)--([xshift=19.05cm,yshift=-11.92cm]current page.north west);
\node[anchor=north west,inner sep=0pt,text=Muted,text width=9.0cm,font=\fontsize{7.7}{9.3}\selectfont] at ([xshift=.95cm,yshift=-12.06cm]current page.north west) {10 · Laboratorio di stima};
\node[anchor=north west,inner sep=0pt,text=Classical,text width=8.2cm,align=right,font=\fontsize{7.3}{9.1}\selectfont] at ([xshift=9.45cm,yshift=-12.06cm]current page.north west) {\href{https://openintrostat.github.io/ims/}{IMS} · \href{https://stat.ethz.ch/R-manual/R-devel/library/stats/html/t.test.html}{RTTEST}};
\node[anchor=north east,inner sep=0pt,text=Muted,font=\fontsize{7.7}{9.3}\selectfont] at ([xshift=19.05cm,yshift=-12.06cm]current page.north west) {102/110};

\end{tikzpicture}
\note{Idea della slide: La conclusione indica stima, intervallo, unità, popolazione di riferimento, procedura e limite principale. Deve evitare frasi come «il 95\% dei fossili cade nell'intervallo della media». L'oggetto della stima va ripetuto quando potrebbe essere confuso.\par\medskip Punto da far emergere: Una buona interpretazione vale più di molte cifre senza condizioni.\par\medskip Fonti della sezione: IMS, RTTEST}
\end{frame}
\begin{lrbox}{\TitleBox}\begin{minipage}[t]{18.2cm}\color{Ink}\raggedright\sloppy\fontsize{20.7}{24}\selectfont
\bfseries\color{Navy}Dal quanto al confronto
\end{minipage}\end{lrbox}
\begin{lrbox}{\BodyBox}\begin{minipage}[t]{10.25cm}\color{Ink}\raggedright\sloppy\fontsize{12.4}{16.4}\selectfont
Il prossimo modulo aggiungerà test di ipotesi, ma la stima non viene sostituita dal p-value. Effetto e intervallo resteranno il centro del risultato. Il test risponde a una domanda complementare sulla compatibilità dei dati con un'ipotesi specificata.
\end{minipage}\end{lrbox}
\begin{lrbox}{\SideBox}\begin{minipage}[t]{7.45cm}\color{Ink}\raggedright\sloppy\fontsize{12.2}{16.1}\selectfont
Stima: quanto grande potrebbe essere l'effetto?\par\smallskip Test: quanto sono incompatibili i dati con l'ipotesi di riferimento?
\end{minipage}\end{lrbox}
\begin{lrbox}{\RecallBox}\begin{minipage}[t]{16.45cm}\color{Ink}\raggedright\sloppy\fontsize{10.1}{13.2}\selectfont
Intervalli, effetti e test vanno interpretati insieme al disegno dello studio.
\end{minipage}\end{lrbox}
\MeasureBox{103-TitleBox}{\TitleBox}
\MeasureBox{103-BodyBox}{\BodyBox}
\MeasureBox{103-SideBox}{\SideBox}
\MeasureBox{103-RecallBox}{\RecallBox}
\begin{frame}[plain]
\hypertarget{slide103}{}\bookmark[level=1,page=103]{Dal quanto al confronto}
\begin{tikzpicture}[remember picture,overlay]
\node[anchor=north east,inner sep=0pt,text=Classical,font=\fontsize{8}{10}\selectfont] at ([xshift=19.05cm,yshift=-.16cm]current page.north west) {\hyperlink{slide3}{Indice}};
\node[anchor=north west,inner sep=0pt,text=Navy,font=\bfseries] at ([xshift=.95cm,yshift=-.62cm]current page.north west) {\usebox{\TitleBox}};
\draw[Classical,line width=1.3pt] ([xshift=.95cm,yshift=-1.68cm]current page.north west)--([xshift=2.25cm,yshift=-1.68cm]current page.north west);
\node[anchor=north west,inner sep=0pt] at ([xshift=.95cm,yshift=-2.12cm]current page.north west) {\usebox{\BodyBox}};
\node[anchor=north west,inner sep=0pt,text=Quantum,font=\bfseries\fontsize{11.5}{14}\selectfont] at ([xshift=11.6cm,yshift=-2.1cm]current page.north west) {Esempio e interpretazione};
\node[anchor=north west,inner sep=0pt] at ([xshift=11.6cm,yshift=-2.75cm]current page.north west) {\usebox{\SideBox}};
\draw[Rule,line width=.5pt] ([xshift=.95cm,yshift=-10.43cm]current page.north west)--([xshift=19.05cm,yshift=-10.43cm]current page.north west);
\node[anchor=north west,inner sep=0pt,text=Classical,font=\bfseries\fontsize{9.4}{12}\selectfont] at ([xshift=.95cm,yshift=-10.66cm]current page.north west) {Richiami};
\node[anchor=north west,inner sep=0pt] at ([xshift=2.6cm,yshift=-10.64cm]current page.north west) {\usebox{\RecallBox}};
\draw[Rule,line width=.35pt] ([xshift=.95cm,yshift=-11.92cm]current page.north west)--([xshift=19.05cm,yshift=-11.92cm]current page.north west);
\node[anchor=north west,inner sep=0pt,text=Muted,text width=9.0cm,font=\fontsize{7.7}{9.3}\selectfont] at ([xshift=.95cm,yshift=-12.06cm]current page.north west) {10 · Laboratorio di stima};
\node[anchor=north west,inner sep=0pt,text=Classical,text width=8.2cm,align=right,font=\fontsize{7.3}{9.1}\selectfont] at ([xshift=9.45cm,yshift=-12.06cm]current page.north west) {\href{https://openintrostat.github.io/ims/}{IMS} · \href{https://stat.ethz.ch/R-manual/R-devel/library/stats/html/t.test.html}{RTTEST}};
\node[anchor=north east,inner sep=0pt,text=Muted,font=\fontsize{7.7}{9.3}\selectfont] at ([xshift=19.05cm,yshift=-12.06cm]current page.north west) {103/110};

\end{tikzpicture}
\note{Idea della slide: Il prossimo modulo aggiungerà test di ipotesi, ma la stima non viene sostituita dal p-value. Effetto e intervallo resteranno il centro del risultato. Il test risponde a una domanda complementare sulla compatibilità dei dati con un'ipotesi specificata.\par\medskip Punto da far emergere: Intervalli, effetti e test vanno interpretati insieme al disegno dello studio.\par\medskip Fonti della sezione: IMS, RTTEST}
\end{frame}
\begin{lrbox}{\TitleBox}\begin{minipage}[t]{18.2cm}\color{Ink}\raggedright\sloppy\fontsize{20.7}{24}\selectfont
\bfseries\color{Navy}Attività 10 · applica il concetto
\end{minipage}\end{lrbox}
\begin{lrbox}{\BodyBox}\begin{minipage}[t]{10.25cm}\color{Ink}\raggedright\sloppy\fontsize{12.4}{16.4}\selectfont
Sul dataset dei fossili confronta l'analisi della miscela con un bootstrap che mantenga 80 individui per morfotipo. Riporta media, dispersione e intervalli. Spiega quale schema imita il generatore e quale modifica il disegno.
\end{minipage}\end{lrbox}
\begin{lrbox}{\SideBox}\begin{minipage}[t]{7.45cm}\color{Ink}\raggedright\sloppy\fontsize{12.2}{16.1}\selectfont
\noindent{\color{Classical}\bfseries 01}\hspace{.35cm}Individua la domanda\par\vspace{.22cm}\noindent{\color{Classical}\bfseries 02}\hspace{.35cm}Esplicita le ipotesi\par\vspace{.22cm}\noindent{\color{Classical}\bfseries 03}\hspace{.35cm}Calcola o rappresenta\par\vspace{.22cm}\noindent{\color{Classical}\bfseries 04}\hspace{.35cm}Scrivi una conclusione\par\vspace{.22cm}\noindent{\color{Classical}\bfseries 05}\hspace{.35cm}Indica un limite\par\vspace{.22cm}
\end{minipage}\end{lrbox}
\begin{lrbox}{\RecallBox}\begin{minipage}[t]{16.45cm}\color{Ink}\raggedright\sloppy\fontsize{10.1}{13.2}\selectfont
Come presentare un numero con la sua incertezza e i suoi limiti?
\end{minipage}\end{lrbox}
\MeasureBox{104-TitleBox}{\TitleBox}
\MeasureBox{104-BodyBox}{\BodyBox}
\MeasureBox{104-SideBox}{\SideBox}
\MeasureBox{104-RecallBox}{\RecallBox}
\begin{frame}[plain]
\hypertarget{slide104}{}\bookmark[level=1,page=104]{Attività 10 · applica il concetto}
\begin{tikzpicture}[remember picture,overlay]
\node[anchor=north east,inner sep=0pt,text=Classical,font=\fontsize{8}{10}\selectfont] at ([xshift=19.05cm,yshift=-.16cm]current page.north west) {\hyperlink{slide3}{Indice}};
\node[anchor=north west,inner sep=0pt,text=Navy,font=\bfseries] at ([xshift=.95cm,yshift=-.62cm]current page.north west) {\usebox{\TitleBox}};
\draw[Classical,line width=1.3pt] ([xshift=.95cm,yshift=-1.68cm]current page.north west)--([xshift=2.25cm,yshift=-1.68cm]current page.north west);
\node[anchor=north west,inner sep=0pt] at ([xshift=.95cm,yshift=-2.12cm]current page.north west) {\usebox{\BodyBox}};
\node[anchor=north west,inner sep=0pt,text=Quantum,font=\bfseries\fontsize{11.5}{14}\selectfont] at ([xshift=11.6cm,yshift=-2.1cm]current page.north west) {Consegna di lavoro};
\node[anchor=north west,inner sep=0pt] at ([xshift=11.6cm,yshift=-2.75cm]current page.north west) {\usebox{\SideBox}};
\draw[Rule,line width=.5pt] ([xshift=.95cm,yshift=-10.43cm]current page.north west)--([xshift=19.05cm,yshift=-10.43cm]current page.north west);
\node[anchor=north west,inner sep=0pt,text=Classical,font=\bfseries\fontsize{9.4}{12}\selectfont] at ([xshift=.95cm,yshift=-10.66cm]current page.north west) {Richiami};
\node[anchor=north west,inner sep=0pt] at ([xshift=2.6cm,yshift=-10.64cm]current page.north west) {\usebox{\RecallBox}};
\draw[Rule,line width=.35pt] ([xshift=.95cm,yshift=-11.92cm]current page.north west)--([xshift=19.05cm,yshift=-11.92cm]current page.north west);
\node[anchor=north west,inner sep=0pt,text=Muted,text width=9.0cm,font=\fontsize{7.7}{9.3}\selectfont] at ([xshift=.95cm,yshift=-12.06cm]current page.north west) {10 · Laboratorio di stima};
\node[anchor=north west,inner sep=0pt,text=Classical,text width=8.2cm,align=right,font=\fontsize{7.3}{9.1}\selectfont] at ([xshift=9.45cm,yshift=-12.06cm]current page.north west) {\href{https://openintrostat.github.io/ims/}{IMS} · \href{https://stat.ethz.ch/R-manual/R-devel/library/stats/html/t.test.html}{RTTEST}};
\node[anchor=north east,inner sep=0pt,text=Muted,font=\fontsize{7.7}{9.3}\selectfont] at ([xshift=19.05cm,yshift=-12.06cm]current page.north west) {104/110};

\end{tikzpicture}
\note{Consegna: Sul dataset dei fossili confronta l'analisi della miscela con un bootstrap che mantenga 80 individui per morfotipo. Riporta media, dispersione e intervalli. Spiega quale schema imita il generatore e quale modifica il disegno.\par\medskip Soluzione per il docente: Il bootstrap stratificato mantiene le quote 80/80 del generatore e stima la media dei due gruppi con pesi 1/2. Il pooled consente quote variabili e rappresenta un diverso campionamento della miscela. La media osservata coincide, ma l'incertezza può differire. Riportare schema, ipotesi, repliche e natura simulata.}
\end{frame}
\begin{lrbox}{\TitleBox}\begin{minipage}[t]{18.2cm}\color{Ink}\raggedright\sloppy\fontsize{20.7}{24}\selectfont
\bfseries\color{Navy}Discussione 10 · soluzione ragionata
\end{minipage}\end{lrbox}
\begin{lrbox}{\BodyBox}\begin{minipage}[t]{10.25cm}\color{Ink}\raggedright\sloppy\fontsize{12.4}{16.4}\selectfont
Il bootstrap stratificato mantiene le quote 80/80 del generatore e stima la media dei due gruppi con pesi 1/2. Il pooled consente quote variabili e rappresenta un diverso campionamento della miscela. La media osservata coincide, ma l'incertezza può differire. Riportare schema, ipotesi, repliche e natura simulata.
\end{minipage}\end{lrbox}
\begin{lrbox}{\SideBox}\begin{minipage}[t]{7.45cm}\color{Ink}\raggedright\sloppy\fontsize{12.2}{16.1}\selectfont
Intervalli, effetti e test vanno interpretati insieme al disegno dello studio.
\end{minipage}\end{lrbox}
\begin{lrbox}{\RecallBox}\begin{minipage}[t]{16.45cm}\color{Ink}\raggedright\sloppy\fontsize{10.1}{13.2}\selectfont
Valuta il collegamento tra domanda, procedimento e conclusione, non soltanto il numero finale.
\end{minipage}\end{lrbox}
\MeasureBox{105-TitleBox}{\TitleBox}
\MeasureBox{105-BodyBox}{\BodyBox}
\MeasureBox{105-SideBox}{\SideBox}
\MeasureBox{105-RecallBox}{\RecallBox}
\begin{frame}[plain]
\hypertarget{slide105}{}\bookmark[level=1,page=105]{Discussione 10 · soluzione ragionata}
\begin{tikzpicture}[remember picture,overlay]
\node[anchor=north east,inner sep=0pt,text=Classical,font=\fontsize{8}{10}\selectfont] at ([xshift=19.05cm,yshift=-.16cm]current page.north west) {\hyperlink{slide3}{Indice}};
\node[anchor=north west,inner sep=0pt,text=Navy,font=\bfseries] at ([xshift=.95cm,yshift=-.62cm]current page.north west) {\usebox{\TitleBox}};
\draw[Classical,line width=1.3pt] ([xshift=.95cm,yshift=-1.68cm]current page.north west)--([xshift=2.25cm,yshift=-1.68cm]current page.north west);
\node[anchor=north west,inner sep=0pt] at ([xshift=.95cm,yshift=-2.12cm]current page.north west) {\usebox{\BodyBox}};
\node[anchor=north west,inner sep=0pt,text=Quantum,font=\bfseries\fontsize{11.5}{14}\selectfont] at ([xshift=11.6cm,yshift=-2.1cm]current page.north west) {Criterio di lettura};
\node[anchor=north west,inner sep=0pt] at ([xshift=11.6cm,yshift=-2.75cm]current page.north west) {\usebox{\SideBox}};
\draw[Rule,line width=.5pt] ([xshift=.95cm,yshift=-10.43cm]current page.north west)--([xshift=19.05cm,yshift=-10.43cm]current page.north west);
\node[anchor=north west,inner sep=0pt,text=Classical,font=\bfseries\fontsize{9.4}{12}\selectfont] at ([xshift=.95cm,yshift=-10.66cm]current page.north west) {Richiami};
\node[anchor=north west,inner sep=0pt] at ([xshift=2.6cm,yshift=-10.64cm]current page.north west) {\usebox{\RecallBox}};
\draw[Rule,line width=.35pt] ([xshift=.95cm,yshift=-11.92cm]current page.north west)--([xshift=19.05cm,yshift=-11.92cm]current page.north west);
\node[anchor=north west,inner sep=0pt,text=Muted,text width=9.0cm,font=\fontsize{7.7}{9.3}\selectfont] at ([xshift=.95cm,yshift=-12.06cm]current page.north west) {10 · Laboratorio di stima};
\node[anchor=north west,inner sep=0pt,text=Classical,text width=8.2cm,align=right,font=\fontsize{7.3}{9.1}\selectfont] at ([xshift=9.45cm,yshift=-12.06cm]current page.north west) {\href{https://openintrostat.github.io/ims/}{IMS} · \href{https://stat.ethz.ch/R-manual/R-devel/library/stats/html/t.test.html}{RTTEST}};
\node[anchor=north east,inner sep=0pt,text=Muted,font=\fontsize{7.7}{9.3}\selectfont] at ([xshift=19.05cm,yshift=-12.06cm]current page.north west) {105/110};

\end{tikzpicture}
\note{Il bootstrap stratificato mantiene le quote 80/80 del generatore e stima la media dei due gruppi con pesi 1/2. Il pooled consente quote variabili e rappresenta un diverso campionamento della miscela. La media osservata coincide, ma l'incertezza può differire. Riportare schema, ipotesi, repliche e natura simulata.}
\end{frame}
\begin{lrbox}{\TitleBox}\begin{minipage}[t]{18.2cm}\color{Ink}\raggedright\sloppy\fontsize{20.7}{24}\selectfont
\bfseries\color{Navy}Sintesi del modulo
\end{minipage}\end{lrbox}
\begin{lrbox}{\BodyBox}\begin{minipage}[t]{10.25cm}\color{Ink}\raggedright\sloppy\fontsize{12.4}{16.4}\selectfont
Definire popolazione, quantità e unità di campionamento.\par\medskip Distinguere deviazione standard ed errore standard.\par\medskip Interpretare intervalli e bootstrap con i loro limiti.
\end{minipage}\end{lrbox}
\begin{lrbox}{\SideBox}\begin{minipage}[t]{7.45cm}\color{Ink}\raggedright\sloppy\fontsize{12.2}{16.1}\selectfont
\noindent{\color{Classical}\bfseries 01}\hspace{.35cm}Definisci la domanda\par\vspace{.22cm}\noindent{\color{Classical}\bfseries 02}\hspace{.35cm}Controlla il dato\par\vspace{.22cm}\noindent{\color{Classical}\bfseries 03}\hspace{.35cm}Scegli il metodo\par\vspace{.22cm}\noindent{\color{Classical}\bfseries 04}\hspace{.35cm}Valuta il risultato\par\vspace{.22cm}\noindent{\color{Classical}\bfseries 05}\hspace{.35cm}Dichiara i limiti\par\vspace{.22cm}
\end{minipage}\end{lrbox}
\begin{lrbox}{\RecallBox}\begin{minipage}[t]{16.45cm}\color{Ink}\raggedright\sloppy\fontsize{10.1}{13.2}\selectfont
La precisione non corregge un campionamento distorto.
\end{minipage}\end{lrbox}
\MeasureBox{106-TitleBox}{\TitleBox}
\MeasureBox{106-BodyBox}{\BodyBox}
\MeasureBox{106-SideBox}{\SideBox}
\MeasureBox{106-RecallBox}{\RecallBox}
\begin{frame}[plain]
\hypertarget{slide106}{}\bookmark[level=1,page=106]{Sintesi del modulo}
\begin{tikzpicture}[remember picture,overlay]
\node[anchor=north east,inner sep=0pt,text=Classical,font=\fontsize{8}{10}\selectfont] at ([xshift=19.05cm,yshift=-.16cm]current page.north west) {\hyperlink{slide3}{Indice}};
\node[anchor=north west,inner sep=0pt,text=Navy,font=\bfseries] at ([xshift=.95cm,yshift=-.62cm]current page.north west) {\usebox{\TitleBox}};
\draw[Classical,line width=1.3pt] ([xshift=.95cm,yshift=-1.68cm]current page.north west)--([xshift=2.25cm,yshift=-1.68cm]current page.north west);
\node[anchor=north west,inner sep=0pt] at ([xshift=.95cm,yshift=-2.12cm]current page.north west) {\usebox{\BodyBox}};
\node[anchor=north west,inner sep=0pt,text=Quantum,font=\bfseries\fontsize{11.5}{14}\selectfont] at ([xshift=11.6cm,yshift=-2.1cm]current page.north west) {Passaggi da seguire};
\node[anchor=north west,inner sep=0pt] at ([xshift=11.6cm,yshift=-2.75cm]current page.north west) {\usebox{\SideBox}};
\draw[Rule,line width=.5pt] ([xshift=.95cm,yshift=-10.43cm]current page.north west)--([xshift=19.05cm,yshift=-10.43cm]current page.north west);
\node[anchor=north west,inner sep=0pt,text=Classical,font=\bfseries\fontsize{9.4}{12}\selectfont] at ([xshift=.95cm,yshift=-10.66cm]current page.north west) {Richiami};
\node[anchor=north west,inner sep=0pt] at ([xshift=2.6cm,yshift=-10.64cm]current page.north west) {\usebox{\RecallBox}};
\draw[Rule,line width=.35pt] ([xshift=.95cm,yshift=-11.92cm]current page.north west)--([xshift=19.05cm,yshift=-11.92cm]current page.north west);
\node[anchor=north west,inner sep=0pt,text=Muted,text width=9.0cm,font=\fontsize{7.7}{9.3}\selectfont] at ([xshift=.95cm,yshift=-12.06cm]current page.north west) {Modulo 05 · Stima e incertezza};
\node[anchor=north west,inner sep=0pt,text=Classical,text width=8.2cm,align=right,font=\fontsize{7.3}{9.1}\selectfont] at ([xshift=9.45cm,yshift=-12.06cm]current page.north west) {};
\node[anchor=north east,inner sep=0pt,text=Muted,font=\fontsize{7.7}{9.3}\selectfont] at ([xshift=19.05cm,yshift=-12.06cm]current page.north west) {106/110};

\end{tikzpicture}
\note{La precisione non corregge un campionamento distorto.}
\end{frame}
\begin{lrbox}{\TitleBox}\begin{minipage}[t]{18.2cm}\color{Ink}\raggedright\sloppy\fontsize{20.7}{24}\selectfont
\bfseries\color{Navy}Verifica e consegna finale
\end{minipage}\end{lrbox}
\begin{lrbox}{\BodyBox}\begin{minipage}[t]{10.25cm}\color{Ink}\raggedright\sloppy\fontsize{12.4}{16.4}\selectfont
Sul dataset dei fossili confronta l'analisi della miscela con un bootstrap che mantenga 80 individui per morfotipo. Riporta media, dispersione e intervalli. Spiega quale schema imita il generatore e quale modifica il disegno.\par\medskip Allega i file necessari a ricostruire i passaggi. Nella conclusione indica una scelta che potresti modificare e come controlleresti il suo effetto.
\end{minipage}\end{lrbox}
\begin{lrbox}{\SideBox}\begin{minipage}[t]{7.45cm}\color{Ink}\raggedright\sloppy\fontsize{12.2}{16.1}\selectfont
Criteri: correttezza della domanda; qualità dei controlli; coerenza del metodo; interpretazione; riproducibilità.\par\medskip Le soluzioni sono nella guida docente e nelle slide di discussione.
\end{minipage}\end{lrbox}
\begin{lrbox}{\RecallBox}\begin{minipage}[t]{16.45cm}\color{Ink}\raggedright\sloppy\fontsize{10.1}{13.2}\selectfont
Una conclusione delimitata e verificabile vale più di un risultato apparentemente definitivo.
\end{minipage}\end{lrbox}
\MeasureBox{107-TitleBox}{\TitleBox}
\MeasureBox{107-BodyBox}{\BodyBox}
\MeasureBox{107-SideBox}{\SideBox}
\MeasureBox{107-RecallBox}{\RecallBox}
\begin{frame}[plain]
\hypertarget{slide107}{}\bookmark[level=1,page=107]{Verifica e consegna finale}
\begin{tikzpicture}[remember picture,overlay]
\node[anchor=north east,inner sep=0pt,text=Classical,font=\fontsize{8}{10}\selectfont] at ([xshift=19.05cm,yshift=-.16cm]current page.north west) {\hyperlink{slide3}{Indice}};
\node[anchor=north west,inner sep=0pt,text=Navy,font=\bfseries] at ([xshift=.95cm,yshift=-.62cm]current page.north west) {\usebox{\TitleBox}};
\draw[Classical,line width=1.3pt] ([xshift=.95cm,yshift=-1.68cm]current page.north west)--([xshift=2.25cm,yshift=-1.68cm]current page.north west);
\node[anchor=north west,inner sep=0pt] at ([xshift=.95cm,yshift=-2.12cm]current page.north west) {\usebox{\BodyBox}};
\node[anchor=north west,inner sep=0pt,text=Quantum,font=\bfseries\fontsize{11.5}{14}\selectfont] at ([xshift=11.6cm,yshift=-2.1cm]current page.north west) {Esempio e interpretazione};
\node[anchor=north west,inner sep=0pt] at ([xshift=11.6cm,yshift=-2.75cm]current page.north west) {\usebox{\SideBox}};
\draw[Rule,line width=.5pt] ([xshift=.95cm,yshift=-10.43cm]current page.north west)--([xshift=19.05cm,yshift=-10.43cm]current page.north west);
\node[anchor=north west,inner sep=0pt,text=Classical,font=\bfseries\fontsize{9.4}{12}\selectfont] at ([xshift=.95cm,yshift=-10.66cm]current page.north west) {Richiami};
\node[anchor=north west,inner sep=0pt] at ([xshift=2.6cm,yshift=-10.64cm]current page.north west) {\usebox{\RecallBox}};
\draw[Rule,line width=.35pt] ([xshift=.95cm,yshift=-11.92cm]current page.north west)--([xshift=19.05cm,yshift=-11.92cm]current page.north west);
\node[anchor=north west,inner sep=0pt,text=Muted,text width=9.0cm,font=\fontsize{7.7}{9.3}\selectfont] at ([xshift=.95cm,yshift=-12.06cm]current page.north west) {Modulo 05 · Stima e incertezza};
\node[anchor=north west,inner sep=0pt,text=Classical,text width=8.2cm,align=right,font=\fontsize{7.3}{9.1}\selectfont] at ([xshift=9.45cm,yshift=-12.06cm]current page.north west) {};
\node[anchor=north east,inner sep=0pt,text=Muted,font=\fontsize{7.7}{9.3}\selectfont] at ([xshift=19.05cm,yshift=-12.06cm]current page.north west) {107/110};

\end{tikzpicture}
\note{Una conclusione delimitata e verificabile vale più di un risultato apparentemente definitivo.}
\end{frame}
\begin{frame}[plain]\hypertarget{slide108}{}\bookmark[level=1,page=108]{Fonti e approfondimenti}\begin{tikzpicture}[remember picture,overlay]
\node[anchor=north west,inner sep=0,text=Navy,font=\bfseries\fontsize{22}{26}\selectfont] at ([xshift=.95cm,yshift=-.65cm]current page.north west) {Fonti e approfondimenti 1/2};
\draw[Classical,line width=1.3pt] ([xshift=.95cm,yshift=-1.68cm]current page.north west)--([xshift=2.25cm,yshift=-1.68cm]current page.north west);
\node[anchor=north west,inner sep=0,text=Classical,text width=2.25cm,font=\bfseries\fontsize{9}{12}\selectfont] at ([xshift=.95cm,yshift=-2.12cm]current page.north west) {\href{https://openintrostat.github.io/ims/}{IMS}};
\node[anchor=north west,inner sep=0,text=Ink,text width=15.6cm,font=\fontsize{11.3}{14}\selectfont] at ([xshift=3.4cm,yshift=-2.12cm]current page.north west) {Introduction to Modern Statistics, seconda edizione\par{\color{Muted}\fontsize{9.4}{12}\selectfont M. Çetinkaya-Rundel e J. Hardin}};
\node[anchor=north west,inner sep=0,text=Classical,text width=2.25cm,font=\bfseries\fontsize{9}{12}\selectfont] at ([xshift=.95cm,yshift=-3.35cm]current page.north west) {CORSO};
\node[anchor=north west,inner sep=0,text=Ink,text width=15.6cm,font=\fontsize{11.3}{14}\selectfont] at ([xshift=3.4cm,yshift=-3.35cm]current page.north west) {Materiali di partenza forniti dal docente\par{\color{Muted}\fontsize{9.4}{12}\selectfont Materiali del corso}};
\node[anchor=north west,inner sep=0,text=Classical,text width=2.25cm,font=\bfseries\fontsize{9}{12}\selectfont] at ([xshift=.95cm,yshift=-4.58cm]current page.north west) {\href{https://stat.ethz.ch/R-manual/R-devel/library/stats/html/00Index.html}{RSTATS}};
\node[anchor=north west,inner sep=0,text=Ink,text width=15.6cm,font=\fontsize{11.3}{14}\selectfont] at ([xshift=3.4cm,yshift=-4.58cm]current page.north west) {Indice del pacchetto stats\par{\color{Muted}\fontsize{9.4}{12}\selectfont R Core Team}};
\node[anchor=north west,inner sep=0,text=Muted,text width=18cm,font=\fontsize{9.8}{12.6}\selectfont] at ([xshift=.95cm,yshift=-10.55cm]current page.north west) {Consultazione: 2 novembre 2026. Riferimenti completi e provenienza in FONTI.md. Testi, esercizi e grafici didattici originali; nessuna figura dei manuali riprodotta.};
\draw[Rule,line width=.35pt] ([xshift=.95cm,yshift=-11.92cm]current page.north west)--([xshift=19.05cm,yshift=-11.92cm]current page.north west);
\node[anchor=north west,inner sep=0pt,text=Muted,text width=9.0cm,font=\fontsize{7.7}{9.3}\selectfont] at ([xshift=.95cm,yshift=-12.06cm]current page.north west) {Fonti primarie e materiali del corso};
\node[anchor=north west,inner sep=0pt,text=Classical,text width=8.2cm,align=right,font=\fontsize{7.3}{9.1}\selectfont] at ([xshift=9.45cm,yshift=-12.06cm]current page.north west) {};
\node[anchor=north east,inner sep=0pt,text=Muted,font=\fontsize{7.7}{9.3}\selectfont] at ([xshift=19.05cm,yshift=-12.06cm]current page.north west) {108/110};
\end{tikzpicture}\end{frame}
\begin{frame}[plain]\hypertarget{slide109}{}\bookmark[level=1,page=109]{Fonti e approfondimenti}\begin{tikzpicture}[remember picture,overlay]
\node[anchor=north west,inner sep=0,text=Navy,font=\bfseries\fontsize{22}{26}\selectfont] at ([xshift=.95cm,yshift=-.65cm]current page.north west) {Fonti e approfondimenti 2/2};
\draw[Classical,line width=1.3pt] ([xshift=.95cm,yshift=-1.68cm]current page.north west)--([xshift=2.25cm,yshift=-1.68cm]current page.north west);
\node[anchor=north west,inner sep=0,text=Classical,text width=2.25cm,font=\bfseries\fontsize{9}{12}\selectfont] at ([xshift=.95cm,yshift=-2.12cm]current page.north west) {\href{https://www.openintro.org/go?id=stat\_os4\_probability\_chapter}{PROB}};
\node[anchor=north west,inner sep=0,text=Ink,text width=15.6cm,font=\fontsize{11.3}{14}\selectfont] at ([xshift=3.4cm,yshift=-2.12cm]current page.north west) {OpenIntro Statistics: Probability\par{\color{Muted}\fontsize{9.4}{12}\selectfont OpenIntro}};
\node[anchor=north west,inner sep=0,text=Classical,text width=2.25cm,font=\bfseries\fontsize{9}{12}\selectfont] at ([xshift=.95cm,yshift=-3.35cm]current page.north west) {\href{https://stat.ethz.ch/R-manual/R-devel/library/stats/html/t.test.html}{RTTEST}};
\node[anchor=north west,inner sep=0,text=Ink,text width=15.6cm,font=\fontsize{11.3}{14}\selectfont] at ([xshift=3.4cm,yshift=-3.35cm]current page.north west) {Student’s t-Test — t.test\par{\color{Muted}\fontsize{9.4}{12}\selectfont R Core Team}};
\node[anchor=north west,inner sep=0,text=Classical,text width=2.25cm,font=\bfseries\fontsize{9}{12}\selectfont] at ([xshift=.95cm,yshift=-4.58cm]current page.north west) {\href{https://cran.r-project.org/doc/manuals/r-release/R-intro.html}{RBASE}};
\node[anchor=north west,inner sep=0,text=Ink,text width=15.6cm,font=\fontsize{11.3}{14}\selectfont] at ([xshift=3.4cm,yshift=-4.58cm]current page.north west) {An Introduction to R / documentazione del linguaggio\par{\color{Muted}\fontsize{9.4}{12}\selectfont R Core Team}};
\node[anchor=north west,inner sep=0,text=Muted,text width=18cm,font=\fontsize{9.8}{12.6}\selectfont] at ([xshift=.95cm,yshift=-10.55cm]current page.north west) {Consultazione: 2 novembre 2026. Riferimenti completi e provenienza in FONTI.md. Testi, esercizi e grafici didattici originali; nessuna figura dei manuali riprodotta.};
\draw[Rule,line width=.35pt] ([xshift=.95cm,yshift=-11.92cm]current page.north west)--([xshift=19.05cm,yshift=-11.92cm]current page.north west);
\node[anchor=north west,inner sep=0pt,text=Muted,text width=9.0cm,font=\fontsize{7.7}{9.3}\selectfont] at ([xshift=.95cm,yshift=-12.06cm]current page.north west) {Fonti primarie e materiali del corso};
\node[anchor=north west,inner sep=0pt,text=Classical,text width=8.2cm,align=right,font=\fontsize{7.3}{9.1}\selectfont] at ([xshift=9.45cm,yshift=-12.06cm]current page.north west) {};
\node[anchor=north east,inner sep=0pt,text=Muted,font=\fontsize{7.7}{9.3}\selectfont] at ([xshift=19.05cm,yshift=-12.06cm]current page.north west) {109/110};
\end{tikzpicture}\end{frame}
\begin{frame}[plain]\hypertarget{slide110}{}
\bookmark[level=1,page=110]{La precisione non corregge un campionamento distorto.}
\begin{tikzpicture}[remember picture,overlay]
\fill[Night] (current page.south west) rectangle (current page.north east);
\draw[Mint,line width=2pt] ([xshift=1.05cm,yshift=-1.35cm]current page.north west)--([xshift=3.0cm,yshift=-1.35cm]current page.north west);
\node[anchor=north west,inner sep=0,text=Mint,font=\bfseries\fontsize{11}{14}\selectfont] at ([xshift=1.05cm,yshift=-1.7cm]current page.north west) {DAL MODULO ALLA COMPETENZA};
\node[anchor=north west,inner sep=0,text=white,text width=15.8cm,align=left,font=\bfseries\fontsize{31}{38}\selectfont] at ([xshift=1.05cm,yshift=-2.65cm]current page.north west) {La precisione non corregge un campionamento distorto.};
\node[anchor=north west,inner sep=0,text=Lilac,text width=15.8cm,align=left,font=\fontsize{15}{21}\selectfont] at ([xshift=1.05cm,yshift=-6.4cm]current page.north west) {Confrontare gruppi e valutare ipotesi};
\node[anchor=north west,inner sep=0,text=white,text width=16.5cm,align=left,font=\fontsize{11}{15}\selectfont] at ([xshift=1.05cm,yshift=-8.6cm]current page.north west) {Riparti dalla domanda, conserva i dati, documenta il metodo e rendi controllabile la conclusione.};
\node[anchor=north east,inner sep=0,text=Mint,font=\fontsize{8}{10}\selectfont] at ([xshift=19.05cm,yshift=-.2cm]current page.north west) {\hyperlink{slide3}{Indice}};
\draw[Lilac!30,line width=.35pt] ([xshift=.95cm,yshift=-11.92cm]current page.north west)--([xshift=19.05cm,yshift=-11.92cm]current page.north west);
\node[anchor=north west,inner sep=0pt,text=Lilac,text width=9.0cm,font=\fontsize{7.7}{9.3}\selectfont] at ([xshift=.95cm,yshift=-12.06cm]current page.north west) {Modulo 05 · Stima e incertezza};
\node[anchor=north west,inner sep=0pt,text=Mint,text width=8.2cm,align=right,font=\fontsize{7.3}{9.1}\selectfont] at ([xshift=9.45cm,yshift=-12.06cm]current page.north west) {};
\node[anchor=north east,inner sep=0pt,text=Lilac,font=\fontsize{7.7}{9.3}\selectfont] at ([xshift=19.05cm,yshift=-12.06cm]current page.north west) {110/110};
\end{tikzpicture}\end{frame}
\end{document}
